
\Data\ID{1003018701}
\Updated{2022-08-25 10:08:50}
\Level{A}
\SubArea{Primitive function}
\LANGen{
\Question{
Evaluate the following integral on \( (0;\pi) \).
\[ \int\left(\pi\,\mathrm{e}^x-\frac3{\sin^2 x}\right)\mathrm{d}x \]}
{\( \pi\,\mathrm{e}^x+3\,\mathrm{cotg}\,x+c\text{, }c\in\mathbb{R} \)}{\( \pi\,\mathrm{e}^x+3\,\mathrm{tg}\,x+c\text{, }c\in\mathbb{R} \)}{\( \pi\,\mathrm{e}^x+\mathrm{cotg}\,x+c\text{, }c\in\mathbb{R} \)}{\( \pi\,\mathrm{e}^x-3\,\mathrm{cotg}\,x+c\text{, }c\in\mathbb{R} \)}{}{}}
\LANGpl{
\Question{
Oblicz całkę na przedziale \( (0;\pi) \).
\[ \int\left(\pi\,\mathrm{e}^x-\frac3{\sin^2 x}\right)\mathrm{d}x \]}
{\( \pi\,\mathrm{e}^x+3\,\mathrm{cotg}\,x+c\text{, }c\in\mathbb{R} \)}{\( \pi\,\mathrm{e}^x+3\,\mathrm{tg}\,x+c\text{, }c\in\mathbb{R} \)}{\( \pi\,\mathrm{e}^x+\mathrm{cotg}\,x+c\text{, }c\in\mathbb{R} \)}{\( \pi\,\mathrm{e}^x-3\,\mathrm{cotg}\,x+c\text{, }c\in\mathbb{R} \)}{}{}}
\LANGcs{
\Question{
Vypočítejte na intervalu \( (0;\pi) \) následující integrál.
\[ \int\left(\pi\,\mathrm{e}^x-\frac3{\sin^2 x}\right)\mathrm{d}x \]}
{\( \pi\,\mathrm{e}^x+3\,\mathrm{cotg}\,x+c\text{, }c\in\mathbb{R} \)}{\( \pi\,\mathrm{e}^x+3\,\mathrm{tg}\,x+c\text{, }c\in\mathbb{R} \)}{\( \pi\,\mathrm{e}^x+\mathrm{cotg}\,x+c\text{, }c\in\mathbb{R} \)}{\( \pi\,\mathrm{e}^x-3\,\mathrm{cotg}\,x+c\text{, }c\in\mathbb{R} \)}{}{}}
\LANGsk{
\Question{
Vypočítajte na intervale \( (0;\pi) \) nasledujúci integrál.
\[ \int\left(\pi\,\mathrm{e}^x-\frac3{\sin^2 x}\right)\mathrm{d}x \]}
{\( \pi\,\mathrm{e}^x+3\,\mathrm{cotg}\,x+c\text{, }c\in\mathbb{R} \)}{\( \pi\,\mathrm{e}^x+3\,\mathrm{tg}\,x+c\text{, }c\in\mathbb{R} \)}{\( \pi\,\mathrm{e}^x+\mathrm{cotg}\,x+c\text{, }c\in\mathbb{R} \)}{\( \pi\,\mathrm{e}^x-3\,\mathrm{cotg}\,x+c\text{, }c\in\mathbb{R} \)}{}{}}
\LANGes{
\Question{
Evalúa la siguiente integral en el intervalo \( (0;\pi) \).
\[ \int\left(\pi\,\mathrm{e}^x-\frac3{\sin^2 x}\right)\mathrm{d}x \]}
{\( \pi\,\mathrm{e}^x+3\,\mathrm{cotg}\,x+c\text{, }c\in\mathbb{R} \)}{\( \pi\,\mathrm{e}^x+3\,\mathrm{tg}\,x+c\text{, }c\in\mathbb{R} \)}{\( \pi\,\mathrm{e}^x+\mathrm{cotg}\,x+c\text{, }c\in\mathbb{R} \)}{\( \pi\,\mathrm{e}^x-3\,\mathrm{cotg}\,x+c\text{, }c\in\mathbb{R} \)}{}{}}
\End

\Data\ID{1003018702}
\Updated{2022-08-25 10:08:50}
\Level{A}
\SubArea{Primitive function}
\LANGen{
\Question{
Evaluate the following integral on \( \mathbb{R} \).
\[ \int\left(3^2+3x^2+3^x-\mathrm{e}^x+2^{\mathrm{e}}\right)\mathrm{d}x \]}
{\( 9x+x^3+\frac{3^x}{\ln3} -\mathrm{e}^x+2^{\mathrm{e}} x+c\text{, }c\in\mathbb{R} \)}{\( x^3+\frac{3^x}{\ln3} -\mathrm{e}^x+c\text{, }c\in\mathbb{R} \)}{\( 9x+3x^3+3^x-\mathrm{e}^x+\frac{2^{\mathrm{e}+1}}{\mathrm{e}+1}+c\text{, }c\in\mathbb{R} \)}{\( 9x+6x+\frac{3^x}{\ln3} -\mathrm{e}^x+2^\mathrm{e} x+c\text{, }c\in\mathbb{R} \)}{}{}}
\LANGpl{
\Question{
Oblicz całkę na przedziale \( \mathbb{R} \).
\[ \int\left(3^2+3x^2+3^x-\mathrm{e}^x+2^{\mathrm{e}}\right)\mathrm{d}x \]}
{\( 9x+x^3+\frac{3^x}{\ln3} -\mathrm{e}^x+2^{\mathrm{e}} x+c\text{, }c\in\mathbb{R} \)}{\( x^3+\frac{3^x}{\ln3} -\mathrm{e}^x+c\text{, }c\in\mathbb{R} \)}{\( 9x+3x^3+3^x-\mathrm{e}^x+\frac{2^{\mathrm{e}+1}}{\mathrm{e}+1}+c\text{, }c\in\mathbb{R} \)}{\( 9x+6x+\frac{3^x}{\ln3} -\mathrm{e}^x+2^\mathrm{e} x+c\text{, }c\in\mathbb{R} \)}{}{}}
\LANGcs{
\Question{
Vypočítejte na množině \( \mathbb{R} \) následující integrál.
\[ \int\left(3^2+3x^2+3^x-\mathrm{e}^x+2^{\mathrm{e}}\right)\mathrm{d}x \]}
{\( 9x+x^3+\frac{3^x}{\ln3} -\mathrm{e}^x+2^{\mathrm{e}} x+c\text{, }c\in\mathbb{R} \)}{\( x^3+\frac{3^x}{\ln3} -\mathrm{e}^x+c\text{, }c\in\mathbb{R} \)}{\( 9x+3x^3+3^x-\mathrm{e}^x+\frac{2^{\mathrm{e}+1}}{\mathrm{e}+1}+c\text{, }c\in\mathbb{R} \)}{\( 9x+6x+\frac{3^x}{\ln3} -\mathrm{e}^x+2^\mathrm{e} x+c\text{, }c\in\mathbb{R} \)}{}{}}
\LANGsk{
\Question{
Vypočítajte na množine \( \mathbb{R} \) nasledujúci integrál.
\[ \int\left(3^2+3x^2+3^x-\mathrm{e}^x+2^{\mathrm{e}}\right)\mathrm{d}x \]}
{\( 9x+x^3+\frac{3^x}{\ln3} -\mathrm{e}^x+2^{\mathrm{e}} x+c\text{, }c\in\mathbb{R} \)}{\( x^3+\frac{3^x}{\ln3} -\mathrm{e}^x+c\text{, }c\in\mathbb{R} \)}{\( 9x+3x^3+3^x-\mathrm{e}^x+\frac{2^{\mathrm{e}+1}}{\mathrm{e}+1}+c\text{, }c\in\mathbb{R} \)}{\( 9x+6x+\frac{3^x}{\ln3} -\mathrm{e}^x+2^\mathrm{e} x+c\text{, }c\in\mathbb{R} \)}{}{}}
\LANGes{
\Question{
Evalúa la siguiente integral en \( \mathbb{R} \).
\[ \int\left(3^2+3x^2+3^x-\mathrm{e}^x+2^{\mathrm{e}}\right)\mathrm{d}x \]}
{\( 9x+x^3+\frac{3^x}{\ln3} -\mathrm{e}^x+2^{\mathrm{e}} x+c\text{, }c\in\mathbb{R} \)}{\( x^3+\frac{3^x}{\ln3} -\mathrm{e}^x+c\text{, }c\in\mathbb{R} \)}{\( 9x+3x^3+3^x-\mathrm{e}^x+\frac{2^{\mathrm{e}+1}}{\mathrm{e}+1}+c\text{, }c\in\mathbb{R} \)}{\( 9x+6x+\frac{3^x}{\ln3} -\mathrm{e}^x+2^\mathrm{e} x+c\text{, }c\in\mathbb{R} \)}{}{}}
\End

\Data\ID{1003018703}
\Updated{2022-08-25 10:08:53}
\Level{A}
\SubArea{Primitive function}
\LANGen{
\Question{
Evaluate the following integral on \( \left(\frac{\pi}2;\pi\right) \).
\[ \int\left(7 \cos x+\frac5{\cos^2x}+\frac3{\sin^2x}\right)\mathrm{d}x \]}
{\( 7\sin x+5\,\mathrm{tg}\,x-3\,\mathrm{cotg}\,x+c\text{, }c\in\mathbb{R} \)}{\( -7\sin x+5\,\mathrm{tg}\,x+3\,\mathrm{cotg}\,x+c\text{, }c\in\mathbb{R} \)}{\( -7\sin x-5\,\mathrm{tg}\,x-3\,\mathrm{cotg}\,x+c\text{, }c\in\mathbb{R} \)}{\( 7\sin x-5\,\mathrm{tg}\,x+3\,\mathrm{cotg}\,x+c\text{, }c\in\mathbb{R} \)}{}{}}
\LANGpl{
\Question{
Oblicz całkę na przedziale \( \left(\frac{\pi}2;\pi\right) \).
\[ \int\left(7 \cos x+\frac5{\cos^2x}+\frac3{\sin^2x}\right)\mathrm{d}x \]}
{\( 7\sin x+5\,\mathrm{tg}\,x-3\,\mathrm{cotg}\,x+c\text{, }c\in\mathbb{R} \)}{\( -7\sin x+5\,\mathrm{tg}\,x+3\,\mathrm{cotg}\,x+c\text{, }c\in\mathbb{R} \)}{\( -7\sin x-5\,\mathrm{tg}\,x-3\,\mathrm{cotg}\,x+c\text{, }c\in\mathbb{R} \)}{\( 7\sin x-5\,\mathrm{tg}\,x+3\,\mathrm{cotg}\,x+c\text{, }c\in\mathbb{R} \)}{}{}}
\LANGcs{
\Question{
Vypočítejte na intervalu \( \left(\frac{\pi}2;\pi\right) \) následující integrál.
\[ \int\left(7 \cos x+\frac5{\cos^2x}+\frac3{\sin^2x}\right)\mathrm{d}x \]}
{\( 7\sin x+5\,\mathrm{tg}\,x-3\,\mathrm{cotg}\,x+c\text{, }c\in\mathbb{R} \)}{\( -7\sin x+5\,\mathrm{tg}\,x+3\,\mathrm{cotg}\,x+c\text{, }c\in\mathbb{R} \)}{\( -7\sin x-5\,\mathrm{tg}\,x-3\,\mathrm{cotg}\,x+c\text{, }c\in\mathbb{R} \)}{\( 7\sin x-5\,\mathrm{tg}\,x+3\,\mathrm{cotg}\,x+c\text{, }c\in\mathbb{R} \)}{}{}}
\LANGsk{
\Question{
Vypočítajte na intervale \( \left(\frac{\pi}2;\pi\right) \) nasledujúci integrál.
\[ \int\left(7 \cos x+\frac5{\cos^2x}+\frac3{\sin^2x}\right)\mathrm{d}x \]}
{\( 7\sin x+5\,\mathrm{tg}\,x-3\,\mathrm{cotg}\,x+c\text{, }c\in\mathbb{R} \)}{\( -7\sin x+5\,\mathrm{tg}\,x+3\,\mathrm{cotg}\,x+c\text{, }c\in\mathbb{R} \)}{\( -7\sin x-5\,\mathrm{tg}\,x-3\,\mathrm{cotg}\,x+c\text{, }c\in\mathbb{R} \)}{\( 7\sin x-5\,\mathrm{tg}\,x+3\,\mathrm{cotg}\,x+c\text{, }c\in\mathbb{R} \)}{}{}}
\LANGes{
\Question{
Evalúa la siguiente integral en el intervalo \( \left(\frac{\pi}2;\pi\right) \).
\[ \int\left(7 \cos x+\frac5{\cos^2x}+\frac3{\sin^2x}\right)\mathrm{d}x \]}
{\( 7\sin x+5\,\mathrm{tg}\,x-3\,\mathrm{cotg}\,x+c\text{, }c\in\mathbb{R} \)}{\( -7\sin x+5\,\mathrm{tg}\,x+3\,\mathrm{cotg}\,x+c\text{, }c\in\mathbb{R} \)}{\( -7\sin x-5\,\mathrm{tg}\,x-3\,\mathrm{cotg}\,x+c\text{, }c\in\mathbb{R} \)}{\( 7\sin x-5\,\mathrm{tg}\,x+3\,\mathrm{cotg}\,x+c\text{, }c\in\mathbb{R} \)}{}{}}
\End

\Data\ID{1003018704}
\Updated{2022-08-25 10:08:53}
\Level{A}
\SubArea{Primitive function}
\LANGen{
\Question{
Evaluate the following integral on \( (0;\infty) \).
\[ \int\left(3\sqrt x+4\sqrt[3]x-35\sqrt[4]{x^3}\right)\mathrm{d}x \]}
{\( 2x\sqrt x+3x\sqrt[3]x-20x\sqrt[4]{x^3}+c\text{, }c\in\mathbb{R} \)}{\( \frac92 x\sqrt x+\frac{16}3 x\sqrt[3]x-\frac{245}4 x\sqrt[4]{x^3}+c\text{, }c\in\mathbb{R} \)}{\( 2\sqrt x+3\sqrt[3]x-20\sqrt[4]{x^3}+c\text{, }c\in\mathbb{R} \)}{\( 2\sqrt[3]x+3\sqrt[4]{x^3}-20\sqrt[7]{x^4}+c\text{, }c\in\mathbb{R} \)}{}{}}
\LANGpl{
\Question{
Oblicz całkę na przedziale \( (0;\infty) \).
\[ \int\left(3\sqrt x+4\sqrt[3]x-35\sqrt[4]{x^3}\right)\mathrm{d}x \]}
{\( 2x\sqrt x+3x\sqrt[3]x-20x\sqrt[4]{x^3}+c\text{, }c\in\mathbb{R} \)}{\( \frac92 x\sqrt x+\frac{16}3 x\sqrt[3]x-\frac{245}4 x\sqrt[4]{x^3}+c\text{, }c\in\mathbb{R} \)}{\( 2\sqrt x+3\sqrt[3]x-20\sqrt[4]{x^3}+c\text{, }c\in\mathbb{R} \)}{\( 2\sqrt[3]x+3\sqrt[4]{x^3}-20\sqrt[7]{x^4}+c\text{, }c\in\mathbb{R} \)}{}{}}
\LANGcs{
\Question{
Vypočítejte na intervalu \( (0;\infty) \) následující integrál.
\[ \int\left(3\sqrt x+4\sqrt[3]x-35\sqrt[4]{x^3}\right)\mathrm{d}x \]}
{\( 2x\sqrt x+3x\sqrt[3]x-20x\sqrt[4]{x^3}+c\text{, }c\in\mathbb{R} \)}{\( \frac92 x\sqrt x+\frac{16}3 x\sqrt[3]x-\frac{245}4 x\sqrt[4]{x^3}+c\text{, }c\in\mathbb{R} \)}{\( 2\sqrt x+3\sqrt[3]x-20\sqrt[4]{x^3}+c\text{, }c\in\mathbb{R} \)}{\( 2\sqrt[3]x+3\sqrt[4]{x^3}-20\sqrt[7]{x^4}+c\text{, }c\in\mathbb{R} \)}{}{}}
\LANGsk{
\Question{
Vypočítajte na intervale \( (0;\infty) \) nasledujúci integrál.
\[ \int\left(3\sqrt x+4\sqrt[3]x-35\sqrt[4]{x^3}\right)\mathrm{d}x \]}
{\( 2x\sqrt x+3x\sqrt[3]x-20x\sqrt[4]{x^3}+c\text{, }c\in\mathbb{R} \)}{\( \frac92 x\sqrt x+\frac{16}3 x\sqrt[3]x-\frac{245}4 x\sqrt[4]{x^3}+c\text{, }c\in\mathbb{R} \)}{\( 2\sqrt x+3\sqrt[3]x-20\sqrt[4]{x^3}+c\text{, }c\in\mathbb{R} \)}{\( 2\sqrt[3]x+3\sqrt[4]{x^3}-20\sqrt[7]{x^4}+c\text{, }c\in\mathbb{R} \)}{}{}}
\LANGes{
\Question{
Evalúa la siguiente integral en el intervalo \( (0;\infty) \).
\[ \int\left(3\sqrt x+4\sqrt[3]x-35\sqrt[4]{x^3}\right)\mathrm{d}x \]}
{\( 2x\sqrt x+3x\sqrt[3]x-20x\sqrt[4]{x^3}+c\text{, }c\in\mathbb{R} \)}{\( \frac92 x\sqrt x+\frac{16}3 x\sqrt[3]x-\frac{245}4 x\sqrt[4]{x^3}+c\text{, }c\in\mathbb{R} \)}{\( 2\sqrt x+3\sqrt[3]x-20\sqrt[4]{x^3}+c\text{, }c\in\mathbb{R} \)}{\( 2\sqrt[3]x+3\sqrt[4]{x^3}-20\sqrt[7]{x^4}+c\text{, }c\in\mathbb{R} \)}{}{}}
\End

\Data\ID{1003018705}
\Updated{2022-08-25 10:08:53}
\Level{A}
\SubArea{Primitive function}
\LANGen{
\Question{
Evaluate the following integral on \( (0;\infty) \).
\[ \int\left(3\sqrt{x^7}-\sqrt[5]{x^4}\right)\mathrm{d}x \]}
{\( \frac23 x^4\sqrt x-\frac59 x\sqrt[5]{x^4}+c\text{, }c\in\mathbb{R} \)}{\( \frac32 x^4\sqrt x-\frac95 x \sqrt[5]{x^4}+c\text{, }c\in\mathbb{R} \)}{\( \frac29 x^4\sqrt x-\frac59 x\sqrt[5]{x^4}+c\text{, }c\in\mathbb{R} \)}{\( \frac23\sqrt[9]{x^2}-\frac59\sqrt[9]{x^5}+c\text{, }c\in\mathbb{R} \)}{}{}}
\LANGpl{
\Question{
Oblicz całkę na przedziale \( (0;\infty) \).
\[ \int\left(3\sqrt{x^7}-\sqrt[5]{x^4}\right)\mathrm{d}x \]}
{\( \frac23 x^4\sqrt x-\frac59 x\sqrt[5]{x^4}+c\text{, }c\in\mathbb{R} \)}{\( \frac32 x^4\sqrt x-\frac95 x \sqrt[5]{x^4}+c\text{, }c\in\mathbb{R} \)}{\( \frac29 x^4\sqrt x-\frac59 x\sqrt[5]{x^4}+c\text{, }c\in\mathbb{R} \)}{\( \frac23\sqrt[9]{x^2}-\frac59\sqrt[9]{x^5}+c\text{, }c\in\mathbb{R} \)}{}{}}
\LANGcs{
\Question{
Vypočítejte na intervalu \( (0;\infty) \) následující integrál.
\[ \int\left(3\sqrt{x^7}-\sqrt[5]{x^4}\right)\mathrm{d}x \]}
{\( \frac23 x^4\sqrt x-\frac59 x\sqrt[5]{x^4}+c\text{, }c\in\mathbb{R} \)}{\( \frac32 x^4\sqrt x-\frac95 x \sqrt[5]{x^4}+c\text{, }c\in\mathbb{R} \)}{\( \frac29 x^4\sqrt x-\frac59 x\sqrt[5]{x^4}+c\text{, }c\in\mathbb{R} \)}{\( \frac23\sqrt[9]{x^2}-\frac59\sqrt[9]{x^5}+c\text{, }c\in\mathbb{R} \)}{}{}}
\LANGsk{
\Question{
Vypočítajte na intervale \( (0;\infty) \) nasledujúci integrál.
\[ \int\left(3\sqrt{x^7}-\sqrt[5]{x^4}\right)\mathrm{d}x \]}
{\( \frac23 x^4\sqrt x-\frac59 x\sqrt[5]{x^4}+c\text{, }c\in\mathbb{R} \)}{\( \frac32 x^4\sqrt x-\frac95 x \sqrt[5]{x^4}+c\text{, }c\in\mathbb{R} \)}{\( \frac29 x^4\sqrt x-\frac59 x\sqrt[5]{x^4}+c\text{, }c\in\mathbb{R} \)}{\( \frac23\sqrt[9]{x^2}-\frac59\sqrt[9]{x^5}+c\text{, }c\in\mathbb{R} \)}{}{}}
\LANGes{
\Question{
Evalúa la siguiente integral en el intervalo \( (0;\infty) \).
\[ \int\left(3\sqrt{x^7}-\sqrt[5]{x^4}\right)\mathrm{d}x \]}
{\( \frac23 x^4\sqrt x-\frac59 x\sqrt[5]{x^4}+c\text{, }c\in\mathbb{R} \)}{\( \frac32 x^4\sqrt x-\frac95 x \sqrt[5]{x^4}+c\text{, }c\in\mathbb{R} \)}{\( \frac29 x^4\sqrt x-\frac59 x\sqrt[5]{x^4}+c\text{, }c\in\mathbb{R} \)}{\( \frac23\sqrt[9]{x^2}-\frac59\sqrt[9]{x^5}+c\text{, }c\in\mathbb{R} \)}{}{}}
\End

\Data\ID{1003018706}
\Updated{2022-08-25 10:08:53}
\Level{A}
\SubArea{Primitive function}
\LANGen{
\Question{
Evaluate the following integral on \( (0;\infty) \).
\[ \int\left(\frac5x-\frac x5\right)\mathrm{d}x \]}
{\( 5\ln|x|-\frac{x^2}{10}+c\text{, }c\in\mathbb{R} \)}{\( 5 x^0+\frac1{10}x^2+c\text{, }c\in\mathbb{R} \)}{\( 5\ln x+\frac{x^2}{10}+c\text{, }c\in\mathbb{R} \)}{\( \ln|x|+\frac{x^2}5+c\text{, }c\in\mathbb{R} \)}{}{}}
\LANGpl{
\Question{
Oblicz całkę na przedziale \( (0;\infty) \).
\[ \int\left(\frac5x-\frac x5\right)\mathrm{d}x \]}
{\( 5\ln|x|-\frac{x^2}{10}+c\text{, }c\in\mathbb{R} \)}{\( 5 x^0+\frac1{10}x^2+c\text{, }c\in\mathbb{R} \)}{\( 5\ln x+\frac{x^2}{10}+c\text{, }c\in\mathbb{R} \)}{\( \ln|x|+\frac{x^2}5+c\text{, }c\in\mathbb{R} \)}{}{}}
\LANGcs{
\Question{
Vypočítejte na intervalu \( (0;\infty) \) následující integrál.
\[ \int\left(\frac5x-\frac x5\right)\mathrm{d}x \]}
{\( 5\ln|x|-\frac{x^2}{10}+c\text{, }c\in\mathbb{R} \)}{\( 5 x^0+\frac1{10}x^2+c\text{, }c\in\mathbb{R} \)}{\( 5\ln x+\frac{x^2}{10}+c\text{, }c\in\mathbb{R} \)}{\( \ln|x|+\frac{x^2}5+c\text{, }c\in\mathbb{R} \)}{}{}}
\LANGsk{
\Question{
Vypočítajte na intervale \( (0;\infty) \) nasledujúci integrál.
\[ \int\left(\frac5x-\frac x5\right)\mathrm{d}x \]}
{\( 5\ln|x|-\frac{x^2}{10}+c\text{, }c\in\mathbb{R} \)}{\( 5 x^0+\frac1{10}x^2+c\text{, }c\in\mathbb{R} \)}{\( 5\ln x+\frac{x^2}{10}+c\text{, }c\in\mathbb{R} \)}{\( \ln|x|+\frac{x^2}5+c\text{, }c\in\mathbb{R} \)}{}{}}
\LANGes{
\Question{
Evalúa la siguiente integral en el intervalo \( (0;\infty) \).
\[ \int\left(\frac5x-\frac x5\right)\mathrm{d}x \]}
{\( 5\ln|x|-\frac{x^2}{10}+c\text{, }c\in\mathbb{R} \)}{\( 5 x^0+\frac1{10}x^2+c\text{, }c\in\mathbb{R} \)}{\( 5\ln x+\frac{x^2}{10}+c\text{, }c\in\mathbb{R} \)}{\( \ln|x|+\frac{x^2}5+c\text{, }c\in\mathbb{R} \)}{}{}}
\End

\Data\ID{1003018707}
\Updated{2022-08-25 10:08:53}
\Level{A}
\SubArea{Primitive function}
\LANGen{
\Question{
Given the function \( F(x)=\frac23\cos x-\frac{x^2}2\cdot\ln4 \), find the function \( f \) such that \( F \) is primitive to \( f \) on \(\mathbb{R} \).}
{\( f(x)=-\frac23\sin x-x\ln4 \)}{\( f(x)=\frac23\sin x-x\ln4 \)}{\( f(x)=\frac23\sin x-2x\ln2 \)}{\( f(x)=-\frac23\sin x-2x \)}{}{}}
\LANGpl{
\Question{
Dana jest funkcja \( F(x)=\frac23\cos x-\frac{x^2}2\cdot\ln4 \), wyznacz funkcję  \( f \) tak, aby \( F \) była funkcją pierwotną   \( f \) w \(\mathbb{R} \).}
{\( f(x)=-\frac23\sin x-x\ln4 \)}{\( f(x)=\frac23\sin x-x\ln4 \)}{\( f(x)=\frac23\sin x-2x\ln2 \)}{\( f(x)=-\frac23\sin x-2x \)}{}{}}
\LANGcs{
\Question{
Je dána funkce \( F(x)=\frac23\cos x-\frac{x^2}2\cdot\ln4 \). Najděte takovou funkci \( f \), k níž je  \( F \) na \(\mathbb{R} \) primitivní.}
{\( f(x)=-\frac23\sin x-x\ln4 \)}{\( f(x)=\frac23\sin x-x\ln4 \)}{\( f(x)=\frac23\sin x-2x\ln2 \)}{\( f(x)=-\frac23\sin x-2x \)}{}{}}
\LANGsk{
\Question{
Je daná funkcia \( F(x)=\frac23\cos x-\frac{x^2}2\cdot\ln4 \). Nájdite takú funkciu \( f \), ku ktorej je  \( F \) na \(\mathbb{R} \) primitívna.}
{\( f(x)=-\frac23\sin x-x\ln4 \)}{\( f(x)=\frac23\sin x-x\ln4 \)}{\( f(x)=\frac23\sin x-2x\ln2 \)}{\( f(x)=-\frac23\sin x-2x \)}{}{}}
\LANGes{
\Question{
Dada la función \( F(x)=\frac23\cos x-\frac{x^2}2\cdot\ln4 \), encuentra la función \( f \) tal que \( F \) es primitiva de \( f \) en \(\mathbb{R} \).}
{\( f(x)=-\frac23\sin x-x\ln4 \)}{\( f(x)=\frac23\sin x-x\ln4 \)}{\( f(x)=\frac23\sin x-2x\ln2 \)}{\( f(x)=-\frac23\sin x-2x \)}{}{}}
\End

\Data\ID{1003018708}
\Updated{2022-08-25 10:08:53}
\Level{A}
\SubArea{Primitive function}
\LANGen{
\Question{
Given the function \( F(x)= 5\,\mathrm{e}^x+2\sqrt x \), find the function \( f \) such that \( F \) is primitive to \( f \) on \( (0;\infty) \).}
{\( f(x)=5\,\mathrm{e}^x+\frac{\sqrt x}x \)}{\( f(x)=5\,\mathrm{e}^x+\sqrt x \)}{\( f(x)=\mathrm{e}^x+\frac{\sqrt x}x \)}{\( f(x)=5\,\mathrm{e}^x-\frac{\sqrt x}x \)}{}{}}
\LANGpl{
\Question{
Dana jest funkcja \( F(x)= 5\,\mathrm{e}^x+2\sqrt x \), wskaż funkcję \( f \) tak, aby \( F \) była funkcją pierwotną \( f \) w \( (0;\infty) \).}
{\( f(x)=5\,\mathrm{e}^x+\frac{\sqrt x}x \)}{\( f(x)=5\,\mathrm{e}^x+\sqrt x \)}{\( f(x)=\mathrm{e}^x+\frac{\sqrt x}x \)}{\( f(x)=5\,\mathrm{e}^x-\frac{\sqrt x}x \)}{}{}}
\LANGcs{
\Question{
Je dána funkce \( F(x)= 5\,\mathrm{e}^x+2\sqrt x \). Najděte takovou funkci \( f \), k níž je \( F \) na \( (0;\infty) \) primitivní.}
{\( f(x)=5\,\mathrm{e}^x+\frac{\sqrt x}x \)}{\( f(x)=5\,\mathrm{e}^x+\sqrt x \)}{\( f(x)=\mathrm{e}^x+\frac{\sqrt x}x \)}{\( f(x)=5\,\mathrm{e}^x-\frac{\sqrt x}x \)}{}{}}
\LANGsk{
\Question{
Je daná funkcia \( F(x)= 5\,\mathrm{e}^x+2\sqrt x \). Nájdite takú funkciu \( f \), ku ktorej je \( F \) na \( (0;\infty) \) primitívna.}
{\( f(x)=5\,\mathrm{e}^x+\frac{\sqrt x}x \)}{\( f(x)=5\,\mathrm{e}^x+\sqrt x \)}{\( f(x)=\mathrm{e}^x+\frac{\sqrt x}x \)}{\( f(x)=5\,\mathrm{e}^x-\frac{\sqrt x}x \)}{}{}}
\LANGes{
\Question{
Dada la función \( F(x)= 5\,\mathrm{e}^x+2\sqrt x \), encuentra la función \( f \) tal que \( F \) es primitiva de \( f \) en \( (0;\infty) \).}
{\( f(x)=5\,\mathrm{e}^x+\frac{\sqrt x}x \)}{\( f(x)=5\,\mathrm{e}^x+\sqrt x \)}{\( f(x)=\mathrm{e}^x+\frac{\sqrt x}x \)}{\( f(x)=5\,\mathrm{e}^x-\frac{\sqrt x}x \)}{}{}}
\End

\Data\ID{1003018709}
\Updated{2022-08-25 10:08:53}
\Level{A}
\SubArea{Primitive function}
\LANGen{
\Question{
Four children evaluated the following integral \( I \) on \( \mathbb{R} \). Who made a mistake?
\[ I =\int\left(5x^4+9x^2-6x\right)\mathrm{d}x \]}
{Paul: \( I =\left(x^3+3\right)\cdot x^2-3x^3+c\text{, }c\in\mathbb{R} \)}{Jane: \( I =x^2\left(x^3+3x-3\right)+c\text{, }c\in\mathbb{R} \)}{Ann: \( I =x^5+3x^3-3x^2+c\text{, }c\in\mathbb{R} \)}{Miky: \( I=x^5+(3x-3)\cdot x^2+c\text{, }c\in\mathbb{R} \)}{}{}}
\LANGpl{
\Question{
Czterech uczniów wyznaczało całkę \( I \) na \( \mathbb{R} \). Kto popełnił błąd?
\[ I =\int\left(5x^4+9x^2-6x\right)\mathrm{d}x \]}
{Paweł: \( I =\left(x^3+3\right)\cdot x^2-3x^3+c\text{, }c\in\mathbb{R} \)}{Jan: \( I =x^2\left(x^3+3x-3\right)+c\text{, }c\in\mathbb{R} \)}{Anna: \( I =x^5+3x^3-3x^2+c\text{, }c\in\mathbb{R} \)}{Michał: \( I=x^5+(3x-3)\cdot x^2+c\text{, }c\in\mathbb{R} \)}{}{}}
\LANGcs{
\Question{
Čtyři studenti počítali následující integrál  \( I \) na \( \mathbb{R} \). Kdo udělal chybu?
\[ I =\int\left(5x^4+9x^2-6x\right)\mathrm{d}x \]}
{Pavel: \( I =\left(x^3+3\right)\cdot x^2-3x^3+c\text{, }c\in\mathbb{R} \)}{Jana: \( I =x^2\left(x^3+3x-3\right)+c\text{, }c\in\mathbb{R} \)}{Anna: \( I =x^5+3x^3-3x^2+c\text{, }c\in\mathbb{R} \)}{Mikuláš: \( I=x^5+(3x-3)\cdot x^2+c\text{, }c\in\mathbb{R} \)}{}{}}
\LANGsk{
\Question{
Štyria študenti počítali následujúci integrál  \( I \) na \( \mathbb{R} \). Kto urobil chybu?
\[ I =\int\left(5x^4+9x^2-6x\right)\mathrm{d}x \]}
{Pavel: \( I =\left(x^3+3\right)\cdot x^2-3x^3+c\text{, }c\in\mathbb{R} \)}{Jana: \( I =x^2\left(x^3+3x-3\right)+c\text{, }c\in\mathbb{R} \)}{Anna: \( I =x^5+3x^3-3x^2+c\text{, }c\in\mathbb{R} \)}{Mikuláš: \( I=x^5+(3x-3)\cdot x^2+c\text{, }c\in\mathbb{R} \)}{}{}}
\LANGes{
\Question{
Cuatro chicos han evaluado la siguiente integral \( I \) en \( \mathbb{R} \). ¿Quién ha cometido un error?
\[ I =\int\left(5x^4+9x^2-6x\right)\mathrm{d}x \]}
{Pablo: \( I =\left(x^3+3\right)\cdot x^2-3x^3+c\text{, }c\in\mathbb{R} \)}{Juana: \( I =x^2\left(x^3+3x-3\right)+c\text{, }c\in\mathbb{R} \)}{Ana: \( I =x^5+3x^3-3x^2+c\text{, }c\in\mathbb{R} \)}{Miguel: \( I=x^5+(3x-3)\cdot x^2+c\text{, }c\in\mathbb{R} \)}{}{}}
\End

\Data\ID{1003018710}
\Updated{2022-08-25 10:08:53}
\Level{A}
\SubArea{Primitive function}
\LANGen{
\Question{
Four children evaluated the following integral \( I \) on \( (0;\infty) \). Who made a mistake?
\[ I =\int\left(\frac18\sqrt[8]{x^3}+\frac12\sqrt{x^9}-\frac15\sqrt[5]{x^6} \right)\mathrm{d}x \]}
{Paul: \( I =\frac1{11}\left(x^3\sqrt[8]x+x\sqrt{x^5}-x\sqrt[5]{x^2}\right)+c\text{, }c\in\mathbb{R} \)}{Jane: \( I  =\frac1{11}\left(x\sqrt[8]{x^3}+x^5\sqrt x-x^2\sqrt[5] x+c\right)\text{, }c\in\mathbb{R} \)}{Ann: \( I =\frac1{11}\left(x\sqrt[8]{x^3}+x^5\sqrt x-x^2\sqrt[5]x\right)+c\text{, }c\in\mathbb{R} \)}{Miky: \( I =\frac x{11}\sqrt[8]{x^3}+\frac{x^5}{11}\sqrt x-\frac{x^2}{11}\sqrt[5]x+c\text{, }c\in\mathbb{R} \)}{}{}}
\LANGpl{
\Question{
Czterech uczniów wyznaczyło całkę \( I \) w \( (0;\infty) \). Kto popełnił błąd?
\[ I =\int\left(\frac18\sqrt[8]{x^3}+\frac12\sqrt{x^9}-\frac15\sqrt[5]{x^6} \right)\mathrm{d}x \]}
{Paweł: \( I =\frac1{11}\left(x^3\sqrt[8]x+x\sqrt{x^5}-x\sqrt[5]{x^2}\right)+c\text{, }c\in\mathbb{R} \)}{Jan: \( I  =\frac1{11}\left(x\sqrt[8]{x^3}+x^5\sqrt x-x^2\sqrt[5] x+c\right)\text{, }c\in\mathbb{R} \)}{Anna: \( I =\frac1{11}\left(x\sqrt[8]{x^3}+x^5\sqrt x-x^2\sqrt[5]x\right)+c\text{, }c\in\mathbb{R} \)}{Michał: \( I =\frac x{11}\sqrt[8]{x^3}+\frac{x^5}{11}\sqrt x-\frac{x^2}{11}\sqrt[5]x+c\text{, }c\in\mathbb{R} \)}{}{}}
\LANGcs{
\Question{
Čtyři studenti počítali následující integrál  \( I \) na  \( (0;\infty) \). Kdo udělal chybu?
\[ I =\int\left(\frac18\sqrt[8]{x^3}+\frac12\sqrt{x^9}-\frac15\sqrt[5]{x^6} \right)\mathrm{d}x \]}
{Pavel: \( I =\frac1{11}\left(x^3\sqrt[8]x+x\sqrt{x^5}-x\sqrt[5]{x^2}\right)+c\text{, }c\in\mathbb{R} \)}{Jana: \( I  =\frac1{11}\left(x\sqrt[8]{x^3}+x^5\sqrt x-x^2\sqrt[5] x+c\right)\text{, }c\in\mathbb{R} \)}{Anna: \( I =\frac1{11}\left(x\sqrt[8]{x^3}+x^5\sqrt x-x^2\sqrt[5]x\right)+c\text{, }c\in\mathbb{R} \)}{Mikuláš: \( I =\frac x{11}\sqrt[8]{x^3}+\frac{x^5}{11}\sqrt x-\frac{x^2}{11}\sqrt[5]x+c\text{, }c\in\mathbb{R} \)}{}{}}
\LANGsk{
\Question{
Štyria študenti počítali následujúci integrál  \( I \) na  \( (0;\infty) \). Kto urobil chybu?
\[ I =\int\left(\frac18\sqrt[8]{x^3}+\frac12\sqrt{x^9}-\frac15\sqrt[5]{x^6} \right)\mathrm{d}x \]}
{Pavol: \( I =\frac1{11}\left(x^3\sqrt[8]x+x\sqrt{x^5}-x\sqrt[5]{x^2}\right)+c\text{, }c\in\mathbb{R} \)}{Jana: \( I  =\frac1{11}\left(x\sqrt[8]{x^3}+x^5\sqrt x-x^2\sqrt[5] x+c\right)\text{, }c\in\mathbb{R} \)}{Anna: \( I =\frac1{11}\left(x\sqrt[8]{x^3}+x^5\sqrt x-x^2\sqrt[5]x\right)+c\text{, }c\in\mathbb{R} \)}{Mikuláš: \( I =\frac x{11}\sqrt[8]{x^3}+\frac{x^5}{11}\sqrt x-\frac{x^2}{11}\sqrt[5]x+c\text{, }c\in\mathbb{R} \)}{}{}}
\LANGes{
\Question{
Cuatro chicos han evaluado la siguiente integral \( I \) en \( (0;\infty) \). ¿Quién ha cometido un error? 
\[ I =\int\left(\frac18\sqrt[8]{x^3}+\frac12\sqrt{x^9}-\frac15\sqrt[5]{x^6} \right)\mathrm{d}x \]}
{Pablo: \( I =\frac1{11}\left(x^3\sqrt[8]x+x\sqrt{x^5}-x\sqrt[5]{x^2}\right)+c\text{, }c\in\mathbb{R} \)}{Juana: \( I  =\frac1{11}\left(x\sqrt[8]{x^3}+x^5\sqrt x-x^2\sqrt[5] x+c\right)\text{, }c\in\mathbb{R} \)}{Ana: \( I =\frac1{11}\left(x\sqrt[8]{x^3}+x^5\sqrt x-x^2\sqrt[5]x\right)+c\text{, }c\in\mathbb{R} \)}{Miguel: \( I =\frac x{11}\sqrt[8]{x^3}+\frac{x^5}{11}\sqrt x-\frac{x^2}{11}\sqrt[5]x+c\text{, }c\in\mathbb{R} \)}{}{}}
\End

\Data\ID{1003027304}
\Updated{2022-08-25 10:08:58}
\Level{A}
\SubArea{Primitive function}
\LANGen{
\Question{
Choose a pair of functions, \( f_1 \) and \( f_2 \), that are antiderivatives of  the same function on \( \mathbb{R} \).}
{\( f_1(x) = 3+\sin x\text{, }f_2(x)=\cos\left(\frac32\pi+x\right) \)}{\( f_1(x) = 5+\sin x\text{, }f_2(x)=-\cos x \)}{\( f_1(x) = \sin(x+\pi)\text{, }f_2(x)=\sin x \)}{\( f_1(x) = \cos x\text{, }f_2(x)=-\cos x \)}{}{}}
\LANGpl{
\Question{
Wybierz parę funkcji, \( f_1 \) i \( f_2 \) tak, aby były funkcjami pierwotnymi tej samej funkcji w \( \mathbb{R} \).}
{\( f_1(x) = 3+\sin x\text{, }f_2(x)=\cos\left(\frac32\pi+x\right) \)}{\( f_1(x) = 5+\sin x\text{, }f_2(x)=-\cos x \)}{\( f_1(x) = \sin(x+\pi)\text{, }f_2(x)=\sin x \)}{\( f_1(x) = \cos x\text{, }f_2(x)=-\cos x \)}{}{}}
\LANGcs{
\Question{
Vyber dvojici funkcí  \( f_1 \) a \( f_2 \) primitivních na  \( \mathbb{R} \) k téže funkci.}
{\( f_1(x) = 3+\sin x\text{, }f_2(x)=\cos\left(\frac32\pi+x\right) \)}{\( f_1(x) = 5+\sin x\text{, }f_2(x)=-\cos x \)}{\( f_1(x) = \sin(x+\pi)\text{, }f_2(x)=\sin x \)}{\( f_1(x) = \cos x\text{, }f_2(x)=-\cos x \)}{}{}}
\LANGsk{
\Question{
Vyber dvojicu funkcií  \( f_1 \) a \( f_2 \) primitívnych na  \( \mathbb{R} \) k tej istej funkcii.}
{\( f_1(x) = 3+\sin x\text{, }f_2(x)=\cos\left(\frac32\pi+x\right) \)}{\( f_1(x) = 5+\sin x\text{, }f_2(x)=-\cos x \)}{\( f_1(x) = \sin(x+\pi)\text{, }f_2(x)=\sin x \)}{\( f_1(x) = \cos x\text{, }f_2(x)=-\cos x \)}{}{}}
\LANGes{
\Question{
Elige la pareja de las funciones \( f_1 \) y \( f_2 \), que son antiderivadas de la misma función en \( \mathbb{R} \).}
{\( f_1(x) = 3+\sin x\text{, }f_2(x)=\cos\left(\frac32\pi+x\right) \)}{\( f_1(x) = 5+\sin x\text{, }f_2(x)=-\cos x \)}{\( f_1(x) = \sin(x+\pi)\text{, }f_2(x)=\sin x \)}{\( f_1(x) = \cos x\text{, }f_2(x)=-\cos x \)}{}{}}
\End

\Data\ID{1003107801}
\Updated{2020-07-15 03:07:12}
\Level{A}
\SubArea{Primitive function}
\LANGen{
\Question{
Identify the difference of the functions $F(x)=\frac{\sin^2x}2$ and $G(x)=-0.25\cdot\cos(2x)$, that are primitive to the same function $f(x)$.}
{$\frac14$}{$\frac18$}{$\frac12$}{$1$}{}{}}
\LANGpl{
\Question{
Określ różnicę funkcji  $F(x)=\frac{\sin^2x}2$ i $G(x)=-0{,}25\cdot\cos(2x)$, funkcje są funkcjami pierwotnymi tej samej funkcji $f(x)$.}
{$\frac14$}{$\frac18$}{$\frac12$}{$1$}{}{}}
\LANGcs{
\Question{
O jakou konstantu se liší dvě funkce $F(x)=\frac{\sin^2x}2$ a $G(x)=-0{,}25\cdot\cos(2x)$, které jsou primitivní k téže funkci $f(x)$?}
{o $\frac14$}{o $\frac18$}{o $\frac12$}{o $1$}{}{}}
\LANGsk{
\Question{
O akú konštantu sa líšia dve funkcie $F(x)=\frac{\sin^2x}2$ a $G(x)=-0{,}25\cdot\cos(2x)$, ktoré sú primitívne k tej istej funkcii $f(x)$?}
{o $\frac14$}{o $\frac18$}{o $\frac12$}{o $1$}{}{}}
\LANGes{
\Question{
Identifica la diferencia de las funciones $F(x)=\frac{\sin^2x}2$ y $G(x)=-0.25\cdot\cos(2x)$, que son primitivas de la misma función $f(x)$.}
{$\frac14$}{$\frac18$}{$\frac12$}{$1$}{}{}}
\End

\Data\ID{1003107805}
\Updated{2020-07-15 03:07:12}
\Level{A}
\SubArea{Primitive function}
\LANGen{
\Question{
Define the function $f(x)$ so that it holds: $f'(x)=x^5-\sqrt[4]x$ on $(0;\infty)$, $f(1)=-1$.}
{$f(x)=\frac{x^6}6-\frac45x\sqrt[4]x-\frac{11}{30}$}{$f(x)=\frac{x^6}6-\frac45\sqrt[4]{x^5}+\frac{11}{30}$}{$f(x)=\frac{x^6}6-\frac54x\sqrt[4]x-\frac{11}{30}$}{$f(x)=\frac{x^6}6-\frac54x\sqrt[4]x+\frac{11}{30}$}{}{}}
\LANGpl{
\Question{
Określ funkcję  $f(x)$ tak, aby: $f'(x)=x^5-\sqrt[4]x$ w $(0;\infty)$, $f(1)=-1$.}
{$f(x)=\frac{x^6}6-\frac45x\sqrt[4]x-\frac{11}{30}$}{$f(x)=\frac{x^6}6-\frac45\sqrt[4]{x^5}+\frac{11}{30}$}{$f(x)=\frac{x^6}6-\frac54x\sqrt[4]x-\frac{11}{30}$}{$f(x)=\frac{x^6}6-\frac54x\sqrt[4]x+\frac{11}{30}$}{}{}}
\LANGcs{
\Question{
Určete funkci $f(x)$ tak, aby platilo: $f'(x)=x^5-\sqrt[4]x$ na $(0;\infty)\wedge f(1)=-1$.}
{$f(x)=\frac{x^6}6-\frac45x\sqrt[4]x-\frac{11}{30}$}{$f(x)=\frac{x^6}6-\frac45\sqrt[4]{x^5}+\frac{11}{30}$}{$f(x)=\frac{x^6}6-\frac54x\sqrt[4]x-\frac{11}{30}$}{$f(x)=\frac{x^6}6-\frac54x\sqrt[4]x+\frac{11}{30}$}{}{}}
\LANGsk{
\Question{
Určte funkciu $f(x)$ tak, aby platilo: $f'(x)=x^5-\sqrt[4]x$ na $(0;\infty)\wedge f(1)=-1$.}
{$f(x)=\frac{x^6}6-\frac45x\sqrt[4]x-\frac{11}{30}$}{$f(x)=\frac{x^6}6-\frac45\sqrt[4]{x^5}+\frac{11}{30}$}{$f(x)=\frac{x^6}6-\frac54x\sqrt[4]x-\frac{11}{30}$}{$f(x)=\frac{x^6}6-\frac54x\sqrt[4]x+\frac{11}{30}$}{}{}}
\LANGes{
\Question{
Define la función $f(x)$ para que valga: $f'(x)=x^5-\sqrt[4]x$ on $(0;\infty)$, $f(1)=-1$.}
{$f(x)=\frac{x^6}6-\frac45x\sqrt[4]x-\frac{11}{30}$}{$f(x)=\frac{x^6}6-\frac45\sqrt[4]{x^5}+\frac{11}{30}$}{$f(x)=\frac{x^6}6-\frac54x\sqrt[4]x-\frac{11}{30}$}{$f(x)=\frac{x^6}6-\frac54x\sqrt[4]x+\frac{11}{30}$}{}{}}
\End

\Data\ID{1003107806}
\Updated{2021-08-04 07:08:43}
\Level{A}
\SubArea{Primitive function}
\LANGen{
\Question{
Define the function $f(x)$ so that the following holds: $f''(x)=\mathrm{e}^x+x^5$ on $\mathbb{R}$, $f(0)=1$, and $f(1)=\frac{43}{42}$.}
{$f(x)=\mathrm{e}^x+\frac{x^7}{42}+(1-\mathrm{e})x$}{$f(x)=\mathrm{e}^x+\frac{x^7}{42}+(-\mathrm{e}-1)x$}{$f(x)=\mathrm{e}^x+\frac{7}{6}x^7+x-\mathrm{e}x$}{$f(x)=\mathrm{e}^x+\frac{x^7}{42}+\frac{43}{42}$}{}{}}
\LANGpl{
\Question{
Określ funkcję  $f(x)$ tak, aby: $f''(x)=\mathrm{e}^x+x^5$ w $\mathbb{R}$, $f(0)=1$, i $f(1)=\frac{43}{42}$.}
{$f(x)=\mathrm{e}^x+\frac{x^7}{42}+(1-\mathrm{e})x$}{$f(x)=\mathrm{e}^x+\frac{x^7}{42}+(-\mathrm{e}-1)x$}{$f(x)=\mathrm{e}^x+\frac{7}{6}x^7+x-\mathrm{e}x$}{$f(x)=\mathrm{e}^x+\frac{x^7}{42}+\frac{43}{42}$}{}{}}
\LANGcs{
\Question{
Určete funkci $f(x)$ tak, aby platilo:  $f''(x)=\mathrm{e}^x+x^5$ na $\mathbb{R} $,  $ f(0)=1$ a $f(1)=\frac{43}{42}$.}
{$f(x)=\mathrm{e}^x+\frac{x^7}{42}+(1-\mathrm{e})x$}{$f(x)=\mathrm{e}^x+\frac{x^7}{42}+(-\mathrm{e}-1)x$}{$f(x)=\mathrm{e}^x+\frac{7}{6}x^7+x-\mathrm{e}x$}{$f(x)=\mathrm{e}^x+\frac{x^7}{42}+\frac{43}{42}$}{}{}}
\LANGsk{
\Question{
Určte funkciu $f(x)$ tak, aby platilo:  $f''(x)=\mathrm{e}^x+x^5$ na $\mathbb{R} $,  $ f(0)=1$ a $f(1)=\frac{43}{42}$.}
{$f(x)=\mathrm{e}^x+\frac{x^7}{42}+(1-\mathrm{e})x$}{$f(x)=\mathrm{e}^x+\frac{x^7}{42}+(-\mathrm{e}-1)x$}{$f(x)=\mathrm{e}^x+\frac{7}{6}x^7+x-\mathrm{e}x$}{$f(x)=\mathrm{e}^x+\frac{x^7}{42}+\frac{43}{42}$}{}{}}
\LANGes{
\Question{
Define la función $f(x)$ para que valga lo siguiente: $f''(x)=\mathrm{e}^x+x^5$ on $\mathbb{R}$, $f(0)=1$, y $f(1)=\frac{43}{42}$.}
{$f(x)=\mathrm{e}^x+\frac{x^7}{42}+(1-\mathrm{e})x$}{$f(x)=\mathrm{e}^x+\frac{x^7}{42}+(-\mathrm{e}-1)x$}{$f(x)=\mathrm{e}^x+\frac{7}{6}x^7+x-\mathrm{e}x$}{$f(x)=\mathrm{e}^x+\frac{x^7}{42}+\frac{43}{42}$}{}{}}
\End

\Data\ID{1003107807}
\Updated{2020-07-15 03:07:12}
\Level{A}
\SubArea{Primitive function}
\LANGen{
\Question{
Find a function $F(x)$, that is primitive to the function $f(x)=2^x\cdot\ln2+4^x\cdot2\ln2+8^x\cdot3\ln2$ on $\mathbb{R}$, and satisfies the condition $F(0)=5$.}
{$F(x)=2^x+4^x+8^x+2$}{$F(x)=\frac{2^x}{\ln 2}+\frac{4^x}{\ln 4}+\frac{8^x}{\ln 8}+2x$}{$F(x)=2^x+4^x+8^x+5$}{$F(x)=2^x\cdot\ln2+2^{x+1}\cdot\ln2+2^{x+3}\cdot\ln2+5$}{}{}}
\LANGpl{
\Question{
Wskaż funkcję  $F(x)$, która jest funkcją pierwotną funkcji $f(x)=2^x\cdot\ln2+4^x\cdot2\ln2+8^x\cdot3\ln2$ dla $\mathbb{R}$ oraz spełnia warunek $F(0)=5$.}
{$F(x)=2^x+4^x+8^x+2$}{$F(x)=\frac{2^x}{\ln 2}+\frac{4^x}{\ln 4}+\frac{8^x}{\ln 8}+2x$}{$F(x)=2^x+4^x+8^x+5$}{$F(x)=2^x\cdot\ln2+2^{x+1}\cdot\ln2+2^{x+3}\cdot\ln2+5$}{}{}}
\LANGcs{
\Question{
Určete takovou funkci $F(x)$, která je primitivní k funkci $f(x)=2^x\cdot\ln2+4^x\cdot2\ln2+8^x\cdot3\ln2$ na $\mathbb{R}$ a splňuje podmínku $F(0)=5$.}
{$F(x)=2^x+4^x+8^x+2$}{$F(x)=\frac{2^x}{\ln 2}+\frac{4^x}{\ln 4}+\frac{8^x}{\ln 8}+2x$}{$F(x)=2^x+4^x+8^x+5$}{$F(x)=2^x\cdot\ln2+2^{x+1}\cdot\ln2+2^{x+3}\cdot\ln2+5$}{}{}}
\LANGsk{
\Question{
Určte takú funkciu $F(x)$, ktorá je primitívna k funkcii $f(x)=2^x\cdot\ln2+4^x\cdot2\ln2+8^x\cdot3\ln2$ na $\mathbb{R}$ a splňuje podmienku $F(0)=5$.}
{$F(x)=2^x+4^x+8^x+2$}{$F(x)=\frac{2^x}{\ln 2}+\frac{4^x}{\ln 4}+\frac{8^x}{\ln 8}+2x$}{$F(x)=2^x+4^x+8^x+5$}{$F(x)=2^x\cdot\ln2+2^{x+1}\cdot\ln2+2^{x+3}\cdot\ln2+5$}{}{}}
\LANGes{
\Question{
Encuentra la función $F(x)$, que es primitiva de la función $f(x)=2^x\cdot\ln2+4^x\cdot2\ln2+8^x\cdot3\ln2$ en $\mathbb{R}$, y cumpla la condición $F(0)=5$.}
{$F(x)=2^x+4^x+8^x+2$}{$F(x)=\frac{2^x}{\ln 2}+\frac{4^x}{\ln 4}+\frac{8^x}{\ln 8}+2x$}{$F(x)=2^x+4^x+8^x+5$}{$F(x)=2^x\cdot\ln2+2^{x+1}\cdot\ln2+2^{x+3}\cdot\ln2+5$}{}{}}
\End

\Data\ID{1103027101}
\Updated{2022-05-29 12:05:45}
\Level{A}
\SubArea{Primitive function}
\LANGen{
\Question{
In which picture can you see a pair of graphs of functions \( f_1 \) and \( f_2 \) that are primitive to the same function?}
{}{}{}{}{}{}}
\LANGpl{
\Question{
Na którym rysunku znajdują się wykresy funkcji \( f_1 \) i \( f_2 \) funkcje te są funkcjami pierwotnymi tej samej funkcji?}
{}{}{}{}{}{}}
\LANGcs{
\Question{
Na kterém obrázku je dvojice grafů funkcí \( f_1 \) a \( f_2 \) primitivních k téže funkci?}
{}{}{}{}{}{}}
\LANGsk{
\Question{
Na ktorom obrázku je dvojica grafov funkcií \( f_1 \) a \( f_2 \) primitívnych k tej istej funkcii?}
{}{}{}{}{}{}}
\LANGes{
\Question{
¿En qué imagen puedes ver las funciones \( f_1 \) y \( f_2 \) que son primitivas de la misma función?}
{}{}{}{}{}{}}

\IMGanswer{1103027101a.pdf}
\IMGanswer{1103027101b.pdf}
\IMGanswer{1103027101c.pdf}
\IMGanswer{1103027101d.pdf}\End

\Data\ID{1103027102}
\Updated{2020-07-15 03:07:10}
\Level{A}
\SubArea{Primitive function}
\LANGen{
\Question{
In which picture can you see a pair of graphs of functions \( f_1 \) and \( f_2 \) that are primitive to the same function?}
{}{}{}{}{}{}}
\LANGpl{
\Question{
Na którym rysunku znajdują się wykresy funkcji \( f_1 \) i \( f_2 \) funkcje te są funkcjami pierwotnymi tej samej funkcji?}
{}{}{}{}{}{}}
\LANGcs{
\Question{
Na kterém obrázku je dvojice grafů funkcí \( f_1 \) a \( f_2 \) primitivních k téže funkci?}
{}{}{}{}{}{}}
\LANGsk{
\Question{
Na ktorom obrázku je dvojica grafov funkcií \( f_1 \) a \( f_2 \) primitívnych k tej istej funkcii?}
{}{}{}{}{}{}}
\LANGes{
\Question{
¿En qué imagen puedes ver gráficas de las funciones \( f_1 \) y \( f_2 \) que son primitivas de la misma función?}
{}{}{}{}{}{}}

\IMGanswer{1103027102a.pdf}
\IMGanswer{1103027102b.pdf}
\IMGanswer{1103027102c.pdf}
\IMGanswer{1103027102d.pdf}\End

\Data\ID{1103027103}
\Updated{2020-07-15 03:07:21}
\Level{A}
\SubArea{Primitive function}
\IMG{1103027103.pdf}
\LANGen{
\Question{
In the picture there are graphs of four functions. Which of the functions is primitive to the function \( f(x)=3 \) on \( \mathbb{R} \)?}
{\( f_1 \)}{\( f_2 \)}{\( f_3 \)}{\( f_4 \)}{}{}}
\LANGpl{
\Question{
Na rysunku znajdują się cztery wykresy funkcji.  Wskaż funkcję, która jest funkcją pierwotną funkcji \( f(x)=3 \) na \( \mathbb{R} \).}
{\( f_1 \)}{\( f_2 \)}{\( f_3 \)}{\( f_4 \)}{}{}}
\LANGcs{
\Question{
Na obrázku jsou grafy čtyř funkcí. Která z nich je na \( \mathbb{R} \) primitivní k funkci \( f(x)=3 \)?}
{\( f_1 \)}{\( f_2 \)}{\( f_3 \)}{\( f_4 \)}{}{}}
\LANGsk{
\Question{
Na obrázku sú grafy štyroch funkcií. Ktorá z nich je na \( \mathbb{R} \) primitívna k funkcii \( f(x)=3 \)?}
{\( f_1 \)}{\( f_2 \)}{\( f_3 \)}{\( f_4 \)}{}{}}
\LANGes{
\Question{
En la imagen hay gráficas de cuatro funciones. ¿Cuál de las funciones son primitivas de la función \( f(x)=3 \) en \( \mathbb{R} \)?}
{\( f_1 \)}{\( f_2 \)}{\( f_3 \)}{\( f_4 \)}{}{}}
\End

\Data\ID{1103027104}
\Updated{2020-07-15 03:07:21}
\Level{A}
\SubArea{Primitive function}
\IMG{1103027104.pdf}
\LANGen{
\Question{
In the picture there are graphs of five functions. Which of the functions is primitive to the function \( f(x)=x^2 \) on \( \mathbb{R} \)?}
{all of \( f_2 \), \( f_3 \) and \( f_4 \)}{only \( f_5 \)}{only \( f_1 \)}{only \( f_3 \) and \( f_4 \)}{}{}}
\LANGpl{
\Question{
Na rysunku znajduje się pięć wykresów funkcji.  Wskaż funkcję, która jest funkcją pierwotną funkcji  \( f(x)=x^2 \) na \( \mathbb{R} \).}
{wszystkie z \( f_2 \), \( f_3 \) i \( f_4 \)}{tylko \( f_5 \)}{tylko \( f_1 \)}{tylko \( f_3 \) i \( f_4 \)}{}{}}
\LANGcs{
\Question{
Na obrázku jsou grafy pěti funkcí. Která z nich je na \( \mathbb{R} \) primitivní k funkci \( f(x)=x^2 \)?}
{\( f_2 \), \( f_3 \) i \( f_4 \)}{jen \( f_5 \)}{jen \( f_1 \)}{jen \( f_3 \) a \( f_4 \)}{}{}}
\LANGsk{
\Question{
Na obrázku sú grafy piatich funkcií. Ktorá z nich je na \( \mathbb{R} \) primitívna k funkcii \( f(x)=x^2 \)?}
{\( f_2 \), \( f_3 \) i \( f_4 \)}{jen \( f_5 \)}{jen \( f_1 \)}{jen \( f_3 \) a \( f_4 \)}{}{}}
\LANGes{
\Question{
En la imagen hay gráficas de cinco funciones. ¿Cuál de las funciones es la primitiva de la función \( f(x)=x^2 \) en \( \mathbb{R} \)?}
{\( f_2 \), \( f_3 \) y \( f_4 \)}{sólo \( f_5 \)}{sólo \( f_1 \)}{sólo \( f_3 \) y \( f_4 \)}{}{}}
\End

\Data\ID{2010005101}
\Updated{2025-06-20 09:06:19}
\Level{A}
\SubArea{Primitive function}
\LANGen{
\Question{
Evaluate the following integral on \( \mathbb{R} \).
\[ \int\left(2^3+2x^3+\mathrm{e}^x-2^x-2^{\mathrm{e}}\right)\mathrm{d}x \]}
{\( 8x+0.5x^4+\mathrm{e}^x-\frac{2^x}{\ln2} -2^{\mathrm{e}} x+c,\ c\in\mathbb{R} \)}{\( -0.5x^4+\mathrm{e}^x-\frac{2^x}{\ln2} +c,\ c\in\mathbb{R} \)}{\( 8x-2x^4+\mathrm{e}^x-2^x-\frac{2^{\mathrm{e}+1}}{\mathrm{e}+1}+c,\ c\in\mathbb{R} \)}{\( 4-6x^4+\mathrm{e}^x -\frac{2^x}{\ln2} -2^\mathrm{e} x+c,\ c\in\mathbb{R} \)}{}{}}
\LANGpl{
\Question{
Oblicz następującą całkę na \( \mathbb{R} \).
\[ \int\left(2^3+2x^3+\mathrm{e}^x-2^x-2^{\mathrm{e}}\right)\mathrm{d}x \]}
{\( 8x+0{,}5x^4+\mathrm{e}^x-\frac{2^x}{\ln2} -2^{\mathrm{e}} x+c,\ c\in\mathbb{R} \)}{\( -0{,}5x^4+\mathrm{e}^x-\frac{2^x}{\ln2} +c,\ c\in\mathbb{R} \)}{\( 8x-2x^4+\mathrm{e}^x-2^x-\frac{2^{\mathrm{e}+1}}{\mathrm{e}+1}+c,\ c\in\mathbb{R} \)}{\( 4-6x^4+\mathrm{e}^x -\frac{2^x}{\ln2} -2^\mathrm{e} x+c,\ c\in\mathbb{R} \)}{}{}}
\LANGcs{
\Question{
Vypočítejte na množině \( \mathbb{R} \) následující integrál.
\[ \int\left(2^3+2x^3+\mathrm{e}^x-2^x-2^{\mathrm{e}}\right)\mathrm{d}x \]}
{\( 8x+0{,}5x^4+\mathrm{e}^x-\frac{2^x}{\ln2} -2^{\mathrm{e}} x+c,\ c\in\mathbb{R} \)}{\( -0{,}5x^4+\mathrm{e}^x-\frac{2^x}{\ln2} +c,\ c\in\mathbb{R} \)}{\( 8x-2x^4+\mathrm{e}^x-2^x-\frac{2^{\mathrm{e}+1}}{\mathrm{e}+1}+c,\ c\in\mathbb{R} \)}{\( 4-6x^4+\mathrm{e}^x -\frac{2^x}{\ln2} -2^\mathrm{e} x+c,\ c\in\mathbb{R} \)}{}{}}
\LANGsk{
\Question{
Vypočítajte na množine \( \mathbb{R} \) nasledujúci integrál. \[ \int\left(2^3+2x^3+\mathrm{e}^x-2^x-2^{\mathrm{e}}\right)\mathrm{d}x \]}
{\( 8x+0{,}5x^4+\mathrm{e}^x-\frac{2^x}{\ln2} -2^{\mathrm{e}} x+c,\ c\in\mathbb{R} \)}{\( -0{,}5x^4+\mathrm{e}^x-\frac{2^x}{\ln2} +c,\ c\in\mathbb{R} \)}{\( 8x-2x^4+\mathrm{e}^x-2^x-\frac{2^{\mathrm{e}+1}}{\mathrm{e}+1}+c,\ c\in\mathbb{R} \)}{\( 4-6x^4+\mathrm{e}^x -\frac{2^x}{\ln2} -2^\mathrm{e} x+c,\ c\in\mathbb{R} \)}{}{}}
\LANGes{
\Question{
Calcula la siguiente integral en \( \mathbb{R} \).
\[ \int\left(2^3+2x^3+\mathrm{e}^x-2^x-2^{\mathrm{e}}\right)\mathrm{d}x \]}
{\( 8x+0.5x^4+\mathrm{e}^x-\frac{2^x}{\ln2} -2^{\mathrm{e}} x+c,\ c\in\mathbb{R} \)}{\( -0.5x^4+\mathrm{e}^x-\frac{2^x}{\ln2} +c,\ c\in\mathbb{R} \)}{\( 8x-2x^4+\mathrm{e}^x-2^x-\frac{2^{\mathrm{e}+1}}{\mathrm{e}+1}+c,\ c\in\mathbb{R} \)}{\( 4-6x^4+\mathrm{e}^x -\frac{2^x}{\ln2} -2^\mathrm{e} x+c,\ c\in\mathbb{R} \)}{}{}}
\End

\Data\ID{2010005102}
\Updated{2022-11-02 09:11:34}
\Level{A}
\SubArea{Primitive function}
\LANGen{
\Question{
Evaluate the following integral on \( \left(\frac{\pi}2;\pi\right) \).
\[ \int\left(5 \sin x-\frac3{\cos^2x}-\frac7{\sin^2x}\right)\mathrm{d}x \]}
{\( -5\cos x-3\,\mathrm{tg}\,x+7\,\mathrm{cotg}\,x+c,\ c\in\mathbb{R} \)}{\( 5\cos x+3\,\mathrm{tg}\,x+7\,\mathrm{cotg}\,x+c,\ c\in\mathbb{R} \)}{\( -5\cos x-3\,\mathrm{tg}\,x-7\,\mathrm{cotg}\,x+c,\ c\in\mathbb{R} \)}{\( 5\cos x+3\,\mathrm{tg}\,x+7\,\mathrm{cotg}\,x+c,\ c\in\mathbb{R} \)}{}{}}
\LANGpl{
\Question{
Oblicz następującą całkę na \( \left(\frac{\pi}2;\pi\right) \).
\[ \int\left(5 \sin x-\frac3{\cos^2x}-\frac7{\sin^2x}\right)\mathrm{d}x \]}
{\( -5\cos x-3\,\mathrm{tg}\,x+7\,\mathrm{cotg}\,x+c,\ c\in\mathbb{R} \)}{\( 5\cos x+3\,\mathrm{tg}\,x+7\,\mathrm{cotg}\,x+c,\ c\in\mathbb{R} \)}{\( -5\cos x-3\,\mathrm{tg}\,x-7\,\mathrm{cotg}\,x+c,\ c\in\mathbb{R} \)}{\( 5\cos x+3\,\mathrm{tg}\,x+7\,\mathrm{cotg}\,x+c,\ c\in\mathbb{R} \)}{}{}}
\LANGcs{
\Question{
Vypočítejte na intervalu \( \left(\frac{\pi}2;\pi\right) \) následující integrál.
\[ \int\left(5 \sin x-\frac3{\cos^2x}-\frac7{\sin^2x}\right)\mathrm{d}x \]}
{\( -5\cos x-3\,\mathrm{tg}\,x+7\,\mathrm{cotg}\,x+c,\ c\in\mathbb{R} \)}{\( 5\cos x+3\,\mathrm{tg}\,x+7\,\mathrm{cotg}\,x+c,\ c\in\mathbb{R} \)}{\( -5\cos x-3\,\mathrm{tg}\,x-7\,\mathrm{cotg}\,x+c,\ c\in\mathbb{R} \)}{\( 5\cos x+3\,\mathrm{tg}\,x+7\,\mathrm{cotg}\,x+c,\ c\in\mathbb{R} \)}{}{}}
\LANGsk{
\Question{
Vypočítajte na intervale \( \left(\frac{\pi}2;\pi\right) \) nasledujúci integrál.
\[ \int\left(5 \sin x-\frac3{\cos^2x}-\frac7{\sin^2x}\right)\mathrm{d}x \]}
{\( -5\cos x-3\,\mathrm{tg}\,x+7\,\mathrm{cotg}\,x+c,\ c\in\mathbb{R} \)}{\( 5\cos x+3\,\mathrm{tg}\,x+7\,\mathrm{cotg}\,x+c,\ c\in\mathbb{R} \)}{\( -5\cos x-3\,\mathrm{tg}\,x-7\,\mathrm{cotg}\,x+c,\ c\in\mathbb{R} \)}{\( 5\cos x+3\,\mathrm{tg}\,x+7\,\mathrm{cotg}\,x+c,\ c\in\mathbb{R} \)}{}{}}
\LANGes{
\Question{
Calcula la siguiente integral en \( \left(\frac{\pi}2;\pi\right) \).
\[ \int\left(5 \sin x-\frac3{\cos^2x}-\frac7{\sin^2x}\right)\mathrm{d}x \]}
{\( -5\cos x-3\,\mathrm{tg}\,x+7\,\mathrm{cotg}\,x+c,\ c\in\mathbb{R} \)}{\( 5\cos x+3\,\mathrm{tg}\,x+7\,\mathrm{cotg}\,x+c,\ c\in\mathbb{R} \)}{\( -5\cos x-3\,\mathrm{tg}\,x-7\,\mathrm{cotg}\,x+c,\ c\in\mathbb{R} \)}{\( 5\cos x+3\,\mathrm{tg}\,x+7\,\mathrm{cotg}\,x+c,\ c\in\mathbb{R} \)}{}{}}
\End

\Data\ID{2010005103}
\Updated{2022-11-02 09:11:34}
\Level{A}
\SubArea{Primitive function}
\LANGen{
\Question{
Evaluate the following integral on \( (0;\infty) \).
\[ \int\left(6\sqrt x-5\sqrt[3]{x^2}+10\sqrt[4]{x}\right)\mathrm{d}x \]}
{\( 4x\sqrt x-3x\sqrt[3]{x^2}+8x\sqrt[4]{x}+c,\ c\in\mathbb{R} \)}{\( 9 x\sqrt x-\frac{25}3 x^3\sqrt[3]{x^2}+\frac{25}2 x\sqrt[4]{x}+c,\ c\in\mathbb{R} \)}{\( 4\sqrt x-3\sqrt[3]{x^2}+8\sqrt[4]{x}+c,\ c\in\mathbb{R} \)}{\(x+20\sqrt{x}+c,\ c\in\mathbb{R} \)}{}{}}
\LANGpl{
\Question{
Oblicz następującą całkę na \( (0;\infty) \).
\[ \int\left(6\sqrt x-5\sqrt[3]{x^2}+10\sqrt[4]{x}\right)\mathrm{d}x \]}
{\( 4x\sqrt x-3x\sqrt[3]{x^2}+8x\sqrt[4]{x}+c,\ c\in\mathbb{R} \)}{\( 9 x\sqrt x-\frac{25}3 x^3\sqrt[3]{x^2}+\frac{25}2 x\sqrt[4]{x}+c,\ c\in\mathbb{R} \)}{\( 4\sqrt x-3\sqrt[3]{x^2}+8\sqrt[4]{x}+c,\ c\in\mathbb{R} \)}{\(x+20\sqrt{x}+c,\ c\in\mathbb{R} \)}{}{}}
\LANGcs{
\Question{
Vypočítejte na intervalu  \( (0;\infty) \) následující integrál.
\[ \int\left(6\sqrt x-5\sqrt[3]{x^2}+10\sqrt[4]{x}\right)\mathrm{d}x \]}
{\( 4x\sqrt x-3x\sqrt[3]{x^2}+8x\sqrt[4]{x}+c,\ c\in\mathbb{R} \)}{\( 9 x\sqrt x-\frac{25}3 x^3\sqrt[3]{x^2}+\frac{25}2 x\sqrt[4]{x}+c,\ c\in\mathbb{R} \)}{\( 4\sqrt x-3\sqrt[3]{x^2}+8\sqrt[4]{x}+c,\ c\in\mathbb{R} \)}{\(x+20\sqrt{x}+c,\ c\in\mathbb{R} \)}{}{}}
\LANGsk{
\Question{
Vypočítajte na intervale \( (0;\infty) \) nasledujúci integrál.
\[ \int\left(6\sqrt x-5\sqrt[3]{x^2}+10\sqrt[4]{x}\right)\mathrm{d}x \]}
{\( 4x\sqrt x-3x\sqrt[3]{x^2}+8x\sqrt[4]{x}+c,\ c\in\mathbb{R} \)}{\( 9 x\sqrt x-\frac{25}3 x^3\sqrt[3]{x^2}+\frac{25}2 x\sqrt[4]{x}+c,\ c\in\mathbb{R} \)}{\( 4\sqrt x-3\sqrt[3]{x^2}+8\sqrt[4]{x}+c,\ c\in\mathbb{R} \)}{\(x+20\sqrt{x}+c,\ c\in\mathbb{R} \)}{}{}}
\LANGes{
\Question{
Calcula la siguiente integral en \( (0;\infty) \).
\[ \int\left(6\sqrt x-5\sqrt[3]{x^2}+10\sqrt[4]{x}\right)\mathrm{d}x \]}
{\( 4x\sqrt x-3x\sqrt[3]{x^2}+8x\sqrt[4]{x}+c,\ c\in\mathbb{R} \)}{\( 9 x\sqrt x-\frac{25}3 x^3\sqrt[3]{x^2}+\frac{25}2 x\sqrt[4]{x}+c,\ c\in\mathbb{R} \)}{\( 4\sqrt x-3\sqrt[3]{x^2}+8\sqrt[4]{x}+c,\ c\in\mathbb{R} \)}{\(x+20\sqrt{x}+c,\ c\in\mathbb{R} \)}{}{}}
\End

\Data\ID{2010005104}
\Updated{2022-11-02 09:11:34}
\Level{A}
\SubArea{Primitive function}
\LANGen{
\Question{
Evaluate the following integral on the interval \((0;+\infty)\). 
\[ 
 \int \left (2x^{-1}+\frac{2}
{x^2} - 3x^{-3} \right )\, \mathrm{d}x
\]}
{\(2\ln |x| - \frac{2} 
{x} +\frac{3} 
{2x^{2}} + c,\ c\in \mathbb{R}\)}{\(2\ln |x| - \frac{2} 
{x} -\frac{3} 
{2x^{2}} + c,\ c\in \mathbb{R}\)}{\(2\ln |x| + \frac{2} 
{x} +\frac{3} 
{2x^{2}} + c,\ c\in \mathbb{R}\)}{\(2\ln |x| + \frac{2} 
{x} -\frac{3} 
{2x^{2}} + c,\ c\in \mathbb{R}\)}{}{}}
\LANGpl{
\Question{
Oblicz następującą całkę na przedziale \((0;+\infty)\). 
\[ 
 \int \left (2x^{-1}+\frac{2}
{x^2} - 3x^{-3} \right )\, \mathrm{d}x
\]}
{\(2\ln |x| - \frac{2} 
{x} +\frac{3} 
{2x^{2}} + c,\ c\in \mathbb{R}\)}{\(2\ln |x| - \frac{2} 
{x} -\frac{3} 
{2x^{2}} + c,\ c\in \mathbb{R}\)}{\(2\ln |x| + \frac{2} 
{x} +\frac{3} 
{2x^{2}} + c,\ c\in \mathbb{R}\)}{\(2\ln |x| + \frac{2} 
{x} -\frac{3} 
{2x^{2}} + c,\ c\in \mathbb{R}\)}{}{}}
\LANGcs{
\Question{
Vypočítejte na intervalu \((0;+\infty)\) následující integrál. 
\[ 
 \int \left (2x^{-1}+\frac{2}
{x^2} - 3x^{-3} \right )\, \mathrm{d}x
\]}
{\(2\ln |x| - \frac{2} 
{x} +\frac{3} 
{2x^{2}} + c,\ c\in \mathbb{R}\)}{\(2\ln |x| - \frac{2} 
{x} -\frac{3} 
{2x^{2}} + c,\ c\in \mathbb{R}\)}{\(2\ln |x| + \frac{2} 
{x} +\frac{3} 
{2x^{2}} + c,\ c\in \mathbb{R}\)}{\(2\ln |x| + \frac{2} 
{x} -\frac{3} 
{2x^{2}} + c,\ c\in \mathbb{R}\)}{}{}}
\LANGsk{
\Question{
Vypočítajte na intervale \((0;+\infty)\) nasledujúci integrál. 
\[ 
 \int \left (2x^{-1}+\frac{2}
{x^2} - 3x^{-3} \right )\, \mathrm{d}x
\]}
{\(2\ln |x| - \frac{2} 
{x} +\frac{3} 
{2x^{2}} + c,\ c\in \mathbb{R}\)}{\(2\ln |x| - \frac{2} 
{x} -\frac{3} 
{2x^{2}} + c,\ c\in \mathbb{R}\)}{\(2\ln |x| + \frac{2} 
{x} +\frac{3} 
{2x^{2}} + c,\ c\in \mathbb{R}\)}{\(2\ln |x| + \frac{2} 
{x} -\frac{3} 
{2x^{2}} + c,\ c\in \mathbb{R}\)}{}{}}
\LANGes{
\Question{
Calcula la siguiente integral en el intervalo \((0;+\infty)\). 
\[ 
 \int \left (2x^{-1}+\frac{2}
{x^2} - 3x^{-3} \right )\, \mathrm{d}x
\]}
{\(2\ln |x| - \frac{2} 
{x} +\frac{3} 
{2x^{2}} + c,\ c\in \mathbb{R}\)}{\(2\ln |x| - \frac{2} 
{x} -\frac{3} 
{2x^{2}} + c,\ c\in \mathbb{R}\)}{\(2\ln |x| + \frac{2} 
{x} +\frac{3} 
{2x^{2}} + c,\ c\in \mathbb{R}\)}{\(2\ln |x| + \frac{2} 
{x} -\frac{3} 
{2x^{2}} + c,\ c\in \mathbb{R}\)}{}{}}
\End

\Data\ID{2010005105}
\Updated{2022-11-02 09:11:36}
\Level{A}
\SubArea{Primitive function}
\LANGen{
\Question{
Given the function \(f(x) =\cos x -\sin x\), find its primitive function F so that the graph of  \(F\) passes through the point \(A = \left [ \pi;3\right ]\).}
{\(F(x) =\sin x +\cos x + 4\)}{\(F(x) =\cos x -\sin x + 4\)}{\(F(x) = \sin x -\cos x + 2\)}{}{}{}}
\LANGpl{
\Question{
Dana jest funkcja \(f(x) =\cos x -\sin x\), wyznacz jej funkcję pierwotną  F taką, by wykres    \(F\) przechodził przez punkt \(A = \left [ \pi;3\right ]\).}
{\(F(x) =\sin x +\cos x + 4\)}{\(F(x) =\cos x -\sin x + 4\)}{\(F(x) = \sin x -\cos x + 2\)}{}{}{}}
\LANGcs{
\Question{
K dané funkci \(f(x) =\cos x -\sin x\) v
\(\mathbb{R}\) určete primitivní funkci \(F\), jejíž graf prochází bodem \(A = \left [ \pi;3\right ]\).}
{\(F(x) =\sin x +\cos x + 4\)}{\(F(x) =\cos x -\sin x + 4\)}{\(F(x) = \sin x -\cos x + 2\)}{}{}{}}
\LANGsk{
\Question{
K danej funkcii \(f(x) =\cos x -\sin x\) v
\(\mathbb{R}\) určte primitívnu
funkciu \(F\), ktorej graf
prechádza bodom \(A = \left [ \pi;3\right ]\).}
{\(F(x) =\sin x +\cos x + 4\)}{\(F(x) =\cos x -\sin x + 4\)}{\(F(x) = \sin x -\cos x + 2\)}{}{}{}}
\LANGes{
\Question{
Dada la función \(f(x) =\cos x -\sin x\), halla su función primitiva F tal que la gráfica de \(F\) pasa por el punto \(A = \left [ \pi;3\right ]\).}
{\(F(x) =\sin x +\cos x + 4\)}{\(F(x) =\cos x -\sin x + 4\)}{\(F(x) = \sin x -\cos x + 2\)}{}{}{}}
\End

\Data\ID{2110005106}
\Updated{2022-11-02 09:11:56}
\Level{A}
\SubArea{Primitive function}
\LANGen{
\Question{
In which picture can you see a pair of graphs of functions \( f_1 \) and \( f_2 \) that are primitive to the same function?}
{}{}{}{}{}{}}
\LANGpl{
\Question{
Na którym obrazku widać parę wykresów funkcji \( f_1 \) i \( f_2 \), które są funkcjami pierwotnymi do tej samej funkcji?}
{}{}{}{}{}{}}
\LANGcs{
\Question{
Na kterém obrázku je dvojice grafů funkcí \( f_1 \) a \( f_2 \) primitivních k téže funkci?}
{}{}{}{}{}{}}
\LANGsk{
\Question{
Na ktorom obrázku je dvojica grafov funkcií \( f_1 \) a \( f_2 \) primitívnych k tej istej funkcii?}
{}{}{}{}{}{}}
\LANGes{
\Question{
¿En qué imagen puedes ver un par de gráficas de funciones \( f_1 \) y \( f_2 \) que son primitivas de la misma función?}
{}{}{}{}{}{}}

\IMGanswer{2110005101-figure0.pdf}
\IMGanswer{2110005101-figure1.pdf}
\IMGanswer{2110005101-figure2.pdf}
\IMGanswer{2110005101-figure3.pdf}\End

\Data\ID{9000065501}
\Updated{2020-07-15 03:07:37}
\Level{A}
\SubArea{Primitive function}
\LANGen{
\Question{
Evaluate the following integral on \(\mathbb{R}\). 
\[ 
 \int (x^{3} + x^{2} - 2x)\, \mathrm{d}x
\]}
{\(\frac{1}
{4}x^{4} + \frac{1} 
{3}x^{3} - x^{2} + c,\ c\in \mathbb{R}\)}{\(\frac{1}
{4}x^{4} -\frac{1} 
{3}x^{3} + x^{2} + c,\ c\in \mathbb{R}\)}{\(3x^{2} + 2x - 2 + c,\ c\in \mathbb{R}\)}{\(3x^{2} - 2x + 2 + c,\ c\in \mathbb{R}\)}{}{}}
\LANGpl{
\Question{
Wyznacz całkę.
\[ 
 \int (x^{3} + x^{2} - 2x)\, \mathrm{d}x
\]}
{\(\frac{1}
{4}x^{4} + \frac{1} 
{3}x^{3} - x^{2} + c,\ c\in \mathbb{R}\)}{\(\frac{1}
{4}x^{4} -\frac{1} 
{3}x^{3} + x^{2} + c,\ c\in \mathbb{R}\)}{\(3x^{2} + 2x - 2 + c,\ c\in \mathbb{R}\)}{\(3x^{2} - 2x + 2 + c,\ c\in \mathbb{R}\)}{}{}}
\LANGcs{
\Question{
Vypočtěte \(\int (x^{3} + x^{2} - 2x)\, \mathrm{d}x\) na \(\mathbb{R}\).}
{\(\frac{1}
{4}x^{4} + \frac{1} 
{3}x^{3} - x^{2} + c,\ c\in\mathbb{R}\)}{\(\frac{1}
{4}x^{4} -\frac{1} 
{3}x^{3} + x^{2} + c, c\in\mathbb{R}\)}{\(3x^{2} + 2x - 2 + c, c\in\mathbb{R}\)}{\(3x^{2} - 2x + 2 + c, c\in\mathbb{R}\)}{}{}}
\LANGsk{
\Question{
Vypočítajte \(\int (x^{3} + x^{2} - 2x)\, \mathrm{d}x\) na \(\mathbb{R}\).}
{\(\frac{1}
{4}x^{4} + \frac{1} 
{3}x^{3} - x^{2} + c,\ c\in\mathbb{R}\)}{\(\frac{1}
{4}x^{4} -\frac{1} 
{3}x^{3} + x^{2} + c, c\in\mathbb{R}\)}{\(3x^{2} + 2x - 2 + c, c\in\mathbb{R}\)}{\(3x^{2} - 2x + 2 + c, c\in\mathbb{R}\)}{}{}}
\LANGes{
\Question{
Evalúa la siguiente integral en \(\mathbb{R}\). 
\[ 
 \int (x^{3} + x^{2} - 2x)\, \mathrm{d}x
\]}
{\(\frac{1}
{4}x^{4} + \frac{1} 
{3}x^{3} - x^{2} + c,\ c\in \mathbb{R}\)}{\(\frac{1}
{4}x^{4} -\frac{1} 
{3}x^{3} + x^{2} + c,\ c\in \mathbb{R}\)}{\(3x^{2} + 2x - 2 + c,\ c\in \mathbb{R}\)}{\(3x^{2} - 2x + 2 + c,\ c\in \mathbb{R}\)}{}{}}
\End

\Data\ID{9000065502}
\Updated{2020-07-15 03:07:37}
\Level{A}
\SubArea{Primitive function}
\LANGen{
\Question{
Evaluate the following integral on \(\mathbb{R}\). 
\[ 
 \int (4x + 7)\, \mathrm{d}x
\]}
{\(2x^{2} + 7x + c,\ c\in \mathbb{R}\)}{\(2x^{2} - 7x + c,\ c\in \mathbb{R}\)}{\(4 + c,\ c\in \mathbb{R}\)}{\(4x^{2} + 7x + c,\ c\in \mathbb{R}\)}{}{}}
\LANGpl{
\Question{
Wyznacz całkę. 
\[ 
 \int (4x + 7)\, \mathrm{d}x
\]}
{\(2x^{2} + 7x + c,\ c\in \mathbb{R}\)}{\(2x^{2} - 7x + c,\ c\in \mathbb{R}\)}{\(4 + c,\ c\in \mathbb{R}\)}{\(4x^{2} + 7x + c,\ c\in \mathbb{R}\)}{}{}}
\LANGcs{
\Question{
Vypočtěte \(\int (4x + 7)\, \mathrm{d}x\) na \(\mathbb{R}\).}
{\(2x^{2} + 7x + c,\ c\in\mathbb{R}\)}{\(2x^{2} - 7x + c,\ c\in\mathbb{R}\)}{\(4 + c,\ c\in\mathbb{R}\)}{\(4x^{2} + 7x + c,\ c\in\mathbb{R}\)}{}{}}
\LANGsk{
\Question{
Vypočítajte \(\int (4x + 7)\, \mathrm{d}x\) na \(\mathbb{R}\).}
{\(2x^{2} + 7x + c,\ c\in\mathbb{R}\)}{\(2x^{2} - 7x + c,\ c\in\mathbb{R}\)}{\(4 + c,\ c\in\mathbb{R}\)}{\(4x^{2} + 7x + c,\ c\in\mathbb{R}\)}{}{}}
\LANGes{
\Question{
Evalúa la siguiente integral en \(\mathbb{R}\). 
\[ 
 \int (4x + 7)\, \mathrm{d}x
\]}
{\(2x^{2} + 7x + c,\ c\in \mathbb{R}\)}{\(2x^{2} - 7x + c,\ c\in \mathbb{R}\)}{\(4 + c,\ c\in \mathbb{R}\)}{\(4x^{2} + 7x + c,\ c\in \mathbb{R}\)}{}{}}
\End

\Data\ID{9000065503}
\Updated{2020-07-15 03:07:37}
\Level{A}
\SubArea{Primitive function}
\LANGen{
\Question{
Evaluate the following integral on the interval \((0;+\infty)\). 
\[ 
 \int (4x^{-3} - x^{-4})\, \mathrm{d}x
\]}
{\(- 2x^{-2} + \frac{1} 
{3}x^{-3} + c,\ c\in \mathbb{R}\)}{\(-\frac{4} 
{3}x^{-2} -\frac{1} 
{3}x^{-3} + c,\ c\in \mathbb{R}\)}{\(-\frac{3} 
{4}x^{-4} -\frac{1} 
{5}x^{-5} + c,\ c\in \mathbb{R}\)}{\(- 12x^{2} + 4x^{-3} + c,\ c\in \mathbb{R}\)}{}{}}
\LANGpl{
\Question{
Wyznacz całkę na przedziale  \((0;+\infty)\). 
\[ 
 \int (4x^{-3} - x^{-4})\, \mathrm{d}x
\]}
{\(- 2x^{-2} + \frac{1} 
{3}x^{-3} + c,\ c\in \mathbb{R}\)}{\(-\frac{4} 
{3}x^{-2} -\frac{1} 
{3}x^{-3} + c,\ c\in \mathbb{R}\)}{\(-\frac{3} 
{4}x^{-4} -\frac{1} 
{5}x^{-5} + c,\ c\in \mathbb{R}\)}{\(- 12x^{2} + 4x^{-3} + c,\ c\in \mathbb{R}\)}{}{}}
\LANGcs{
\Question{
Vypočtěte \(\int (4x^{-3} - x^{-4})\, \mathrm{d}x\) na intervalu \((0;+\infty)\).}
{\(- 2x^{-2} + \frac{1} 
{3}x^{-3} + c,\ c\in\mathbb{R}\)}{\(-\frac{4} 
{3}x^{-2} -\frac{1} 
{3}x^{-3} + c,\ c\in\mathbb{R}\)}{\(-\frac{3} 
{4}x^{-4} -\frac{1} 
{5}x^{-5} + c,\ c\in\mathbb{R}\)}{\(- 12x^{2} + 4x^{-3} + c,\ c\in\mathbb{R}\)}{}{}}
\LANGsk{
\Question{
Vypočítajte \(\int (4x^{-3} - x^{-4})\, \mathrm{d}x\) na intervale \((0;+\infty)\).}
{\(- 2x^{-2} + \frac{1} 
{3}x^{-3} + c,\ c\in\mathbb{R}\)}{\(-\frac{4} 
{3}x^{-2} -\frac{1} 
{3}x^{-3} + c,\ c\in\mathbb{R}\)}{\(-\frac{3} 
{4}x^{-4} -\frac{1} 
{5}x^{-5} + c,\ c\in\mathbb{R}\)}{\(- 12x^{2} + 4x^{-3} + c,\ c\in\mathbb{R}\)}{}{}}
\LANGes{
\Question{
Evalúa la siguiente integral en el intervalo \((0;+\infty)\). 
\[ 
 \int (4x^{-3} - x^{-4})\, \mathrm{d}x
\]}
{\(- 2x^{-2} + \frac{1} 
{3}x^{-3} + c,\ c\in \mathbb{R}\)}{\(-\frac{4} 
{3}x^{-2} -\frac{1} 
{3}x^{-3} + c,\ c\in \mathbb{R}\)}{\(-\frac{3} 
{4}x^{-4} -\frac{1} 
{5}x^{-5} + c,\ c\in \mathbb{R}\)}{\(- 12x^{2} + 4x^{-3} + c,\ c\in \mathbb{R}\)}{}{}}
\End

\Data\ID{9000065507}
\Updated{2020-07-15 03:07:37}
\Level{A}
\SubArea{Primitive function}
\LANGen{
\Question{
Given the function 
\[ 
 F(x) = \frac{1} 
{4}x^{4} -\frac{2} 
{3}x^{3},
\]
find the function \(f\) 
such that \(F\) is
primitive to \(f\) 
on \(\mathbb{R}\).}
{\(f(x) = x^{3} - 2x^{2}\)}{\(f(x) = x^{5} - 2x^{4}\)}{\(f(x) = x^{5} - 3x^{2}\)}{\(f(x) = -4x^{-4} - 3x^{2}\)}{}{}}
\LANGpl{
\Question{
Dana jest funkcja 
\[ 
 F(x) = \frac{1} 
{4}x^{4} -\frac{2} 
{3}x^{3}
\] Wybierz funkcję $f$, dla której funkcja $F$ jest pierwotna.}
{\(f(x) = x^{3} - 2x^{2}\)}{\(f(x) = x^{5} - 2x^{4}\)}{\(f(x) = x^{5} - 3x^{2}\)}{\(f(x) = -4x^{-4} - 3x^{2}\)}{}{}}
\LANGcs{
\Question{
Je dána funkce \(F\) 
předpisem: \(F(x) = \frac{1} 
{4}x^{4} -\frac{2} 
{3}x^{3}\).
Vyberte funkci \(f\),
k níž je \(F\) 
funkcí primitivní na \(\mathbb{R}\).}
{\(f(x) = x^{3} - 2x^{2}\)}{\(f(x) = x^{5} - 2x^{4}\)}{\(f(x) = x^{5} - 3x^{2}\)}{\(f(x) = -4x^{-4} - 3x^{2}\)}{}{}}
\LANGsk{
\Question{
Je daná funkcia \(F\) 
predpisom: \(F(x) = \frac{1} 
{4}x^{4} -\frac{2} 
{3}x^{3}\).
Vyberte funkciu \(f\),
v  ktorej je \(F\) 
funkcia primitivná na \(\mathbb{R}\).}
{\(f(x) = x^{3} - 2x^{2}\)}{\(f(x) = x^{5} - 2x^{4}\)}{\(f(x) = x^{5} - 3x^{2}\)}{\(f(x) = -4x^{-4} - 3x^{2}\)}{}{}}
\LANGes{
\Question{
Dada la función
\[ 
 F(x) = \frac{1} 
{4}x^{4} -\frac{2} 
{3}x^{3},
\]
encuentra la función \(f\) 
tal que \(F\) es
primitiva de \(f\) 
en \(\mathbb{R}\).}
{\(f(x) = x^{3} - 2x^{2}\)}{\(f(x) = x^{5} - 2x^{4}\)}{\(f(x) = x^{5} - 3x^{2}\)}{\(f(x) = -4x^{-4} - 3x^{2}\)}{}{}}
\End

\Data\ID{9000065508}
\Updated{2020-07-15 03:07:37}
\Level{A}
\SubArea{Primitive function}
\LANGen{
\Question{
Given the function 
\[ 
 F(x) = \frac{1} 
{4}x^{4} -\frac{5} 
{2}x^{2},
\]
find the function \(f\) 
such that \(F\) is
primitive to \(f\) 
on \(\mathbb{R}\).}
{\(f(x) = x(x^{2} - 5)\)}{\(f(x) = x^{3} - 5x^{2}\)}{\(f(x) = x^{5} - 5x^{3}\)}{\(f(x) = x^{5} - 2x\)}{}{}}
\LANGpl{
\Question{
Dana jest funkcja 
\[ 
 F(x) = \frac{1} 
{4}x^{4} -\frac{5} 
{2}x^{2}
\]
Wybierz funkcję $f$, dla której funkcja $F$ jest pierwotna.}
{\(f(x) = x(x^{2} - 5)\)}{\(f(x) = x^{3} - 5x^{2}\)}{\(f(x) = x^{5} - 5x^{3}\)}{\(f(x) = x^{5} - 2x\)}{}{}}
\LANGcs{
\Question{
Je dána funkce \(F\) 
předpisem: \(F(x) = \frac{1} 
{4}x^{4} -\frac{5} 
{2}x^{2}\).
Vyberte funkci \(f\),
k níž je \(F\) 
funkcí primitivní na \(\mathbb{R}\).}
{\(f(x) = x(x^{2} - 5)\)}{\(f(x) = x^{3} - 5x^{2}\)}{\(f(x) = x^{5} - 5x^{3}\)}{\(f(x) = x^{5} - 2x\)}{}{}}
\LANGsk{
\Question{
Je daná funkcia \(F\) 
predpisom: \(F(x) = \frac{1} 
{4}x^{4} -\frac{5} 
{2}x^{2}\).
Vyberte funkciu \(f\),
v ktorej je \(F\) 
funkcia primitívna na \(\mathbb{R}\).}
{\(f(x) = x(x^{2} - 5)\)}{\(f(x) = x^{3} - 5x^{2}\)}{\(f(x) = x^{5} - 5x^{3}\)}{\(f(x) = x^{5} - 2x\)}{}{}}
\LANGes{
\Question{
Dada la función
\[ 
 F(x) = \frac{1} 
{4}x^{4} -\frac{5} 
{2}x^{2},
\]
encuentra la función \(f\) 
tal que \(F\) es
primitiva de \(f\) 
en \(\mathbb{R}\).}
{\(f(x) = x(x^{2} - 5)\)}{\(f(x) = x^{3} - 5x^{2}\)}{\(f(x) = x^{5} - 5x^{3}\)}{\(f(x) = x^{5} - 2x\)}{}{}}
\End

\Data\ID{9000065509}
\Updated{2020-07-15 03:07:37}
\Level{A}
\SubArea{Primitive function}
\LANGen{
\Question{
Given the function 
\[ 
 F(x) = x + \frac{9} 
{2}x^{2} + 9x^{3} + \frac{27} 
 {4} x^{4},
\]
find the function \(f\) 
such that \(F\) is
primitive to \(f\) 
on \(\mathbb{R}\).}
{\(f(x) = (1 + 3x)^{3}\)}{\(f(x) = (1 + 3x)^{2}\)}{\(f(x) = 1 + 3x + 3x^{2} + 3x^{3}\)}{\(f(x) = (1 + 3x)^{4}\)}{}{}}
\LANGpl{
\Question{
Dana jest funkcja \[ 
 F(x) = x + \frac{9} 
{2}x^{2} + 9x^{3} + \frac{27} 
 {4} x^{4}
\]
Wybierz funkcję $f$, dla której funkcja $F$ jest pierwotna.}
{\(f(x) = (1 + 3x)^{3}\)}{\(f(x) = (1 + 3x)^{2}\)}{\(f(x) = 1 + 3x + 3x^{2} + 3x^{3}\)}{\(f(x) = (1 + 3x)^{4}\)}{}{}}
\LANGcs{
\Question{
Je dána funkce \(F\) 
předpisem: \(F(x) = x + \frac{9} 
{2}x^{2} + 9x^{3} + \frac{27} 
 {4} x^{4}\).
Vyberte funkci \(f\),
k níž je \(F\) 
funkcí primitivní na \(\mathbb{R}\).}
{\(f(x) = (1 + 3x)^{3}\)}{\(f(x) = (1 + 3x)^{2}\)}{\(f(x) = 1 + 3x + 3x^{2} + 3x^{3}\)}{\(f(x) = (1 + 3x)^{4}\)}{}{}}
\LANGsk{
\Question{
Je daná funkcia \(F\) 
predpisom: \(F(x) = x + \frac{9} 
{2}x^{2} + 9x^{3} + \frac{27} 
 {4} x^{4}\).
Vyberte funkciu \(f\),
v ktorej je \(F\) 
funkcia primitívna na \(\mathbb{R}\).}
{\(f(x) = (1 + 3x)^{3}\)}{\(f(x) = (1 + 3x)^{2}\)}{\(f(x) = 1 + 3x + 3x^{2} + 3x^{3}\)}{\(f(x) = (1 + 3x)^{4}\)}{}{}}
\LANGes{
\Question{
Dada la función
\[ 
 F(x) = x + \frac{9} 
{2}x^{2} + 9x^{3} + \frac{27} 
 {4} x^{4},
\]
encuentra la función \(f\) 
tal que \(F\) es
primitiva de \(f\) 
en \(\mathbb{R}\).}
{\(f(x) = (1 + 3x)^{3}\)}{\(f(x) = (1 + 3x)^{2}\)}{\(f(x) = 1 + 3x + 3x^{2} + 3x^{3}\)}{\(f(x) = (1 + 3x)^{4}\)}{}{}}
\End

\Data\ID{9000065510}
\Updated{2020-07-15 03:07:37}
\Level{A}
\SubArea{Primitive function}
\LANGen{
\Question{
Given the function 
\[ 
 F(x) = \frac{6} 
{7}x^{3}\sqrt{x},
\]
find the function \(f\) 
such that \(F\) is
primitive to \(f\) 
on the interval \((0;+\infty)\).}
{\(f(x) = 3x^{2}\sqrt{x}\)}{\(f(x) = 3x\sqrt{x}\)}{\(f(x) = 3x^{3}\sqrt{x}\)}{\(f(x) = 7x\sqrt{x}\)}{}{}}
\LANGpl{
\Question{
Dana jest funkcja \[ 
 F(x) = \frac{6} 
{7}x^{3}\sqrt{x},
\]
Wybierz funkcję $f$, dla której funkcja $F$ jest pierwotna w przedziale
 \((0;+\infty)\).}
{\(f(x) = 3x^{2}\sqrt{x}\)}{\(f(x) = 3x\sqrt{x}\)}{\(f(x) = 3x^{3}\sqrt{x}\)}{\(f(x) = 7x\sqrt{x}\)}{}{}}
\LANGcs{
\Question{
Je dána funkce \(F\) 
předpisem: \(F(x) = \frac{6} 
{7}x^{3}\sqrt{x}\).
Vyberte funkci \(f\),
k níž je \(F\) 
funkcí primitivní na intervalu \((0;+\infty)\).}
{\(f(x) = 3x^{2}\sqrt{x}\)}{\(f(x) = 3x\sqrt{x}\)}{\(f(x) = 3x^{3}\sqrt{x}\)}{\(f(x) = 7x\sqrt{x}\)}{}{}}
\LANGsk{
\Question{
Je daná funkcia \(F\) 
predpisom: \(F(x) = \frac{6} 
{7}x^{3}\sqrt{x}\).
Vyberte funkciu \(f\),
v ktorej je \(F\) 
funkcia primitívna na intervale \((0;+\infty)\).}
{\(f(x) = 3x^{2}\sqrt{x}\)}{\(f(x) = 3x\sqrt{x}\)}{\(f(x) = 3x^{3}\sqrt{x}\)}{\(f(x) = 7x\sqrt{x}\)}{}{}}
\LANGes{
\Question{
Dada la función
\[ 
 F(x) = \frac{6} 
{7}x^{3}\sqrt{x},
\]
encuentra la función \(f\) 
tal que \(F\) es
primitiva de \(f\) 
en el intervalo  \((0;+\infty)\).}
{\(f(x) = 3x^{2}\sqrt{x}\)}{\(f(x) = 3x\sqrt{x}\)}{\(f(x) = 3x^{3}\sqrt{x}\)}{\(f(x) = 7x\sqrt{x}\)}{}{}}
\End

\Data\ID{9000065902}
\Updated{2020-07-15 03:07:37}
\Level{A}
\SubArea{Primitive function}
\LANGen{
\Question{
Evaluate the following integral on the interval \((0;+\infty)\). 
\[ 
 \int \left (2 + \frac{1} 
{x}\right )\, \text{d}x
\]}
{\(2x +\ln |x| + c,\ c\in \mathbb{R}\)}{\(\ln |x| + c,\ c\in \mathbb{R}\)}{\(2 +\ln |x| + c,\ c\in \mathbb{R}\)}{\(2x^{2} +\ln |x| + c,\ c\in \mathbb{R}\)}{}{}}
\LANGpl{
\Question{
Wyznacz całkę na przedziale  \((0;+\infty)\). 
\[ 
 \int \left (2 + \frac{1} 
{x}\right )\, \text{d}x
\]}
{\(2x +\ln |x| + c,\ c\in \mathbb{R}\)}{\(\ln |x| + c,\ c\in \mathbb{R}\)}{\(2 +\ln |x| + c,\ c\in \mathbb{R}\)}{\(2x^{2} +\ln |x| + c,\ c\in \mathbb{R}\)}{}{}}
\LANGcs{
\Question{
Vypočtěte \(\int \left (2 + \frac{1} 
{x}\right )\, \text{d}x\) na intervalu \((0;+\infty)\).}
{\(2x +\ln |x| + c,\ c\in\mathbb{R}\)}{\(\ln |x| + c,\ c\in\mathbb{R}\)}{\(2 +\ln |x| + c,\ c\in\mathbb{R}\)}{\(2x^{2} +\ln |x| + c,\ c\in\mathbb{R}\)}{}{}}
\LANGsk{
\Question{
Vypočítajte \(\int \left (2 + \frac{1} 
{x}\right )\, \text{d}x\) na intervale \((0;+\infty)\).}
{\(2x +\ln |x| + c,\ c\in\mathbb{R}\)}{\(\ln |x| + c,\ c\in\mathbb{R}\)}{\(2 +\ln |x| + c,\ c\in\mathbb{R}\)}{\(2x^{2} +\ln |x| + c,\ c\in\mathbb{R}\)}{}{}}
\LANGes{
\Question{
Evalúa la siguiente integral en el intervalo\((0;+\infty)\). 
\[ 
 \int \left (2 + \frac{1} 
{x}\right )\, \text{d}x
\]}
{\(2x +\ln |x| + c,\ c\in \mathbb{R}\)}{\(\ln |x| + c,\ c\in \mathbb{R}\)}{\(2 +\ln |x| + c,\ c\in \mathbb{R}\)}{\(2x^{2} +\ln |x| + c,\ c\in \mathbb{R}\)}{}{}}
\End

\Data\ID{9000065908}
\Updated{2020-07-15 03:07:37}
\Level{A}
\SubArea{Primitive function}
\LANGen{
\Question{
Given the function 
\[ 
 F(x) = \frac{1} 
{2}x^{2} - x,
\]
find the function \(f\) 
such that \(F\) is
primitive to \(f\) 
on \((1;+\infty )\).}
{\(f(x) = \frac{x^{2}-1} 
 {x+1} \)}{\(f(x) = \frac{x^{2}-1} 
 {x-1} \)}{\(f(x) = \frac{x+1} 
{x^{2}-1}\)}{\(f(x) = \frac{x-1} 
{x^{2}-1}\)}{}{}}
\LANGpl{
\Question{
Dana jest funkcja 
\[ 
 F(x) = \frac{1} 
{2}x^{2} - x
\]
Wybierz funkcję $f$, dla której funkcja $F$ jest pierwotna
w przedziale $(1;+\infty)$.}
{\(f(x) = \frac{x^{2}-1} 
 {x+1} \)}{\(f(x) = \frac{x^{2}-1} 
 {x-1} \)}{\(f(x) = \frac{x+1} 
{x^{2}-1}\)}{\(f(x) = \frac{x-1} 
{x^{2}-1}\)}{}{}}
\LANGcs{
\Question{
Je dána funkce \(F(x) = \frac{1} 
{2}x^{2} - x\).
Vyberte funkci \(f\),
pro niž je \(F\) funkcí
primitivní na intervalu \((1;+\infty )\).}
{\(f(x) = \frac{x^{2}-1} 
 {x+1} \)}{\(f(x) = \frac{x^{2}-1} 
 {x-1} \)}{\(f(x) = \frac{x+1} 
{x^{2}-1}\)}{\(f(x) = \frac{x-1} 
{x^{2}-1}\)}{}{}}
\LANGsk{
\Question{
Je daná funkcia \(F(x) = \frac{1} 
{2}x^{2} - x\).
Vyberte funkciu \(f\),
pre ktorú je \(F\) funkciou
primitívnou na intervale \((1;+\infty )\).}
{\(f(x) = \frac{x^{2}-1} 
 {x+1} \)}{\(f(x) = \frac{x^{2}-1} 
 {x-1} \)}{\(f(x) = \frac{x+1} 
{x^{2}-1}\)}{\(f(x) = \frac{x-1} 
{x^{2}-1}\)}{}{}}
\LANGes{
\Question{
Dada la función
\[ 
 F(x) = \frac{1} 
{2}x^{2} - x,
\]
encuentra la función \(f\) 
tal que \(F\) es
primitiva de \(f\) 
en \((1;+\infty )\).}
{\(f(x) = \frac{x^{2}-1} 
 {x+1} \)}{\(f(x) = \frac{x^{2}-1} 
 {x-1} \)}{\(f(x) = \frac{x+1} 
{x^{2}-1}\)}{\(f(x) = \frac{x-1} 
{x^{2}-1}\)}{}{}}
\End

\Data\ID{9000065909}
\Updated{2020-07-15 03:07:37}
\Level{A}
\SubArea{Primitive function}
\LANGen{
\Question{
Given function 
\[ 
 F(x) = 2\ln |x + 1|,
\]
find the function \(f\) 
such that \(F\) is
primitive to \(f\) 
on \((-1;+\infty )\).}
{\(f(x) = \frac{2} 
{x+1}\)}{\(f(x) = 2\mathrm{e}^{x+1}\)}{\(f(x) = \frac{1} 
{2(x+1)}\)}{\(f(x) = \frac{2} 
{2x+2}\)}{}{}}
\LANGpl{
\Question{
Dana jest funkcja 
\[ 
 F(x) = 2\ln |x + 1|,
\]
Wybierz funkcję $f$, dla której funkcja $F$ jest pierwotna
w przedziale \((-1;+\infty )\).}
{\(f(x) = \frac{2} 
{x+1}\)}{\(f(x) = 2\mathrm{e}^{x+1}\)}{\(f(x) = \frac{1} 
{2(x+1)}\)}{\(f(x) = \frac{2} 
{2x+2}\)}{}{}}
\LANGcs{
\Question{
Je dána funkce \(F(x) = 2\ln |x + 1|\).
Vyberte funkci \(f\),
pro niž je \(F\) funkcí
primitivní na intervalu \((-1;+\infty )\).}
{\(f(x) = \frac{2} 
{x+1}\)}{\(f(x) = 2\mathrm{e}^{x+1}\)}{\(f(x) = \frac{1} 
{2(x+1)}\)}{\(f(x) = \frac{2} 
{2x+2}\)}{}{}}
\LANGsk{
\Question{
Je daná funkcia \(F(x) = 2\ln |x + 1|\).
Vyberte funkciu \(f\),
pre ktorú je \(F\) funkcia
primitívna na intervale \((-1;+\infty )\).}
{\(f(x) = \frac{2} 
{x+1}\)}{\(f(x) = 2\mathrm{e}^{x+1}\)}{\(f(x) = \frac{1} 
{2(x+1)}\)}{\(f(x) = \frac{2} 
{2x+2}\)}{}{}}
\LANGes{
\Question{
Dada la función
\[ 
 F(x) = 2\ln |x + 1|,
\]
encuentra la función \(f\) 
tal que \(F\) es
primitiva de \(f\) 
en \((-1;+\infty )\).}
{\(f(x) = \frac{2} 
{x+1}\)}{\(f(x) = 2\mathrm{e}^{x+1}\)}{\(f(x) = \frac{1} 
{2(x+1)}\)}{\(f(x) = \frac{2} 
{2x+2}\)}{}{}}
\End

\Data\ID{9000065910}
\Updated{2020-07-15 03:07:37}
\Level{A}
\SubArea{Primitive function}
\LANGen{
\Question{
Given function 
\[ 
 F(x) = x + 2\ln |x|-\frac{1} 
{x},
\]
find the function \(f\) 
such that \(F\) is
primitive to \(f\) 
on \((0;+\infty )\).}
{\(f(x) = \frac{x^{2}+2x+1} 
 {x^{2}} \)}{\(f(x) = \frac{x^{2}} 
{(x+1)^{2}} \)}{\(f(x) = \frac{x^{2}-1} 
 {x^{2}} \)}{\(f(x) = \frac{x^{2}} 
{(x-1)^{2}} \)}{}{}}
\LANGpl{
\Question{
Dana jest funkcja 
\[ 
 F(x) = x + 2\ln |x|-\frac{1} 
{x},
\]
Wybierz funkcję $f$, dla której funkcja $F$ jest pierwotna
w przedziale \((0;+\infty )\).}
{\(f(x) = \frac{x^{2}+2x+1} 
 {x^{2}} \)}{\(f(x) = \frac{x^{2}} 
{(x+1)^{2}} \)}{\(f(x) = \frac{x^{2}-1} 
 {x^{2}} \)}{\(f(x) = \frac{x^{2}} 
{(x-1)^{2}} \)}{}{}}
\LANGcs{
\Question{
Je dána funkce \(F(x) = x + 2\ln |x|-\frac{1} 
{x}\).
Vyberte funkci \(f\),
pro niž je \(F\) funkcí
primitivní na intervalu \((0;+\infty )\).}
{\(f(x) = \frac{x^{2}+2x+1} 
 {x^{2}} \)}{\(f(x) = \frac{x^{2}} 
{(x+1)^{2}} \)}{\(f(x) = \frac{x^{2}-1} 
 {x^{2}} \)}{\(f(x) = \frac{x^{2}} 
{(x-1)^{2}} \)}{}{}}
\LANGsk{
\Question{
Je daná funkcia \(F(x) = x + 2\ln |x|-\frac{1} 
{x}\).
Vyberte funkciu \(f\),
pre ktorú je \(F\) funkcia
primitívna na intervale \((0;+\infty )\).}
{\(f(x) = \frac{x^{2}+2x+1} 
 {x^{2}} \)}{\(f(x) = \frac{x^{2}} 
{(x+1)^{2}} \)}{\(f(x) = \frac{x^{2}-1} 
 {x^{2}} \)}{\(f(x) = \frac{x^{2}} 
{(x-1)^{2}} \)}{}{}}
\LANGes{
\Question{
Dada la función
\[ 
 F(x) = x + 2\ln |x|-\frac{1} 
{x},
\]
encuentra la función \(f\) 
tal que \(F\) es
primitiva de \(f\) 
en \((0;+\infty )\).}
{\(f(x) = \frac{x^{2}+2x+1} 
 {x^{2}} \)}{\(f(x) = \frac{x^{2}} 
{(x+1)^{2}} \)}{\(f(x) = \frac{x^{2}-1} 
 {x^{2}} \)}{\(f(x) = \frac{x^{2}} 
{(x-1)^{2}} \)}{}{}}
\End

\Data\ID{9000071204}
\Updated{2020-07-15 03:07:37}
\Level{A}
\SubArea{Primitive function}
\LANGen{
\Question{
Evaluate the following integral on the interval \((0;+\infty)\). 
\[ 
 \int \left (2e^{x} -\frac{3} 
{x}\right )\, \mathrm{d}x
\]}
{\(2e^{x} - 3\ln \left |x\right | + c,\ c\in \mathbb{R}\)}{\(2\ln \left |x\right |- \frac{3} 
{2x^{2}} + c,\ c\in \mathbb{R}\)}{\(2e^{x} - 3 + c,\ c\in \mathbb{R}\)}{}{}{}}
\LANGpl{
\Question{
Wyznacz całkę na przedziale  \((0;+\infty)\). 
\[ 
 \int \left (2e^{x} -\frac{3} 
{x}\right )\, \mathrm{d}x
\]}
{\(2e^{x} - 3\ln \left |x\right | + c,\ c\in \mathbb{R}\)}{\(2\ln \left |x\right |- \frac{3} 
{2x^{2}} + c,\ c\in \mathbb{R}\)}{\(2e^{x} - 3 + c,\ c\in \mathbb{R}\)}{}{}{}}
\LANGcs{
\Question{
Vypočtěte \(\int \left (2e^{x} -\frac{3} 
{x}\right )\, \mathrm{d}x\) na intervalu \((0;+\infty)\).}
{\(2e^{x} - 3\ln \left |x\right | + c,\ c\in \mathbb{R}\)}{\(2\ln \left |x\right |- \frac{3} 
{2x^{2}} + c,\ c\in \mathbb{R}\)}{\(2e^{x} - 3 + c,\ c\in \mathbb{R}\)}{}{}{}}
\LANGsk{
\Question{
Vypočítajte \(\int \left (2e^{x} -\frac{3} 
{x}\right )\, \mathrm{d}x\) na intervale \((0;+\infty)\).}
{\(2e^{x} - 3\ln \left |x\right | + c,\ c\in \mathbb{R}\)}{\(2\ln \left |x\right |- \frac{3} 
{2x^{2}} + c,\ c\in \mathbb{R}\)}{\(2e^{x} - 3 + c,\ c\in \mathbb{R}\)}{}{}{}}
\LANGes{
\Question{
Evalúa la siguiente integral en el intervalo \((0;+\infty)\). 
\[ 
 \int \left (2e^{x} -\frac{3} 
{x}\right )\, \mathrm{d}x
\]}
{\(2e^{x} - 3\ln \left |x\right | + c,\ c\in \mathbb{R}\)}{\(2\ln \left |x\right |- \frac{3} 
{2x^{2}} + c,\ c\in \mathbb{R}\)}{\(2e^{x} - 3 + c,\ c\in \mathbb{R}\)}{}{}{}}
\End

\Data\ID{9000071205}
\Updated{2020-07-15 03:07:37}
\Level{A}
\SubArea{Primitive function}
\LANGen{
\Question{
Evaluate the following integral on \(\mathbb{R}\). 
\[ 
 \int \left (x^{2} + 2^{x}\right )\, \mathrm{d}x
\]}
{\(\frac{x^{3}} 
 {3} + \frac{2^{x}} 
 {\ln 2} + c,\ c\in \mathbb{R}\)}{\(\frac{x^{3}} 
 {3} + \frac{2^{x+1}} 
{x+1} + c,\ c\in \mathbb{R}\)}{\(2x + \frac{2^{x}} 
{\ln \left |x\right |} + c,\ c\in \mathbb{R}\)}{}{}{}}
\LANGpl{
\Question{
Wyznacz całkę.
\[ 
 \int \left (x^{2} + 2^{x}\right )\, \mathrm{d}x
\]}
{\(\frac{x^{3}} 
 {3} + \frac{2^{x}} 
 {\ln 2} + c,\ c\in \mathbb{R}\)}{\(\frac{x^{3}} 
 {3} + \frac{2^{x+1}} 
{x+1} + c,\ c\in \mathbb{R}\)}{\(2x + \frac{2^{x}} 
{\ln \left |x\right |} + c,\ c\in \mathbb{R}\)}{}{}{}}
\LANGcs{
\Question{
Vypočtěte \(\int \left (x^{2} + 2^{x}\right )\, \mathrm{d}x\) na \(\mathbb{R}\).}
{\(\frac{x^{3}} 
 {3} + \frac{2^{x}} 
 {\ln 2} + c,\ c\in \mathbb{R}\)}{\(\frac{x^{3}} 
 {3} + \frac{2^{x+1}} 
{x+1} + c,\ c\in \mathbb{R}\)}{\(2x + \frac{2^{x}} 
{\ln \left |x\right |} + c,\ c\in \mathbb{R}\)}{}{}{}}
\LANGsk{
\Question{
Vypočítajte \(\int \left (x^{2} + 2^{x}\right )\, \mathrm{d}x\) na \(\mathbb{R}\).}
{\(\frac{x^{3}} 
 {3} + \frac{2^{x}} 
 {\ln 2} + c,\ c\in \mathbb{R}\)}{\(\frac{x^{3}} 
 {3} + \frac{2^{x+1}} 
{x+1} + c,\ c\in \mathbb{R}\)}{\(2x + \frac{2^{x}} 
{\ln \left |x\right |} + c,\ c\in \mathbb{R}\)}{}{}{}}
\LANGes{
\Question{
Evalúa la siguiente integral en \(\mathbb{R}\). 
\[ 
 \int \left (x^{2} + 2^{x}\right )\, \mathrm{d}x
\]}
{\(\frac{x^{3}} 
 {3} + \frac{2^{x}} 
 {\ln 2} + c,\ c\in \mathbb{R}\)}{\(\frac{x^{3}} 
 {3} + \frac{2^{x+1}} 
{x+1} + c,\ c\in \mathbb{R}\)}{\(2x + \frac{2^{x}} 
{\ln \left |x\right |} + c,\ c\in \mathbb{R}\)}{}{}{}}
\End

\Data\ID{9000071206}
\Updated{2022-04-23 05:04:14}
\Level{A}
\SubArea{Primitive function}
\LANGen{
\Question{
Given the function \(f(x) =\sin x +\cos x\), find its
primitive function \(F\) so that the graph of \(F\) passes through the point \(A = \left [ \frac{\pi }{2};3\right ]\).}
{\(F(x) =\sin x -\cos x + 2\)}{\(F(x) =\cos x -\sin x + 4\)}{\(F(x) = -\cos x +\sin x + 4\)}{}{}{}}
\LANGpl{
\Question{
Wyznacz pochodną funkcji  \(f\colon y =\sin x +\cos x\), 
 której wykres przechodzi przez punkt
 \(A = \left [ \frac{\pi }{2};3\right ]\).}
{\(F\colon y =\sin x -\cos x + 2\)}{\(F\colon y =\cos x -\sin x + 4\)}{\(F\colon y = -\cos x +\sin x + 4\)}{}{}{}}
\LANGcs{
\Question{
K dané funkci \(f\colon y =\sin x +\cos x\) v
\(\mathbb{R}\) určete primitivní
funkci \(F\), jejíž graf
prochází bodem \(A = \left [ \frac{\pi }{2};3\right ]\).}
{\(F\colon y =\sin x -\cos x + 2\)}{\(F\colon y =\cos x -\sin x + 4\)}{\(F\colon y = -\cos x +\sin x + 4\)}{}{}{}}
\LANGsk{
\Question{
K danej funkcii \(f\colon y =\sin x +\cos x\) v
\(\mathbb{R}\) určte primitívnu
funkciu \(F\), ktorej graf
prechádza bodom \(A = \left [ \frac{\pi }{2};3\right ]\).}
{\(F\colon y =\sin x -\cos x + 2\)}{\(F\colon y =\cos x -\sin x + 4\)}{\(F\colon y = -\cos x +\sin x + 4\)}{}{}{}}
\LANGes{
\Question{
Dada la función \(f\colon y =\sin x +\cos x\), encuentra su función primitiva \(F\) cuya gráfica pasa por el punto \(A = \left [ \frac{\pi }{2};3\right ]\).}
{\(F\colon y =\sin x -\cos x + 2\)}{\(F\colon y =\cos x -\sin x + 4\)}{\(F\colon y = -\cos x +\sin x + 4\)}{}{}{}}
\End

\Data\ID{9000150101}
\Updated{2020-07-15 03:07:37}
\Level{A}
\SubArea{Primitive function}
\LANGen{
\Question{
Evaluate the following integral on \(\mathbb{R}\). 
\[ 
 \int \left (\cos x -\sin x\right )\, \mathrm{d}x
\]}
{\(\sin x +\cos x + c,\ c\in \mathbb{R}\)}{\(\sin x -\cos x + c,\ c\in \mathbb{R}\)}{\(-\sin x +\cos x + c,\ c\in \mathbb{R}\)}{\(-\sin x -\cos x + c,\ c\in \mathbb{R}\)}{}{}}
\LANGpl{
\Question{
Wyznacz całkę.
\[ 
 \int \left (\cos x -\sin x\right )\, \mathrm{d}x
\]}
{\(\sin x +\cos x + c,\ c\in \mathbb{R}\)}{\(\sin x -\cos x + c,\ c\in \mathbb{R}\)}{\(-\sin x +\cos x + c,\ c\in \mathbb{R}\)}{\(-\sin x -\cos x + c,\ c\in \mathbb{R}\)}{}{}}
\LANGcs{
\Question{
Vypočtěte \(\int \left (\cos x -\sin x\right )\, \mathrm{d}x\) na \(\mathbb{R}\).}
{\(\sin x +\cos x + c,\ c\in \mathbb{R}\)}{\(\sin x -\cos x + c,\ c\in \mathbb{R}\)}{\(-\sin x +\cos x + c,\ c\in \mathbb{R}\)}{\(-\sin x -\cos x + c,\ c\in \mathbb{R}\)}{}{}}
\LANGsk{
\Question{
Vypočítajte \(\int \left (\cos x -\sin x\right )\, \mathrm{d}x\) na \(\mathbb{R}\).}
{\(\sin x +\cos x + c,\ c\in \mathbb{R}\)}{\(\sin x -\cos x + c,\ c\in \mathbb{R}\)}{\(-\sin x +\cos x + c,\ c\in \mathbb{R}\)}{\(-\sin x -\cos x + c,\ c\in \mathbb{R}\)}{}{}}
\LANGes{
\Question{
Evalúa la siguiente integral en \(\mathbb{R}\). 
\[ 
 \int \left (\cos x -\sin x\right )\, \mathrm{d}x
\]}
{\(\sin x +\cos x + c,\ c\in \mathbb{R}\)}{\(\sin x -\cos x + c,\ c\in \mathbb{R}\)}{\(-\sin x +\cos x + c,\ c\in \mathbb{R}\)}{\(-\sin x -\cos x + c,\ c\in \mathbb{R}\)}{}{}}
\End

\Data\ID{9000150103}
\Updated{2020-07-15 04:07:14}
\Level{A}
\SubArea{Primitive function}
\LANGen{
\Question{
Evaluate the following integral on the interval \(\left(-\frac{\pi}2;\frac{\pi}2\right)\). 
\[ 
 \int \left ( \frac{3}
{\cos ^{2}x} - 3\mathrm{e}^{x}\right )\, \mathrm{d}x
\]}
{\(3\mathop{\mathrm{tg}}\nolimits x - 3\mathrm{e}^{x} + c,\ c\in \mathbb{R}\)}{\(- 3\mathop{\mathrm{tg}}\nolimits x - 3\mathrm{e}^{x} + c,\ c\in \mathbb{R}\)}{\(- 3\mathop{\mathrm{tg}}\nolimits x + 3\mathrm{e}^{x} + c,\ c\in \mathbb{R}\)}{\(3\mathop{\mathrm{tg}}\nolimits x + 3\mathrm{e}^{x} + c,\ c\in \mathbb{R}\)}{}{}}
\LANGpl{
\Question{
Wyznacz całkę na przedziale  \(\left(-\frac{\pi}2;\frac{\pi}2\right)\). 
\[ 
 \int \left ( \frac{3}
{\cos ^{2}x} - 3\mathrm{e}^{x}\right )\, \mathrm{d}x
\]}
{\(3\mathop{\mathrm{tg}}\nolimits x - 3\mathrm{e}^{x} + c,\ c\in \mathbb{R}\)}{\(- 3\mathop{\mathrm{tg}}\nolimits x - 3\mathrm{e}^{x} + c,\ c\in \mathbb{R}\)}{\(- 3\mathop{\mathrm{tg}}\nolimits x + 3\mathrm{e}^{x} + c,\ c\in \mathbb{R}\)}{\(3\mathop{\mathrm{tg}}\nolimits x + 3\mathrm{e}^{x} + c,\ c\in \mathbb{R}\)}{}{}}
\LANGcs{
\Question{
Vypočtěte \(\int \left ( \frac{3}
{\cos ^{2}x} - 3\mathrm{e}^{x}\right )\, \mathrm{d}x\) na intervalu \(\left(-\frac{\pi}2,\frac{\pi}2\right)\).}
{\(3\mathop{\mathrm{tg}}\nolimits x - 3\mathrm{e}^{x} + c,\ c\in \mathbb{R}\)}{\(- 3\mathop{\mathrm{tg}}\nolimits x - 3\mathrm{e}^{x} + c,\ c\in \mathbb{R}\)}{\(- 3\mathop{\mathrm{tg}}\nolimits x + 3\mathrm{e}^{x} + c,\ c\in \mathbb{R}\)}{\(3\mathop{\mathrm{tg}}\nolimits x + 3\mathrm{e}^{x} + c,\ c\in \mathbb{R}\)}{}{}}
\LANGsk{
\Question{
Vypočítajte \(\int \left ( \frac{3}
{\cos ^{2}x} - 3\mathrm{e}^{x}\right )\, \mathrm{d}x\) na intervale \(\left(-\frac{\pi}2,\frac{\pi}2\right)\).}
{\(3\mathop{\mathrm{tg}}\nolimits x - 3\mathrm{e}^{x} + c,\ c\in \mathbb{R}\)}{\(- 3\mathop{\mathrm{tg}}\nolimits x - 3\mathrm{e}^{x} + c,\ c\in \mathbb{R}\)}{\(- 3\mathop{\mathrm{tg}}\nolimits x + 3\mathrm{e}^{x} + c,\ c\in \mathbb{R}\)}{\(3\mathop{\mathrm{tg}}\nolimits x + 3\mathrm{e}^{x} + c,\ c\in \mathbb{R}\)}{}{}}
\LANGes{
\Question{
Evalúa la siguiente integral en el intervalo \(\left(-\frac{\pi}2;\frac{\pi}2\right)\). 
\[ 
 \int \left ( \frac{3}
{\cos ^{2}x} - 3\mathrm{e}^{x}\right )\, \mathrm{d}x
\]}
{\(3\mathop{\mathrm{tg}}\nolimits x - 3\mathrm{e}^{x} + c,\ c\in \mathbb{R}\)}{\(- 3\mathop{\mathrm{tg}}\nolimits x - 3\mathrm{e}^{x} + c,\ c\in \mathbb{R}\)}{\(- 3\mathop{\mathrm{tg}}\nolimits x + 3\mathrm{e}^{x} + c,\ c\in \mathbb{R}\)}{\(3\mathop{\mathrm{tg}}\nolimits x + 3\mathrm{e}^{x} + c,\ c\in \mathbb{R}\)}{}{}}
\End

\Data\ID{9000150105}
\Updated{2020-07-15 03:07:37}
\Level{A}
\SubArea{Primitive function}
\LANGen{
\Question{
Evaluate the following integral on \(\mathbb{R}\). 
\[ 
 \int \left (6^{x} - 6x^{6}\right )\, \mathrm{d}x
\]}
{\(\frac{6^{x}}
 {\ln 6} -\frac{6x^{7}} 
 {7} + c,\ c\in \mathbb{R}\)}{\(6^{x}\ln 6 - 6x^{7} + c,\ c\in \mathbb{R}\)}{\(6^{x}\ln 6 -\frac{6x^{7}} 
 {7} + c,\ c\in \mathbb{R}\)}{\(\frac{6^{x}}
 {\ln 6} - 6x^{7} + c,\ c\in \mathbb{R}\)}{}{}}
\LANGpl{
\Question{
Wyznacz całkę.
\[ 
 \int \left (6^{x} - 6x^{6}\right )\, \mathrm{d}x
\]}
{\(\frac{6^{x}}
 {\ln 6} -\frac{6x^{7}} 
 {7} + c,\ c\in \mathbb{R}\)}{\(6^{x}\ln 6 - 6x^{7} + c,\ c\in \mathbb{R}\)}{\(6^{x}\ln 6 -\frac{6x^{7}} 
 {7} + c,\ c\in \mathbb{R}\)}{\(\frac{6^{x}}
 {\ln 6} - 6x^{7} + c,\ c\in \mathbb{R}\)}{}{}}
\LANGcs{
\Question{
Vypočtěte \(\int \left (6^{x} - 6x^{6}\right )\, \mathrm{d}x\) na \(\mathbb{R}\).}
{\(\frac{6^{x}}
 {\ln 6} -\frac{6x^{7}} 
 {7} + c,\ c\in \mathbb{R}\)}{\(6^{x}\ln 6 - 6x^{7} + c,\ c\in \mathbb{R}\)}{\(6^{x}\ln 6 -\frac{6x^{7}} 
 {7} + c,\ c\in \mathbb{R}\)}{\(\frac{6^{x}}
 {\ln 6} - 6x^{7} + c,\ c\in \mathbb{R}\)}{}{}}
\LANGsk{
\Question{
Vypočítajte \(\int \left (6^{x} - 6x^{6}\right )\, \mathrm{d}x\) na \(\mathbb{R}\).}
{\(\frac{6^{x}}
 {\ln 6} -\frac{6x^{7}} 
 {7} + c,\ c\in \mathbb{R}\)}{\(6^{x}\ln 6 - 6x^{7} + c,\ c\in \mathbb{R}\)}{\(6^{x}\ln 6 -\frac{6x^{7}} 
 {7} + c,\ c\in \mathbb{R}\)}{\(\frac{6^{x}}
 {\ln 6} - 6x^{7} + c,\ c\in \mathbb{R}\)}{}{}}
\LANGes{
\Question{
Evalúa la siguiente integral en \(\mathbb{R}\). 
\[ 
 \int \left (6^{x} - 6x^{6}\right )\, \mathrm{d}x
\]}
{\(\frac{6^{x}}
 {\ln 6} -\frac{6x^{7}} 
 {7} + c,\ c\in \mathbb{R}\)}{\(6^{x}\ln 6 - 6x^{7} + c,\ c\in \mathbb{R}\)}{\(6^{x}\ln 6 -\frac{6x^{7}} 
 {7} + c,\ c\in \mathbb{R}\)}{\(\frac{6^{x}}
 {\ln 6} - 6x^{7} + c,\ c\in \mathbb{R}\)}{}{}}
\End

\Data\ID{9000150108}
\Updated{2022-04-23 05:04:14}
\Level{A}
\SubArea{Primitive function}
\LANGen{
\Question{
Evaluate the following integral on the interval \((0;+\infty)\). 
\[ 
 \int \left (\frac{3}
{x} - 3x^{-2} + \frac{2} 
{x^{3}}\right )\, \mathrm{d}x
\]}
{\(3\ln |x| + \frac{3} 
{x} - \frac{1} 
{x^{2}} + c,\ c\in \mathbb{R}\)}{\(3\ln |x|-\frac{3} 
{x} - \frac{1} 
{x^{2}} + c,\ c\in \mathbb{R}\)}{\(3\ln |x| + \frac{3} 
{x} + \frac{1} 
{x^{2}} + c,\ c\in \mathbb{R}\)}{\(3\ln |x|-\frac{3} 
{x} + \frac{1} 
{x^{2}} + c,\ c\in \mathbb{R}\)}{}{}}
\LANGpl{
\Question{
Wyznacz całkę na przedziale  \((0;+\infty)\). 
\[ 
 \int \left (\frac{3}
{x} - 3x^{-2} + \frac{2} 
{x^{3}}\right )\, \mathrm{d}x
\]}
{\(3\ln |x| + \frac{3} 
{x} - \frac{1} 
{x^{2}} + c,\ c\in \mathbb{R}\)}{\(3\ln |x|-\frac{3} 
{x} - \frac{1} 
{x^{2}} + c,\ c\in \mathbb{R}\)}{\(3\ln |x| + \frac{3} 
{x} + \frac{1} 
{x^{2}} + c,\ c\in \mathbb{R}\)}{\(3\ln |x|-\frac{3} 
{x} + \frac{1} 
{x^{2}} + c,\ c\in \mathbb{R}\)}{}{}}
\LANGcs{
\Question{
Vypočtěte \(\int \left (\frac{3}
{x} - 3x^{-2} + \frac{2} 
{x^{3}} \right )\, \mathrm{d}x\) na intervalu \((0;+\infty)\).}
{\(3\ln |x| + \frac{3} 
{x} - \frac{1} 
{x^{2}} + c,\ c\in \mathbb{R}\)}{\(3\ln |x|-\frac{3} 
{x} - \frac{1} 
{x^{2}} + c,\ c\in \mathbb{R}\)}{\(3\ln |x| + \frac{3} 
{x} + \frac{1} 
{x^{2}} + c,\ c\in \mathbb{R}\)}{\(3\ln |x|-\frac{3} 
{x} + \frac{1} 
{x^{2}} + c,\ c\in \mathbb{R}\)}{}{}}
\LANGsk{
\Question{
Vypočítajte\(\int \left (\frac{3}
{x} - 3x^{-2} + \frac{2} 
{x^{3}} \right )\, \mathrm{d}x\) na intervale \((0;+\infty)\).}
{\(3\ln |x| + \frac{3} 
{x} - \frac{1} 
{x^{2}} + c,\ c\in \mathbb{R}\)}{\(3\ln |x|-\frac{3} 
{x} - \frac{1} 
{x^{2}} + c,\ c\in \mathbb{R}\)}{\(3\ln |x| + \frac{3} 
{x} + \frac{1} 
{x^{2}} + c,\ c\in \mathbb{R}\)}{\(3\ln |x|-\frac{3} 
{x} + \frac{1} 
{x^{2}} + c,\ c\in \mathbb{R}\)}{}{}}
\LANGes{
\Question{
Evalúa la siguiente integral en el intervalo \((0;+\infty)\). 
\[ 
 \int \left (\frac{3}
{x} - 3x^{-2} + \frac{2} 
{x^{3}}\right )\, \mathrm{d}x
\]}
{\(3\ln |x| + \frac{3} 
{x} - \frac{1} 
{x^{2}} + c,\ c\in \mathbb{R}\)}{\(3\ln |x|-\frac{3} 
{x} - \frac{1} 
{x^{2}} + c,\ c\in \mathbb{R}\)}{\(3\ln |x| + \frac{3} 
{x} + \frac{1} 
{x^{2}} + c,\ c\in \mathbb{R}\)}{\(3\ln |x|-\frac{3} 
{x} + \frac{1} 
{x^{2}} + c,\ c\in \mathbb{R}\)}{}{}}
\End

\Data\ID{9000150301}
\Updated{2020-07-15 03:07:37}
\Level{A}
\SubArea{Primitive function}
\LANGen{
\Question{
Evaluate the following integral on \(\mathbb{R}\). 
\[ 
 \int 9x^{8}\, \text{d}x
\]}
{\(x^{9} + c,\ c\in \mathbb{R}\)}{\(9x^{9} + c,\ c\in \mathbb{R}\)}{\(\frac{x^{9}} 
 {9} + c,\ c\in \mathbb{R}\)}{\(9x^{7} + c,\ c\in \mathbb{R}\)}{}{}}
\LANGpl{
\Question{
Wyznacz całkę.
\[ 
 \int 9x^{8}\, \text{d}x
\]}
{\(x^{9} + c,\ c\in \mathbb{R}\)}{\(9x^{9} + c,\ c\in \mathbb{R}\)}{\(\frac{x^{9}} 
 {9} + c,\ c\in \mathbb{R}\)}{\(9x^{7} + c,\ c\in \mathbb{R}\)}{}{}}
\LANGcs{
\Question{
Vypočtěte \(\int 9x^{8}\, \text{d}x\) na \(\mathbb{R}\).}
{\(x^{9} + c,\ c\in \mathbb{R}\)}{\(9x^{9} + c,\ c\in \mathbb{R}\)}{\(\frac{x^{9}} 
 {9} + c,\ c\in \mathbb{R}\)}{\(9x^{7} + c,\ c\in \mathbb{R}\)}{}{}}
\LANGsk{
\Question{
Vypočítajte \(\int 9x^{8}\, \text{d}x\) na \(\mathbb{R}\).}
{\(x^{9} + c,\ c\in \mathbb{R}\)}{\(9x^{9} + c,\ c\in \mathbb{R}\)}{\(\frac{x^{9}} 
 {9} + c,\ c\in \mathbb{R}\)}{\(9x^{7} + c,\ c\in \mathbb{R}\)}{}{}}
\LANGes{
\Question{
Evalúa la siguiente integral en \(\mathbb{R}\). 
\[ 
 \int 9x^{8}\, \text{d}x
\]}
{\(x^{9} + c,\ c\in \mathbb{R}\)}{\(9x^{9} + c,\ c\in \mathbb{R}\)}{\(\frac{x^{9}} 
 {9} + c,\ c\in \mathbb{R}\)}{\(9x^{7} + c,\ c\in \mathbb{R}\)}{}{}}
\End

\Data\ID{9000150302}
\Updated{2020-07-15 03:07:37}
\Level{A}
\SubArea{Primitive function}
\LANGen{
\Question{
Evaluate the following integral on \(\mathbb{R}\). 
\[ 
 \int 8\sin x\, \text{d}x
\]}
{\(- 8\cos x + c,\ c\in \mathbb{R}\)}{\(8\cos x + c,\ c\in \mathbb{R}\)}{\(8\sin x + c,\ c\in \mathbb{R}\)}{\(- 8\sin x + c,\ c\in \mathbb{R}\)}{}{}}
\LANGpl{
\Question{
Wyznacz całkę.
\[ 
 \int 8\sin x\, \text{d}x
\]}
{\(- 8\cos x + c,\ c\in \mathbb{R}\)}{\(8\cos x + c,\ c\in \mathbb{R}\)}{\(8\sin x + c,\ c\in \mathbb{R}\)}{\(- 8\sin x + c,\ c\in \mathbb{R}\)}{}{}}
\LANGcs{
\Question{
Vypočtěte \(\int 8\sin x\, \text{d}x\) na \(\mathbb{R}\).}
{\(- 8\cos x + c,\ c\in \mathbb{R}\)}{\(8\cos x + c,\ c\in \mathbb{R}\)}{\(8\sin x + c,\ c\in \mathbb{R}\)}{\(- 8\sin x + c,\ c\in \mathbb{R}\)}{}{}}
\LANGsk{
\Question{
Vypočítajte \(\int 8\sin x\, \text{d}x\) na \(\mathbb{R}\).}
{\(- 8\cos x + c,\ c\in \mathbb{R}\)}{\(8\cos x + c,\ c\in \mathbb{R}\)}{\(8\sin x + c,\ c\in \mathbb{R}\)}{\(- 8\sin x + c,\ c\in \mathbb{R}\)}{}{}}
\LANGes{
\Question{
Evalúa la siguiente integral en \(\mathbb{R}\). 
\[ 
 \int 8\sin x\, \text{d}x
\]}
{\(- 8\cos x + c,\ c\in \mathbb{R}\)}{\(8\cos x + c,\ c\in \mathbb{R}\)}{\(8\sin x + c,\ c\in \mathbb{R}\)}{\(- 8\sin x + c,\ c\in \mathbb{R}\)}{}{}}
\End

\Data\ID{9000150303}
\Updated{2020-07-15 03:07:37}
\Level{A}
\SubArea{Primitive function}
\LANGen{
\Question{
Evaluate the following integral on \(\mathbb{R}\). 
\[ 
 \int 9\mathrm{e}^{x}\, \text{d}x
\]}
{\(9\mathrm{e}^{x} + c,\ c\in \mathbb{R}\)}{\(9 -\mathrm{e}^{x} + c,\ c\in \mathbb{R}\)}{\(9 +\mathrm{e} ^{x} + c,\ c\in \mathbb{R}\)}{\(- 9\mathrm{e}^{x} + c,\ c\in \mathbb{R}\)}{}{}}
\LANGpl{
\Question{
Wyznacz całkę.
\[ 
 \int 9\mathrm{e}^{x}\, \text{d}x
\]}
{\(9\mathrm{e}^{x} + c,\ c\in \mathbb{R}\)}{\(9 -\mathrm{e}^{x} + c,\ c\in \mathbb{R}\)}{\(9 +\mathrm{e} ^{x} + c,\ c\in \mathbb{R}\)}{\(- 9\mathrm{e}^{x} + c,\ c\in \mathbb{R}\)}{}{}}
\LANGcs{
\Question{
Vypočtěte \(\int 9\mathrm{e}^{x}\, \text{d}x\) na \(\mathbb{R}\).}
{\(9\mathrm{e}^{x} + c,\ c\in \mathbb{R}\)}{\(9 -\mathrm{e}^{x} + c,\ c\in \mathbb{R}\)}{\(9 +\mathrm{e} ^{x} + c,\ c\in \mathbb{R}\)}{\(- 9\mathrm{e}^{x} + c,\ c\in \mathbb{R}\)}{}{}}
\LANGsk{
\Question{
Vypočítajte \(\int 9\mathrm{e}^{x}\, \text{d}x\) na \(\mathbb{R}\).}
{\(9\mathrm{e}^{x} + c,\ c\in \mathbb{R}\)}{\(9 -\mathrm{e}^{x} + c,\ c\in \mathbb{R}\)}{\(9 +\mathrm{e} ^{x} + c,\ c\in \mathbb{R}\)}{\(- 9\mathrm{e}^{x} + c,\ c\in \mathbb{R}\)}{}{}}
\LANGes{
\Question{
Evalúa la siguiente integral en \(\mathbb{R}\). 
\[ 
 \int 9\mathrm{e}^{x}\, \text{d}x
\]}
{\(9\mathrm{e}^{x} + c,\ c\in \mathbb{R}\)}{\(9 -\mathrm{e}^{x} + c,\ c\in \mathbb{R}\)}{\(9 +\mathrm{e} ^{x} + c,\ c\in \mathbb{R}\)}{\(- 9\mathrm{e}^{x} + c,\ c\in \mathbb{R}\)}{}{}}
\End

\Data\ID{9000150304}
\Updated{2020-07-15 03:07:37}
\Level{A}
\SubArea{Primitive function}
\LANGen{
\Question{
Evaluate the following integral on the interval \((0;+\infty)\). 
\[ 
 \int \frac{5}
{x}\, \text{d}x
\]}
{\(5\ln |x| + c,\ c\in \mathbb{R}\)}{\(5x^{2} + c,\ c\in \mathbb{R}\)}{\(\frac{\ln |x|}
{5} + c,\ c\in \mathbb{R}\)}{\(\frac{5}
{x^{2}} + c,\ c\in \mathbb{R}\)}{}{}}
\LANGpl{
\Question{
Wyznacz całkę na przedziale  \((0;+\infty)\). 
\[ 
 \int \frac{5}
{x}\, \text{d}x
\]}
{\(5\ln |x| + c,\ c\in \mathbb{R}\)}{\(5x^{2} + c,\ c\in \mathbb{R}\)}{\(\frac{\ln |x|}
{5} + c,\ c\in \mathbb{R}\)}{\(\frac{5}
{x^{2}} + c,\ c\in \mathbb{R}\)}{}{}}
\LANGcs{
\Question{
Vypočtěte \(\int \frac{5}
{x}\, \text{d}x\) na intervalu \((0;+\infty)\).}
{\(5\ln |x| + c,\ c\in \mathbb{R}\)}{\(5x^{2} + c,\ c\in \mathbb{R}\)}{\(\frac{\ln |x|}
{5} + c,\ c\in \mathbb{R}\)}{\(\frac{5}
{x^{2}} + c,\ c\in \mathbb{R}\)}{}{}}
\LANGsk{
\Question{
Vypočítajte \(\int \frac{5}
{x}\, \text{d}x\) na intervalu \((0;+\infty)\).}
{\(5\ln |x| + c,\ c\in \mathbb{R}\)}{\(5x^{2} + c,\ c\in \mathbb{R}\)}{\(\frac{\ln |x|}
{5} + c,\ c\in \mathbb{R}\)}{\(\frac{5}
{x^{2}} + c,\ c\in \mathbb{R}\)}{}{}}
\LANGes{
\Question{
Evalúa la siguiente integral en el intervalo \((0;+\infty)\). 
\[ 
 \int \frac{5}
{x}\, \text{d}x
\]}
{\(5\ln |x| + c,\ c\in \mathbb{R}\)}{\(5x^{2} + c,\ c\in \mathbb{R}\)}{\(\frac{\ln |x|}
{5} + c,\ c\in \mathbb{R}\)}{\(\frac{5}
{x^{2}} + c,\ c\in \mathbb{R}\)}{}{}}
\End

\Data\ID{9000150305}
\Updated{2020-07-15 03:07:37}
\Level{A}
\SubArea{Primitive function}
\LANGen{
\Question{
Evaluate the following integral on the interval \(\left(0;\frac{\pi}2\right)\). 
\[ 
 \int \frac{8}
{\cos ^{2}x}\, \text{d}x
\]}
{\(8\mathop{\mathrm{tg}}\nolimits x + c,\ c\in \mathbb{R}\)}{\(- 8\mathop{\mathrm{cotg}}\nolimits x + c,\ c\in \mathbb{R}\)}{\(8\mathop{\mathrm{cotg}}\nolimits x + c,\ c\in \mathbb{R}\)}{\(- 8\mathop{\mathrm{tg}}\nolimits x + c,\ c\in \mathbb{R}\)}{}{}}
\LANGpl{
\Question{
Wyznacz całkę na przedziale \(\left(0;\frac{\pi}2\right)\). 
\[ 
 \int \frac{8}
{\cos ^{2}x}\, \text{d}x
\]}
{\(8\mathop{\mathrm{tg}}\nolimits x + c,\ c\in \mathbb{R}\)}{\(- 8\mathop{\mathrm{cotg}}\nolimits x + c,\ c\in \mathbb{R}\)}{\(8\mathop{\mathrm{cotg}}\nolimits x + c,\ c\in \mathbb{R}\)}{\(- 8\mathop{\mathrm{tg}}\nolimits x + c,\ c\in \mathbb{R}\)}{}{}}
\LANGcs{
\Question{
Vypočtěte \(\int \frac{8}
{\cos ^{2}x}\, \text{d}x\) na intervalu \(\left(0;\frac{\pi}2\right)\).}
{\(8\mathop{\mathrm{tg}}\nolimits x + c,\ c\in \mathbb{R}\)}{\(- 8\mathop{\mathrm{cotg}}\nolimits x + c,\ c\in \mathbb{R}\)}{\(8\mathop{\mathrm{cotg}}\nolimits x + c,\ c\in \mathbb{R}\)}{\(- 8\mathop{\mathrm{tg}}\nolimits x + c,\ c\in \mathbb{R}\)}{}{}}
\LANGsk{
\Question{
Vypočítajte \(\int \frac{8}
{\cos ^{2}x}\, \text{d}x\) na intervale \(\left(0;\frac{\pi}2\right)\).}
{\(8\mathop{\mathrm{tg}}\nolimits x + c,\ c\in \mathbb{R}\)}{\(- 8\mathop{\mathrm{cotg}}\nolimits x + c,\ c\in \mathbb{R}\)}{\(8\mathop{\mathrm{cotg}}\nolimits x + c,\ c\in \mathbb{R}\)}{\(- 8\mathop{\mathrm{tg}}\nolimits x + c,\ c\in \mathbb{R}\)}{}{}}
\LANGes{
\Question{
Evalúa la siguiente integral en el intervalo \(\left(0;\frac{\pi}2\right)\). 
\[ 
 \int \frac{8}
{\cos ^{2}x}\, \text{d}x
\]}
{\(8\mathop{\mathrm{tg}}\nolimits x + c,\ c\in \mathbb{R}\)}{\(- 8\mathop{\mathrm{cotg}}\nolimits x + c,\ c\in \mathbb{R}\)}{\(8\mathop{\mathrm{cotg}}\nolimits x + c,\ c\in \mathbb{R}\)}{\(- 8\mathop{\mathrm{tg}}\nolimits x + c,\ c\in \mathbb{R}\)}{}{}}
\End

\Data\ID{9000150306}
\Updated{2020-07-15 03:07:37}
\Level{A}
\SubArea{Primitive function}
\LANGen{
\Question{
Evaluate the following integral on the interval \((0;+\infty)\). 
\[ 
 \int \frac{9}
{x^{5}}\, \text{d}x
\]}
{\(- \frac{9} 
{4x^{4}} + c,\ c\in \mathbb{R}\)}{\(\frac{9}
{x^{6}} + c,\ c\in \mathbb{R}\)}{\(- \frac{3} 
{2x^{6}} + c,\ c\in \mathbb{R}\)}{\(\frac{9}
{x^{4}} + c,\ c\in \mathbb{R}\)}{}{}}
\LANGpl{
\Question{
Wyznacz całkę na przedziale  \((0;+\infty)\). 
\[ 
 \int \frac{9}
{x^{5}}\, \text{d}x
\]}
{\(- \frac{9} 
{4x^{4}} + c,\ c\in \mathbb{R}\)}{\(\frac{9}
{x^{6}} + c,\ c\in \mathbb{R}\)}{\(- \frac{3} 
{2x^{6}} + c,\ c\in \mathbb{R}\)}{\(\frac{9}
{x^{4}} + c,\ c\in \mathbb{R}\)}{}{}}
\LANGcs{
\Question{
Vypočtěte \(\int \frac{9}
{x^{5}} \, \text{d}x\) na intervalu \((0;+\infty)\).}
{\(- \frac{9} 
{4x^{4}} + c,\ c\in \mathbb{R}\)}{\(\frac{9}
{x^{6}} + c,\ c\in \mathbb{R}\)}{\(- \frac{3} 
{2x^{6}} + c,\ c\in \mathbb{R}\)}{\(\frac{9}
{x^{4}} + c,\ c\in \mathbb{R}\)}{}{}}
\LANGsk{
\Question{
Vypočítajte \(\int \frac{9}
{x^{5}} \, \text{d}x\) na intervale \((0;+\infty)\).}
{\(- \frac{9} 
{4x^{4}} + c,\ c\in \mathbb{R}\)}{\(\frac{9}
{x^{6}} + c,\ c\in \mathbb{R}\)}{\(- \frac{3} 
{2x^{6}} + c,\ c\in \mathbb{R}\)}{\(\frac{9}
{x^{4}} + c,\ c\in \mathbb{R}\)}{}{}}
\LANGes{
\Question{
Evalúa la siguiente integral en el intervalo \((0;+\infty)\). 
\[ 
 \int \frac{9}
{x^{5}}\, \text{d}x
\]}
{\(- \frac{9} 
{4x^{4}} + c,\ c\in \mathbb{R}\)}{\(\frac{9}
{x^{6}} + c,\ c\in \mathbb{R}\)}{\(- \frac{3} 
{2x^{6}} + c,\ c\in \mathbb{R}\)}{\(\frac{9}
{x^{4}} + c,\ c\in \mathbb{R}\)}{}{}}
\End

\Data\ID{9000150307}
\Updated{2020-07-15 03:07:37}
\Level{A}
\SubArea{Primitive function}
\LANGen{
\Question{
Evaluate the following integral on \(\mathbb{R}\). 
\[ 
 \int 8\cdot 5^{x}\, \text{d}x
\]}
{\(\frac{8\cdot 5^{x}}
 {\ln 5} + c,\ c\in \mathbb{R}\)}{\(\frac{8\cdot 5^{x}}
 {\ln x} + c,\ c\in \mathbb{R}\)}{\(8\cdot 5^{x}\cdot \ln 5 + c,\ c\in \mathbb{R}\)}{\(8\cdot 5^{x}\cdot \ln x + c,\ c\in \mathbb{R}\)}{}{}}
\LANGpl{
\Question{
Wyznacz całkę.
\[ 
 \int 8\cdot 5^{x}\, \text{d}x
\]}
{\(\frac{8\cdot 5^{x}}
 {\ln 5} + c,\ c\in \mathbb{R}\)}{\(\frac{8\cdot 5^{x}}
 {\ln x} + c,\ c\in \mathbb{R}\)}{\(8\cdot 5^{x}\cdot \ln 5 + c,\ c\in \mathbb{R}\)}{\(8\cdot 5^{x}\cdot \ln x + c,\ c\in \mathbb{R}\)}{}{}}
\LANGcs{
\Question{
Vypočtěte \(\int 8\cdot 5^{x}\, \text{d}x\) na \(\mathbb{R}\).}
{\(\frac{8\cdot 5^{x}}
 {\ln 5} + c,\ c\in \mathbb{R}\)}{\(\frac{8\cdot 5^{x}}
 {\ln x} + c,\ c\in \mathbb{R}\)}{\(8\cdot 5^{x}\cdot \ln 5 + c,\ c\in \mathbb{R}\)}{\(8\cdot 5^{x}\cdot \ln x + c,\ c\in \mathbb{R}\)}{}{}}
\LANGsk{
\Question{
Vypočítajte \(\int 8\cdot 5^{x}\, \text{d}x\) na \(\mathbb{R}\).}
{\(\frac{8\cdot 5^{x}}
 {\ln 5} + c,\ c\in \mathbb{R}\)}{\(\frac{8\cdot 5^{x}}
 {\ln x} + c,\ c\in \mathbb{R}\)}{\(8\cdot 5^{x}\cdot \ln 5 + c,\ c\in \mathbb{R}\)}{\(8\cdot 5^{x}\cdot \ln x + c,\ c\in \mathbb{R}\)}{}{}}
\LANGes{
\Question{
Evalúa la siguiente integral en \(\mathbb{R}\). 
\[ 
 \int 8\cdot 5^{x}\, \text{d}x
\]}
{\(\frac{8\cdot 5^{x}}
 {\ln 5} + c,\ c\in \mathbb{R}\)}{\(\frac{8\cdot 5^{x}}
 {\ln x} + c,\ c\in \mathbb{R}\)}{\(8\cdot 5^{x}\cdot \ln 5 + c,\ c\in \mathbb{R}\)}{\(8\cdot 5^{x}\cdot \ln x + c,\ c\in \mathbb{R}\)}{}{}}
\End

\Data\ID{9000150308}
\Updated{2020-07-15 03:07:37}
\Level{A}
\SubArea{Primitive function}
\LANGen{
\Question{
Evaluate the following integral on \(\mathbb{R}\). 
\[ 
 \int \frac{\cos x}
{8}\, \text{d}x
\]}
{\(\frac{\sin x}
{8} + c,\ c\in \mathbb{R}\)}{\(-\frac{\sin x} 
{8} + c,\ c\in \mathbb{R}\)}{\(\sin x + c,\ c\in \mathbb{R}\)}{\(-\sin x + c,\ c\in \mathbb{R}\)}{}{}}
\LANGpl{
\Question{
Wyznacz całkę.
\[ 
 \int \frac{\cos x}
{8}\, \text{d}x
\]}
{\(\frac{\sin x}
{8} + c,\ c\in \mathbb{R}\)}{\(-\frac{\sin x} 
{8} + c,\ c\in \mathbb{R}\)}{\(\sin x + c,\ c\in \mathbb{R}\)}{\(-\sin x + c,\ c\in \mathbb{R}\)}{}{}}
\LANGcs{
\Question{
Vypočtěte \(\int \frac{\cos x}
{8} \, \text{d}x\) na \(\mathbb{R}\).}
{\(\frac{\sin x}
{8} + c,\ c\in \mathbb{R}\)}{\(-\frac{\sin x} 
{8} + c,\ c\in \mathbb{R}\)}{\(\sin x + c,\ c\in \mathbb{R}\)}{\(-\sin x + c,\ c\in \mathbb{R}\)}{}{}}
\LANGsk{
\Question{
Vypočítajte \(\int \frac{\cos x}
{8} \, \text{d}x\) na \(\mathbb{R}\).}
{\(\frac{\sin x}
{8} + c,\ c\in \mathbb{R}\)}{\(-\frac{\sin x} 
{8} + c,\ c\in \mathbb{R}\)}{\(\sin x + c,\ c\in \mathbb{R}\)}{\(-\sin x + c,\ c\in \mathbb{R}\)}{}{}}
\LANGes{
\Question{
Evalúa la siguiente integral en \(\mathbb{R}\). 
\[ 
 \int \frac{\cos x}
{8}\, \text{d}x
\]}
{\(\frac{\sin x}
{8} + c,\ c\in \mathbb{R}\)}{\(-\frac{\sin x} 
{8} + c,\ c\in \mathbb{R}\)}{\(\sin x + c,\ c\in \mathbb{R}\)}{\(-\sin x + c,\ c\in \mathbb{R}\)}{}{}}
\End

\Data\ID{1003027301}
\Updated{2022-08-25 10:08:54}
\Level{B}
\SubArea{Primitive function}
\LANGen{
\Question{
Evaluate the following indefinite integral on \( \left(0;\frac{\pi}2 \right) \).
\[ \int\frac{(\sin x+\cos x)^2-1}{\sin x\cos x}\,\mathrm{d}x \]}
{\( 2x+c\text{, }c\in\mathbb{R} \)}{\( x+c\text{, }c\in\mathbb{R} \)}{\( \mathrm{tg}\,x+c\text{, }c\in\mathbb{R} \)}{\( c\text{, }c\in\mathbb{R} \)}{}{}}
\LANGpl{
\Question{
Wyznacz całkę nieoznaczoną na przedziale \( \left(0;\frac{\pi}2 \right) \).
\[ \int\frac{(\sin x+\cos x)^2-1}{\sin x\cos x}\,\mathrm{d}x \]}
{\( 2x+c\text{, }c\in\mathbb{R} \)}{\( x+c\text{, }c\in\mathbb{R} \)}{\( \mathrm{tg}\,x+c\text{, }c\in\mathbb{R} \)}{\( c\text{, }c\in\mathbb{R} \)}{}{}}
\LANGcs{
\Question{
Vypočítejte na intervalu \( \left(0;\frac{\pi}2 \right) \) následující integrál.
\[ \int\frac{(\sin x+\cos x)^2-1}{\sin x\cos x}\,\mathrm{d}x \]}
{\( 2x+c\text{, }c\in\mathbb{R} \)}{\( x+c\text{, }c\in\mathbb{R} \)}{\( \mathrm{tg}\,x+c\text{, }c\in\mathbb{R} \)}{\( c\text{, }c\in\mathbb{R} \)}{}{}}
\LANGsk{
\Question{
Vypočítajte na intervale \( \left(0;\frac{\pi}2 \right) \) následujúci integrál.
\[ \int\frac{(\sin x+\cos x)^2-1}{\sin x\cos x}\,\mathrm{d}x \]}
{\( 2x+c\text{, }c\in\mathbb{R} \)}{\( x+c\text{, }c\in\mathbb{R} \)}{\( \mathrm{tg}\,x+c\text{, }c\in\mathbb{R} \)}{\( c\text{, }c\in\mathbb{R} \)}{}{}}
\LANGes{
\Question{
Evalúa la siguiente integral en el intervalo \( \left(0;\frac{\pi}2 \right) \).
\[ \int\frac{(\sin x+\cos x)^2-1}{\sin x\cos x}\,\mathrm{d}x \]}
{\( 2x+c\text{, }c\in\mathbb{R} \)}{\( x+c\text{, }c\in\mathbb{R} \)}{\( \mathrm{tg}\,x+c\text{, }c\in\mathbb{R} \)}{\( c\text{, }c\in\mathbb{R} \)}{}{}}
\End

\Data\ID{1003027302}
\Updated{2022-08-25 10:08:54}
\Level{B}
\SubArea{Primitive function}
\LANGen{
\Question{
Evaluate the following indefinite integral on \( \left(-\frac{\pi}2;\frac{\pi}2 \right) \).
\[ \int \left(\frac1{\cos x}-\sin x\cdot\mathrm{tg}\,x\right)\,\mathrm{d}x \]}
{\( \sin x+c\text{, }c\in\mathbb{R} \)}{\( \cos x+c\text{, }c\in\mathbb{R} \)}{\( -\sin x+c\text{, }c\in\mathbb{R} \)}{\( -\cos x+c\text{, }c\in\mathbb{R} \)}{}{}}
\LANGpl{
\Question{
Wyznacz całkę nieoznaczoną na przedziale \( \left(-\frac{\pi}2;\frac{\pi}2 \right) \).
\[ \int \left(\frac1{\cos x}-\sin x\cdot\mathrm{tg}\,x\right)\,\mathrm{d}x \]}
{\( \sin x+c\text{, }c\in\mathbb{R} \)}{\( \cos x+c\text{, }c\in\mathbb{R} \)}{\( -\sin x+c\text{, }c\in\mathbb{R} \)}{\( -\cos x+c\text{, }c\in\mathbb{R} \)}{}{}}
\LANGcs{
\Question{
Vypočítejte na intervalu \( \left(-\frac{\pi}2;\frac{\pi}2 \right) \) následující integrál.
\[ \int \left(\frac1{\cos x}-\sin x\cdot\mathrm{tg}\,x\right)\,\mathrm{d}x \]}
{\( \sin x+c\text{, }c\in\mathbb{R} \)}{\( \cos x+c\text{, }c\in\mathbb{R} \)}{\( -\sin x+c\text{, }c\in\mathbb{R} \)}{\( -\cos x+c\text{, }c\in\mathbb{R} \)}{}{}}
\LANGsk{
\Question{
Vypočítajte na intervale \( \left(-\frac{\pi}2;\frac{\pi}2 \right) \) nasledujúci integrál.
\[ \int \left(\frac1{\cos x}-\sin x\cdot\mathrm{tg}\,x\right)\,\mathrm{d}x \]}
{\( \sin x+c\text{, }c\in\mathbb{R} \)}{\( \cos x+c\text{, }c\in\mathbb{R} \)}{\( -\sin x+c\text{, }c\in\mathbb{R} \)}{\( -\cos x+c\text{, }c\in\mathbb{R} \)}{}{}}
\LANGes{
\Question{
Evalúa la siguiente integral en el intervalo \( \left(-\frac{\pi}2;\frac{\pi}2 \right) \).
\[ \int \left(\frac1{\cos x}-\sin x\cdot\mathrm{tg}\,x\right)\,\mathrm{d}x \]}
{\( \sin x+c\text{, }c\in\mathbb{R} \)}{\( \cos x+c\text{, }c\in\mathbb{R} \)}{\( -\sin x+c\text{, }c\in\mathbb{R} \)}{\( -\cos x+c\text{, }c\in\mathbb{R} \)}{}{}}
\End

\Data\ID{1003027303}
\Updated{2022-08-25 10:08:58}
\Level{B}
\SubArea{Primitive function}
\LANGen{
\Question{
Choose the function \( g \) whose derivative on \( (-\pi;0) \)  is the following function:
\[ f(x)=\sin^2x\cdot\mathrm{cotg}^2\,x+\sin^2x-1\]}
{\( g(x)=5 \)}{\( g(x)=x+3 \)}{\( g(x)=\mathrm{tg}\,x \)}{\( g(x)=\sin\,x \)}{}{}}
\LANGpl{
\Question{
Wybierz funkcję  \( g \), której pochodną na przedziale  \( (-\pi;0) \)  jest funkcja:
\[ f(x)=\sin^2x\cdot\mathrm{cotg}^2\,x+\sin^2x-1\]}
{\( g(x)=5 \)}{\( g(x)=x+3 \)}{\( g(x)=\mathrm{tg}\,x \)}{\( g(x)=\sin\,x \)}{}{}}
\LANGcs{
\Question{
Vyber takovou funkci \( g \), jejíž derivací na  \( (-\pi;0) \)  je funkce:
\[ f(x)=\sin^2x\cdot\mathrm{cotg}^2\,x+\sin^2x-1\]}
{\( g(x)=5 \)}{\( g(x)=x+3 \)}{\( g(x)=\mathrm{tg}\,x \)}{\( g(x)=\sin\,x \)}{}{}}
\LANGsk{
\Question{
Vyber takú funkciu \( g \), ktorej derivácia na  \( (-\pi;0) \)  je funkcia:
\[ f(x)=\sin^2x\cdot\mathrm{cotg}^2\,x+\sin^2x-1\]}
{\( g(x)=5 \)}{\( g(x)=x+3 \)}{\( g(x)=\mathrm{tg}\,x \)}{\( g(x)=\sin\,x \)}{}{}}
\LANGes{
\Question{
Elige la función \( g \) cuya derivada en \( (-\pi;0) \) es la siguiente función:
\[ f(x)=\sin^2x\cdot\mathrm{cotg}^2\,x+\sin^2x-1\]}
{\( g(x)=5 \)}{\( g(x)=x+3 \)}{\( g(x)=\mathrm{tg}\,x \)}{\( g(x)=\sin\,x \)}{}{}}
\End

\Data\ID{1003027305}
\Updated{2022-08-25 10:08:58}
\Level{B}
\SubArea{Primitive function}
\LANGen{
\Question{
Choose an incorrect evaluation of the following indefinite integral on \( (0;\infty) \).
\[ \int\left(\sqrt{x}-1\right)\left(\sqrt{x}+1\right)\,\mathrm{d}x \]}
{\( x^2-2x+c\text{, }c\in\mathbb{R} \)}{\( \frac{x^2}2-x+c\text{, }c\in\mathbb{R} \)}{\( \frac{x^2-2x}2+c\text{, }c\in\mathbb{R} \)}{\( \frac{x^4-4x^2}{2x(x+2)}+c\text{, }c\in\mathbb{R} \)}{}{}}
\LANGpl{
\Question{
Wskaż niepoprawne obliczenie całki  w \( (0;\infty) \).
\[ \int\left(\sqrt{x}-1\right)\left(\sqrt{x}+1\right)\,\mathrm{d}x \]}
{\( x^2-2x+c\text{, }c\in\mathbb{R} \)}{\( \frac{x^2}2-x+c\text{, }c\in\mathbb{R} \)}{\( \frac{x^2-2x}2+c\text{, }c\in\mathbb{R} \)}{\( \frac{x^4-4x^2}{2x(x+2)}+c\text{, }c\in\mathbb{R} \)}{}{}}
\LANGcs{
\Question{
Který výpočet daného integrálu na  \( (0;\infty) \) není správný?
\[ \int\left(\sqrt{x}-1\right)\left(\sqrt{x}+1\right)\,\mathrm{d}x \]}
{\( x^2-2x+c\text{, }c\in\mathbb{R} \)}{\( \frac{x^2}2-x+c\text{, }c\in\mathbb{R} \)}{\( \frac{x^2-2x}2+c\text{, }c\in\mathbb{R} \)}{\( \frac{x^4-4x^2}{2x(x+2)}+c\text{, }c\in\mathbb{R} \)}{}{}}
\LANGsk{
\Question{
Ktorý výpočet daného integrálu na  \( (0;\infty) \) nie je správny?
\[ \int\left(\sqrt{x}-1\right)\left(\sqrt{x}+1\right)\,\mathrm{d}x \]}
{\( x^2-2x+c\text{, }c\in\mathbb{R} \)}{\( \frac{x^2}2-x+c\text{, }c\in\mathbb{R} \)}{\( \frac{x^2-2x}2+c\text{, }c\in\mathbb{R} \)}{\( \frac{x^4-4x^2}{2x(x+2)}+c\text{, }c\in\mathbb{R} \)}{}{}}
\LANGes{
\Question{
Elige el resultado incorrecto de la siguiente integral indefinida en \( (0;\infty) \).
\[ \int\left(\sqrt{x}-1\right)\left(\sqrt{x}+1\right)\,\mathrm{d}x \]}
{\( x^2-2x+c\text{, }c\in\mathbb{R} \)}{\( \frac{x^2}2-x+c\text{, }c\in\mathbb{R} \)}{\( \frac{x^2-2x}2+c\text{, }c\in\mathbb{R} \)}{\( \frac{x^4-4x^2}{2x(x+2)}+c\text{, }c\in\mathbb{R} \)}{}{}}
\End

\Data\ID{1003027306}
\Updated{2022-08-25 10:08:58}
\Level{B}
\SubArea{Primitive function}
\LANGen{
\Question{
Evaluate the following indefinite integral on \( \left(3;\infty \right) \).
\[ \int\frac{x^2-5x+6}{x-3}\,\mathrm{d}x \]}
{\( \frac{x^2}2-2x+c\text{, }c\in\mathbb{R} \)}{\( x-2+c\text{, }c\in\mathbb{R} \)}{\( \frac{x^2}2+2x+c\text{, }c\in\mathbb{R} \)}{\( \frac{x^2}2-x+c\text{, }c\in\mathbb{R} \)}{}{}}
\LANGpl{
\Question{
Wyznacz całkę nieoznaczoną na przedziale \( \left(3;\infty \right) \).
\[ \int\frac{x^2-5x+6}{x-3}\,\mathrm{d}x \]}
{\( \frac{x^2}2-2x+c\text{, }c\in\mathbb{R} \)}{\( x-2+c\text{, }c\in\mathbb{R} \)}{\( \frac{x^2}2+2x+c\text{, }c\in\mathbb{R} \)}{\( \frac{x^2}2-x+c\text{, }c\in\mathbb{R} \)}{}{}}
\LANGcs{
\Question{
Vyber správný výpočet následujícího integrálu na \( \left(3;\infty \right) \).
\[ \int\frac{x^2-5x+6}{x-3}\,\mathrm{d}x \]}
{\( \frac{x^2}2-2x+c\text{, }c\in\mathbb{R} \)}{\( x-2+c\text{, }c\in\mathbb{R} \)}{\( \frac{x^2}2+2x+c\text{, }c\in\mathbb{R} \)}{\( \frac{x^2}2-x+c\text{, }c\in\mathbb{R} \)}{}{}}
\LANGsk{
\Question{
Vyber správny výpočet nasledujúceho integrálu na \( \left(3;\infty \right) \).
\[ \int\frac{x^2-5x+6}{x-3}\,\mathrm{d}x \]}
{\( \frac{x^2}2-2x+c\text{, }c\in\mathbb{R} \)}{\( x-2+c\text{, }c\in\mathbb{R} \)}{\( \frac{x^2}2+2x+c\text{, }c\in\mathbb{R} \)}{\( \frac{x^2}2-x+c\text{, }c\in\mathbb{R} \)}{}{}}
\LANGes{
\Question{
Evalúa la siguiente integral en el intervalo \( \left(3;\infty \right) \).
\[ \int\frac{x^2-5x+6}{x-3}\,\mathrm{d}x \]}
{\( \frac{x^2}2-2x+c\text{, }c\in\mathbb{R} \)}{\( x-2+c\text{, }c\in\mathbb{R} \)}{\( \frac{x^2}2+2x+c\text{, }c\in\mathbb{R} \)}{\( \frac{x^2}2-x+c\text{, }c\in\mathbb{R} \)}{}{}}
\End

\Data\ID{1003107804}
\Updated{2021-04-08 07:04:39}
\Level{B}
\SubArea{Primitive function}
\LANGen{
\Question{
Four girls evaluated the integral $I=\int\sin x\cdot\cos x\,\mathrm{d}x$ on $\mathbb{R}$. Ann started to integrate by parts like this: $I=\int\sin x\cdot\cos x\,\mathrm{d}x=\sin^2x-\int\cos x\cdot\sin x\,\mathrm{d}x$. Beth started to integrate by parts like this: $I=\int\sin x\cdot\cos x\,\mathrm{d}x=-\cos^2 x-\int\sin x\cdot\cos x\,\mathrm{d}x$. Claire used substitution $a=\sin x$ like this: $I=\int\sin x\cdot\cos x\,\mathrm{d}x=\int a\,\mathrm{d}a$. Diana integrated $\int\sin x\cdot\cos x\,\mathrm{d}x=-\cos x\cdot\sin x+c$, $c\in\mathbb{R}$. Which of the girls made a mistake?}
{Diana}{Ann}{Beth}{Claire}{}{}}
\LANGpl{
\Question{
Cztery dziewczyny obliczyły całkę  $I=\int\sin x\cdot\cos x\,\mathrm{d}x$ w $\mathbb{R}$. Ania rozpoczęła całkowanie przez części w następujący sposób: $I=\int\sin x\cdot\cos x\,\mathrm{d}x=\sin^2x-\int\cos x\cdot\sin x\,\mathrm{d}x$. Beata rozpoczęła całkowanie przez części w następujący sposób: $I=\int\sin x\cdot\cos x\,\mathrm{d}x=-\cos^2 x-\int\sin x\cdot\cos x\,\mathrm{d}x$. Kasia użyła podstawienia $a=\sin x$ w następujący sposób: $I=\int\sin x\cdot\cos x\,\mathrm{d}x=\int a\,\mathrm{d}a$. Natomiast Daria wykonała  $\int\sin x\cdot\cos x\,\mathrm{d}x=-\cos x\cdot\sin x+c$, $c\in\mathbb{R}$. Która z dziewczyn popełniła błąd?}
{Diana}{Ann}{Beth}{Claire}{}{}}
\LANGcs{
\Question{
Čtyři dívky počítaly neurčitý integrál $I=\int\sin x\cdot\cos x\,\mathrm{d}x$ na $\mathbb{R}$. Anička začala integrovat metodou per partes takto: $I=\int\sin x\cdot\cos x\,\mathrm{d}x=\sin^2x-\int\cos x\cdot\sin x\,\mathrm{d}x$. Bětka integrovala rovněž metodou per partes, ale takto: $I=\int\sin x\cdot\cos x\,\mathrm{d}x=-\cos^2 x-\int\sin x\cdot\cos x\,\mathrm{d}x$. Klára použila substituci $a=\sin x$ takto: $I=\int\sin x\cdot\cos x\,\mathrm{d}x=\int a\,\mathrm{d}a$. Diana integrovala přímo $\int\sin x\cdot\cos x\,\mathrm{d}x=-\cos x\cdot\sin x+c$, $c\in\mathbb{R}$. Která z dívek udělala chybu?}
{Diana}{Anička}{Bětka}{Klára}{}{}}
\LANGsk{
\Question{
Štyri dievčatá počítali neurčitý integrál $I=\int\sin x\cdot\cos x\,\mathrm{d}x$ na $\mathbb{R}$. Anička začala integrovať metódou per partes takto: $I=\int\sin x\cdot\cos x\,\mathrm{d}x=\sin^2x-\int\cos x\cdot\sin x\,\mathrm{d}x$. Betka integrovala rovnako metódou per partes, ale takto: $I=\int\sin x\cdot\cos x\,\mathrm{d}x=-\cos^2 x-\int\sin x\cdot\cos x\,\mathrm{d}x$. Klára použila substitúciu $a=\sin x$ takto: $I=\int\sin x\cdot\cos x\,\mathrm{d}x=\int a\,\mathrm{d}a$. Diana integrovala priamo $\int\sin x\cdot\cos x\,\mathrm{d}x=-\cos x\cdot\sin x+c$, $c\in\mathbb{R}$. Ktorá z dievčat urobila chybu?}
{Diana}{Anička}{Betka}{Klára}{}{}}
\LANGes{
\Question{
Cuatro chicas evaluaron la integral $I=\int\sin x\cdot\cos x\,\mathrm{d}x$ en $\mathbb{R}$. Ana empezó a integrar por partes de esta manera: $I=\int\sin x\cdot\cos x\,\mathrm{d}x=\sin^2x-\int\cos x\cdot\sin x\,\mathrm{d}x$. Beatriz empezó a integrar por partes de esta manera: $I=\int\sin x\cdot\cos x\,\mathrm{d}x=-\cos^2 x-\int\sin x\cdot\cos x\,\mathrm{d}x$. Clara usó la sustitución $a=\sin x$ de la manera siguiente: $I=\int\sin x\cdot\cos x\,\mathrm{d}x=\int a\,\mathrm{d}a$. Diana integró $\int\sin x\cdot\cos x\,\mathrm{d}x=-\cos x\cdot\sin x+c$, $c\in\mathbb{R}$. ¿Cuál de las cuatro chicas cometió un error?}
{Diana}{Ana}{Beatriz}{Clara}{}{}}
\End

\Data\ID{1003107808}
\Updated{2022-04-23 05:04:08}
\Level{B}
\SubArea{Primitive function}
\LANGen{
\Question{
Find the indefinite integral  $\int\frac{\ln^5x}x\,\mathrm{d}x$ on $(0;\infty)$. Use a substitution $a=\ln x$.}
{$\frac{\ln^6x}6+c$, $c\in\mathbb{R}$}{$5\ln^4x+c$, $c\in\mathbb{R}$}{$\frac{\ln^2x}2+c$, $c\in\mathbb{R}$}{$\frac12\ln^5x+c$, $c\in\mathbb{R}$}{}{}}
\LANGpl{
\Question{
Wyznacz całkę nieoznaczoną  $\int\frac{\ln^5x}x\,\mathrm{d}x$ w $(0;\infty)$. Użyj podstawienia $a=\ln x$.}
{$\frac{\ln^6x}6+c$, $c\in\mathbb{R}$}{$5\ln^4x+c$, $c\in\mathbb{R}$}{$\frac{\ln^2x}2+c$, $c\in\mathbb{R}$}{$\frac12\ln^5x+c$, $c\in\mathbb{R}$}{}{}}
\LANGcs{
\Question{
Řešte užitím substituce $a=\ln x$ neurčitý integrál $\int\frac{\ln^5x}x\,\mathrm{d}x$ na $(0;\infty)$.}
{$\frac{\ln^6x}6+c$, $c\in\mathbb{R}$}{$5\ln^4x+c$, $c\in\mathbb{R}$}{$\frac{\ln^2x}2+c$, $c\in\mathbb{R}$}{$\frac12\ln^5x+c$, $c\in\mathbb{R}$}{}{}}
\LANGsk{
\Question{
Riešte použitím substitúcie $a=\ln x$ neurčitý integrál $\int\frac{\ln^5x}x\,\mathrm{d}x$ na $(0;\infty)$.}
{$\frac{\ln^6x}6+c$, $c\in\mathbb{R}$}{$5\ln^4x+c$, $c\in\mathbb{R}$}{$\frac{\ln^2x}2+c$, $c\in\mathbb{R}$}{$\frac12\ln^5x+c$, $c\in\mathbb{R}$}{}{}}
\LANGes{
\Question{
Resuelve la integral indefinida $\int\frac{\ln^5x}x\,\mathrm{d}x$ en $(0;\infty)$ por sustitución $a=\ln x$.}
{$\frac{\ln^6x}6+c$, $c\in\mathbb{R}$}{$5\ln^4x+c$, $c\in\mathbb{R}$}{$\frac{\ln^2x}2+c$, $c\in\mathbb{R}$}{$\frac12\ln^5x+c$, $c\in\mathbb{R}$}{}{}}
\End

\Data\ID{1003107809}
\Updated{2020-07-15 03:07:12}
\Level{B}
\SubArea{Primitive function}
\LANGen{
\Question{
Solve an indefinite integral $\int\frac{8x^7-30x^5}{x^8-5x^6+2}\,\mathrm{d}x$ on $(3;\infty)$.}
{$\ln\left|x^8-5x^6+2\right|+c$, $c\in\mathbb{R}$}{$\ln\left|8x^7-30x^5+2x\right|+c$, $c\in\mathbb{R}$}{$\log\left|x^8-5x^6+2\right|+c$, $c\in\mathbb{R}$}{$\log\left|8x^7-30x^5+2x\right|+c$, $c\in\mathbb{R}$}{}{}}
\LANGpl{
\Question{
Oblicz całkę nieoznaczoną  $\int\frac{8x^7-30x^5}{x^8-5x^6+2}\,\mathrm{d}x$ w $(3;\infty)$.}
{$\ln\left|x^8-5x^6+2\right|+c$, $c\in\mathbb{R}$}{$\ln\left|8x^7-30x^5+2x\right|+c$, $c\in\mathbb{R}$}{$\log\left|x^8-5x^6+2\right|+c$, $c\in\mathbb{R}$}{$\log\left|8x^7-30x^5+2x\right|+c$, $c\in\mathbb{R}$}{}{}}
\LANGcs{
\Question{
Řešte neurčitý integrál $\int\frac{8x^7-30x^5}{x^8-5x^6+2}\,\mathrm{d}x$ na $(3;\infty)$.}
{$\ln\left|x^8-5x^6+2\right|+c$, $c\in\mathbb{R}$}{$\ln\left|8x^7-30x^5+2x\right|+c$, $c\in\mathbb{R}$}{$\log\left|x^8-5x^6+2\right|+c$, $c\in\mathbb{R}$}{$\log\left|8x^7-30x^5+2x\right|+c$, $c\in\mathbb{R}$}{}{}}
\LANGsk{
\Question{
Riešte neurčitý integrál $\int\frac{8x^7-30x^5}{x^8-5x^6+2}\,\mathrm{d}x$ na $(3;\infty)$.}
{$\ln\left|x^8-5x^6+2\right|+c$, $c\in\mathbb{R}$}{$\ln\left|8x^7-30x^5+2x\right|+c$, $c\in\mathbb{R}$}{$\log\left|x^8-5x^6+2\right|+c$, $c\in\mathbb{R}$}{$\log\left|8x^7-30x^5+2x\right|+c$, $c\in\mathbb{R}$}{}{}}
\LANGes{
\Question{
Resuelve la integral indefinida $\int\frac{8x^7-30x^5}{x^8-5x^6+2}\,\mathrm{d}x$ en $(3;\infty)$.}
{$\ln\left|x^8-5x^6+2\right|+c$, $c\in\mathbb{R}$}{$\ln\left|8x^7-30x^5+2x\right|+c$, $c\in\mathbb{R}$}{$\log\left|x^8-5x^6+2\right|+c$, $c\in\mathbb{R}$}{$\log\left|8x^7-30x^5+2x\right|+c$, $c\in\mathbb{R}$}{}{}}
\End

\Data\ID{1103107802}
\Updated{2020-07-15 03:07:12}
\Level{B}
\SubArea{Primitive function}
\IMG{1103107802.pdf}
\LANGen{
\Question{
What is the color of the graph of a function $F(x)$, that is primitive to the function $f(x)=\frac2{x+1}$ on $(-1;\infty)$? (See the picture.)}
{yellow}{green}{blue}{red}{}{}}
\LANGpl{
\Question{
Który kolor wykresu funkcji $F(x)$,  jest funkcją pierwotną funkcji $f(x)=\frac2{x+1}$ w $(-1;\infty)$? (Spójrz na rysunek.)}
{żółty}{zielony}{niebieski}{czerwony}{}{}}
\LANGcs{
\Question{
Jakou barvu má na obrázku graf funkce $F(x)$, která je primitivní k funkci $f(x)=\frac2{x+1}$  na $(-1;\infty)$?}
{žlutou}{zelenou}{modrou}{červenou}{}{}}
\LANGsk{
\Question{
Akú farbu má na obrázku graf funkcie $F(x)$, ktorá je primitívna k funkcii $f(x)=\frac2{x+1}$  na $(-1;\infty)$?}
{žltú}{zelenú}{modrú}{červenú}{}{}}
\LANGes{
\Question{
¿Qué color tiene la gráfica de la función $F(x)$, que es primitiva de la función $f(x)=\frac2{x+1}$ en $(-1;\infty)$? (Mira la imagen.)}
{amarillo}{verde}{azul}{rojo}{}{}}
\End

\Data\ID{1103107803}
\Updated{2020-07-15 03:07:10}
\Level{B}
\SubArea{Primitive function}
\LANGen{
\Question{
Which of the following pictures shows the  graph of a function $F(x)$, that is primitive to the function $f(x)=\frac{x^2-11x+28}{x-7}$ on $(-\infty;7)$?}
{}{}{}{}{}{}}
\LANGpl{
\Question{
Na którym rysunku znajduje się wykres funkcji $F(x)$, funkcji pierwotnej do funkcji  $f(x)=\frac{x^2-11x+28}{x-7}$ w $(-\infty;7)$?}
{}{}{}{}{}{}}
\LANGcs{
\Question{
Na kterém z uvedených obrázků je znázorněn graf funkce $F(x)$ primitivní k $f(x)=\frac{x^2-11x+28}{x-7}$ na $(-\infty;7)$?}
{}{}{}{}{}{}}
\LANGsk{
\Question{
Na ktorom z uvedených obrázkov je znázornený graf funkcie $F(x)$ primitívny k $f(x)=\frac{x^2-11x+28}{x-7}$ na $(-\infty;7)$?}
{}{}{}{}{}{}}
\LANGes{
\Question{
¿Cuál de las siguientes imágenes representa la gráfica de la función $F(x)$, que es primitiva de la función $f(x)=\frac{x^2-11x+28}{x-7}$ en $(-\infty;7)$?}
{}{}{}{}{}{}}

\IMGanswer{1103107803a_0.pdf}
\IMGanswer{1103107803b_0.pdf}
\IMGanswer{1103107803c_0.pdf}
\IMGanswer{1103107803d_0.pdf}\End

\Data\ID{2010000301}
\Updated{2022-08-25 10:08:55}
\Level{B}
\SubArea{Primitive function}
\LANGen{
\Question{
Evaluate the following integral on the interval \(\left(0;\frac{\pi}2\right)\). 
\[ 
 \int \frac{\cos 2x}
{\cos ^{2}x}\, \mathrm{d}x
\]}
{\(2x -\mathop{\mathrm{tg}}\nolimits x + c,\ c\in \mathbb{R}\)}{\(\frac{\sin 2x}
{\frac{1} 
{3} \sin ^{3}x} + c,\ c\in \mathbb{R}\)}{\(2x +\mathop{\mathrm{cotg}}\nolimits x  + c,\ c\in \mathbb{R}\)}{}{}{}}
\LANGpl{
\Question{
Oblicz następującą całkę na przedziale \(\left(0;\frac{\pi}2\right)\). 
\[ 
 \int \frac{\cos 2x}
{\cos ^{2}x}\, \mathrm{d}x
\]}
{\(2x -\mathop{\mathrm{tg}}\nolimits x + c,\ c\in \mathbb{R}\)}{\(\frac{\sin 2x}
{\frac{1} 
{3} \sin ^{3}x} + c,\ c\in \mathbb{R}\)}{\(2x +\mathop{\mathrm{cotg}}\nolimits x  + c,\ c\in \mathbb{R}\)}{}{}{}}
\LANGcs{
\Question{
Vypočtěte \[ \int \frac{\cos 2x}
{\cos ^{2}x}\, \mathrm{d}x
\] na intervalu \(\left(0;\frac{\pi}2\right)\).}
{\(2x -\mathop{\mathrm{tg}}\nolimits x + c,\ c\in \mathbb{R}\)}{\(\frac{\sin 2x}
{\frac{1} 
{3} \sin ^{3}x} + c,\ c\in \mathbb{R}\)}{\(2x +\mathop{\mathrm{cotg}}\nolimits x  + c,\ c\in \mathbb{R}\)}{}{}{}}
\LANGsk{
\Question{
Vypočítajte  
\[ 
 \int \frac{\cos 2x}
{\cos ^{2}x}\, \mathrm{d}x
\]
na intervale \(\left(0;\frac{\pi}2\right)\).}
{\(2x -\mathop{\mathrm{tg}}\nolimits x + c,\ c\in \mathbb{R}\)}{\(\frac{\sin 2x}
{\frac{1} 
{3} \sin ^{3}x} + c,\ c\in \mathbb{R}\)}{\(2x +\mathop{\mathrm{cotg}}\nolimits x  + c,\ c\in \mathbb{R}\)}{}{}{}}
\LANGes{
\Question{
Evalúa la siguiente integral en el intervalo \(\left(0;\frac{\pi}2\right)\). 
\[ 
 \int \frac{\cos 2x}
{\cos ^{2}x}\, \mathrm{d}x
\]}
{\(2x -\mathop{\mathrm{tg}}\nolimits x + c,\ c\in \mathbb{R}\)}{\(\frac{\sin 2x}
{\frac{1} 
{3} \sin ^{3}x} + c,\ c\in \mathbb{R}\)}{\(2x +\mathop{\mathrm{cotg}}\nolimits x  + c,\ c\in \mathbb{R}\)}{}{}{}}
\End

\Data\ID{2010000302}
\Updated{2022-08-25 10:08:55}
\Level{B}
\SubArea{Primitive function}
\LANGen{
\Question{
Evaluate the following integral on the interval \(\left(\sqrt{\frac34};+\infty\right)\). 
\[ 
 \int \frac{8x}
{(4x^{2} - 3)^{2}}\, \mathrm{d}x
\]}
{\(\frac{1}
{3-4x^{2}} + c,\ c\in \mathbb{R}\)}{\(\frac{4x^{2}} 
{\frac{16}{5}x^{5}-8x^{3}+9x} + c,\ c\in \mathbb{R}\)}{\(\frac{1}
{4x^{2}-3} + c,\ c\in \mathbb{R}\)}{}{}{}}
\LANGpl{
\Question{
Oblicz następującą całkę na przedziale \(\left(\sqrt{\frac34};+\infty\right)\). 
\[ 
 \int \frac{8x}
{(4x^{2} - 3)^{2}}\, \mathrm{d}x
\]}
{\(\frac{1}
{3-4x^{2}} + c,\ c\in \mathbb{R}\)}{\(\frac{4x^{2}} 
{\frac{16}{5}x^{5}-8x^{3}+9x} + c,\ c\in \mathbb{R}\)}{\(\frac{1}
{4x^{2}-3} + c,\ c\in \mathbb{R}\)}{}{}{}}
\LANGcs{
\Question{
Vypočtěte
\[ 
 \int \frac{8x}
{(4x^{2} - 3)^{2}}\, \mathrm{d}x
\] na intervalu
\(\left(\sqrt{\frac34};+\infty\right)\).}
{\(\frac{1}
{3-4x^{2}} + c,\ c\in \mathbb{R}\)}{\(\frac{4x^{2}} 
{\frac{16}{5}x^{5}-8x^{3}+9x} + c,\ c\in \mathbb{R}\)}{\(\frac{1}
{4x^{2}-3} + c,\ c\in \mathbb{R}\)}{}{}{}}
\LANGsk{
\Question{
Vypočítajte 
\[ 
 \int \frac{8x}
{(4x^{2} - 3)^{2}}\, \mathrm{d}x
\]
na intervale \(\left(\sqrt{\frac34};+\infty\right)\).}
{\(\frac{1}
{3-4x^{2}} + c,\ c\in \mathbb{R}\)}{\(\frac{4x^{2}} 
{\frac{16}{5}x^{5}-8x^{3}+9x} + c,\ c\in \mathbb{R}\)}{\(\frac{1}
{4x^{2}-3} + c,\ c\in \mathbb{R}\)}{}{}{}}
\LANGes{
\Question{
Evalúa la siguiente integral en el interval \(\left(\sqrt{\frac34};+\infty\right)\). 
\[ 
 \int \frac{8x}
{(4x^{2} - 3)^{2}}\, \mathrm{d}x
\]}
{\(\frac{1}
{3-4x^{2}} + c,\ c\in \mathbb{R}\)}{\(\frac{4x^{2}} 
{\frac{16}{5}x^{5}-8x^{3}+9x} + c,\ c\in \mathbb{R}\)}{\(\frac{1}
{4x^{2}-3} + c,\ c\in \mathbb{R}\)}{}{}{}}
\End

\Data\ID{2010000303}
\Updated{2022-08-25 10:08:55}
\Level{B}
\SubArea{Primitive function}
\LANGen{
\Question{
Evaluate the following integral on the interval \(\left(\frac54;+\infty\right)\). 
\[ 
 \int \frac{3}
{5 - 4x}\, \mathrm{d}x
\]}
{\(-\frac{3} 
{4}\ln |5 - 4x| + c,\ c\in \mathbb{R}\)}{\(-\frac{3}
{4\cdot \ln |5-4x|} + c,\ c\in \mathbb{R}\)}{\(\frac{3} 
{4}\ln |5 - 4x| + c,\ c\in \mathbb{R}\)}{\( \frac{3} 
{4\cdot \ln |5-4x|} + c,\ c\in \mathbb{R}\)}{}{}}
\LANGpl{
\Question{
Oblicz następującą całkę na przedziale \(\left(\frac54;+\infty\right)\). 
\[ 
 \int \frac{3}
{5 - 4x}\, \mathrm{d}x
\]}
{\(-\frac{3} 
{4}\ln |5 - 4x| + c,\ c\in \mathbb{R}\)}{\(-\frac{3}
{4\cdot \ln |5-4x|} + c,\ c\in \mathbb{R}\)}{\(\frac{3} 
{4}\ln |5 - 4x| + c,\ c\in \mathbb{R}\)}{\( \frac{3} 
{4\cdot \ln |5-4x|} + c,\ c\in \mathbb{R}\)}{}{}}
\LANGcs{
\Question{
Vypočtěte 
\[ 
 \int \frac{3}
{5 - 4x}\, \mathrm{d}x
\] na intervalu \(\left(\frac54;+\infty\right)\).}
{\(-\frac{3} 
{4}\ln |5 - 4x| + c,\ c\in \mathbb{R}\)}{\(-\frac{3}
{4\cdot \ln |5-4x|} + c,\ c\in \mathbb{R}\)}{\(\frac{3} 
{4}\ln |5 - 4x| + c,\ c\in \mathbb{R}\)}{\( \frac{3} 
{4\cdot \ln |5-4x|} + c,\ c\in \mathbb{R}\)}{}{}}
\LANGsk{
\Question{
Vypočítajte 
\[ 
 \int \frac{3}
{5 - 4x}\, \mathrm{d}x
\]
na intervale \(\left(\frac54;+\infty\right)\).}
{\(-\frac{3} 
{4}\ln |5 - 4x| + c,\ c\in \mathbb{R}\)}{\(-\frac{3}
{4\cdot \ln |5-4x|} + c,\ c\in \mathbb{R}\)}{\(\frac{3} 
{4}\ln |5 - 4x| + c,\ c\in \mathbb{R}\)}{\( \frac{3} 
{4\cdot \ln |5-4x|} + c,\ c\in \mathbb{R}\)}{}{}}
\LANGes{
\Question{
Evalúa la siguiente integral en el intervalo \(\left(\frac54;+\infty\right)\). 
\[ 
 \int \frac{3}
{5 - 4x}\, \mathrm{d}x
\]}
{\(-\frac{3} 
{4}\ln |5 - 4x| + c,\ c\in \mathbb{R}\)}{\(-\frac{3}
{4\cdot \ln |5-4x|} + c,\ c\in \mathbb{R}\)}{\(\frac{3} 
{4}\ln |5 - 4x| + c,\ c\in \mathbb{R}\)}{\( \frac{3} 
{4\cdot \ln |5-4x|} + c,\ c\in \mathbb{R}\)}{}{}}
\End

\Data\ID{2010001501}
\Updated{2022-08-25 10:08:23}
\Level{B}
\SubArea{Primitive function}
\LANGen{
\Question{
Find the indefinite integral  $\int\frac{4\ln^3x}x\,\mathrm{d}x$ on $(0;\infty)$. Use a substitution $a=\ln x$.}
{$\ln^4x+c$, $c\in\mathbb{R}$}{$12\ln^2x+c$, $c\in\mathbb{R}$}{$2\ln^2x+c$, $c\in\mathbb{R}$}{$16\ln^4x+c$, $c\in\mathbb{R}$}{}{}}
\LANGpl{
\Question{
Znajdź całkę nieoznaczoną $\int\frac{4\ln^3x}x\,\mathrm{d}x$ na  $(0;\infty)$. Podstaw $a=\ln x$.}
{$\ln^4x+c$, $c\in\mathbb{R}$}{$12\ln^2x+c$, $c\in\mathbb{R}$}{$2\ln^2x+c$, $c\in\mathbb{R}$}{$16\ln^4x+c$, $c\in\mathbb{R}$}{}{}}
\LANGcs{
\Question{
Řešte užitím substituce  $a=\ln x$ neurčitý integrál $\int\frac{4\ln^3x}x\ \mathrm{d}x$ na $(0;\infty)$.}
{$\ln^4x+c$, $c\in\mathbb{R}$}{$12\ln^2x+c$, $c\in\mathbb{R}$}{$2\ln^2x+c$, $c\in\mathbb{R}$}{$16\ln^4x+c$, $c\in\mathbb{R}$}{}{}}
\LANGsk{
\Question{
Využitím substitúcie $a=\ln x$ riešte neurčitý integrál $\int\frac{4\ln^3x}x\,\mathrm{d}x$ na $(0;\infty)$.}
{$\ln^4x+c$, $c\in\mathbb{R}$}{$12\ln^2x+c$, $c\in\mathbb{R}$}{$2\ln^2x+c$, $c\in\mathbb{R}$}{$16\ln^4x+c$, $c\in\mathbb{R}$}{}{}}
\LANGes{
\Question{
Encuentra la integral indefinida $\int\frac{4\ln^3x}x\,\mathrm{d}x$ en $(0;\infty)$. Usa la sustitución $a=\ln x$.}
{$\ln^4x+c$, $c\in\mathbb{R}$}{$12\ln^2x+c$, $c\in\mathbb{R}$}{$2\ln^2x+c$, $c\in\mathbb{R}$}{$16\ln^4x+c$, $c\in\mathbb{R}$}{}{}}
\End

\Data\ID{2010001502}
\Updated{2022-08-25 10:08:23}
\Level{B}
\SubArea{Primitive function}
\LANGen{
\Question{
Evaluate the following integral on \(\mathbb{R}\). 
\[ 
 \int (5x^{2} -2)(2-x^{2} )\, \mathrm{d}x
\]}
{\(-x^5 +4x^3-4x + c,\ c\in \mathbb{R}\)}{\( \left(\frac{5}{3}x^{3}  -2x \right)\left(2x -\frac{x^3}{3}\right)
+ c,\ c\in \mathbb{R}\)}{\(-20x^3 +24x + c,\ c\in \mathbb{R}\)}{\(\frac{4x^3}{3}+ c,\ c\in \mathbb{R}\)}{}{}}
\LANGpl{
\Question{
Oblicz następującą całkę na \(\mathbb{R}\). 
\[ 
 \int (5x^{2} -2)(2-x^{2} )\, \mathrm{d}x
\]}
{\(-x^5 +4x^3-4x + c,\ c\in \mathbb{R}\)}{\( \left(\frac{5}{3}x^{3}  -2x \right)\left(2x -\frac{x^3}{3}\right)
+ c,\ c\in \mathbb{R}\)}{\(-20x^3 +24x + c,\ c\in \mathbb{R}\)}{\(\frac{4x^3}{3}+ c,\ c\in \mathbb{R}\)}{}{}}
\LANGcs{
\Question{
Určete  \[ \int (5x^{2} -2)(2-x^{2} )\, \mathrm{d}x \] na \(\mathbb{R}\).}
{\(-x^5 +4x^3-4x + c,\ c\in \mathbb{R}\)}{\( \left(\frac{5}{3}x^{3}  -2x \right)\left(2x -\frac{x^3}{3}\right)
+ c,\ c\in \mathbb{R}\)}{\(-20x^3 +24x + c,\ c\in \mathbb{R}\)}{\(\frac{4x^3}{3}+ c,\ c\in \mathbb{R}\)}{}{}}
\LANGsk{
\Question{
Vypočítajte na množine \(\mathbb{R}\) nasledujúci integrál.
\[ 
 \int (5x^{2} -2)(2-x^{2} )\, \mathrm{d}x
\]}
{\(-x^5 +4x^3-4x + c,\ c\in \mathbb{R}\)}{\( \left(\frac{5}{3}x^{3}  -2x \right)\left(2x -\frac{x^3}{3}\right)
+ c,\ c\in \mathbb{R}\)}{\(-20x^3 +24x + c,\ c\in \mathbb{R}\)}{\(\frac{4x^3}{3}+ c,\ c\in \mathbb{R}\)}{}{}}
\LANGes{
\Question{
Evalúa la siguiente integral en \(\mathbb{R}\). 
\[ 
 \int (5x^{2} -2)(2-x^{2} )\, \mathrm{d}x
\]}
{\(-x^5 +4x^3-4x + c,\ c\in \mathbb{R}\)}{\( \left(\frac{5}{3}x^{3}  -2x \right)\left(2x -\frac{x^3}{3}\right)
+ c,\ c\in \mathbb{R}\)}{\(-20x^3 +24x + c,\ c\in \mathbb{R}\)}{\(\frac{4x^3}{3}+ c,\ c\in \mathbb{R}\)}{}{}}
\End

\Data\ID{2010001503}
\Updated{2022-08-25 10:08:23}
\Level{B}
\SubArea{Primitive function}
\LANGen{
\Question{
Evaluate the following integral on the interval \((0;+\infty)\). 
\[ 
 \int \frac{2x^{4} -x^2}
 {x^{3}} \, \text{d}x
\]}
{\(x^2 -\ln |x| + c,\ c\in \mathbb{R}\)}{\(\frac{\frac{2}{5}x^5-\frac{x^3}{3}}{\frac{x^4}{4}}
+ c,\ c\in \mathbb{R}\)}{\(2-\ln |x| + c,\ c\in \mathbb{R}\)}{\(4x^2 -\ln |x| + c,\ c\in \mathbb{R}\)}{}{}}
\LANGpl{
\Question{
Oblicz następującą całkę na \((0;+\infty)\). 
\[ 
 \int \frac{2x^{4} -x^2}
 {x^{3}} \, \text{d}x
\]}
{\(x^2 -\ln |x| + c,\ c\in \mathbb{R}\)}{\(\frac{\frac{2}{5}x^5-\frac{x^3}{3}}{\frac{x^4}{4}}
+ c,\ c\in \mathbb{R}\)}{\(2-\ln |x| + c,\ c\in \mathbb{R}\)}{\(4x^2 -\ln |x| + c,\ c\in \mathbb{R}\)}{}{}}
\LANGcs{
\Question{
Vypočtěte  
\[ 
 \int \frac{2x^{4} -x^2}
 {x^{3}} \, \text{d}x
\] na intervalu \((0;+\infty)\).}
{\(x^2 -\ln |x| + c,\ c\in \mathbb{R}\)}{\(\frac{\frac{2}{5}x^5-\frac{x^3}{3}}{\frac{x^4}{4}}
+ c,\ c\in \mathbb{R}\)}{\(2-\ln |x| + c,\ c\in \mathbb{R}\)}{\(4x^2 -\ln |x| + c,\ c\in \mathbb{R}\)}{}{}}
\LANGsk{
\Question{
Vypočítajte na intervale \((0;+\infty)\) nasledujúci integrál. 
\[ 
 \int \frac{2x^{4} -x^2}
 {x^{3}} \, \text{d}x
\]}
{\(x^2 -\ln |x| + c,\ c\in \mathbb{R}\)}{\(\frac{\frac{2}{5}x^5-\frac{x^3}{3}}{\frac{x^4}{4}}
+ c,\ c\in \mathbb{R}\)}{\(2-\ln |x| + c,\ c\in \mathbb{R}\)}{\(4x^2 -\ln |x| + c,\ c\in \mathbb{R}\)}{}{}}
\LANGes{
\Question{
Evalúa la siguiente integral en el intervalo \((0;+\infty)\). 
\[ 
 \int \frac{2x^{4} -x^2}
 {x^{3}} \, \text{d}x
\]}
{\(x^2 -\ln |x| + c,\ c\in \mathbb{R}\)}{\(\frac{\frac{2}{5}x^5-\frac{x^3}{3}}{\frac{x^4}{4}}
+ c,\ c\in \mathbb{R}\)}{\(2-\ln |x| + c,\ c\in \mathbb{R}\)}{\(4x^2 -\ln |x| + c,\ c\in \mathbb{R}\)}{}{}}
\End

\Data\ID{2010001504}
\Updated{2022-08-25 10:08:23}
\Level{B}
\SubArea{Primitive function}
\LANGen{
\Question{
Evaluate the following integral on \(\mathbb{R}\). 
\[ 
 \int x\cos x\, \mathrm{d}x
\]}
{\( x\sin x +\cos x + c,\ c\in \mathbb{R}\)}{\( x\sin x -\cos x + c,\ c\in \mathbb{R}\)}{\( -x\sin x +\cos x + c,\ c\in \mathbb{R}\)}{\(- x\sin x -\cos x + c,\ c\in \mathbb{R}\)}{}{}}
\LANGpl{
\Question{
Oblicz następującą całkę na \(\mathbb{R}\). 
\[ 
 \int x\cos x\, \mathrm{d}x
\]}
{\( x\sin x +\cos x + c,\ c\in \mathbb{R}\)}{\( x\sin x -\cos x + c,\ c\in \mathbb{R}\)}{\( -x\sin x +\cos x + c,\ c\in \mathbb{R}\)}{\(- x\sin x -\cos x + c,\ c\in \mathbb{R}\)}{}{}}
\LANGcs{
\Question{
Vypočtěte  
\[ 
 \int x\cos x\, \mathrm{d}x
\] na \(\mathbb{R}\).}
{\( x\sin x +\cos x + c,\ c\in \mathbb{R}\)}{\( x\sin x -\cos x + c,\ c\in \mathbb{R}\)}{\( -x\sin x +\cos x + c,\ c\in \mathbb{R}\)}{\(- x\sin x -\cos x + c,\ c\in \mathbb{R}\)}{}{}}
\LANGsk{
\Question{
Vypočítajte na množine \(\mathbb{R}\) nasledujúci integrál.
\[ 
 \int x\cos x\, \mathrm{d}x
\]}
{\( x\sin x +\cos x + c,\ c\in \mathbb{R}\)}{\( x\sin x -\cos x + c,\ c\in \mathbb{R}\)}{\( -x\sin x +\cos x + c,\ c\in \mathbb{R}\)}{\(- x\sin x -\cos x + c,\ c\in \mathbb{R}\)}{}{}}
\LANGes{
\Question{
Evalúa la siguiente integral en \(\mathbb{R}\). 
\[ 
 \int x\cos x\, \mathrm{d}x
\]}
{\( x\sin x +\cos x + c,\ c\in \mathbb{R}\)}{\( x\sin x -\cos x + c,\ c\in \mathbb{R}\)}{\( -x\sin x +\cos x + c,\ c\in \mathbb{R}\)}{\(- x\sin x -\cos x + c,\ c\in \mathbb{R}\)}{}{}}
\End

\Data\ID{2010001505}
\Updated{2022-08-25 10:08:23}
\Level{B}
\SubArea{Primitive function}
\LANGen{
\Question{
Evaluate the following integral on \(\mathbb{R}\). 
\[ 
 \int x2^{x}\, \mathrm{d}x
\]}
{\(\frac{x2^x}{\ln 2} - \frac{2^x}{\ln^2 2} + c,\ c\in \mathbb{R}\)}{\(x2^x -2^x + c,\ c\in \mathbb{R}\)}{\(\frac{x^22^x}{2\ln 2}+ c,\ c\in \mathbb{R}\)}{\( \frac{2^x(x-1)}{\ln 2} + c,\ c\in \mathbb{R}\)}{}{}}
\LANGpl{
\Question{
Oblicz następującą całkę na \(\mathbb{R}\). 
\[ 
 \int x2^{x}\, \mathrm{d}x
\]}
{\(\frac{x2^x}{\ln 2} - \frac{2^x}{\ln^2 2} + c,\ c\in \mathbb{R}\)}{\(x2^x -2^x + c,\ c\in \mathbb{R}\)}{\(\frac{x^22^x}{2\ln 2}+ c,\ c\in \mathbb{R}\)}{\( \frac{2^x(x-1)}{\ln 2} + c,\ c\in \mathbb{R}\)}{}{}}
\LANGcs{
\Question{
Vypočtěte 
\[ 
 \int x2^{x}\, \mathrm{d}x
\]  na \(\mathbb{R}\).}
{\(\frac{x2^x}{\ln 2} - \frac{2^x}{\ln^2 2} + c,\ c\in \mathbb{R}\)}{\(x2^x -2^x + c,\ c\in \mathbb{R}\)}{\(\frac{x^22^x}{2\ln 2}+ c,\ c\in \mathbb{R}\)}{\( \frac{2^x(x-1)}{\ln 2} + c,\ c\in \mathbb{R}\)}{}{}}
\LANGsk{
\Question{
Vypočítajte na množine \(\mathbb{R}\) nasledujúci integrál.
\[ 
 \int x2^{x}\, \mathrm{d}x
\]}
{\(\frac{x2^x}{\ln 2} - \frac{2^x}{\ln^2 2} + c,\ c\in \mathbb{R}\)}{\(x2^x -2^x + c,\ c\in \mathbb{R}\)}{\(\frac{x^22^x}{2\ln 2}+ c,\ c\in \mathbb{R}\)}{\( \frac{2^x(x-1)}{\ln 2} + c,\ c\in \mathbb{R}\)}{}{}}
\LANGes{
\Question{
Evalúa la siguiente integral en \(\mathbb{R}\). 
\[ 
 \int x2^{x}\, \mathrm{d}x
\]}
{\(\frac{x2^x}{\ln 2} - \frac{2^x}{\ln^2 2} + c,\ c\in \mathbb{R}\)}{\(x2^x -2^x + c,\ c\in \mathbb{R}\)}{\(\frac{x^22^x}{2\ln 2}+ c,\ c\in \mathbb{R}\)}{\( \frac{2^x(x-1)}{\ln 2} + c,\ c\in \mathbb{R}\)}{}{}}
\End

\Data\ID{2010001506}
\Updated{2022-08-25 10:08:23}
\Level{B}
\SubArea{Primitive function}
\LANGen{
\Question{
Evaluate the following integral on \(\mathbb{R}\). 
\[ 
 \int (x^{4} - 5)^{2}\, \mathrm{d}x
\]}
{\(\frac{x^{9}} 
 {9} - 2x^{5} + 25x + c,\ c\in \mathbb{R}\)}{\(\frac{(x^{4}-5)^{3}} 
 {3} + c,\ c\in \mathbb{R}\)}{\(8x^{9} - 40x^{5} + 25x + c,\ c\in \mathbb{R}\)}{}{}{}}
\LANGpl{
\Question{
Oblicz następującą całkę na \(\mathbb{R}\). 
\[ 
 \int (x^{4} - 5)^{2}\, \mathrm{d}x
\]}
{\(\frac{x^{9}} 
 {9} - 2x^{5} + 25x + c,\ c\in \mathbb{R}\)}{\(\frac{(x^{4}-5)^{3}} 
 {3} + c,\ c\in \mathbb{R}\)}{\(8x^{9} - 40x^{5} + 25x + c,\ c\in \mathbb{R}\)}{}{}{}}
\LANGcs{
\Question{
Vypočtěte
\[ 
 \int (x^{4} - 5)^{2}\, \mathrm{d}x
\]
na  \(\mathbb{R}\).}
{\(\frac{x^{9}} 
 {9} - 2x^{5} + 25x + c,\ c\in \mathbb{R}\)}{\(\frac{(x^{4}-5)^{3}} 
 {3} + c,\ c\in \mathbb{R}\)}{\(8x^{9} - 40x^{5} + 25x + c,\ c\in \mathbb{R}\)}{}{}{}}
\LANGsk{
\Question{
Vypočítajte na množine \(\mathbb{R}\) nasledujúci integrál.
\[ 
 \int (x^{4} - 5)^{2}\, \mathrm{d}x
\]}
{\(\frac{x^{9}} 
 {9} - 2x^{5} + 25x + c,\ c\in \mathbb{R}\)}{\(\frac{(x^{4}-5)^{3}} 
 {3} + c,\ c\in \mathbb{R}\)}{\(8x^{9} - 40x^{5} + 25x + c,\ c\in \mathbb{R}\)}{}{}{}}
\LANGes{
\Question{
Evalúa la siguiente integral en \(\mathbb{R}\). 
\[ 
 \int (x^{4} - 5)^{2}\, \mathrm{d}x
\]}
{\(\frac{x^{9}} 
 {9} - 2x^{5} + 25x + c,\ c\in \mathbb{R}\)}{\(\frac{(x^{4}-5)^{3}} 
 {3} + c,\ c\in \mathbb{R}\)}{\(8x^{9} - 40x^{5} + 25x + c,\ c\in \mathbb{R}\)}{}{}{}}
\End

\Data\ID{2010001507}
\Updated{2022-08-25 10:08:23}
\Level{B}
\SubArea{Primitive function}
\LANGen{
\Question{
Evaluate the following integral on the interval \((0;+\infty)\). 
\[ 
 \int \frac{19\root{3}\of{x^{4}} - 3}
 {\root{4}\of{x^{3}}} \, \mathrm{d}x
\]}
{\(12(x\root{12}\of{x^{7}} -\root{4}\of{x}) + c,\ c\in \mathbb{R}\)}{\(\frac{\frac{57} 
 {7} \root{3}\of{x^{7}}-3x}
 {\frac{4} 
{7} \root{4}\of{x^{7}}} + c,\ c\in \mathbb{R}\)}{\(\frac{361}
 {12} \root{12}\of{x^{19}} -\frac{3} 
{4}\root{4}\of{x} + c,\ c\in \mathbb{R}\)}{}{}{}}
\LANGpl{
\Question{
Oblicz następującą całkę na przedziale \((0;+\infty)\). 
\[ 
 \int \frac{19\root{3}\of{x^{4}} - 3}
 {\root{4}\of{x^{3}}} \, \mathrm{d}x
\]}
{\(12(x\root{12}\of{x^{7}} -\root{4}\of{x}) + c,\ c\in \mathbb{R}\)}{\(\frac{\frac{57} 
 {7} \root{3}\of{x^{7}}-3x}
 {\frac{4} 
{7} \root{4}\of{x^{7}}} + c,\ c\in \mathbb{R}\)}{\(\frac{361}
 {12} \root{12}\of{x^{19}} -\frac{3} 
{4}\root{4}\of{x} + c,\ c\in \mathbb{R}\)}{}{}{}}
\LANGcs{
\Question{
Vypočtěte  
\[ 
 \int \frac{19\root{3}\of{x^{4}} - 3}
 {\root{4}\of{x^{3}}} \, \mathrm{d}x
\]
na intervalu \((0;+\infty)\).}
{\(12(x\root{12}\of{x^{7}} -\root{4}\of{x}) + c,\ c\in \mathbb{R}\)}{\(\frac{\frac{57} 
 {7} \root{3}\of{x^{7}}-3x}
 {\frac{4} 
{7} \root{4}\of{x^{7}}} + c,\ c\in \mathbb{R}\)}{\(\frac{361}
 {12} \root{12}\of{x^{19}} -\frac{3} 
{4}\root{4}\of{x} + c,\ c\in \mathbb{R}\)}{}{}{}}
\LANGsk{
\Question{
Vypočítajte na intervale \((0;+\infty)\) nasledujúci integrál. 
\[ 
 \int \frac{19\root{3}\of{x^{4}} - 3}
 {\root{4}\of{x^{3}}} \, \mathrm{d}x
\]}
{\(12(x\root{12}\of{x^{7}} -\root{4}\of{x}) + c,\ c\in \mathbb{R}\)}{\(\frac{\frac{57} 
 {7} \root{3}\of{x^{7}}-3x}
 {\frac{4} 
{7} \root{4}\of{x^{7}}} + c,\ c\in \mathbb{R}\)}{\(\frac{361}
 {12} \root{12}\of{x^{19}} -\frac{3} 
{4}\root{4}\of{x} + c,\ c\in \mathbb{R}\)}{}{}{}}
\LANGes{
\Question{
Evalúa la siguiente integral en el intervalo \((0;+\infty)\). 
\[ 
 \int \frac{19\root{3}\of{x^{4}} - 3}
 {\root{4}\of{x^{3}}} \, \mathrm{d}x
\]}
{\(12(x\root{12}\of{x^{7}} -\root{4}\of{x}) + c,\ c\in \mathbb{R}\)}{\(\frac{\frac{57} 
 {7} \root{3}\of{x^{7}}-3x}
 {\frac{4} 
{7} \root{4}\of{x^{7}}} + c,\ c\in \mathbb{R}\)}{\(\frac{361}
 {12} \root{12}\of{x^{19}} -\frac{3} 
{4}\root{4}\of{x} + c,\ c\in \mathbb{R}\)}{}{}{}}
\End

\Data\ID{2010008101}
\Updated{2022-08-25 10:08:51}
\Level{B}
\SubArea{Primitive function}
\LANGen{
\Question{
Evaluate the following integral on the interval \( (0;+\infty)\).
\[
\int \left( \frac{1}{\sqrt{x}}+\frac{2}{x} + \frac{3}{x^2} \right) \mathrm{d}x
\]}
{\( 2\sqrt{x} +2 \ln x -\frac{3}{x} +c;~c \in \mathbb{R}\)}{\( \frac{3}{2\sqrt{x^3}}+\frac{4}{x^2} + \frac{9}{x^3}+c;~c \in \mathbb{R}\)}{\( \frac{2}{3\sqrt{x^3}}+\frac{1}{x^2} + \frac{1}{x^3}+c;~c \in \mathbb{R}\)}{\( \frac{\sqrt{x}}{2} +2\ln x - \frac{3}{x} +c;~c \in \mathbb{R}\)}{}{}}
\LANGpl{
\Question{
Oblicz następującą całkę na przedziale \( (0;+\infty)\).
\[
\int \left( \frac{1}{\sqrt{x}}+\frac{2}{x} + \frac{3}{x^2} \right) \mathrm{d}x
\]}
{\( 2\sqrt{x} +2 \ln x -\frac{3}{x} +c;~c \in \mathbb{R}\)}{\( \frac{3}{2\sqrt{x^3}}+\frac{4}{x^2} + \frac{9}{x^3}+c;~c \in \mathbb{R}\)}{\( \frac{2}{3\sqrt{x^3}}+\frac{1}{x^2} + \frac{1}{x^3}+c;~c \in \mathbb{R}\)}{\( \frac{\sqrt{x}}{2} +2\ln x - \frac{3}{x} +c;~c \in \mathbb{R}\)}{}{}}
\LANGcs{
\Question{
Vypočtěte 
\[
\int \left( \frac{1}{\sqrt{x}}+\frac{2}{x} + \frac{3}{x^2} \right) \mathrm{d}x
\] na intervalu \( (0;+\infty)\).}
{\( 2\sqrt{x} +2 \ln x -\frac{3}{x} +c;~c \in \mathbb{R}\)}{\( \frac{3}{2\sqrt{x^3}}+\frac{4}{x^2} + \frac{9}{x^3}+c;~c \in \mathbb{R}\)}{\( \frac{2}{3\sqrt{x^3}}+\frac{1}{x^2} + \frac{1}{x^3}+c;~c \in \mathbb{R}\)}{\( \frac{\sqrt{x}}{2} +2\ln x - \frac{3}{x} +c;~c \in \mathbb{R}\)}{}{}}
\LANGsk{
\Question{
Vypočítajte nasledujúci integrál na intervale \( (0;+\infty)\).
\[
\int \left( \frac{1}{\sqrt{x}}+\frac{2}{x} + \frac{3}{x^2} \right) \mathrm{d}x
\]}
{\( 2\sqrt{x} +2 \ln x -\frac{3}{x} +c;~c \in \mathbb{R}\)}{\( \frac{3}{2\sqrt{x^3}}+\frac{4}{x^2} + \frac{9}{x^3}+c;~c \in \mathbb{R}\)}{\( \frac{2}{3\sqrt{x^3}}+\frac{1}{x^2} + \frac{1}{x^3}+c;~c \in \mathbb{R}\)}{\( \frac{\sqrt{x}}{2} +2\ln x - \frac{3}{x} +c;~c \in \mathbb{R}\)}{}{}}
\LANGes{
\Question{
Evalúa la siguiente integral en el intervalo \( (0;+\infty)\).
\[
\int \left( \frac{1}{\sqrt{x}}+\frac{2}{x} + \frac{3}{x^2} \right) \mathrm{d}x
\]}
{\( 2\sqrt{x} +2 \ln x -\frac{3}{x} +c;~c \in \mathbb{R}\)}{\( \frac{3}{2\sqrt{x^3}}+\frac{4}{x^2} + \frac{9}{x^3}+c;~c \in \mathbb{R}\)}{\( \frac{2}{3\sqrt{x^3}}+\frac{1}{x^2} + \frac{1}{x^3}+c;~c \in \mathbb{R}\)}{\( \frac{\sqrt{x}}{2} +2\ln x - \frac{3}{x} +c;~c \in \mathbb{R}\)}{}{}}
\End

\Data\ID{2010008102}
\Updated{2022-08-25 10:08:51}
\Level{B}
\SubArea{Primitive function}
\LANGen{
\Question{
Evaluate the following integral on the interval \( (0;+\infty)\).
\[
\int \left( \frac{1}{x} - \frac{2}{x^2} +\frac{3}{\sqrt{x}}\right) \mathrm{d}x
\]}
{\( \ln x +\frac{2}{x}+6\sqrt{x}+c;~c \in \mathbb{R}\)}{\( \frac{2}{x^2}-\frac{6}{x^3}+\frac{9}{2\sqrt{x^3}}+c;~c \in \mathbb{R}\)}{\( \frac{1}{2x^2}-\frac{2}{3x^3}+\frac{2}{\sqrt{x^3}}+c;~c \in \mathbb{R}\)}{\( \ln x +\frac{2}{x}+\frac{3\sqrt{x}}{2}+c;~c \in \mathbb{R}\)}{}{}}
\LANGpl{
\Question{
Oblicz następującą całkę na przedziale \( (0;+\infty)\).
\[
\int \left( \frac{1}{x} - \frac{2}{x^2} +\frac{3}{\sqrt{x}}\right) \mathrm{d}x
\]}
{\( \ln x +\frac{2}{x}+6\sqrt{x}+c;~c \in \mathbb{R}\)}{\( \frac{2}{x^2}-\frac{6}{x^3}+\frac{9}{2\sqrt{x^3}}+c;~c \in \mathbb{R}\)}{\( \frac{1}{2x^2}-\frac{2}{3x^3}+\frac{2}{\sqrt{x^3}}+c;~c \in \mathbb{R}\)}{\( \ln x +\frac{2}{x}+\frac{3\sqrt{x}}{2}+c;~c \in \mathbb{R}\)}{}{}}
\LANGcs{
\Question{
Vypočtěte 
\[
\int \left( \frac{1}{x} - \frac{2}{x^2} +\frac{3}{\sqrt{x}}\right) \mathrm{d}x
\]
na intervalu \( (0;+\infty)\).}
{\( \ln x +\frac{2}{x}+6\sqrt{x}+c;~c \in \mathbb{R}\)}{\( \frac{2}{x^2}-\frac{6}{x^3}+\frac{9}{2\sqrt{x^3}}+c;~c \in \mathbb{R}\)}{\( \frac{1}{2x^2}-\frac{2}{3x^3}+\frac{2}{\sqrt{x^3}}+c;~c \in \mathbb{R}\)}{\( \ln x +\frac{2}{x}+\frac{3\sqrt{x}}{2}+c;~c \in \mathbb{R}\)}{}{}}
\LANGsk{
\Question{
Vypočítajte nasledujúci integrál na intervale \( (0;+\infty)\).
\[
\int \left( \frac{1}{x} - \frac{2}{x^2} +\frac{3}{\sqrt{x}}\right) \mathrm{d}x
\]}
{\( \ln x +\frac{2}{x}+6\sqrt{x}+c;~c \in \mathbb{R}\)}{\( \frac{2}{x^2}-\frac{6}{x^3}+\frac{9}{2\sqrt{x^3}}+c;~c \in \mathbb{R}\)}{\( \frac{1}{2x^2}-\frac{2}{3x^3}+\frac{2}{\sqrt{x^3}}+c;~c \in \mathbb{R}\)}{\( \ln x +\frac{2}{x}+\frac{3\sqrt{x}}{2}+c;~c \in \mathbb{R}\)}{}{}}
\LANGes{
\Question{
Evalúa la siguiente integral en el intervalo \( (0;+\infty)\).
\[
\int \left( \frac{1}{x} - \frac{2}{x^2} +\frac{3}{\sqrt{x}}\right) \mathrm{d}x
\]}
{\( \ln x +\frac{2}{x}+6\sqrt{x}+c;~c \in \mathbb{R}\)}{\( \frac{2}{x^2}-\frac{6}{x^3}+\frac{9}{2\sqrt{x^3}}+c;~c \in \mathbb{R}\)}{\( \frac{1}{2x^2}-\frac{2}{3x^3}+\frac{2}{\sqrt{x^3}}+c;~c \in \mathbb{R}\)}{\( \ln x +\frac{2}{x}+\frac{3\sqrt{x}}{2}+c;~c \in \mathbb{R}\)}{}{}}
\End

\Data\ID{2010008103}
\Updated{2022-08-25 10:08:51}
\Level{B}
\SubArea{Primitive function}
\LANGen{
\Question{
Evaluate the following integral on \( \mathbb{R}\).
\[
\int \left(  x^2+2\sin^2 x + 3 \mathrm{e}^{2x}\right) \mathrm{d}x
\]}
{\( \frac{x^3}{3}+x-\sin x \cos x + \frac32 \mathrm{e}^{2x}+c;~c \in \mathbb{R}\)}{\( 3x^3+x-\sin x \cos x + 6\mathrm{e}^{2x}+c;~c \in \mathbb{R}\)}{\( \frac{x^3}{3}+\frac23\sin^3 x +\mathrm{e}^{3x}+c;~c \in \mathbb{R}\)}{\( \frac{x^3}{3}-2\cos^2 x + 3\mathrm{e}^{2x}+c;~c \in \mathbb{R}\)}{}{}}
\LANGpl{
\Question{
Oblicz następującą całkę na \( \mathbb{R}\).
\[
\int \left(  x^2+2\sin^2 x + 3 \mathrm{e}^{2x}\right) \mathrm{d}x
\]}
{\( \frac{x^3}{3}+x-\sin x \cos x + \frac32 \mathrm{e}^{2x}+c;~c \in \mathbb{R}\)}{\( 3x^3+x-\sin x \cos x + 6\mathrm{e}^{2x}+c;~c \in \mathbb{R}\)}{\( \frac{x^3}{3}+\frac23\sin^3 x +\mathrm{e}^{3x}+c;~c \in \mathbb{R}\)}{\( \frac{x^3}{3}-2\cos^2 x + 3\mathrm{e}^{2x}+c;~c \in \mathbb{R}\)}{}{}}
\LANGcs{
\Question{
Vypočtěte 
\[
\int \left(  x^2+2\sin^2 x + 3 \mathrm{e}^{2x}\right) \mathrm{d}x
\]
na \( \mathbb{R}\).}
{\( \frac{x^3}{3}+x-\sin x \cos x + \frac32 \mathrm{e}^{2x}+c;~c \in \mathbb{R}\)}{\( 3x^3+x-\sin x \cos x + 6\mathrm{e}^{2x}+c;~c \in \mathbb{R}\)}{\( \frac{x^3}{3}+\frac23\sin^3 x +\mathrm{e}^{3x}+c;~c \in \mathbb{R}\)}{\( \frac{x^3}{3}-2\cos^2 x + 3\mathrm{e}^{2x}+c;~c \in \mathbb{R}\)}{}{}}
\LANGsk{
\Question{
Vypočítajte nasledujúci integrál na \( \mathbb{R}\).
\[
\int \left(  x^2+2\sin^2 x + 3 \mathrm{e}^{2x}\right) \mathrm{d}x
\]}
{\( \frac{x^3}{3}+x-\sin x \cos x + \frac32 \mathrm{e}^{2x}+c;~c \in \mathbb{R}\)}{\( 3x^3+x-\sin x \cos x + 6\mathrm{e}^{2x}+c;~c \in \mathbb{R}\)}{\( \frac{x^3}{3}+\frac23\sin^3 x +\mathrm{e}^{3x}+c;~c \in \mathbb{R}\)}{\( \frac{x^3}{3}-2\cos^2 x + 3\mathrm{e}^{2x}+c;~c \in \mathbb{R}\)}{}{}}
\LANGes{
\Question{
Evalúa la siguiente integral en \( \mathbb{R}\).
\[
\int \left(  x^2+2\sin^2 x + 3 \mathrm{e}^{2x}\right) \mathrm{d}x
\]}
{\( \frac{x^3}{3}+x-\sin x \cos x + \frac32 \mathrm{e}^{2x}+c;~c \in \mathbb{R}\)}{\( 3x^3+x-\sin x \cos x + 6\mathrm{e}^{2x}+c;~c \in \mathbb{R}\)}{\( \frac{x^3}{3}+\frac23\sin^3 x +\mathrm{e}^{3x}+c;~c \in \mathbb{R}\)}{\( \frac{x^3}{3}-2\cos^2 x + 3\mathrm{e}^{2x}+c;~c \in \mathbb{R}\)}{}{}}
\End

\Data\ID{2010008104}
\Updated{2022-08-25 10:08:51}
\Level{B}
\SubArea{Primitive function}
\LANGen{
\Question{
Evaluate the following integral on \( \mathbb{R}\).
\[
\int \left(  3x^3 + \mathrm{e}^{2x}-\cos^2 x\right) \mathrm{d}x
\]}
{\( \frac{3x^4}{4}+\frac{\mathrm{e}^{2x}}{2}-\frac{x}2-\frac12\sin x \cos x+c;~c \in \mathbb{R}\)}{\( 12x^4+2\mathrm{e}^{2x}-\frac{x}2-\frac12\sin x \cos x+c;~c \in \mathbb{R}\)}{\( \frac{3x^4}{4}+\frac{\mathrm{e}^{3x}}{3}-\frac13\cos^3 x +c;~c \in \mathbb{R}\)}{\( \frac{3x^4}{4}+\mathrm{e}^{2x}-\sin^2 x +c;~c \in \mathbb{R}\)}{}{}}
\LANGpl{
\Question{
Oblicz następującą całkę na \( \mathbb{R}\).
\[
\int \left(  3x^3 + \mathrm{e}^{2x}-\cos^2 x\right) \mathrm{d}x
\]}
{\( \frac{3x^4}{4}+\frac{\mathrm{e}^{2x}}{2}-\frac{x}2-\frac12\sin x \cos x+c;~c \in \mathbb{R}\)}{\( 12x^4+2\mathrm{e}^{2x}-\frac{x}2-\frac12\sin x \cos x+c;~c \in \mathbb{R}\)}{\( \frac{3x^4}{4}+\frac{\mathrm{e}^{3x}}{3}-\frac13\cos^3 x +c;~c \in \mathbb{R}\)}{\( \frac{3x^4}{4}+\mathrm{e}^{2x}-\sin^2 x +c;~c \in \mathbb{R}\)}{}{}}
\LANGcs{
\Question{
Vypočtěte 
\[
\int \left(  3x^3 + \mathrm{e}^{2x}-\cos^2 x\right) \mathrm{d}x
\]
na \( \mathbb{R}\).}
{\( \frac{3x^4}{4}+\frac{\mathrm{e}^{2x}}{2}-\frac{x}2-\frac12\sin x \cos x+c;~c \in \mathbb{R}\)}{\( 12x^4+2\mathrm{e}^{2x}-\frac{x}2-\frac12\sin x \cos x+c;~c \in \mathbb{R}\)}{\( \frac{3x^4}{4}+\frac{\mathrm{e}^{3x}}{3}-\frac13\cos^3 x +c;~c \in \mathbb{R}\)}{\( \frac{3x^4}{4}+\mathrm{e}^{2x}-\sin^2 x +c;~c \in \mathbb{R}\)}{}{}}
\LANGsk{
\Question{
Vypočítajte nasledujúci integrál na \( \mathbb{R}\).
\[
\int \left(  3x^3 + \mathrm{e}^{2x}-\cos^2 x\right) \mathrm{d}x
\]}
{\( \frac{3x^4}{4}+\frac{\mathrm{e}^{2x}}{2}-\frac{x}2-\frac12\sin x \cos x+c;~c \in \mathbb{R}\)}{\( 12x^4+2\mathrm{e}^{2x}-\frac{x}2-\frac12\sin x \cos x+c;~c \in \mathbb{R}\)}{\( \frac{3x^4}{4}+\frac{\mathrm{e}^{3x}}{3}-\frac13\cos^3 x +c;~c \in \mathbb{R}\)}{\( \frac{3x^4}{4}+\mathrm{e}^{2x}-\sin^2 x +c;~c \in \mathbb{R}\)}{}{}}
\LANGes{
\Question{
Evalúa la siguiente integral en \( \mathbb{R}\).
\[
\int \left(  3x^3 + \mathrm{e}^{2x}-\cos^2 x\right) \mathrm{d}x
\]}
{\( \frac{3x^4}{4}+\frac{\mathrm{e}^{2x}}{2}-\frac{x}2-\frac12\sin x \cos x+c;~c \in \mathbb{R}\)}{\( 12x^4+2\mathrm{e}^{2x}-\frac{x}2-\frac12\sin x \cos x+c;~c \in \mathbb{R}\)}{\( \frac{3x^4}{4}+\frac{\mathrm{e}^{3x}}{3}-\frac13\cos^3 x +c;~c \in \mathbb{R}\)}{\( \frac{3x^4}{4}+\mathrm{e}^{2x}-\sin^2 x +c;~c \in \mathbb{R}\)}{}{}}
\End

\Data\ID{2010008105}
\Updated{2022-08-25 10:08:51}
\Level{B}
\SubArea{Primitive function}
\LANGen{
\Question{
Evaluate the following integral on the interval \((0;+\infty)\).
\[
\int \left(   \frac{1}{x} + \frac{1}{x+1}+\frac{x}{x^2+2}\right) \mathrm{d}x
\]}
{\( \ln x + \ln(x+1) + \frac12 \ln(x^2+2)+c;~c \in \mathbb{R}\)}{\( \ln x + \ln(x+1) + \frac12 x^2\ln(x^2+2)+c;~c \in \mathbb{R}\)}{\( \frac{2}{x^2} + \frac{1}{\frac12 x^2+x} + \frac{\frac12x^2}{\frac13x^3+2x}+c;~c \in \mathbb{R}\)}{\( \ln x + \ln(x+1) + \frac{3x}{2(x^2+6)}+c;~c \in \mathbb{R}\)}{}{}}
\LANGpl{
\Question{
Oblicz następującą całkę na przedziale \((0;+\infty)\).
\[
\int \left(   \frac{1}{x} + \frac{1}{x+1}+\frac{x}{x^2+2}\right) \mathrm{d}x
\]}
{\( \ln x + \ln(x+1) + \frac12 \ln(x^2+2)+c;~c \in \mathbb{R}\)}{\( \ln x + \ln(x+1) + \frac12 x^2\ln(x^2+2)+c;~c \in \mathbb{R}\)}{\( \frac{2}{x^2} + \frac{1}{\frac12 x^2+x} + \frac{\frac12x^2}{\frac13x^3+2x}+c;~c \in \mathbb{R}\)}{\( \ln x + \ln(x+1) + \frac{3x}{2(x^2+6)}+c;~c \in \mathbb{R}\)}{}{}}
\LANGcs{
\Question{
Vypočtěte 
\[
\int \left(   \frac{1}{x} + \frac{1}{x+1}+\frac{x}{x^2+2}\right) \mathrm{d}x
\]
na intervalu \((0;+\infty)\).}
{\( \ln x + \ln(x+1) + \frac12 \ln(x^2+2)+c;~c \in \mathbb{R}\)}{\( \ln x + \ln(x+1) + \frac12 x^2\ln(x^2+2)+c;~c \in \mathbb{R}\)}{\( \frac{2}{x^2} + \frac{1}{\frac12 x^2+x} + \frac{\frac12x^2}{\frac13x^3+2x}+c;~c \in \mathbb{R}\)}{\( \ln x + \ln(x+1) + \frac{3x}{2(x^2+6)}+c;~c \in \mathbb{R}\)}{}{}}
\LANGsk{
\Question{
Vypočítajte nasledujúci integrál na intervale \((0;+\infty)\).
\[
\int \left(   \frac{1}{x} + \frac{1}{x+1}+\frac{x}{x^2+2}\right) \mathrm{d}x
\]}
{\( \ln x + \ln(x+1) + \frac12 \ln(x^2+2)+c;~c \in \mathbb{R}\)}{\( \ln x + \ln(x+1) + \frac12 x^2\ln(x^2+2)+c;~c \in \mathbb{R}\)}{\( \frac{2}{x^2} + \frac{1}{\frac12 x^2+x} + \frac{\frac12x^2}{\frac13x^3+2x}+c;~c \in \mathbb{R}\)}{\( \ln x + \ln(x+1) + \frac{3x}{2(x^2+6)}+c;~c \in \mathbb{R}\)}{}{}}
\LANGes{
\Question{
Evalúa la siguiente integral en el intervalo \((0;+\infty)\).
\[
\int \left(   \frac{1}{x} + \frac{1}{x+1}+\frac{x}{x^2+2}\right) \mathrm{d}x
\]}
{\( \ln x + \ln(x+1) + \frac12 \ln(x^2+2)+c;~c \in \mathbb{R}\)}{\( \ln x + \ln(x+1) + \frac12 x^2\ln(x^2+2)+c;~c \in \mathbb{R}\)}{\( \frac{2}{x^2} + \frac{1}{\frac12 x^2+x} + \frac{\frac12x^2}{\frac13x^3+2x}+c;~c \in \mathbb{R}\)}{\( \ln x + \ln(x+1) + \frac{3x}{2(x^2+6)}+c;~c \in \mathbb{R}\)}{}{}}
\End

\Data\ID{2010008106}
\Updated{2022-08-25 10:08:51}
\Level{B}
\SubArea{Primitive function}
\LANGen{
\Question{
Evaluate the following integral on the interval \((0;+\infty)\).
\[
\int \left(   \frac{2}{x} + \frac{1}{x+2}+\frac{x}{x^2+1}\right) \mathrm{d}x
\]}
{\( 2\ln x + \ln(x+2) + \frac12 \ln(x^2+1)+c;~c \in \mathbb{R}\)}{\( 2\ln x + \ln(x+2) + \frac12 x^2\ln(x^2+1)+c;~c \in \mathbb{R}\)}{\( \frac{4}{x^2} + \frac{1}{\frac12 x^2+2x} + \frac{\frac12x^2}{\frac13x^3+x}+c;~c \in \mathbb{R}\)}{\( 2\ln x + \ln(x+2) + \frac{3x}{2(x^2+3)}+c;~c \in \mathbb{R}\)}{}{}}
\LANGpl{
\Question{
Oblicz następującą całkę na przedziale  \((0;+\infty)\).
\[
\int \left(   \frac{2}{x} + \frac{1}{x+2}+\frac{x}{x^2+1}\right) \mathrm{d}x
\]}
{\( 2\ln x + \ln(x+2) + \frac12 \ln(x^2+1)+c;~c \in \mathbb{R}\)}{\( 2\ln x + \ln(x+2) + \frac12 x^2\ln(x^2+1)+c;~c \in \mathbb{R}\)}{\( \frac{4}{x^2} + \frac{1}{\frac12 x^2+2x} + \frac{\frac12x^2}{\frac13x^3+x}+c;~c \in \mathbb{R}\)}{\( 2\ln x + \ln(x+2) + \frac{3x}{2(x^2+3)}+c;~c \in \mathbb{R}\)}{}{}}
\LANGcs{
\Question{
Vypočtěte 
\[
\int \left(   \frac{2}{x} + \frac{1}{x+2}+\frac{x}{x^2+1}\right) \mathrm{d}x
\]
na intervalu \((0;+\infty)\).}
{\( 2\ln x + \ln(x+2) + \frac12 \ln(x^2+1)+c;~c \in \mathbb{R}\)}{\( 2\ln x + \ln(x+2) + \frac12 x^2\ln(x^2+1)+c;~c \in \mathbb{R}\)}{\( \frac{4}{x^2} + \frac{1}{\frac12 x^2+2x} + \frac{\frac12x^2}{\frac13x^3+x}+c;~c \in \mathbb{R}\)}{\( 2\ln x + \ln(x+2) + \frac{3x}{2(x^2+3)}+c;~c \in \mathbb{R}\)}{}{}}
\LANGsk{
\Question{
Vypočítajte nasledujúci integrál na intervale \((0;+\infty)\).
\[
\int \left(   \frac{2}{x} + \frac{1}{x+2}+\frac{x}{x^2+1}\right) \mathrm{d}x
\]}
{\( 2\ln x + \ln(x+2) + \frac12 \ln(x^2+1)+c;~c \in \mathbb{R}\)}{\( 2\ln x + \ln(x+2) + \frac12 x^2\ln(x^2+1)+c;~c \in \mathbb{R}\)}{\( \frac{4}{x^2} + \frac{1}{\frac12 x^2+2x} + \frac{\frac12x^2}{\frac13x^3+x}+c;~c \in \mathbb{R}\)}{\( 2\ln x + \ln(x+2) + \frac{3x}{2(x^2+3)}+c;~c \in \mathbb{R}\)}{}{}}
\LANGes{
\Question{
Evalúa la siguiente integral en el intervalo \((0;+\infty)\).
\[
\int \left(   \frac{2}{x} + \frac{1}{x+2}+\frac{x}{x^2+1}\right) \mathrm{d}x
\]}
{\( 2\ln x + \ln(x+2) + \frac12 \ln(x^2+1)+c;~c \in \mathbb{R}\)}{\( 2\ln x + \ln(x+2) + \frac12 x^2\ln(x^2+1)+c;~c \in \mathbb{R}\)}{\( \frac{4}{x^2} + \frac{1}{\frac12 x^2+2x} + \frac{\frac12x^2}{\frac13x^3+x}+c;~c \in \mathbb{R}\)}{\( 2\ln x + \ln(x+2) + \frac{3x}{2(x^2+3)}+c;~c \in \mathbb{R}\)}{}{}}
\End

\Data\ID{2010008107}
\Updated{2022-08-25 10:08:51}
\Level{B}
\SubArea{Primitive function}
\LANGen{
\Question{
Evaluate the following integral on the interval \((0;+\infty)\).
\[
\int \left( x\sqrt{x}+x\cos x - x\mathrm{e}^x\right) \mathrm{d}x
\]}
{\( \frac25x^2\sqrt{x}+x\sin x+\cos x-x\mathrm{e}^x+\mathrm{e}^x+c;~c \in \mathbb{R}\)}{\( \frac{x^2}2\left(\frac23x^{\frac32}+\sin x-\mathrm{e}^x\right)+c;~c \in \mathbb{R}\)}{\( \frac25x^2\sqrt{x}+x\sin x-\cos x-x\mathrm{e}^x+\mathrm{e}^x+c;~c \in \mathbb{R}\)}{\( \frac25x^2\sqrt{x}+x\sin x + \cos x-x\mathrm{e}^x-\mathrm{e}^x+c;~c \in \mathbb{R}\)}{}{}}
\LANGpl{
\Question{
Oblicz następującą całkę na przedziale  \((0;+\infty)\).
\[
\int \left( x\sqrt{x}+x\cos x - x\mathrm{e}^x\right) \mathrm{d}x
\]}
{\( \frac25x^2\sqrt{x}+x\sin x+\cos x-x\mathrm{e}^x+\mathrm{e}^x+c;~c \in \mathbb{R}\)}{\( \frac{x^2}2\left(\frac23x^{\frac32}+\sin x-\mathrm{e}^x\right)+c;~c \in \mathbb{R}\)}{\( \frac25x^2\sqrt{x}+x\sin x-\cos x-x\mathrm{e}^x+\mathrm{e}^x+c;~c \in \mathbb{R}\)}{\( \frac25x^2\sqrt{x}+x\sin x + \cos x-x\mathrm{e}^x-\mathrm{e}^x+c;~c \in \mathbb{R}\)}{}{}}
\LANGcs{
\Question{
Vypočtěte
\[
\int \left( x\sqrt{x}+x\cos x - x\mathrm{e}^x\right) \mathrm{d}x
\]
na intervalu \((0;+\infty)\).}
{\( \frac25x^2\sqrt{x}+x\sin x+\cos x-x\mathrm{e}^x+\mathrm{e}^x+c;~c \in \mathbb{R}\)}{\( \frac{x^2}2\left(\frac23x^{\frac32}+\sin x-\mathrm{e}^x\right)+c;~c \in \mathbb{R}\)}{\( \frac25x^2\sqrt{x}+x\sin x-\cos x-x\mathrm{e}^x+\mathrm{e}^x+c;~c \in \mathbb{R}\)}{\( \frac25x^2\sqrt{x}+x\sin x + \cos x-x\mathrm{e}^x-\mathrm{e}^x+c;~c \in \mathbb{R}\)}{}{}}
\LANGsk{
\Question{
Vypočítajte nasledujúci integrál na intervale \((0;+\infty)\).
\[
\int \left( x\sqrt{x}+x\cos x - x\mathrm{e}^x\right) \mathrm{d}x
\]}
{\( \frac25x^2\sqrt{x}+x\sin x+\cos x-x\mathrm{e}^x+\mathrm{e}^x+c;~c \in \mathbb{R}\)}{\( \frac{x^2}2\left(\frac23x^{\frac32}+\sin x-\mathrm{e}^x\right)+c;~c \in \mathbb{R}\)}{\( \frac25x^2\sqrt{x}+x\sin x-\cos x-x\mathrm{e}^x+\mathrm{e}^x+c;~c \in \mathbb{R}\)}{\( \frac25x^2\sqrt{x}+x\sin x + \cos x-x\mathrm{e}^x-\mathrm{e}^x+c;~c \in \mathbb{R}\)}{}{}}
\LANGes{
\Question{
Evalúa la siguiente integral en el intervalo \((0;+\infty)\).
\[
\int \left( x\sqrt{x}+x\cos x - x\mathrm{e}^x\right) \mathrm{d}x
\]}
{\( \frac25x^2\sqrt{x}+x\sin x+\cos x-x\mathrm{e}^x+\mathrm{e}^x+c;~c \in \mathbb{R}\)}{\( \frac{x^2}2\left(\frac23x^{\frac32}+\sin x-\mathrm{e}^x\right)+c;~c \in \mathbb{R}\)}{\( \frac25x^2\sqrt{x}+x\sin x-\cos x-x\mathrm{e}^x+\mathrm{e}^x+c;~c \in \mathbb{R}\)}{\( \frac25x^2\sqrt{x}+x\sin x + \cos x-x\mathrm{e}^x-\mathrm{e}^x+c;~c \in \mathbb{R}\)}{}{}}
\End

\Data\ID{2010008108}
\Updated{2022-08-25 10:08:51}
\Level{B}
\SubArea{Primitive function}
\LANGen{
\Question{
Evaluate the following integral on the interval \((0;+\infty)\).
\[
\int \left( x\sqrt[3]{x}+x\sin x + x\mathrm{e}^x\right) \mathrm{d}x
\]}
{\( \frac37x^2\sqrt[3]{x}+x\cos x+\sin x+x\mathrm{e}^x-\mathrm{e}^x+c;~c \in \mathbb{R}\)}{\( \frac{x^2}2\left(\frac34x^{\frac43}-\cos x+\mathrm{e}^x\right)+c;~c \in \mathbb{R}\)}{\( \frac37x^2\sqrt[3]{x}+x\cos x-\sin x+x\mathrm{e}^x-\mathrm{e}^x+c;~c \in \mathbb{R}\)}{\( \frac37x^2\sqrt[3]{x}-x\cos x + \sin x+x\mathrm{e}^x+\mathrm{e}^x+c;~c \in \mathbb{R}\)}{}{}}
\LANGpl{
\Question{
Oblicz następującą całkę na przedziale  \((0;+\infty)\).
\[
\int \left( x\sqrt[3]{x}+x\sin x + x\mathrm{e}^x\right) \mathrm{d}x
\]}
{\( \frac37x^2\sqrt[3]{x}+x\cos x+\sin x+x\mathrm{e}^x-\mathrm{e}^x+c;~c \in \mathbb{R}\)}{\( \frac{x^2}2\left(\frac34x^{\frac43}-\cos x+\mathrm{e}^x\right)+c;~c \in \mathbb{R}\)}{\( \frac37x^2\sqrt[3]{x}+x\cos x-\sin x+x\mathrm{e}^x-\mathrm{e}^x+c;~c \in \mathbb{R}\)}{\( \frac37x^2\sqrt[3]{x}-x\cos x + \sin x+x\mathrm{e}^x+\mathrm{e}^x+c;~c \in \mathbb{R}\)}{}{}}
\LANGcs{
\Question{
Vypočtěte
\[
\int \left( x\sqrt[3]{x}+x\sin x + x\mathrm{e}^x\right) \mathrm{d}x
\]
na intervalu \((0;+\infty)\).}
{\( \frac37x^2\sqrt[3]{x}+x\cos x+\sin x+x\mathrm{e}^x-\mathrm{e}^x+c;~c \in \mathbb{R}\)}{\( \frac{x^2}2\left(\frac34x^{\frac43}-\cos x+\mathrm{e}^x\right)+c;~c \in \mathbb{R}\)}{\( \frac37x^2\sqrt[3]{x}+x\cos x-\sin x+x\mathrm{e}^x-\mathrm{e}^x+c;~c \in \mathbb{R}\)}{\( \frac37x^2\sqrt[3]{x}-x\cos x + \sin x+x\mathrm{e}^x+\mathrm{e}^x+c;~c \in \mathbb{R}\)}{}{}}
\LANGsk{
\Question{
Vypočítajte nasledujúci integrál na intervale \((0;+\infty)\).
\[
\int \left( x\sqrt[3]{x}+x\sin x + x\mathrm{e}^x\right) \mathrm{d}x
\]}
{\( \frac37x^2\sqrt[3]{x}+x\cos x+\sin x+x\mathrm{e}^x-\mathrm{e}^x+c;~c \in \mathbb{R}\)}{\( \frac{x^2}2\left(\frac34x^{\frac43}-\cos x+\mathrm{e}^x\right)+c;~c \in \mathbb{R}\)}{\( \frac37x^2\sqrt[3]{x}+x\cos x-\sin x+x\mathrm{e}^x-\mathrm{e}^x+c;~c \in \mathbb{R}\)}{\( \frac37x^2\sqrt[3]{x}-x\cos x + \sin x+x\mathrm{e}^x+\mathrm{e}^x+c;~c \in \mathbb{R}\)}{}{}}
\LANGes{
\Question{
Evalúa la siguiente integral en el intervalo \((0;+\infty)\).
\[
\int \left( x\sqrt[3]{x}+x\sin x + x\mathrm{e}^x\right) \mathrm{d}x
\]}
{\( \frac37x^2\sqrt[3]{x}+x\cos x+\sin x+x\mathrm{e}^x-\mathrm{e}^x+c;~c \in \mathbb{R}\)}{\( \frac{x^2}2\left(\frac34x^{\frac43}-\cos x+\mathrm{e}^x\right)+c;~c \in \mathbb{R}\)}{\( \frac37x^2\sqrt[3]{x}+x\cos x-\sin x+x\mathrm{e}^x-\mathrm{e}^x+c;~c \in \mathbb{R}\)}{\( \frac37x^2\sqrt[3]{x}-x\cos x + \sin x+x\mathrm{e}^x+\mathrm{e}^x+c;~c \in \mathbb{R}\)}{}{}}
\End

\Data\ID{2010008109}
\Updated{2022-08-25 10:08:51}
\Level{B}
\SubArea{Primitive function}
\LANGen{
\Question{
Evaluate the following integral on the interval \(\left(\frac{\pi}{4};\frac{3\pi}{4}\right)\).
\[
\int \left(\frac{1}{2x}+\sin 2x - \frac{1}{\cos^2 2x}\right) \mathrm{d}x
\]}
{\( \frac12\left(\ln x - \cos 2x- \mathrm{tg}\, 2x\right) +c;~c \in \mathbb{R}\)}{\( \frac12\left(\ln(2x) - \cos 2x- \mathrm{tg }\, 2x \right)+c;~c \in \mathbb{R}\)}{\( \ln(2x) - \cos 2x- \mathrm{tg }\, 2x +c;~c \in \mathbb{R}\)}{\( \ln(2x) + \cos 2x- \mathrm{cotg }\, 2x +c;~c \in \mathbb{R}\)}{}{}}
\LANGpl{
\Question{
Oblicz następującą całkę na przedziale  \(\left(\frac{\pi}{4};\frac{3\pi}{4}\right)\).
\[
\int \left(\frac{1}{2x}+\sin 2x - \frac{1}{\cos^2 2x}\right) \mathrm{d}x
\]}
{\( \frac12\left(\ln x - \cos 2x- \mathrm{tg}\, 2x\right) +c;~c \in \mathbb{R}\)}{\( \frac12\left(\ln(2x) - \cos 2x- \mathrm{tg }\, 2x \right)+c;~c \in \mathbb{R}\)}{\( \ln(2x) - \cos 2x- \mathrm{tg }\, 2x +c;~c \in \mathbb{R}\)}{\( \ln(2x) + \cos 2x- \mathrm{cotg }\, 2x +c;~c \in \mathbb{R}\)}{}{}}
\LANGcs{
\Question{
Vypočtěte
\[
\int \left(\frac{1}{2x}+\sin 2x - \frac{1}{\cos^2 2x}\right) \mathrm{d}x
\]
na intervalu \(\left(\frac{\pi}{4};\frac{3\pi}{4}\right)\).}
{\( \frac12\left(\ln x - \cos 2x- \mathrm{tg}\, 2x\right) +c;~c \in \mathbb{R}\)}{\( \frac12\left(\ln(2x) - \cos 2x- \mathrm{tg }\, 2x \right)+c;~c \in \mathbb{R}\)}{\( \ln(2x) - \cos 2x- \mathrm{tg }\, 2x +c;~c \in \mathbb{R}\)}{\( \ln(2x) + \cos 2x- \mathrm{cotg }\, 2x +c;~c \in \mathbb{R}\)}{}{}}
\LANGsk{
\Question{
Vypočítajte nasledujúci integrál na intervale \(\left(\frac{\pi}{4};\frac{3\pi}{4}\right)\).
\[
\int \left(\frac{1}{2x}+\sin 2x - \frac{1}{\cos^2 2x}\right) \mathrm{d}x
\]}
{\( \frac12\left(\ln x - \cos 2x- \mathrm{tg}\, 2x\right) +c;~c \in \mathbb{R}\)}{\( \frac12\left(\ln(2x) - \cos 2x- \mathrm{tg }\, 2x \right)+c;~c \in \mathbb{R}\)}{\( \ln(2x) - \cos 2x- \mathrm{tg }\, 2x +c;~c \in \mathbb{R}\)}{\( \ln(2x) + \cos 2x- \mathrm{cotg }\, 2x +c;~c \in \mathbb{R}\)}{}{}}
\LANGes{
\Question{
Evalúa la siguiente integral en el intervalo \(\left(\frac{\pi}{4};\frac{3\pi}{4}\right)\).
\[
\int \left(\frac{1}{2x}+\sin 2x - \frac{1}{\cos^2 2x}\right) \mathrm{d}x
\]}
{\( \frac12\left(\ln x - \cos 2x- \mathrm{tg}\, 2x\right) +c;~c \in \mathbb{R}\)}{\( \frac12\left(\ln(2x) - \cos 2x- \mathrm{tg }\, 2x \right)+c;~c \in \mathbb{R}\)}{\( \ln(2x) - \cos 2x- \mathrm{tg }\, 2x +c;~c \in \mathbb{R}\)}{\( \ln(2x) + \cos 2x- \mathrm{cotg }\, 2x +c;~c \in \mathbb{R}\)}{}{}}
\End

\Data\ID{2010008110}
\Updated{2022-08-25 10:08:51}
\Level{B}
\SubArea{Primitive function}
\LANGen{
\Question{
Evaluate the following integral on the interval \(\left(0;\frac{\pi}{2}\right)\).
\[
\int \left(\cos 2x+ \frac{1}{\sin^2 2x}-\frac{1}{2x} \right) \mathrm{d}x
\]}
{\( \frac12\left(\sin 2x- \mathrm{cotg}\, 2x-\ln x \right) +c;~c \in \mathbb{R}\)}{\( \frac12\left( \sin 2x- \mathrm{cotg }\, 2x -\ln 2x\right)+c;~c \in \mathbb{R}\)}{\( \sin 2x- \mathrm{cotg }\, 2x - \ln 2x +c;~c \in \mathbb{R}\)}{\( \sin 2x+ \mathrm{cotg }\, 2x +\ln 2x +c;~c \in \mathbb{R}\)}{}{}}
\LANGpl{
\Question{
Oblicz następującą całkę na przedziale \(\left(0;\frac{\pi}{2}\right)\).
\[
\int \left(\cos 2x+ \frac{1}{\sin^2 2x}-\frac{1}{2x} \right) \mathrm{d}x
\]}
{\( \frac12\left(\sin 2x- \mathrm{cotg}\, 2x-\ln x \right) +c;~c \in \mathbb{R}\)}{\( \frac12\left( \sin 2x- \mathrm{cotg }\, 2x -\ln 2x\right)+c;~c \in \mathbb{R}\)}{\( \sin 2x- \mathrm{cotg }\, 2x - \ln 2x +c;~c \in \mathbb{R}\)}{\( \sin 2x+ \mathrm{cotg }\, 2x +\ln 2x +c;~c \in \mathbb{R}\)}{}{}}
\LANGcs{
\Question{
Vypočtěte
\[
\int \left(\cos 2x+ \frac{1}{\sin^2 2x}-\frac{1}{2x} \right) \mathrm{d}x
\]
na intervalu \(\left(0;\frac{\pi}{2}\right)\).}
{\( \frac12\left(\sin 2x- \mathrm{cotg}\, 2x-\ln x \right) +c;~c \in \mathbb{R}\)}{\( \frac12\left( \sin 2x- \mathrm{cotg }\, 2x -\ln 2x\right)+c;~c \in \mathbb{R}\)}{\( \sin 2x- \mathrm{cotg }\, 2x - \ln 2x +c;~c \in \mathbb{R}\)}{\( \sin 2x+ \mathrm{cotg }\, 2x +\ln 2x +c;~c \in \mathbb{R}\)}{}{}}
\LANGsk{
\Question{
Vypočítajte nasledujúci integrál na intervale \(\left(0;\frac{\pi}{2}\right)\).
\[
\int \left(\cos 2x+ \frac{1}{\sin^2 2x}-\frac{1}{2x} \right) \mathrm{d}x
\]}
{\( \frac12\left(\sin 2x- \mathrm{cotg}\, 2x-\ln x \right) +c;~c \in \mathbb{R}\)}{\( \frac12\left( \sin 2x- \mathrm{cotg }\, 2x -\ln 2x\right)+c;~c \in \mathbb{R}\)}{\( \sin 2x- \mathrm{cotg }\, 2x - \ln 2x +c;~c \in \mathbb{R}\)}{\( \sin 2x+ \mathrm{cotg }\, 2x +\ln 2x +c;~c \in \mathbb{R}\)}{}{}}
\LANGes{
\Question{
Evalúa la siguiente integral en el intervalo \(\left(0;\frac{\pi}{2}\right)\).
\[
\int \left(\cos 2x+ \frac{1}{\sin^2 2x}-\frac{1}{2x} \right) \mathrm{d}x
\]}
{\( \frac12\left(\sin 2x- \mathrm{cotg}\, 2x-\ln x \right) +c;~c \in \mathbb{R}\)}{\( \frac12\left( \sin 2x- \mathrm{cotg }\, 2x -\ln 2x\right)+c;~c \in \mathbb{R}\)}{\( \sin 2x- \mathrm{cotg }\, 2x - \ln 2x +c;~c \in \mathbb{R}\)}{\( \sin 2x+ \mathrm{cotg }\, 2x +\ln 2x +c;~c \in \mathbb{R}\)}{}{}}
\End

\Data\ID{9000065504}
\Updated{2020-07-15 03:07:37}
\Level{B}
\SubArea{Primitive function}
\LANGen{
\Question{
Evaluate the following integral on the interval \((0;+\infty)\). 
\[ 
 \int (1 -\sqrt{x})(1 + \sqrt{x})\, \mathrm{d}x
\]}
{\(x -\frac{1} 
{2}x^{2} + c,\ c\in \mathbb{R}\)}{\((x -\frac{1} 
{2}x^{2})(x + \frac{1} 
{2}x^{2}) + c,\ c\in \mathbb{R}\)}{\(x -\frac{1} 
{2}x^{\frac{1} 
{2} } + c,\ c\in \mathbb{R}\)}{\((x -\frac{1} 
{2}x^{-\frac{1} 
{2} })(x + \frac{1} {2}x^{-\frac{1} 
{2} }) + c,\ c\in \mathbb{R}\)}{}{}}
\LANGpl{
\Question{
Wyznacz całkę na przedziale  \((0;+\infty)\). 
\[ 
 \int (1 -\sqrt{x})(1 + \sqrt{x})\, \mathrm{d}x
\]}
{\(x -\frac{1} 
{2}x^{2} + c,\ c\in \mathbb{R}\)}{\((x -\frac{1} 
{2}x^{2})(x + \frac{1} 
{2}x^{2}) + c,\ c\in \mathbb{R}\)}{\(x -\frac{1} 
{2}x^{\frac{1} 
{2} } + c,\ c\in \mathbb{R}\)}{\((x -\frac{1} 
{2}x^{-\frac{1} 
{2} })(x + \frac{1} {2}x^{-\frac{1} 
{2} }) + c,\ c\in \mathbb{R}\)}{}{}}
\LANGcs{
\Question{
Vypočtěte \(\int (1 -\sqrt{x})(1 + \sqrt{x})\, \mathrm{d}x\) na intervalu \((0;+\infty)\).}
{\(x -\frac{1} 
{2}x^{2} + c,\ c\in\mathbb{R}\)}{\((x -\frac{1} 
{2}x^{2})(x + \frac{1} 
{2}x^{2}) + c,\ c\in\mathbb{R}\)}{\(x -\frac{1} 
{2}x^{\frac{1} 
{2} } + c,\ c\in\mathbb{R}\)}{\((x -\frac{1} 
{2}x^{-\frac{1} 
{2} })(x + \frac{1} {2}x^{-\frac{1} 
{2} }) + c,\ c\in\mathbb{R}\)}{}{}}
\LANGsk{
\Question{
Vypočítajte \(\int (1 -\sqrt{x})(1 + \sqrt{x})\, \mathrm{d}x\) na intervale \((0;+\infty)\).}
{\(x -\frac{1} 
{2}x^{2} + c,\ c\in\mathbb{R}\)}{\((x -\frac{1} 
{2}x^{2})(x + \frac{1} 
{2}x^{2}) + c,\ c\in\mathbb{R}\)}{\(x -\frac{1} 
{2}x^{\frac{1} 
{2} } + c,\ c\in\mathbb{R}\)}{\((x -\frac{1} 
{2}x^{-\frac{1} 
{2} })(x + \frac{1} {2}x^{-\frac{1} 
{2} }) + c,\ c\in\mathbb{R}\)}{}{}}
\LANGes{
\Question{
Evalúa la siguiente integral en el intervalo \((0;+\infty)\). 
\[ 
 \int (1 -\sqrt{x})(1 + \sqrt{x})\, \mathrm{d}x
\]}
{\(x -\frac{1} 
{2}x^{2} + c,\ c\in \mathbb{R}\)}{\((x -\frac{1} 
{2}x^{2})(x + \frac{1} 
{2}x^{2}) + c,\ c\in \mathbb{R}\)}{\(x -\frac{1} 
{2}x^{\frac{1} 
{2} } + c,\ c\in \mathbb{R}\)}{\((x -\frac{1} 
{2}x^{-\frac{1} 
{2} })(x + \frac{1} {2}x^{-\frac{1} 
{2} }) + c,\ c\in \mathbb{R}\)}{}{}}
\End

\Data\ID{9000065505}
\Updated{2022-04-23 05:04:08}
\Level{B}
\SubArea{Primitive function}
\LANGen{
\Question{
Evaluate the following integral on \(\mathbb{R}\). 
\[ 
 \int (x^{2} + 3)(x^{2} - 1)\, \mathrm{d}x
\]}
{\(\frac{1}
{5}x^{5} + \frac{2} 
{3}x^{3} - 3x + c,\ c\in \mathbb{R}\)}{\((\frac{1}
{3}x^{3} + 3x)(\frac{1}
{3}x^{3} - x) + c,\ c\in \mathbb{R}\)}{\(4x^{2} + c,\ c\in \mathbb{R}\)}{\(4x^{3} + 4x + c,\ c\in \mathbb{R}\)}{}{}}
\LANGpl{
\Question{
Wyznacz całkę.
\[ 
 \int (x^{2} + 3)(x^{2} - 1)\, \mathrm{d}x
\]}
{\(\frac{1}
{5}x^{5} + \frac{2} 
{3}x^{3} - 3x + c,\ c\in \mathbb{R}\)}{\((\frac{1}
{3}x^{3} + 3x)(\frac{1}
{3}x^{3} - x) + c,\ c\in \mathbb{R}\)}{\(4x^{2} + c,\ c\in \mathbb{R}\)}{\(4x^{3} + 4x + c,\ c\in \mathbb{R}\)}{}{}}
\LANGcs{
\Question{
Určete \(\int (x^{2} + 3)(x^{2} - 1)\, \mathrm{d}x\) na \(\mathbb{R}\).}
{\(\frac{1}
{5}x^{5} + \frac{2} 
{3}x^{3} - 3x + c,\ c\in\mathbb{R}\)}{\((\frac{1}
{3}x^{3} + 3x)(\frac{1}
{3}x^{3} - x) + c,\ c\in\mathbb{R}\)}{\(4x^{2} + c,\ c\in\mathbb{R}\)}{\(4x^{3} + 4x + c,\ c\in\mathbb{R}\)}{}{}}
\LANGsk{
\Question{
Určte \(\int (x^{2} + 3)(x^{2} - 1)\, \mathrm{d}x\) na \(\mathbb{R}\).}
{\(\frac{1}
{5}x^{5} + \frac{2} 
{3}x^{3} - 3x + c,\ c\in\mathbb{R}\)}{\((\frac{1}
{3}x^{3} + 3x)(\frac{1}
{3}x^{3} - x) + c,\ c\in\mathbb{R}\)}{\(4x^{2} + c,\ c\in\mathbb{R}\)}{\(4x^{3} + 4x + c,\ c\in\mathbb{R}\)}{}{}}
\LANGes{
\Question{
Evalúa la siguiente integral en \(\mathbb{R}\). 
\[ 
 \int (x^{2} + 3)(x^{2} - 1)\, \mathrm{d}x
\]}
{\(\frac{1}
{5}x^{5} + \frac{2} 
{3}x^{3} - 3x + c,\ c\in \mathbb{R}\)}{\((\frac{1}
{3}x^{3} + 3x)(\frac{1}
{3}x^{3} - x) + c,\ c\in \mathbb{R}\)}{\(4x^{2} + c,\ c\in \mathbb{R}\)}{\(4x^{3} + 4x + c,\ c\in \mathbb{R}\)}{}{}}
\End

\Data\ID{9000065506}
\Updated{2020-07-15 03:07:37}
\Level{B}
\SubArea{Primitive function}
\LANGen{
\Question{
Evaluate the following integral on the interval \((0;+\infty)\). 
\[ 
 \int \frac{x^{2}}
{\sqrt{x}}\, \mathrm{d}x
\]}
{\(\frac{2}
{5}x^{2}\sqrt{x} + c,\ c\in \mathbb{R}\)}{\(\frac{2\sqrt{x}}
 {x} + c,\ c\in \mathbb{R}\)}{\(\frac{2}
{5}x\sqrt{x} + c,\ c\in \mathbb{R}\)}{\(\frac{\sqrt{x}}
 {x} + c,\ c\in \mathbb{R}\)}{}{}}
\LANGpl{
\Question{
Wyznacz całkę na przedziale  \((0;+\infty)\). 
\[ 
 \int \frac{x^{2}}
{\sqrt{x}}\, \mathrm{d}x
\]}
{\(\frac{2}
{5}x^{2}\sqrt{x} + c,\ c\in \mathbb{R}\)}{\(\frac{2\sqrt{x}}
 {x} + c,\ c\in \mathbb{R}\)}{\(\frac{2}
{5}x\sqrt{x} + c,\ c\in \mathbb{R}\)}{\(\frac{\sqrt{x}}
 {x} + c,\ c\in \mathbb{R}\)}{}{}}
\LANGcs{
\Question{
Vypočtěte \(\int \frac{x^{2}} 
{\sqrt{x}}\, \mathrm{d}x\) na intervalu \((0;+\infty)\).}
{\(\frac{2}
{5}x^{2}\sqrt{x} + c,\ c\in\mathbb{R}\)}{\(\frac{2\sqrt{x}}
 {x} + c,\ c\in\mathbb{R}\)}{\(\frac{2}
{5}x\sqrt{x} + c,\ c\in\mathbb{R}\)}{\(\frac{\sqrt{x}}
 {x} + c,\ c\in\mathbb{R}\)}{}{}}
\LANGsk{
\Question{
Vypočítajte \(\int \frac{x^{2}} 
{\sqrt{x}}\, \mathrm{d}x\) na intervale \((0;+\infty)\).}
{\(\frac{2}
{5}x^{2}\sqrt{x} + c,\ c\in\mathbb{R}\)}{\(\frac{2\sqrt{x}}
 {x} + c,\ c\in\mathbb{R}\)}{\(\frac{2}
{5}x\sqrt{x} + c,\ c\in\mathbb{R}\)}{\(\frac{\sqrt{x}}
 {x} + c,\ c\in\mathbb{R}\)}{}{}}
\LANGes{
\Question{
Evalúa la siguiente integral en el intervalo \((0;+\infty)\). 
\[ 
 \int \frac{x^{2}}
{\sqrt{x}}\, \mathrm{d}x
\]}
{\(\frac{2}
{5}x^{2}\sqrt{x} + c,\ c\in \mathbb{R}\)}{\(\frac{2\sqrt{x}}
 {x} + c,\ c\in \mathbb{R}\)}{\(\frac{2}
{5}x\sqrt{x} + c,\ c\in \mathbb{R}\)}{\(\frac{\sqrt{x}}
 {x} + c,\ c\in \mathbb{R}\)}{}{}}
\End

\Data\ID{9000065901}
\Updated{2020-07-15 03:07:37}
\Level{B}
\SubArea{Primitive function}
\LANGen{
\Question{
Evaluate the following integral on the interval \((-1;+\infty)\). 
\[ 
 \int \frac{1}
{x + 1}\, \text{d}x
\]}
{\(\ln |x + 1| + c,\ c\in \mathbb{R}\)}{\(\ln |x| + c,\ c\in \mathbb{R}\)}{\(\frac{1}
{x} + c,\ c\in \mathbb{R}\)}{\(-\frac{1} 
{2}(x + 1)^{-2} + c,\ c\in \mathbb{R}\)}{}{}}
\LANGpl{
\Question{
Wyznacz całkę na przedziale \((-1;+\infty)\). 
\[ 
 \int \frac{1}
{x + 1}\, \text{d}x
\]}
{\(\ln |x + 1| + c,\ c\in \mathbb{R}\)}{\(\ln |x| + c,\ c\in \mathbb{R}\)}{\(\frac{1}
{x} + c,\ c\in \mathbb{R}\)}{\(-\frac{1} 
{2}(x + 1)^{-2} + c,\ c\in \mathbb{R}\)}{}{}}
\LANGcs{
\Question{
Vypočtěte \(\int \frac{1}
{x+1}\, \text{d}x\) na intervalu \((-1;+\infty)\).}
{\(\ln |x + 1| + c,\ c\in\mathbb{R}\)}{\(\ln |x| + c,\ c\in\mathbb{R}\)}{\(\frac{1}
{x} + c,\ c\in\mathbb{R}\)}{\(-\frac{1} 
{2}(x + 1)^{-2} + c,\ c\in\mathbb{R}\)}{}{}}
\LANGsk{
\Question{
Vypočítajte \(\int \frac{1}
{x+1}\, \text{d}x\) na intervale \((-1;+\infty)\).}
{\(\ln |x + 1| + c,\ c\in\mathbb{R}\)}{\(\ln |x| + c,\ c\in\mathbb{R}\)}{\(\frac{1}
{x} + c,\ c\in\mathbb{R}\)}{\(-\frac{1} 
{2}(x + 1)^{-2} + c,\ c\in\mathbb{R}\)}{}{}}
\LANGes{
\Question{
Evalúa la siguiente integral en el intervalo \((-1;+\infty)\). 
\[ 
 \int \frac{1}
{x + 1}\, \text{d}x
\]}
{\(\ln |x + 1| + c,\ c\in \mathbb{R}\)}{\(\ln |x| + c,\ c\in \mathbb{R}\)}{\(\frac{1}
{x} + c,\ c\in \mathbb{R}\)}{\(-\frac{1} 
{2}(x + 1)^{-2} + c,\ c\in \mathbb{R}\)}{}{}}
\End

\Data\ID{9000065903}
\Updated{2020-07-15 03:07:37}
\Level{B}
\SubArea{Primitive function}
\LANGen{
\Question{
Evaluate the following integral on the interval \((-6;+\infty)\). 
\[ 
 \int \frac{1}
{6x + 36}\, \text{d}x
\]}
{\(\frac{1}
{6}\ln |x + 6| + c,\ c\in \mathbb{R}\)}{\(-\frac{1} 
{2}(6x + 36)^{-2} + c,\ c\in \mathbb{R}\)}{\(6\ln |x + 6| + c,\ c\in \mathbb{R}\)}{\(12x^{2} + 36x + c,\ c\in \mathbb{R}\)}{}{}}
\LANGpl{
\Question{
Wyznacz całkę na przedziale  \((-6;+\infty)\). 
\[ 
 \int \frac{1}
{6x + 36}\, \text{d}x
\]}
{\(\frac{1}
{6}\ln |x + 6| + c,\ c\in \mathbb{R}\)}{\(-\frac{1} 
{2}(6x + 36)^{-2} + c,\ c\in \mathbb{R}\)}{\(6\ln |x + 6| + c,\ c\in \mathbb{R}\)}{\(12x^{2} + 36x + c,\ c\in \mathbb{R}\)}{}{}}
\LANGcs{
\Question{
Vypočtěte \(\int \frac{1}
{6x+36}\, \text{d}x\) na intervalu \((-6;+\infty)\).}
{\(\frac{1}
{6}\ln |x + 6| + c,\ c\in\mathbb{R}\)}{\(-\frac{1} 
{2}(6x + 36)^{-2} + c,\ c\in\mathbb{R}\)}{\(6\ln |x + 6| + c,\ c\in\mathbb{R}\)}{\(12x^{2} + 36x + c,\ c\in\mathbb{R}\)}{}{}}
\LANGsk{
\Question{
Vypočítajte \(\int \frac{1}
{6x+36}\, \text{d}x\) na intervale \((-6;+\infty)\).}
{\(\frac{1}
{6}\ln |x + 6| + c,\ c\in\mathbb{R}\)}{\(-\frac{1} 
{2}(6x + 36)^{-2} + c,\ c\in\mathbb{R}\)}{\(6\ln |x + 6| + c,\ c\in\mathbb{R}\)}{\(12x^{2} + 36x + c,\ c\in\mathbb{R}\)}{}{}}
\LANGes{
\Question{
Evalúa la siguiente integral en el intervalo \((-6;+\infty)\). 
\[ 
 \int \frac{1}
{6x + 36}\, \text{d}x
\]}
{\(\frac{1}
{6}\ln |x + 6| + c,\ c\in \mathbb{R}\)}{\(-\frac{1} 
{2}(6x + 36)^{-2} + c,\ c\in \mathbb{R}\)}{\(6\ln |x + 6| + c,\ c\in \mathbb{R}\)}{\(12x^{2} + 36x + c,\ c\in \mathbb{R}\)}{}{}}
\End

\Data\ID{9000065904}
\Updated{2022-04-23 05:04:08}
\Level{B}
\SubArea{Primitive function}
\LANGen{
\Question{
Evaluate the following integral on the interval \((0;+\infty)\). 
\[ 
 \int \frac{x^{3} + 2x}
 {x^{2}} \, \text{d}x
\]}
{\(\frac{1}
{2}x^{2} + 2\ln |x| + c,\ c\in \mathbb{R}\)}{\(x +\ln |x| + c,\ c\in \mathbb{R}\)}{\(\frac{1}
{4}x^{4} + 4x^{2} +\ln |x^{2}| + c,\ c\in \mathbb{R}\)}{\(2x^{2} + 2 +\ln |x^{2}| + c,\ c\in \mathbb{R}\)}{}{}}
\LANGpl{
\Question{
Wyznacz całkę na przedziale  \((0;+\infty)\). 
\[ 
 \int \frac{x^{3} + 2x}
 {x^{2}} \, \text{d}x
\]}
{\(\frac{1}
{2}x^{2} + 2\ln |x| + c,\ c\in \mathbb{R}\)}{\(x +\ln |x| + c,\ c\in \mathbb{R}\)}{\(\frac{1}
{4}x^{4} + 4x^{2} +\ln |x^{2}| + c,\ c\in \mathbb{R}\)}{\(2x^{2} + 2 +\ln |x^{2}| + c,\ c\in \mathbb{R}\)}{}{}}
\LANGcs{
\Question{
Vypočtěte \(\int \frac{x^{3}+2x}
 {x^{2}} \, \text{d}x\) na intervalu \((0;+\infty)\).}
{\(\frac{1}
{2}x^{2} + 2\ln |x| + c,\ c\in\mathbb{R}\)}{\(x +\ln |x| + c,\ c\in\mathbb{R}\)}{\(\frac{1}
{4}x^{4} + 4x^{2} +\ln |x^{2}| + c,\ c\in\mathbb{R}\)}{\(2x^{2} + 2 +\ln |x^{2}| + c,\ c\in\mathbb{R}\)}{}{}}
\LANGsk{
\Question{
Vypočítajte \(\int \frac{x^{3}+2x}
 {x^{2}} \, \text{d}x\) na intervale \((0;+\infty)\).}
{\(\frac{1}
{2}x^{2} + 2\ln |x| + c,\ c\in\mathbb{R}\)}{\(x +\ln |x| + c,\ c\in\mathbb{R}\)}{\(\frac{1}
{4}x^{4} + 4x^{2} +\ln |x^{2}| + c,\ c\in\mathbb{R}\)}{\(2x^{2} + 2 +\ln |x^{2}| + c,\ c\in\mathbb{R}\)}{}{}}
\LANGes{
\Question{
Evalúa la siguiente integral en el intervalo \((0;+\infty)\). 
\[ 
 \int \frac{x^{3} + 2x}
 {x^{2}} \, \text{d}x
\]}
{\(\frac{1}
{2}x^{2} + 2\ln |x| + c,\ c\in \mathbb{R}\)}{\(x +\ln |x| + c,\ c\in \mathbb{R}\)}{\(\frac{1}
{4}x^{4} + 4x^{2} +\ln |x^{2}| + c,\ c\in \mathbb{R}\)}{\(2x^{2} + 2 +\ln |x^{2}| + c,\ c\in \mathbb{R}\)}{}{}}
\End

\Data\ID{9000065905}
\Updated{2020-07-15 03:07:37}
\Level{B}
\SubArea{Primitive function}
\LANGen{
\Question{
Evaluate the following integral on the interval \((0;+\infty)\). 
\[ 
 \int \frac{\left (\sqrt{x} + 2\right )^{2}}
 {x} \, \text{d}x
\]}
{\(x + 8\sqrt{x} + 4\ln |x| + c,\ c\in \mathbb{R}\)}{\(\sqrt{x} + 8x + 4\ln |x| + c,\ c\in \mathbb{R}\)}{\(\frac{1}
{2}x^{-\frac{1} 
{2} } + 2x +\ln |x| + c,\ c\in \mathbb{R}\)}{\(1 + 8\sqrt{x} + 4\ln |x| + c,\ c\in \mathbb{R}\)}{}{}}
\LANGpl{
\Question{
Wyznacz całkę na przedziale  \((0;+\infty)\). 
\[ 
 \int \frac{\left (\sqrt{x} + 2\right )^{2}}
 {x} \, \text{d}x
\]}
{\(x + 8\sqrt{x} + 4\ln |x| + c,\ c\in \mathbb{R}\)}{\(\sqrt{x} + 8x + 4\ln |x| + c,\ c\in \mathbb{R}\)}{\(\frac{1}
{2}x^{-\frac{1} 
{2} } + 2x +\ln |x| + c,\ c\in \mathbb{R}\)}{\(1 + 8\sqrt{x} + 4\ln |x| + c,\ c\in \mathbb{R}\)}{}{}}
\LANGcs{
\Question{
Vypočtěte \(\int \frac{\left (\sqrt{x}+2\right )^{2}} 
 {x} \, \text{d}x\) na intervalu \((0;+\infty)\).}
{\(x + 8\sqrt{x} + 4\ln |x| + c,\ c\in\mathbb{R}\)}{\(\sqrt{x} + 8x + 4\ln |x| + c,\ c\in\mathbb{R}\)}{\(\frac{1}
{2}x^{-\frac{1} 
{2} } + 2x +\ln |x| + c,\ c\in\mathbb{R}\)}{\(1 + 8\sqrt{x} + 4\ln |x| + c,\ c\in\mathbb{R}\)}{}{}}
\LANGsk{
\Question{
Vypočítajte \(\int \frac{\left (\sqrt{x}+2\right )^{2}} 
 {x} \, \text{d}x\) na intervale \((0;+\infty)\).}
{\(x + 8\sqrt{x} + 4\ln |x| + c,\ c\in\mathbb{R}\)}{\(\sqrt{x} + 8x + 4\ln |x| + c,\ c\in\mathbb{R}\)}{\(\frac{1}
{2}x^{-\frac{1} 
{2} } + 2x +\ln |x| + c,\ c\in\mathbb{R}\)}{\(1 + 8\sqrt{x} + 4\ln |x| + c,\ c\in\mathbb{R}\)}{}{}}
\LANGes{
\Question{
Evalúa la siguiente integral en el intervalo \((0;+\infty)\). 
\[ 
 \int \frac{\left (\sqrt{x} + 2\right )^{2}}
 {x} \, \text{d}x
\]}
{\(x + 8\sqrt{x} + 4\ln |x| + c,\ c\in \mathbb{R}\)}{\(\sqrt{x} + 8x + 4\ln |x| + c,\ c\in \mathbb{R}\)}{\(\frac{1}
{2}x^{-\frac{1} 
{2} } + 2x +\ln |x| + c,\ c\in \mathbb{R}\)}{\(1 + 8\sqrt{x} + 4\ln |x| + c,\ c\in \mathbb{R}\)}{}{}}
\End

\Data\ID{9000065906}
\Updated{2020-07-15 03:07:37}
\Level{B}
\SubArea{Primitive function}
\LANGen{
\Question{
Evaluate the following integral on the interval \((-3;+\infty)\). 
\[ 
 \int \frac{x^{2} - 9}
 {x + 3} \, \text{d}x
\]}
{\(\frac{1}
{2}x^{2} - 3x + c,\ c\in \mathbb{R}\)}{\(\frac{1}
{3}x^{3} - 9x +\ln |x + 3| + c,\ c\in \mathbb{R}\)}{\(2x - x^{-2} + c,\ c\in \mathbb{R}\)}{\(\frac{1}
{2}x^{2} + 3x + c,\ c\in \mathbb{R}\)}{}{}}
\LANGpl{
\Question{
Wyznacz całkę na przedziale  \((-3;+\infty)\). 
\[ 
 \int \frac{x^{2} - 9}
 {x + 3} \, \text{d}x
\]}
{\(\frac{1}
{2}x^{2} - 3x + c,\ c\in \mathbb{R}\)}{\(\frac{1}
{3}x^{3} - 9x +\ln |x + 3| + c,\ c\in \mathbb{R}\)}{\(2x - x^{-2} + c,\ c\in \mathbb{R}\)}{\(\frac{1}
{2}x^{2} + 3x + c,\ c\in \mathbb{R}\)}{}{}}
\LANGcs{
\Question{
Vypočtěte \(\int \frac{x^{2}-9}
 {x+3} \, \text{d}x\) na intervalu \((-3;+\infty)\).}
{\(\frac{1}
{2}x^{2} - 3x + c,\ c\in\mathbb{R}\)}{\(\frac{1}
{3}x^{3} - 9x +\ln |x + 3| + c,\ c\in\mathbb{R}\)}{\(2x - x^{-2} + c,\ c\in\mathbb{R}\)}{\(\frac{1}
{2}x^{2} + 3x + c,\ c\in\mathbb{R}\)}{}{}}
\LANGsk{
\Question{
Vypočítajte \(\int \frac{x^{2}-9}
 {x+3} \, \text{d}x\) na intervale \((-3;+\infty)\).}
{\(\frac{1}
{2}x^{2} - 3x + c,\ c\in\mathbb{R}\)}{\(\frac{1}
{3}x^{3} - 9x +\ln |x + 3| + c,\ c\in\mathbb{R}\)}{\(2x - x^{-2} + c,\ c\in\mathbb{R}\)}{\(\frac{1}
{2}x^{2} + 3x + c,\ c\in\mathbb{R}\)}{}{}}
\LANGes{
\Question{
Evalúa la siguiente integral en el intervalo \((-3;+\infty)\). 
\[ 
 \int \frac{x^{2} - 9}
 {x + 3} \, \text{d}x
\]}
{\(\frac{1}
{2}x^{2} - 3x + c,\ c\in \mathbb{R}\)}{\(\frac{1}
{3}x^{3} - 9x +\ln |x + 3| + c,\ c\in \mathbb{R}\)}{\(2x - x^{-2} + c,\ c\in \mathbb{R}\)}{\(\frac{1}
{2}x^{2} + 3x + c,\ c\in \mathbb{R}\)}{}{}}
\End

\Data\ID{9000065907}
\Updated{2020-07-15 03:07:37}
\Level{B}
\SubArea{Primitive function}
\LANGen{
\Question{
Evaluate the following integral on \(\mathbb{R}\). 
\[ 
 \int \frac{x^{4} - 1}
{x^{2} + 1}\, \text{d}x
\]}
{\(\frac{1}
{3}x^{3} - x + c,\ c\in \mathbb{R}\)}{\(\frac{1}
{3}x^{3} + x + c,\ c\in \mathbb{R}\)}{\(\frac{1}
{5}x^{5} - x +\ln |x^{2} - 1| + c,\ c\in \mathbb{R}\)}{\(3x^{2} -\ln |x^{2} - 1| + c,\ c\in \mathbb{R}\)}{}{}}
\LANGpl{
\Question{
Wyznacz całkę.
\[ 
 \int \frac{x^{4} - 1}
{x^{2} + 1}\, \text{d}x
\]}
{\(\frac{1}
{3}x^{3} - x + c,\ c\in \mathbb{R}\)}{\(\frac{1}
{3}x^{3} + x + c,\ c\in \mathbb{R}\)}{\(\frac{1}
{5}x^{5} - x +\ln |x^{2} - 1| + c,\ c\in \mathbb{R}\)}{\(3x^{2} -\ln |x^{2} - 1| + c,\ c\in \mathbb{R}\)}{}{}}
\LANGcs{
\Question{
Vypočtěte \(\int \frac{x^{4}-1}
{x^{2}+1}\, \text{d}x\) na \(\mathbb{R}\).}
{\(\frac{1}
{3}x^{3} - x + c,\ c\in\mathbb{R}\)}{\(\frac{1}
{3}x^{3} + x + c,\ c\in\mathbb{R}\)}{\(\frac{1}
{5}x^{5} - x +\ln |x^{2} - 1| + c,\ c\in\mathbb{R}\)}{\(3x^{2} -\ln |x^{2} - 1| + c,\ c\in\mathbb{R}\)}{}{}}
\LANGsk{
\Question{
Vypočítajte \(\int \frac{x^{4}-1}
{x^{2}+1}\, \text{d}x\) na \(\mathbb{R}\).}
{\(\frac{1}
{3}x^{3} - x + c,\ c\in\mathbb{R}\)}{\(\frac{1}
{3}x^{3} + x + c,\ c\in\mathbb{R}\)}{\(\frac{1}
{5}x^{5} - x +\ln |x^{2} - 1| + c,\ c\in\mathbb{R}\)}{\(3x^{2} -\ln |x^{2} - 1| + c,\ c\in\mathbb{R}\)}{}{}}
\LANGes{
\Question{
Evalúa la siguiente integral en \(\mathbb{R}\). 
\[ 
 \int \frac{x^{4} - 1}
{x^{2} + 1}\, \text{d}x
\]}
{\(\frac{1}
{3}x^{3} - x + c,\ c\in \mathbb{R}\)}{\(\frac{1}
{3}x^{3} + x + c,\ c\in \mathbb{R}\)}{\(\frac{1}
{5}x^{5} - x +\ln |x^{2} - 1| + c,\ c\in \mathbb{R}\)}{\(3x^{2} -\ln |x^{2} - 1| + c,\ c\in \mathbb{R}\)}{}{}}
\End

\Data\ID{9000066002}
\Updated{2022-04-23 05:04:08}
\Level{B}
\SubArea{Primitive function}
\LANGen{
\Question{
Evaluate the following integral on \(\mathbb{R}\). 
\[ 
 \int x\sin x\, \mathrm{d}x
\]}
{\(- x\cos x +\sin x + c,\ c\in \mathbb{R}\)}{\(- x\cos x -\sin x + c,\ c\in \mathbb{R}\)}{\(x\cos x +\sin x + c,\ c\in \mathbb{R}\)}{\(x\cos x -\sin x + c,\ c\in \mathbb{R}\)}{}{}}
\LANGpl{
\Question{
Wyznacz całkę.
\[ 
 \int x\sin x\, \mathrm{d}x
\]}
{\(- x\cos x +\sin x + c,\ c\in \mathbb{R}\)}{\(- x\cos x -\sin x + c,\ c\in \mathbb{R}\)}{\(x\cos x +\sin x + c,\ c\in \mathbb{R}\)}{\(x\cos x -\sin x + c,\ c\in \mathbb{R}\)}{}{}}
\LANGcs{
\Question{
Vypočtěte \(\int x\sin x\, \mathrm{d}x\) na \(\mathbb{R}\).}
{\(- x\cos x +\sin x + c,\ c\in\mathbb{R}\)}{\(- x\cos x -\sin x + c,\ c\in\mathbb{R}\)}{\(x\cos x +\sin x + c,\ c\in\mathbb{R}\)}{\(x\cos x -\sin x + c,\ c\in\mathbb{R}\)}{}{}}
\LANGsk{
\Question{
Vypočítajte \(\int x\sin x\, \mathrm{d}x\) na \(\mathbb{R}\).}
{\(- x\cos x +\sin x + c,\ c\in\mathbb{R}\)}{\(- x\cos x -\sin x + c,\ c\in\mathbb{R}\)}{\(x\cos x +\sin x + c,\ c\in\mathbb{R}\)}{\(x\cos x -\sin x + c,\ c\in\mathbb{R}\)}{}{}}
\LANGes{
\Question{
Evalúa la siguiente integral en \(\mathbb{R}\). 
\[ 
 \int x\sin x\, \mathrm{d}x
\]}
{\(- x\cos x +\sin x + c,\ c\in \mathbb{R}\)}{\(- x\cos x -\sin x + c,\ c\in \mathbb{R}\)}{\(x\cos x +\sin x + c,\ c\in \mathbb{R}\)}{\(x\cos x -\sin x + c,\ c\in \mathbb{R}\)}{}{}}
\End

\Data\ID{9000066003}
\Updated{2020-11-12 01:11:11}
\Level{B}
\SubArea{Primitive function}
\LANGen{
\Question{
Evaluate the following integral on \(\mathbb{R}\). 
\[ 
 \int x\cos x\, \mathrm{d}x
\]}
{\(x\sin x +\cos x + c,\ c\in \mathbb{R}\)}{\(- x\cos x +\sin x + c,\ c\in \mathbb{R}\)}{\(x\cos x -\sin x + c,\ c\in \mathbb{R}\)}{\(x\sin x -\cos x + c,\ c\in \mathbb{R}\)}{}{}}
\LANGpl{
\Question{
Wyznacz całkę.
\[ 
 \int x\cos x\, \mathrm{d}x
\]}
{\(x\sin x +\cos x + c,\ c\in \mathbb{R}\)}{\(- x\cos x +\sin x + c,\ c\in \mathbb{R}\)}{\(x\cos x -\sin x + c,\ c\in \mathbb{R}\)}{\(x\sin x -\cos x + c,\ c\in \mathbb{R}\)}{}{}}
\LANGcs{
\Question{
Vypočtěte \(\int x\cos x\, \mathrm{d}x\) na \(\mathbb{R}\).}
{\(x\sin x +\cos x + c,\ c\in\mathbb{R}\)}{\(- x\cos x +\sin x + c,\ c\in\mathbb{R}\)}{\(x\cos x -\sin x + c,\ c\in\mathbb{R}\)}{\(x\sin x -\cos x + c,\ c\in\mathbb{R}\)}{}{}}
\LANGsk{
\Question{
Vypočítajte \(\int x\cos x\, \mathrm{d}x\) na \(\mathbb{R}\).}
{\(x\sin x +\cos x + c,\ c\in\mathbb{R}\)}{\(- x\cos x +\sin x + c,\ c\in\mathbb{R}\)}{\(x\cos x -\sin x + c,\ c\in\mathbb{R}\)}{\(x\sin x -\cos x + c,\ c\in\mathbb{R}\)}{}{}}
\LANGes{
\Question{
Resuelve la siguiente integral en \(\mathbb{R}\).  
\[ 
 \int x\cos x\, \mathrm{d}x
\]}
{\(x\sin x +\cos x + c,\ c\in \mathbb{R}\)}{\(- x\cos x +\sin x + c,\ c\in \mathbb{R}\)}{\(x\cos x -\sin x + c,\ c\in \mathbb{R}\)}{\(x\sin x -\cos x + c,\ c\in \mathbb{R}\)}{}{}}
\End

\Data\ID{9000066008}
\Updated{2022-04-23 05:04:08}
\Level{B}
\SubArea{Primitive function}
\LANGen{
\Question{
Evaluate the following integral on \(\mathbb{R}\). 
\[ 
 \int x\mathrm{e}^{x}\, \mathrm{d}x
\]}
{\(x\mathrm{e}^{x} -\mathrm{e}^{x} + c,\ c\in \mathbb{R}\)}{\(x^{2}\mathrm{e}^{x} - 2x\mathrm{e}^{x} + 2\mathrm{e}^{x} + c,\ c\in \mathbb{R}\)}{\(2x^{3}\mathrm{e}^{x} - x\mathrm{e}^{x} + c,\ c\in \mathbb{R}\)}{\(\frac{1}
{2}x^{2}\mathrm{e}^{x} + c,\ c\in \mathbb{R}\)}{}{}}
\LANGpl{
\Question{
Wyznacz całkę.
\[ 
 \int x\mathrm{e}^{x}\, \mathrm{d}x
\]}
{\(x\mathrm{e}^{x} -\mathrm{e}^{x} + c,\ c\in \mathbb{R}\)}{\(x^{2}\mathrm{e}^{x} - 2x\mathrm{e}^{x} + 2\mathrm{e}^{x} + c,\ c\in \mathbb{R}\)}{\(2x^{3}\mathrm{e}^{x} - x\mathrm{e}^{x} + c,\ c\in \mathbb{R}\)}{\(\frac{1}
{2}x^{2}\mathrm{e}^{x} + c,\ c\in \mathbb{R}\)}{}{}}
\LANGcs{
\Question{
Vypočtěte \(\int x\mathrm{e}^{x}\, \mathrm{d}x\) na \(\mathbb{R}\).}
{\(x\mathrm{e}^{x} -\mathrm{e}^{x} + c,\ c\in\mathbb{R}\)}{\(x^{2}\mathrm{e}^{x} - 2x\mathrm{e}^{x} + 2\mathrm{e}^{x} + c,\ c\in\mathbb{R}\)}{\(2x^{3}\mathrm{e}^{x} - x\mathrm{e}^{x} + c,\ c\in\mathbb{R}\)}{\(\frac{1}
{2}x^{2}\mathrm{e}^{x} + c,\ c\in\mathbb{R}\)}{}{}}
\LANGsk{
\Question{
Vypočítajte \(\int x\mathrm{e}^{x}\, \mathrm{d}x\) na \(\mathbb{R}\).}
{\(x\mathrm{e}^{x} -\mathrm{e}^{x} + c,\ c\in\mathbb{R}\)}{\(x^{2}\mathrm{e}^{x} - 2x\mathrm{e}^{x} + 2\mathrm{e}^{x} + c,\ c\in\mathbb{R}\)}{\(2x^{3}\mathrm{e}^{x} - x\mathrm{e}^{x} + c,\ c\in\mathbb{R}\)}{\(\frac{1}
{2}x^{2}\mathrm{e}^{x} + c,\ c\in\mathbb{R}\)}{}{}}
\LANGes{
\Question{
Resuelve la siguiente integral en \(\mathbb{R}\).  
\[ 
 \int x\mathrm{e}^{x}\, \mathrm{d}x
\]}
{\(x\mathrm{e}^{x} -\mathrm{e}^{x} + c,\ c\in \mathbb{R}\)}{\(x^{2}\mathrm{e}^{x} - 2x\mathrm{e}^{x} + 2\mathrm{e}^{x} + c,\ c\in \mathbb{R}\)}{\(2x^{3}\mathrm{e}^{x} - x\mathrm{e}^{x} + c,\ c\in \mathbb{R}\)}{\(\frac{1}
{2}x^{2}\mathrm{e}^{x} + c,\ c\in \mathbb{R}\)}{}{}}
\End

\Data\ID{9000071201}
\Updated{2022-04-23 05:04:08}
\Level{B}
\SubArea{Primitive function}
\LANGen{
\Question{
Evaluate the following integral on \(\mathbb{R}\). 
\[ 
 \int (x^{3} - 2)^{2}\, \mathrm{d}x
\]}
{\(\frac{x^{7}} 
 {7} - x^{4} + 4x + c,\ c\in \mathbb{R}\)}{\(\frac{(x^{3}-2)^{3}} 
 {3} + c,\ c\in \mathbb{R}\)}{\(6x^{7} - 12x^{4} + 4x + c,\ c\in \mathbb{R}\)}{}{}{}}
\LANGpl{
\Question{
Wyznacz całkę.
\[ 
 \int (x^{3} - 2)^{2}\, \mathrm{d}x
\]}
{\(\frac{x^{7}} 
 {7} - x^{4} + 4x + c,\ c\in \mathbb{R}\)}{\(\frac{(x^{3}-2)^{3}} 
 {3} + c,\ c\in \mathbb{R}\)}{\(6x^{7} - 12x^{4} + 4x + c,\ c\in \mathbb{R}\)}{}{}{}}
\LANGcs{
\Question{
Vypočtěte \(\int (x^{3} - 2)^{2}\, \mathrm{d}x\) na \(\mathbb{R}\).}
{\(\frac{x^{7}} 
 {7} - x^{4} + 4x + c,\ c\in \mathbb{R}\)}{\(\frac{(x^{3}-2)^{3}} 
 {3} + c,\ c\in \mathbb{R}\)}{\(6x^{7} - 12x^{4} + 4x + c,\ c\in \mathbb{R}\)}{}{}{}}
\LANGsk{
\Question{
Vypočítajte \(\int (x^{3} - 2)^{2}\, \mathrm{d}x\) na \(\mathbb{R}\).}
{\(\frac{x^{7}} 
 {7} - x^{4} + 4x + c,\ c\in \mathbb{R}\)}{\(\frac{(x^{3}-2)^{3}} 
 {3} + c,\ c\in \mathbb{R}\)}{\(6x^{7} - 12x^{4} + 4x + c,\ c\in \mathbb{R}\)}{}{}{}}
\LANGes{
\Question{
Evalúa la siguiente integral en \(\mathbb{R}\). 
\[ 
 \int (x^{3} - 2)^{2}\, \mathrm{d}x
\]}
{\(\frac{x^{7}} 
 {7} - x^{4} + 4x + c,\ c\in \mathbb{R}\)}{\(\frac{(x^{3}-2)^{3}} 
 {3} + c,\ c\in \mathbb{R}\)}{\(6x^{7} - 12x^{4} + 4x + c,\ c\in \mathbb{R}\)}{}{}{}}
\End

\Data\ID{9000071202}
\Updated{2022-04-23 05:04:08}
\Level{B}
\SubArea{Primitive function}
\LANGen{
\Question{
Evaluate the following integral on the interval \((0;+\infty)\). 
\[ 
 \int \frac{11\sqrt{x^{3}} - 2}
 {\root{3}\of{x^{2}}} \, \mathrm{d}x
\]}
{\(6(x\root{6}\of{x^{5}} -\root{3}\of{x}) + c,\ c\in \mathbb{R}\)}{\(\frac{\frac{22} 
 {5} \sqrt{x^{5}}-2x}
 {\frac{3} 
{5} \root{3}\of{x^{5}}} + c,\ c\in \mathbb{R}\)}{\(\frac{121}
 {6} \root{6}\of{x^{11}} -\frac{2} 
{3}\root{3}\of{x} + c,\ c\in \mathbb{R}\)}{}{}{}}
\LANGpl{
\Question{
Wyznacz całkę na przedziale \((0;+\infty)\). 
\[ 
 \int \frac{11\sqrt{x^{3}} - 2}
 {\root{3}\of{x^{2}}} \, \mathrm{d}x
\]}
{\(6(x\root{6}\of{x^{5}} -\root{3}\of{x}) + c,\ c\in \mathbb{R}\)}{\(\frac{\frac{22} 
 {5} \sqrt{x^{5}}-2x}
 {\frac{3} 
{5} \root{3}\of{x^{5}}} + c,\ c\in \mathbb{R}\)}{\(\frac{121}
 {6} \root{6}\of{x^{11}} -\frac{2} 
{3}\root{3}\of{x} + c,\ c\in \mathbb{R}\)}{}{}{}}
\LANGcs{
\Question{
Vypočtěte \(\int \frac{11\sqrt{x^{3}}-2}
 {\root{3}\of{x^{2}}} \, \mathrm{d}x\) na intervalu \((0;+\infty)\).}
{\(6(x\root{6}\of{x^{5}} -\root{3}\of{x}) + c,\ c\in \mathbb{R}\)}{\(\frac{\frac{22} 
 {5} \sqrt{x^{5}}-2x}
 {\frac{3} 
{5} \root{3}\of{x^{5}}} + c,\ c\in \mathbb{R}\)}{\(\frac{121}
 {6} \root{6}\of{x^{11}} -\frac{2} 
{3}\root{3}\of{x} + c,\ c\in \mathbb{R}\)}{}{}{}}
\LANGsk{
\Question{
Vypočítajte  \(\int \frac{11\sqrt{x^{3}}-2}
 {\root{3}\of{x^{2}}} \, \mathrm{d}x\) na intervale \((0;+\infty)\).}
{\(6(x\root{6}\of{x^{5}} -\root{3}\of{x}) + c,\ c\in \mathbb{R}\)}{\(\frac{\frac{22} 
 {5} \sqrt{x^{5}}-2x}
 {\frac{3} 
{5} \root{3}\of{x^{5}}} + c,\ c\in \mathbb{R}\)}{\(\frac{121}
 {6} \root{6}\of{x^{11}} -\frac{2} 
{3}\root{3}\of{x} + c,\ c\in \mathbb{R}\)}{}{}{}}
\LANGes{
\Question{
Evalúa la siguiente integral en el intervalo \((0;+\infty)\). 
\[ 
 \int \frac{11\sqrt{x^{3}} - 2}
 {\root{3}\of{x^{2}}} \, \mathrm{d}x
\]}
{\(6(x\root{6}\of{x^{5}} -\root{3}\of{x}) + c,\ c\in \mathbb{R}\)}{\(\frac{\frac{22} 
 {5} \sqrt{x^{5}}-2x}
 {\frac{3} 
{5} \root{3}\of{x^{5}}} + c,\ c\in \mathbb{R}\)}{\(\frac{121}
 {6} \root{6}\of{x^{11}} -\frac{2} 
{3}\root{3}\of{x} + c,\ c\in \mathbb{R}\)}{}{}{}}
\End

\Data\ID{9000071203}
\Updated{2022-04-23 05:04:02}
\Level{B}
\SubArea{Primitive function}
\LANGen{
\Question{
Evaluate the following integral on the interval \((0;\frac{\pi}2)\). 
\[ 
 \int \frac{\cos 2x}
{\sin ^{2}x}\, \mathrm{d}x
\]}
{\(- 2x -\mathop{\mathrm{cotg}}\nolimits x + c,\ c\in \mathbb{R}\)}{\(\frac{\sin 2x}
{-\frac{1} 
{3} \cos ^{3}x} + c,\ c\in \mathbb{R}\)}{\(\mathop{\mathrm{tg}}\nolimits x - 2x + c,\ c\in \mathbb{R}\)}{}{}{}}
\LANGpl{
\Question{
Wyznacz całkę na przedziale \((0;\frac{\pi}2)\). 
\[ 
 \int \frac{\cos 2x}
{\sin ^{2}x}\, \mathrm{d}x
\]}
{\(- 2x -\mathop{\mathrm{cotg}}\nolimits x + c,\ c\in \mathbb{R}\)}{\(\frac{\sin 2x}
{-\frac{1} 
{3} \cos ^{3}x} + c,\ c\in \mathbb{R}\)}{\(\mathop{\mathrm{tg}}\nolimits x - 2x + c,\ c\in \mathbb{R}\)}{}{}{}}
\LANGcs{
\Question{
Vypočtěte \(\int \frac{\cos 2x}
{\sin ^{2}x}\, \mathrm{d}x\) na intervalu \((0;\frac{\pi}2)\).}
{\(- 2x -\mathop{\mathrm{cotg}}\nolimits x + c,\ c\in \mathbb{R}\)}{\(\frac{\sin 2x}
{-\frac{1} 
{3} \cos ^{3}x} + c,\ c\in \mathbb{R}\)}{\(\mathop{\mathrm{tg}}\nolimits x - 2x + c,\ c\in \mathbb{R}\)}{}{}{}}
\LANGsk{
\Question{
Vypočítajte \(\int \frac{\cos 2x}
{\sin ^{2}x}\, \mathrm{d}x\) na intervale \((0;\frac{\pi}2)\).}
{\(- 2x -\mathop{\mathrm{cotg}}\nolimits x + c,\ c\in \mathbb{R}\)}{\(\frac{\sin 2x}
{-\frac{1} 
{3} \cos ^{3}x} + c,\ c\in \mathbb{R}\)}{\(\mathop{\mathrm{tg}}\nolimits x - 2x + c,\ c\in \mathbb{R}\)}{}{}{}}
\LANGes{
\Question{
Evalúa la siguiente integral en el intervalo \((0;\frac{\pi}2)\). 
\[ 
 \int \frac{\cos 2x}
{\sin ^{2}x}\, \mathrm{d}x
\]}
{\(- 2x -\mathop{\mathrm{cotg}}\nolimits x + c,\ c\in \mathbb{R}\)}{\(\frac{\sin 2x}
{-\frac{1} 
{3} \cos ^{3}x} + c,\ c\in \mathbb{R}\)}{\(\mathop{\mathrm{tg}}\nolimits x - 2x + c,\ c\in \mathbb{R}\)}{}{}{}}
\End

\Data\ID{9000071207}
\Updated{2022-04-23 05:04:02}
\Level{B}
\SubArea{Primitive function}
\LANGen{
\Question{
Evaluate the following integral on the interval \(\left(\sqrt{\frac43};+\infty\right)\). 
\[ 
 \int \frac{6x}
{(3x^{2} - 4)^{2}}\, \mathrm{d}x
\]}
{\(\frac{1}
{4-3x^{2}} + c,\ c\in \mathbb{R}\)}{\(\frac{3x^{2}} 
{x^{3}-12x^{2}+16x} + c,\ c\in \mathbb{R}\)}{\(\frac{1}
{(3x^{2}-4)^{2}} + c,\ c\in \mathbb{R}\)}{}{}{}}
\LANGpl{
\Question{
Wyznacz całkę na przedziale  \(\left(\sqrt{\frac43};+\infty\right)\). 
\[ 
 \int \frac{6x}
{(3x^{2} - 4)^{2}}\, \mathrm{d}x
\]}
{\(\frac{1}
{4-3x^{2}} + c,\ c\in \mathbb{R}\)}{\(\frac{3x^{2}} 
{x^{3}-12x^{2}+16x} + c,\ c\in \mathbb{R}\)}{\(\frac{1}
{(3x^{2}-4)^{2}} + c,\ c\in \mathbb{R}\)}{}{}{}}
\LANGcs{
\Question{
Vypočtěte \(\int \frac{6x}
{(3x^{2}-4)^{2}} \, \mathrm{d}x\) na intervalu \(\left(\sqrt{\frac43};+\infty\right)\).}
{\(\frac{1}
{4-3x^{2}} + c,\ c\in \mathbb{R}\)}{\(\frac{3x^{2}} 
{x^{3}-12x^{2}+16x} + c,\ c\in \mathbb{R}\)}{\(\frac{1}
{(3x^{2}-4)^{2}} + c,\ c\in \mathbb{R}\)}{}{}{}}
\LANGsk{
\Question{
Vypočítajte \(\int \frac{6x}
{(3x^{2}-4)^{2}} \, \mathrm{d}x\) na intervale \(\left(\sqrt{\frac43};+\infty\right)\).}
{\(\frac{1}
{4-3x^{2}} + c,\ c\in \mathbb{R}\)}{\(\frac{3x^{2}} 
{x^{3}-12x^{2}+16x} + c,\ c\in \mathbb{R}\)}{\(\frac{1}
{(3x^{2}-4)^{2}} + c,\ c\in \mathbb{R}\)}{}{}{}}
\LANGes{
\Question{
Evalúa la siguiente integral en el intervalo \(\left(\sqrt{\frac43};+\infty\right)\). 
\[ 
 \int \frac{6x}
{(3x^{2} - 4)^{2}}\, \mathrm{d}x
\]}
{\(\frac{1}
{4-3x^{2}} + c,\ c\in \mathbb{R}\)}{\(\frac{3x^{2}} 
{x^{3}-12x^{2}+16x} + c,\ c\in \mathbb{R}\)}{\(\frac{1}
{(3x^{2}-4)^{2}} + c,\ c\in \mathbb{R}\)}{}{}{}}
\End

\Data\ID{9000150102}
\Updated{2020-07-15 03:07:37}
\Level{B}
\SubArea{Primitive function}
\LANGen{
\Question{
Evaluate the following integral on \(\mathbb{R}\). 
\[ 
 \int 2\sin 2x\, \mathrm{d}x
\]}
{\(-\cos 2x + c,\ c\in \mathbb{R}\)}{\(\cos 2x + c,\ c\in \mathbb{R}\)}{\(- 4\cos 2x + c,\ c\in \mathbb{R}\)}{\(4\cos 2x + c,\ c\in \mathbb{R}\)}{}{}}
\LANGpl{
\Question{
Wyznacz całkę.
\[ 
 \int 2\sin 2x\, \mathrm{d}x
\]}
{\(-\cos 2x + c,\ c\in \mathbb{R}\)}{\(\cos 2x + c,\ c\in \mathbb{R}\)}{\(- 4\cos 2x + c,\ c\in \mathbb{R}\)}{\(4\cos 2x + c,\ c\in \mathbb{R}\)}{}{}}
\LANGcs{
\Question{
Vypočtěte \(\int 2\sin 2x\, \mathrm{d}x\) na \(\mathbb{R}\).}
{\(-\cos 2x + c,\ c\in \mathbb{R}\)}{\(\cos 2x + c,\ c\in \mathbb{R}\)}{\(- 4\cos 2x + c,\ c\in \mathbb{R}\)}{\(4\cos 2x + c,\ c\in \mathbb{R}\)}{}{}}
\LANGsk{
\Question{
Vypočítajte \(\int 2\sin 2x\, \mathrm{d}x\) na \(\mathbb{R}\).}
{\(-\cos 2x + c,\ c\in \mathbb{R}\)}{\(\cos 2x + c,\ c\in \mathbb{R}\)}{\(- 4\cos 2x + c,\ c\in \mathbb{R}\)}{\(4\cos 2x + c,\ c\in \mathbb{R}\)}{}{}}
\LANGes{
\Question{
Evalúa la siguiente integral en \(\mathbb{R}\). 
\[ 
 \int 2\sin 2x\, \mathrm{d}x
\]}
{\(-\cos 2x + c,\ c\in \mathbb{R}\)}{\(\cos 2x + c,\ c\in \mathbb{R}\)}{\(- 4\cos 2x + c,\ c\in \mathbb{R}\)}{\(4\cos 2x + c,\ c\in \mathbb{R}\)}{}{}}
\End

\Data\ID{9000150106}
\Updated{2022-04-23 05:04:02}
\Level{B}
\SubArea{Primitive function}
\LANGen{
\Question{
Evaluate the following integral on the interval \(\left(\frac25;+\infty\right)\). 
\[ 
 \int \frac{7}
{2 - 5x}\, \mathrm{d}x
\]}
{\(-\frac{7} 
{5}\ln |2 - 5x| + c,\ c\in \mathbb{R}\)}{\(- \frac{7} 
{5\cdot \ln |2-5x|} + c,\ c\in \mathbb{R}\)}{\(\frac{7}
{5}\ln |2 - 5x| + c,\ c\in \mathbb{R}\)}{\(\frac{7}
{5\cdot \ln |2-5x|} + c,\ c\in \mathbb{R}\)}{}{}}
\LANGpl{
\Question{
Wyznacz całkę na przedziale \(\left(\frac25;+\infty\right)\). 
\[ 
 \int \frac{7}
{2 - 5x}\, \mathrm{d}x
\]}
{\(-\frac{7} 
{5}\ln |2 - 5x| + c,\ c\in \mathbb{R}\)}{\(- \frac{7} 
{5\cdot \ln |2-5x|} + c,\ c\in \mathbb{R}\)}{\(\frac{7}
{5}\ln |2 - 5x| + c,\ c\in \mathbb{R}\)}{\(\frac{7}
{5\cdot \ln |2-5x|} + c,\ c\in \mathbb{R}\)}{}{}}
\LANGcs{
\Question{
Vypočtěte \(\int \frac{7}
{2-5x}\, \mathrm{d}x\) na intervalu \(\left(\frac25;+\infty\right)\).}
{\(-\frac{7} 
{5}\ln |2 - 5x| + c,\ c\in \mathbb{R}\)}{\(- \frac{7} 
{5\cdot \ln |2-5x|} + c,\ c\in \mathbb{R}\)}{\(\frac{7}
{5}\ln |2 - 5x| + c,\ c\in \mathbb{R}\)}{\(\frac{7}
{5\cdot \ln |2-5x|} + c,\ c\in \mathbb{R}\)}{}{}}
\LANGsk{
\Question{
Vypočítajte \(\int \frac{7}
{2-5x}\, \mathrm{d}x\) na intervale \(\left(\frac25;+\infty\right)\).}
{\(-\frac{7} 
{5}\ln |2 - 5x| + c,\ c\in \mathbb{R}\)}{\(- \frac{7} 
{5\cdot \ln |2-5x|} + c,\ c\in \mathbb{R}\)}{\(\frac{7}
{5}\ln |2 - 5x| + c,\ c\in \mathbb{R}\)}{\(\frac{7}
{5\cdot \ln |2-5x|} + c,\ c\in \mathbb{R}\)}{}{}}
\LANGes{
\Question{
Evalúa la siguiente integral en el intervalo \(\left(\frac25;+\infty\right)\). 
\[ 
 \int \frac{7}
{2 - 5x}\, \mathrm{d}x
\]}
{\(-\frac{7} 
{5}\ln |2 - 5x| + c,\ c\in \mathbb{R}\)}{\(- \frac{7} 
{5\cdot \ln |2-5x|} + c,\ c\in \mathbb{R}\)}{\(\frac{7}
{5}\ln |2 - 5x| + c,\ c\in \mathbb{R}\)}{\(\frac{7}
{5\cdot \ln |2-5x|} + c,\ c\in \mathbb{R}\)}{}{}}
\End

\Data\ID{9000150107}
\Updated{2020-07-15 03:07:37}
\Level{B}
\SubArea{Primitive function}
\LANGen{
\Question{
Evaluate the following integral on the interval \((3;+\infty)\). 
\[ 
 \int \frac{x^{3} - 27}
 {x - 3} \, \mathrm{d}x
\]}
{\(\frac{x^{3}} 
 {3} + \frac{3x^{2}} 
 {2} + 9x + c,\ c\in \mathbb{R}\)}{\(\frac{x^{3}} 
 {3} -\frac{3x^{2}} 
 {2} + 9x + c,\ c\in \mathbb{R}\)}{\(\frac{x^{3}} 
 {3} -\frac{3x^{2}} 
 {2} - 9x + c,\ c\in \mathbb{R}\)}{\(\frac{x^{3}} 
 {3} + \frac{3x^{2}} 
 {2} - 9x + c,\ c\in \mathbb{R}\)}{}{}}
\LANGpl{
\Question{
Wyznacz całkę na przedziale  \((3;+\infty)\). 
\[ 
 \int \frac{x^{3} - 27}
 {x - 3} \, \mathrm{d}x
\]}
{\(\frac{x^{3}} 
 {3} + \frac{3x^{2}} 
 {2} + 9x + c,\ c\in \mathbb{R}\)}{\(\frac{x^{3}} 
 {3} -\frac{3x^{2}} 
 {2} + 9x + c,\ c\in \mathbb{R}\)}{\(\frac{x^{3}} 
 {3} -\frac{3x^{2}} 
 {2} - 9x + c,\ c\in \mathbb{R}\)}{\(\frac{x^{3}} 
 {3} + \frac{3x^{2}} 
 {2} - 9x + c,\ c\in \mathbb{R}\)}{}{}}
\LANGcs{
\Question{
Vypočtěte \(\int \frac{x^{3}-27}
 {x-3} \, \mathrm{d}x\) na intervalu \((3;+\infty)\).}
{\(\frac{x^{3}} 
 {3} + \frac{3x^{2}} 
 {2} + 9x + c,\ c\in \mathbb{R}\)}{\(\frac{x^{3}} 
 {3} -\frac{3x^{2}} 
 {2} + 9x + c,\ c\in \mathbb{R}\)}{\(\frac{x^{3}} 
 {3} -\frac{3x^{2}} 
 {2} - 9x + c,\ c\in \mathbb{R}\)}{\(\frac{x^{3}} 
 {3} + \frac{3x^{2}} 
 {2} - 9x + c,\ c\in \mathbb{R}\)}{}{}}
\LANGsk{
\Question{
Vypočítajte \(\int \frac{x^{3}-27}
 {x-3} \, \mathrm{d}x\) na intervale \((3;+\infty)\).}
{\(\frac{x^{3}} 
 {3} + \frac{3x^{2}} 
 {2} + 9x + c,\ c\in \mathbb{R}\)}{\(\frac{x^{3}} 
 {3} -\frac{3x^{2}} 
 {2} + 9x + c,\ c\in \mathbb{R}\)}{\(\frac{x^{3}} 
 {3} -\frac{3x^{2}} 
 {2} - 9x + c,\ c\in \mathbb{R}\)}{\(\frac{x^{3}} 
 {3} + \frac{3x^{2}} 
 {2} - 9x + c,\ c\in \mathbb{R}\)}{}{}}
\LANGes{
\Question{
Evalúa la siguiente integral en el intervalo \((3;+\infty)\). 
\[ 
 \int \frac{x^{3} - 27}
 {x - 3} \, \mathrm{d}x
\]}
{\(\frac{x^{3}} 
 {3} + \frac{3x^{2}} 
 {2} + 9x + c,\ c\in \mathbb{R}\)}{\(\frac{x^{3}} 
 {3} -\frac{3x^{2}} 
 {2} + 9x + c,\ c\in \mathbb{R}\)}{\(\frac{x^{3}} 
 {3} -\frac{3x^{2}} 
 {2} - 9x + c,\ c\in \mathbb{R}\)}{\(\frac{x^{3}} 
 {3} + \frac{3x^{2}} 
 {2} - 9x + c,\ c\in \mathbb{R}\)}{}{}}
\End

\Data\ID{1003107601}
\Updated{2020-07-15 03:07:12}
\Level{C}
\SubArea{Primitive function}
\LANGen{
\Question{
Evaluate the following integral on \( (3;\infty) \).
\[ \int\frac{5x-3}{x^2-2x-3}\mathrm{d}x \]}
{\( \ln\!\left[(x-3)^3\cdot(x+1)^2\right]+c\text{, }c\in\mathbb{R} \)}{\( \ln\!\left[(x-3)^2\cdot(x+1)^3\right]+c\text{, }c\in\mathbb{R} \)}{\( \ln\frac{(x-3)^2}{(x+1)^3}+c\text{, }c\in\mathbb{R} \)}{\( \ln\frac{(x-3)^3}{(x+1)^2}+c\text{, }c\in\mathbb{R} \)}{}{}}
\LANGpl{
\Question{
Oblicz całkę na przedziale \( (3;\infty) \).
\[ \int\frac{5x-3}{x^2-2x-3}\mathrm{d}x \]}
{\( \ln\!\left[(x-3)^3\cdot(x+1)^2\right]+c\text{, }c\in\mathbb{R} \)}{\( \ln\!\left[(x-3)^2\cdot(x+1)^3\right]+c\text{, }c\in\mathbb{R} \)}{\( \ln\frac{(x-3)^2}{(x+1)^3}+c\text{, }c\in\mathbb{R} \)}{\( \ln\frac{(x-3)^3}{(x+1)^2}+c\text{, }c\in\mathbb{R} \)}{}{}}
\LANGcs{
\Question{
Vypočítejte na intervalu \( (3;\infty) \) následující integrál.
\[ \int\frac{5x-3}{x^2-2x-3}\mathrm{d}x \]}
{\( \ln\!\left[(x-3)^3\cdot(x+1)^2\right]+c\text{, }c\in\mathbb{R} \)}{\( \ln\!\left[(x-3)^2\cdot(x+1)^3\right]+c\text{, }c\in\mathbb{R} \)}{\( \ln\frac{(x-3)^2}{(x+1)^3}+c\text{, }c\in\mathbb{R} \)}{\( \ln\frac{(x-3)^3}{(x+1)^2}+c\text{, }c\in\mathbb{R} \)}{}{}}
\LANGsk{
\Question{
Vypočítajte na intervale \( (3;\infty) \) nasledujúci integrál.
\[ \int\frac{5x-3}{x^2-2x-3}\mathrm{d}x \]}
{\( \ln\!\left[(x-3)^3\cdot(x+1)^2\right]+c\text{, }c\in\mathbb{R} \)}{\( \ln\!\left[(x-3)^2\cdot(x+1)^3\right]+c\text{, }c\in\mathbb{R} \)}{\( \ln\frac{(x-3)^2}{(x+1)^3}+c\text{, }c\in\mathbb{R} \)}{\( \ln\frac{(x-3)^3}{(x+1)^2}+c\text{, }c\in\mathbb{R} \)}{}{}}
\LANGes{
\Question{
Evalúa la siguiente integral en el intervalo \( (3;\infty) \).
\[ \int\frac{5x-3}{x^2-2x-3}\mathrm{d}x \]}
{\( \ln\!\left[(x-3)^3\cdot(x+1)^2\right]+c\text{, }c\in\mathbb{R} \)}{\( \ln\!\left[(x-3)^2\cdot(x+1)^3\right]+c\text{, }c\in\mathbb{R} \)}{\( \ln\frac{(x-3)^2}{(x+1)^3}+c\text{, }c\in\mathbb{R} \)}{\( \ln\frac{(x-3)^3}{(x+1)^2}+c\text{, }c\in\mathbb{R} \)}{}{}}
\End

\Data\ID{1003107602}
\Updated{2020-07-15 03:07:12}
\Level{C}
\SubArea{Primitive function}
\LANGen{
\Question{
Evaluate the following integral on \( (3;\infty) \).
\[ \int\frac7{x^2+x-12}\mathrm{d}x \]}
{\( \ln\frac{x-3}{x+4}+c\text{, }c\in\mathbb{R} \)}{\( \ln\!\left[(x-3)(x+4)\right]+c\text{, }c\in\mathbb{R} \)}{\( \ln\frac{x+4}{x-3}+c\text{, }c\in\mathbb{R} \)}{\( \ln\frac{x+6}{x-2}+c\text{, }c\in\mathbb{R} \)}{}{}}
\LANGpl{
\Question{
Oblicz całkę na przedziale \( (3;\infty) \).
\[ \int\frac7{x^2+x-12}\mathrm{d}x \]}
{\( \ln\frac{x-3}{x+4}+c\text{, }c\in\mathbb{R} \)}{\( \ln\!\left[(x-3)(x+4)\right]+c\text{, }c\in\mathbb{R} \)}{\( \ln\frac{x+4}{x-3}+c\text{, }c\in\mathbb{R} \)}{\( \ln\frac{x+6}{x-2}+c\text{, }c\in\mathbb{R} \)}{}{}}
\LANGcs{
\Question{
Vypočítejte na intervalu \( (3;\infty) \) následující integrál.
\[ \int\frac7{x^2+x-12}\mathrm{d}x \]}
{\( \ln\frac{x-3}{x+4}+c\text{, }c\in\mathbb{R} \)}{\( \ln\!\left[(x-3)(x+4)\right]+c\text{, }c\in\mathbb{R} \)}{\( \ln\frac{x+4}{x-3}+c\text{, }c\in\mathbb{R} \)}{\( \ln\frac{x+6}{x-2}+c\text{, }c\in\mathbb{R} \)}{}{}}
\LANGsk{
\Question{
Vypočítajte na intervale \( (3;\infty) \) nasledujúci integrál.
\[ \int\frac7{x^2+x-12}\mathrm{d}x \]}
{\( \ln\frac{x-3}{x+4}+c\text{, }c\in\mathbb{R} \)}{\( \ln\!\left[(x-3)(x+4)\right]+c\text{, }c\in\mathbb{R} \)}{\( \ln\frac{x+4}{x-3}+c\text{, }c\in\mathbb{R} \)}{\( \ln\frac{x+6}{x-2}+c\text{, }c\in\mathbb{R} \)}{}{}}
\LANGes{
\Question{
Evalúa la siguiente integral en el intervalo \( (3;\infty) \).
\[ \int\frac7{x^2+x-12}\mathrm{d}x \]}
{\( \ln\frac{x-3}{x+4}+c\text{, }c\in\mathbb{R} \)}{\( \ln\!\left[(x-3)(x+4)\right]+c\text{, }c\in\mathbb{R} \)}{\( \ln\frac{x+4}{x-3}+c\text{, }c\in\mathbb{R} \)}{\( \ln\frac{x+6}{x-2}+c\text{, }c\in\mathbb{R} \)}{}{}}
\End

\Data\ID{1003107603}
\Updated{2020-07-15 03:07:12}
\Level{C}
\SubArea{Primitive function}
\LANGen{
\Question{
Evaluate the following integral on \( (0;\infty) \).
\[ \int\frac{3x^3+3x^2-x+1}{x^2+x}\mathrm{d}x \]}
{\( 1.5x^2+\ln\frac x{x^2+2x+1}+c\text{, }c\in\mathbb{R} \)}{\( \ln\frac x{(x+1)^2}+3x+c\text{, }c\in\mathbb{R} \)}{\( \ln\!\left[x\cdot(x+1)^2\right]+3x+c\text{, }c\in\mathbb{R} \)}{\( \ln\frac{(x+1)^2}x+c\text{, }c\in\mathbb{R} \)}{}{}}
\LANGpl{
\Question{
Oblicz całkę na przedziale \( (0;\infty) \).
\[ \int\frac{3x^3+3x^2-x+1}{x^2+x}\mathrm{d}x \]}
{\( 1{,}5x^2+\ln\frac x{x^2+2x+1}+c\text{, }c\in\mathbb{R} \)}{\( \ln\frac x{(x+1)^2}+3x+c\text{, }c\in\mathbb{R} \)}{\( \ln\!\left[x\cdot(x+1)^2\right]+3x+c\text{, }c\in\mathbb{R} \)}{\( \ln\frac{(x+1)^2}x+c\text{, }c\in\mathbb{R} \)}{}{}}
\LANGcs{
\Question{
Vypočítejte na intervalu \( (0;\infty) \) následující integrál.
\[ \int\frac{3x^3+3x^2-x+1}{x^2+x}\mathrm{d}x \]}
{\( 1{,}5x^2+\ln\frac x{x^2+2x+1}+c\text{, }c\in\mathbb{R} \)}{\( \ln\frac x{(x+1)^2}+3x+c\text{, }c\in\mathbb{R} \)}{\( \ln\!\left[x\cdot(x+1)^2\right]+3x+c\text{, }c\in\mathbb{R} \)}{\( \ln\frac{(x+1)^2}x+c\text{, }c\in\mathbb{R} \)}{}{}}
\LANGsk{
\Question{
Vypočítajte na intervale \( (0;\infty) \) nasledujúci integrál.
\[ \int\frac{3x^3+3x^2-x+1}{x^2+x}\mathrm{d}x \]}
{\( 1{,}5x^2+\ln\frac x{x^2+2x+1}+c\text{, }c\in\mathbb{R} \)}{\( \ln\frac x{(x+1)^2}+3x+c\text{, }c\in\mathbb{R} \)}{\( \ln\!\left[x\cdot(x+1)^2\right]+3x+c\text{, }c\in\mathbb{R} \)}{\( \ln\frac{(x+1)^2}x+c\text{, }c\in\mathbb{R} \)}{}{}}
\LANGes{
\Question{
Evalúa la siguiente integral en el intervalo \( (0;\infty) \).
\[ \int\frac{3x^3+3x^2-x+1}{x^2+x}\mathrm{d}x \]}
{\( 1.5x^2+\ln\frac x{x^2+2x+1}+c\text{, }c\in\mathbb{R} \)}{\( \ln\frac x{(x+1)^2}+3x+c\text{, }c\in\mathbb{R} \)}{\( \ln\!\left[x\cdot(x+1)^2\right]+3x+c\text{, }c\in\mathbb{R} \)}{\( \ln\frac{(x+1)^2}x+c\text{, }c\in\mathbb{R} \)}{}{}}
\End

\Data\ID{1003107604}
\Updated{2020-07-15 03:07:12}
\Level{C}
\SubArea{Primitive function}
\LANGen{
\Question{
Evaluate the following integral on \( (0;\infty) \).
\[ \int\frac{-2x^2+3x+2}{x^3+x^2}\mathrm{d}x \]}
{\( \ln\frac x{x^3+3x^2+3x+1}-\frac2x+c\text{, }c\in\mathbb{R} \)}{\( \ln x-\ln(x+1)^3+c\text{, }c\in\mathbb{R} \)}{\( \ln\frac{x\cdot x^2}{(x+1)^3}+c\text{, }c\in\mathbb{R} \)}{\(\ln\frac{x^3}{x^3+3x^2+3x+1}-\frac2x+c\text{, }c\in\mathbb{R} \)}{}{}}
\LANGpl{
\Question{
Oblicz całkę na przedziale \( (0;\infty) \).
\[ \int\frac{-2x^2+3x+2}{x^3+x^2}\mathrm{d}x \]}
{\( \ln\frac x{x^3+3x^2+3x+1}-\frac2x+c\text{, }c\in\mathbb{R} \)}{\( \ln x-\ln(x+1)^3+c\text{, }c\in\mathbb{R} \)}{\( \ln\frac{x\cdot x^2}{(x+1)^3}+c\text{, }c\in\mathbb{R} \)}{\(\ln\frac{x^3}{x^3+3x^2+3x+1}-\frac2x+c\text{, }c\in\mathbb{R} \)}{}{}}
\LANGcs{
\Question{
Vypočítejte na intervalu \( (0;\infty) \) následující integrál.
\[ \int\frac{-2x^2+3x+2}{x^3+x^2}\mathrm{d}x \]}
{\( \ln\frac x{x^3+3x^2+3x+1}-\frac2x+c\text{, }c\in\mathbb{R} \)}{\( \ln x-\ln(x+1)^3+c\text{, }c\in\mathbb{R} \)}{\( \ln\frac{x\cdot x^2}{(x+1)^3}+c\text{, }c\in\mathbb{R} \)}{\(\ln\frac{x^3}{x^3+3x^2+3x+1}-\frac2x+c\text{, }c\in\mathbb{R} \)}{}{}}
\LANGsk{
\Question{
Vypočítajte na intervale \( (0;\infty) \) nasledujúci integrál.
\[ \int\frac{-2x^2+3x+2}{x^3+x^2}\mathrm{d}x \]}
{\( \ln\frac x{x^3+3x^2+3x+1}-\frac2x+c\text{, }c\in\mathbb{R} \)}{\( \ln x-\ln(x+1)^3+c\text{, }c\in\mathbb{R} \)}{\( \ln\frac{x\cdot x^2}{(x+1)^3}+c\text{, }c\in\mathbb{R} \)}{\(\ln\frac{x^3}{x^3+3x^2+3x+1}-\frac2x+c\text{, }c\in\mathbb{R} \)}{}{}}
\LANGes{
\Question{
Evalúa la siguiente integral en el intervalo \( (0;\infty) \).
\[ \int\frac{-2x^2+3x+2}{x^3+x^2}\mathrm{d}x \]}
{\( \ln\frac x{x^3+3x^2+3x+1}-\frac2x+c\text{, }c\in\mathbb{R} \)}{\( \ln x-\ln(x+1)^3+c\text{, }c\in\mathbb{R} \)}{\( \ln\frac{x\cdot x^2}{(x+1)^3}+c\text{, }c\in\mathbb{R} \)}{\(\ln\frac{x^3}{x^3+3x^2+3x+1}-\frac2x+c\text{, }c\in\mathbb{R} \)}{}{}}
\End

\Data\ID{1003107605}
\Updated{2020-07-15 03:07:12}
\Level{C}
\SubArea{Primitive function}
\LANGen{
\Question{
Evaluate the following integral on \( (0;\infty) \).
\[ \int\frac{7x+2}{15x^2+3x}\mathrm{d}x \]}
{\( \frac23\ln x-\frac15\ln\left(x+\frac15\right)+c\text{, }c\in\mathbb{R} \)}{\( \ln\frac{\sqrt[3]{x^2}}{\sqrt[5]{x+1}}+c\text{, }c\in\mathbb{R} \)}{\( \frac23\ln x-\ln(5x+1)+c\text{, }c\in\mathbb{R} \)}{\(\ln\frac{\sqrt[3]{x^2}}{\sqrt[5]{x+5}}+c\text{, }c\in\mathbb{R} \)}{}{}}
\LANGpl{
\Question{
Oblicz całkę na przedziale \( (0;\infty) \).
\[ \int\frac{7x+2}{15x^2+3x}\mathrm{d}x \]}
{\( \frac23\ln x-\frac15\ln\left(x+\frac15\right)+c\text{, }c\in\mathbb{R} \)}{\( \ln\frac{\sqrt[3]{x^2}}{\sqrt[5]{x+1}}+c\text{, }c\in\mathbb{R} \)}{\( \frac23\ln x-\ln(5x+1)+c\text{, }c\in\mathbb{R} \)}{\(\ln\frac{\sqrt[3]{x^2}}{\sqrt[5]{x+5}}+c\text{, }c\in\mathbb{R} \)}{}{}}
\LANGcs{
\Question{
Vypočítejte na intervalu \( (0;\infty) \) následující integrál.
\[ \int\frac{7x+2}{15x^2+3x}\mathrm{d}x \]}
{\( \frac23\ln x-\frac15\ln\left(x+\frac15\right)+c\text{, }c\in\mathbb{R} \)}{\( \ln\frac{\sqrt[3]{x^2}}{\sqrt[5]{x+1}}+c\text{, }c\in\mathbb{R} \)}{\( \frac23\ln x-\ln(5x+1)+c\text{, }c\in\mathbb{R} \)}{\(\ln\frac{\sqrt[3]{x^2}}{\sqrt[5]{x+5}}+c\text{, }c\in\mathbb{R} \)}{}{}}
\LANGsk{
\Question{
Vypočítajte na intervale \( (0;\infty) \) nasledujúci integrál.
\[ \int\frac{7x+2}{15x^2+3x}\mathrm{d}x \]}
{\( \frac23\ln x-\frac15\ln\left(x+\frac15\right)+c\text{, }c\in\mathbb{R} \)}{\( \ln\frac{\sqrt[3]{x^2}}{\sqrt[5]{x+1}}+c\text{, }c\in\mathbb{R} \)}{\( \frac23\ln x-\ln(5x+1)+c\text{, }c\in\mathbb{R} \)}{\(\ln\frac{\sqrt[3]{x^2}}{\sqrt[5]{x+5}}+c\text{, }c\in\mathbb{R} \)}{}{}}
\LANGes{
\Question{
Evalúa la siguiente integral en el intervalo \( (0;\infty) \).
\[ \int\frac{7x+2}{15x^2+3x}\mathrm{d}x \]}
{\( \frac23\ln x-\frac15\ln\left(x+\frac15\right)+c\text{, }c\in\mathbb{R} \)}{\( \ln\frac{\sqrt[3]{x^2}}{\sqrt[5]{x+1}}+c\text{, }c\in\mathbb{R} \)}{\( \frac23\ln x-\ln(5x+1)+c\text{, }c\in\mathbb{R} \)}{\(\ln\frac{\sqrt[3]{x^2}}{\sqrt[5]{x+5}}+c\text{, }c\in\mathbb{R} \)}{}{}}
\End

\Data\ID{1003107901}
\Updated{2020-07-15 03:07:19}
\Level{C}
\SubArea{Primitive function}
\LANGen{
\Question{
Solve the indefinite integral 
\[ \int\sin^3 x\cos^2x\,\mathrm{d}x \]
of a real-valued function using a suitable substitution.}
{\( \frac{\cos^5x}5-\frac{\cos^3x}3+c \)}{\( -\frac{\cos^5x}5+\frac{\cos^3x}3+c \)}{\( \frac{\cos^5x}5+\frac{\cos^3x}3+c \)}{\( -\frac{\cos^5x}5-\frac{\cos^3x}3+c \)}{}{}}
\LANGpl{
\Question{
Oblicz całkę nieoznaczoną
\[ \int\sin^3 x\cos^2x\,\mathrm{d}x \]
funkcji rzeczywistej używając właściwego podstawienia.}
{\( \frac{\cos^5x}5-\frac{\cos^3x}3+c \)}{\( -\frac{\cos^5x}5+\frac{\cos^3x}3+c \)}{\( \frac{\cos^5x}5+\frac{\cos^3x}3+c \)}{\( -\frac{\cos^5x}5-\frac{\cos^3x}3+c \)}{}{}}
\LANGcs{
\Question{
Užitím vhodné substituce řešte v oboru reálných čísel neurčitý integrál.
\[ \int\sin^3 x\cos^2x\,\mathrm{d}x \]}
{\( \frac{\cos^5x}5-\frac{\cos^3x}3+c \)}{\( -\frac{\cos^5x}5+\frac{\cos^3x}3+c \)}{\( \frac{\cos^5x}5+\frac{\cos^3x}3+c \)}{\( -\frac{\cos^5x}5-\frac{\cos^3x}3+c \)}{}{}}
\LANGsk{
\Question{
Použitím vhodnej substitúcie riešte v obore reálnych čísel neurčitý integrál.
\[ \int\sin^3 x\cos^2x\,\mathrm{d}x \]}
{\( \frac{\cos^5x}5-\frac{\cos^3x}3+c \)}{\( -\frac{\cos^5x}5+\frac{\cos^3x}3+c \)}{\( \frac{\cos^5x}5+\frac{\cos^3x}3+c \)}{\( -\frac{\cos^5x}5-\frac{\cos^3x}3+c \)}{}{}}
\LANGes{
\Question{
Resuelve la integral indefinida  
\[ \int\sin^3 x\cos^2x\,\mathrm{d}x \]
en el conjunto de números reales usando una sustitución adecuada.}
{\( \frac{\cos^5x}5-\frac{\cos^3x}3+c \)}{\( -\frac{\cos^5x}5+\frac{\cos^3x}3+c \)}{\( \frac{\cos^5x}5+\frac{\cos^3x}3+c \)}{\( -\frac{\cos^5x}5-\frac{\cos^3x}3+c \)}{}{}}
\End

\Data\ID{1003107902}
\Updated{2020-07-15 03:07:19}
\Level{C}
\SubArea{Primitive function}
\LANGen{
\Question{
Solve the indefinite integral
\[ \int\left(ab\mathrm{e}^c-bx^2+5^b-\sin c \right)\mathrm{d}x \]
of a real-valued function, where \( a \), \( b \), \( c \) are real numbers.}
{\( ab\mathrm{e}^c x-b\frac{x^3}3+5^b x-x \sin c+k \), \( k\in\mathbb{R} \)}{\( -b\frac{x^3}3+k \), \( k\in\mathbb{R} \)}{\( ab\mathrm{e}^c-2b+5^b-\sin c+k \), \( k\in\mathbb{R} \)}{\( ab\mathrm{e}^c x-2bx+5^b x-x \sin c+k \), \( k\in\mathbb{R} \)}{}{}}
\LANGpl{
\Question{
Oblicz całkę nieoznaczoną
\[ \int\left(ab\mathrm{e}^c-bx^2+5^b-\sin c \right)\mathrm{d}x \]
funkcji rzeczywistej, gdzie \( a \), \( b \), \( c \) to liczby rzeczywiste.}
{\( ab\mathrm{e}^c x-b\frac{x^3}3+5^b x-x \sin c+k \), \( k\in\mathbb{R} \)}{\( -b\frac{x^3}3+k \), \( k\in\mathbb{R} \)}{\( ab\mathrm{e}^c-2b+5^b-\sin c+k \), \( k\in\mathbb{R} \)}{\( ab\mathrm{e}^c x-2bx+5^b x-x \sin c+k \), \( k\in\mathbb{R} \)}{}{}}
\LANGcs{
\Question{
Řešte v oboru reálných čísel neurčitý integrál
\[ \int\left(ab\mathrm{e}^c-bx^2+5^b-\sin c \right)\mathrm{d}x, \]
 \( a \), \( b \), \( c \) jsou reálná čísla.}
{\( ab\mathrm{e}^c x-b\frac{x^3}3+5^b x-x \sin c+k \), \( k\in\mathbb{R} \)}{\( -b\frac{x^3}3+k \), \( k\in\mathbb{R} \)}{\( ab\mathrm{e}^c-2b+5^b-\sin c+k \), \( k\in\mathbb{R} \)}{\( ab\mathrm{e}^c x-2bx+5^b x-x \sin c+k \), \( k\in\mathbb{R} \)}{}{}}
\LANGsk{
\Question{
Riešte v obore reálnych čísel neurčitý integrál
\[ \int\left(ab\mathrm{e}^c-bx^2+5^b-\sin c \right)\mathrm{d}x, \]
 \( a \), \( b \), \( c \) sú reálne čísla.}
{\( ab\mathrm{e}^c x-b\frac{x^3}3+5^b x-x \sin c+k \), \( k\in\mathbb{R} \)}{\( -b\frac{x^3}3+k \), \( k\in\mathbb{R} \)}{\( ab\mathrm{e}^c-2b+5^b-\sin c+k \), \( k\in\mathbb{R} \)}{\( ab\mathrm{e}^c x-2bx+5^b x-x \sin c+k \), \( k\in\mathbb{R} \)}{}{}}
\LANGes{
\Question{
Resuelve la integral indefinida 
\[ \int\left(ab\mathrm{e}^c-bx^2+5^b-\sin c \right)\mathrm{d}x \]
en el conjunto de números reales, donde \( a \), \( b \), \( c \) son números reales.}
{\( ab\mathrm{e}^c x-b\frac{x^3}3+5^b x-x \sin c+k \), \( k\in\mathbb{R} \)}{\( -b\frac{x^3}3+k \), \( k\in\mathbb{R} \)}{\( ab\mathrm{e}^c-2b+5^b-\sin c+k \), \( k\in\mathbb{R} \)}{\( ab\mathrm{e}^c x-2bx+5^b x-x \sin c+k \), \( k\in\mathbb{R} \)}{}{}}
\End

\Data\ID{1003107903}
\Updated{2020-07-15 03:07:19}
\Level{C}
\SubArea{Primitive function}
\LANGen{
\Question{
Solve the indefinite integral 
\[ \int\left( ab\mathrm{e}^c-bx^2+5^b-\sin c\right)\mathrm{d}a \]
of a real-valued function, where \( x \), \( b \), \( c \) are real numbers.}
{\( \frac{a^2 b\mathrm{e}^c}2-abx^2+a5^b-a\sin c+k \), \( k\in\mathbb{R} \)}{\( \frac{a^2}2b\mathrm{e}^c+k \), \( k\in\mathbb{R} \)}{\( b\mathrm{e}^c-bx^2+5^b-\sin c+k \), \( k\in\mathbb{R} \)}{\( ab\mathrm{e}^c-b\frac{x^3}3+k \), \( k\in\mathbb{R} \)}{}{}}
\LANGpl{
\Question{
Oblicz całkę nieoznaczoną
\[ \int\left( ab\mathrm{e}^c-bx^2+5^b-\sin c\right)\mathrm{d}a \]
funkcji rzeczywistej, gdzie \( x \), \( b \), \( c \) to liczby rzeczywiste.}
{\( \frac{a^2 b\mathrm{e}^c}2-abx^2+a5^b-a\sin c+k \), \( k\in\mathbb{R} \)}{\( \frac{a^2}2b\mathrm{e}^c+k \), \( k\in\mathbb{R} \)}{\( b\mathrm{e}^c-bx^2+5^b-\sin c+k \), \( k\in\mathbb{R} \)}{\( ab\mathrm{e}^c-b\frac{x^3}3+k \), \( k\in\mathbb{R} \)}{}{}}
\LANGcs{
\Question{
Řešte v oboru reálných čísel neurčitý integrál
\[ \int\left( ab\mathrm{e}^c-bx^2+5^b-\sin c\right)\mathrm{d}a, \]
 \( x \), \( b \), \( c \) jsou reálná čísla.}
{\( \frac{a^2 b\mathrm{e}^c}2-abx^2+a5^b-a\sin c+k \), \( k\in\mathbb{R} \)}{\( \frac{a^2}2b\mathrm{e}^c+k \), \( k\in\mathbb{R} \)}{\( b\mathrm{e}^c-bx^2+5^b-\sin c+k \), \( k\in\mathbb{R} \)}{\( ab\mathrm{e}^c-b\frac{x^3}3+k \), \( k\in\mathbb{R} \)}{}{}}
\LANGsk{
\Question{
Riešte v obore reálnych čísel neurčitý integrál
\[ \int\left( ab\mathrm{e}^c-bx^2+5^b-\sin c\right)\mathrm{d}a, \]
 \( x \), \( b \), \( c \) sú reálne čísla.}
{\( \frac{a^2 b\mathrm{e}^c}2-abx^2+a5^b-a\sin c+k \), \( k\in\mathbb{R} \)}{\( \frac{a^2}2b\mathrm{e}^c+k \), \( k\in\mathbb{R} \)}{\( b\mathrm{e}^c-bx^2+5^b-\sin c+k \), \( k\in\mathbb{R} \)}{\( ab\mathrm{e}^c-b\frac{x^3}3+k \), \( k\in\mathbb{R} \)}{}{}}
\LANGes{
\Question{
Resuelve la integral indefinida  
\[ \int\left( ab\mathrm{e}^c-bx^2+5^b-\sin c\right)\mathrm{d}a \]
en el conjunto de números reales, donde \( x \), \( b \), \( c \) son números reales.}
{\( \frac{a^2 b\mathrm{e}^c}2-abx^2+a5^b-a\sin c+k \), \( k\in\mathbb{R} \)}{\( \frac{a^2}2b\mathrm{e}^c+k \), \( k\in\mathbb{R} \)}{\( b\mathrm{e}^c-bx^2+5^b-\sin c+k \), \( k\in\mathbb{R} \)}{\( ab\mathrm{e}^c-b\frac{x^3}3+k \), \( k\in\mathbb{R} \)}{}{}}
\End

\Data\ID{1003107904}
\Updated{2020-07-15 03:07:19}
\Level{C}
\SubArea{Primitive function}
\LANGen{
\Question{
Solve the indefinite integral 
\[ \int\left(ab\mathrm{e}^c-bx^2+5^b-\sin c\right) \mathrm{d}b \]
of a real-valued function, where \( x \), \( a \), \( c \) are real numbers.}
{\( 0.5a\mathrm{e}^cb^2-\frac{b^2}2 x^2+\frac{5^b}{\ln5} -b\sin c+k \), \( k\in\mathbb{R} \)}{\( 0.5a\mathrm{e}^cb^2-\frac{b^2}2\cdot\frac{x^3}3+\frac{5^b}{\ln5} -b \sin c+k \), \( k\in\mathbb{R} \)}{\( a\mathrm{e}^c-2bx+5^b-\sin c+k \), \( k\in\mathbb{R} \)}{\( a\mathrm{e}^c-bx^2+5^b-b\sin c+k \), \( k\in\mathbb{R} \)}{}{}}
\LANGpl{
\Question{
Oblicz całkę nieoznaczoną
\[ \int\left(ab\mathrm{e}^c-bx^2+5^b-\sin c\right) \mathrm{d}b \]
funkcji rzeczywistej, gdzie \( x \), \( a \), \( c \) to liczby rzeczywiste.}
{\( 0{,}5a\mathrm{e}^cb^2-\frac{b^2}2 x^2+\frac{5^b}{\ln5} -b\sin c+k \), \( k\in\mathbb{R} \)}{\( 0{,}5a\mathrm{e}^cb^2-\frac{b^2}2\cdot\frac{x^3}3+\frac{5^b}{\ln5} -b \sin c+k \), \( k\in\mathbb{R} \)}{\( a\mathrm{e}^c-2bx+5^b-\sin c+k \), \( k\in\mathbb{R} \)}{\( a\mathrm{e}^c-bx^2+5^b-b\sin c+k \), \( k\in\mathbb{R} \)}{}{}}
\LANGcs{
\Question{
Řešte v oboru reálných čísel neurčitý integrál
\[ \int\left(ab\mathrm{e}^c-bx^2+5^b-\sin c\right) \mathrm{d}b, \]
\( x \), \( a \), \( c \) jsou reálná čísla.}
{\( 0{,}5a\mathrm{e}^cb^2-\frac{b^2}2 x^2+\frac{5^b}{\ln5} -b\sin c+k \), \( k\in\mathbb{R} \)}{\( 0{,}5a\mathrm{e}^cb^2-\frac{b^2}2\cdot\frac{x^3}3+\frac{5^b}{\ln5} -b \sin c+k \), \( k\in\mathbb{R} \)}{\( a\mathrm{e}^c-2bx+5^b-\sin c+k \), \( k\in\mathbb{R} \)}{\( a\mathrm{e}^c-bx^2+5^b-b\sin c+k \), \( k\in\mathbb{R} \)}{}{}}
\LANGsk{
\Question{
Riešte v obore reálnych čísel neurčitý integrál
\[ \int\left(ab\mathrm{e}^c-bx^2+5^b-\sin c\right) \mathrm{d}b, \]
\( x \), \( a \), \( c \) sú reálne čísla.}
{\( 0{,}5a\mathrm{e}^cb^2-\frac{b^2}2 x^2+\frac{5^b}{\ln5} -b\sin c+k \), \( k\in\mathbb{R} \)}{\( 0{,}5a\mathrm{e}^cb^2-\frac{b^2}2\cdot\frac{x^3}3+\frac{5^b}{\ln5} -b \sin c+k \), \( k\in\mathbb{R} \)}{\( a\mathrm{e}^c-2bx+5^b-\sin c+k \), \( k\in\mathbb{R} \)}{\( a\mathrm{e}^c-bx^2+5^b-b\sin c+k \), \( k\in\mathbb{R} \)}{}{}}
\LANGes{
\Question{
Resuelve la integral indefinida  
\[ \int\left(ab\mathrm{e}^c-bx^2+5^b-\sin c\right) \mathrm{d}b \]
en el conjunto real, donde \( x \), \( a \), \( c \) son números reales.}
{\( 0.5a\mathrm{e}^cb^2-\frac{b^2}2 x^2+\frac{5^b}{\ln5} -b\sin c+k \), \( k\in\mathbb{R} \)}{\( 0.5a\mathrm{e}^cb^2-\frac{b^2}2\cdot\frac{x^3}3+\frac{5^b}{\ln5} -b \sin c+k \), \( k\in\mathbb{R} \)}{\( a\mathrm{e}^c-2bx+5^b-\sin c+k \), \( k\in\mathbb{R} \)}{\( a\mathrm{e}^c-bx^2+5^b-b\sin c+k \), \( k\in\mathbb{R} \)}{}{}}
\End

\Data\ID{1003107905}
\Updated{2020-07-15 03:07:19}
\Level{C}
\SubArea{Primitive function}
\LANGen{
\Question{
Solve the indefinite integral 
\[ \int\left(ab\mathrm{e}^c-bx^2+5^b-\sin c\right) \mathrm{d}c \]
of a real-valued function, where \( a \), \( b \), \( x \) are real numbers.}
{\( ab\mathrm{e}^c-bcx^2+c\cdot5^b+\cos c+k \), \( k\in\mathbb{R} \)}{\( \mathrm{e}^c-\cos c+k \), \( k\in\mathbb{R} \)}{\( ab\mathrm{e}^c-bcx^2+c\cdot5^b-\cos c+k \), \( k\in\mathbb{R} \)}{\( ab\mathrm{e}^c-bcx^2+5^b\cdot c-\sin c+k \), \( k\in\mathbb{R} \)}{}{}}
\LANGpl{
\Question{
Oblicz całkę nieznaczoną
\[ \int\left(ab\mathrm{e}^c-bx^2+5^b-\sin c\right) \mathrm{d}c \]
funkcji rzeczywistej, gdzie \( a \), \( b \), \( x \) to liczby rzeczywiste.}
{\( ab\mathrm{e}^c-bcx^2+c\cdot5^b+\cos c+k \), \( k\in\mathbb{R} \)}{\( \mathrm{e}^c-\cos c+k \), \( k\in\mathbb{R} \)}{\( ab\mathrm{e}^c-bcx^2+c\cdot5^b-\cos c+k \), \( k\in\mathbb{R} \)}{\( ab\mathrm{e}^c-bcx^2+5^b\cdot c-\sin c+k \), \( k\in\mathbb{R} \)}{}{}}
\LANGcs{
\Question{
Řešte v oboru reálných čísel
\[ \int\left(ab\mathrm{e}^c-bx^2+5^b-\sin c\right) \mathrm{d}c, \]
\( a \), \( b \), \( x \) jsou reálná čísla.}
{\( ab\mathrm{e}^c-bcx^2+c\cdot5^b+\cos c+k \), \( k\in\mathbb{R} \)}{\( \mathrm{e}^c-\cos c+k \), \( k\in\mathbb{R} \)}{\( ab\mathrm{e}^c-bcx^2+c\cdot5^b-\cos c+k \), \( k\in\mathbb{R} \)}{\( ab\mathrm{e}^c-bcx^2+5^b\cdot c-\sin c+k \), \( k\in\mathbb{R} \)}{}{}}
\LANGsk{
\Question{
Riešte v obore reálnych čísel
\[ \int\left(ab\mathrm{e}^c-bx^2+5^b-\sin c\right) \mathrm{d}c, \]
\( a \), \( b \), \( x \) sú reálne čísla.}
{\( ab\mathrm{e}^c-bcx^2+c\cdot5^b+\cos c+k \), \( k\in\mathbb{R} \)}{\( \mathrm{e}^c-\cos c+k \), \( k\in\mathbb{R} \)}{\( ab\mathrm{e}^c-bcx^2+c\cdot5^b-\cos c+k \), \( k\in\mathbb{R} \)}{\( ab\mathrm{e}^c-bcx^2+5^b\cdot c-\sin c+k \), \( k\in\mathbb{R} \)}{}{}}
\LANGes{
\Question{
Resuelve la integral indefinida  
\[ \int\left(ab\mathrm{e}^c-bx^2+5^b-\sin c\right) \mathrm{d}c \]
en el conjunto de números reales, donde \( a \), \( b \), \( x \) son números reales.}
{\( ab\mathrm{e}^c-bcx^2+c\cdot5^b+\cos c+k \), \( k\in\mathbb{R} \)}{\( \mathrm{e}^c-\cos c+k \), \( k\in\mathbb{R} \)}{\( ab\mathrm{e}^c-bcx^2+c\cdot5^b-\cos c+k \), \( k\in\mathbb{R} \)}{\( ab\mathrm{e}^c-bcx^2+5^b\cdot c-\sin c+k \), \( k\in\mathbb{R} \)}{}{}}
\End

\Data\ID{1003107906}
\Updated{2020-07-15 03:07:19}
\Level{C}
\SubArea{Primitive function}
\LANGen{
\Question{
Solve the indefinite integral 
\[\int\frac{\mathrm{d}x}{\mathrm{cotg}\,x\cdot\sin2x} \]
in the range \( \left(0;\frac{\pi}2\right) \).}
{\( \frac12\mathrm{tg}\,x+c \), \( c\in\mathbb{R} \)}{\( \frac12\mathrm{cotg}\,x+c \), \( c\in\mathbb{R} \)}{\( \ln|\sin2x |+c \), \( c\in\mathbb{R} \)}{\( -\frac12\mathrm{tg}\,x+c \), \( c\in\mathbb{R} \)}{}{}}
\LANGpl{
\Question{
Oblicz całkę nieznaczoną
\[\int\frac{\mathrm{d}x}{\mathrm{cotg}\,x\cdot\sin2x} \]
w  \( \left(0;\frac{\pi}2\right) \).}
{\( \frac12\mathrm{tg}\,x+c \), \( c\in\mathbb{R} \)}{\( \frac12\mathrm{cotg}\,x+c \), \( c\in\mathbb{R} \)}{\( \ln|\sin2x |+c \), \( c\in\mathbb{R} \)}{\( -\frac12\mathrm{tg}\,x+c \), \( c\in\mathbb{R} \)}{}{}}
\LANGcs{
\Question{
Řešte neurčitý integrál
\[\int\frac{\mathrm{d}x}{\mathrm{cotg}\,x\cdot\sin2x} \]
na intervalu \( \left(0;\frac{\pi}2\right) \).}
{\( \frac12\mathrm{tg}\,x+c \), \( c\in\mathbb{R} \)}{\( \frac12\mathrm{cotg}\,x+c \), \( c\in\mathbb{R} \)}{\( \ln|\sin2x |+c \), \( c\in\mathbb{R} \)}{\( -\frac12\mathrm{tg}\,x+c \), \( c\in\mathbb{R} \)}{}{}}
\LANGsk{
\Question{
Riešte neurčitý integrál
\[\int\frac{\mathrm{d}x}{\mathrm{cotg}\,x\cdot\sin2x} \]
na intervale \( \left(0;\frac{\pi}2\right) \).}
{\( \frac12\mathrm{tg}\,x+c \), \( c\in\mathbb{R} \)}{\( \frac12\mathrm{cotg}\,x+c \), \( c\in\mathbb{R} \)}{\( \ln|\sin2x |+c \), \( c\in\mathbb{R} \)}{\( -\frac12\mathrm{tg}\,x+c \), \( c\in\mathbb{R} \)}{}{}}
\LANGes{
\Question{
Resuelve la integral indefinida 
\[\int\frac{\mathrm{d}x}{\mathrm{cotg}\,x\cdot\sin2x} \]
en el rango \( \left(0;\frac{\pi}2\right) \).}
{\( \frac12\mathrm{tg}\,x+c \), \( c\in\mathbb{R} \)}{\( \frac12\mathrm{cotg}\,x+c \), \( c\in\mathbb{R} \)}{\( \ln|\sin2x |+c \), \( c\in\mathbb{R} \)}{\( -\frac12\mathrm{tg}\,x+c \), \( c\in\mathbb{R} \)}{}{}}
\End

\Data\ID{1003107907}
\Updated{2020-07-15 03:07:19}
\Level{C}
\SubArea{Primitive function}
\LANGen{
\Question{
Solve the indefinite integral 
\[ \int\mathrm{cotg}^2x\,\mathrm{d}x \]
in the range \( (\pi; 2\pi) \).}
{\( -\mathrm{cotg}\,x-x+c \), \( c\in\mathbb{R} \)}{\( \mathrm{cotg}\,x-x+c \), \( c\in\mathbb{R} \)}{\( \mathrm{tg}\,x-x+c \), \( c\in\mathbb{R} \)}{\( -\mathrm{tg}^2x+c \), \( c\in\mathbb{R} \)}{}{}}
\LANGpl{
\Question{
Oblicz całkę nieznaczoną
\[ \int\mathrm{cotg}^2x\,\mathrm{d}x \]
w przedziale \( (\pi; 2\pi) \).}
{\( -\mathrm{cotg}\,x-x+c \), \( c\in\mathbb{R} \)}{\( \mathrm{cotg}\,x-x+c \), \( c\in\mathbb{R} \)}{\( \mathrm{tg}\,x-x+c \), \( c\in\mathbb{R} \)}{\( -\mathrm{tg}^2x+c \), \( c\in\mathbb{R} \)}{}{}}
\LANGcs{
\Question{
Řešte neurčitý integrál
\[ \int\mathrm{cotg}^2x\,\mathrm{d}x \]
na intervalu \( (\pi; 2\pi) \).}
{\( -\mathrm{cotg}\,x-x+c \), \( c\in\mathbb{R} \)}{\( \mathrm{cotg}\,x-x+c \), \( c\in\mathbb{R} \)}{\( \mathrm{tg}\,x-x+c \), \( c\in\mathbb{R} \)}{\( -\mathrm{tg}^2x+c \), \( c\in\mathbb{R} \)}{}{}}
\LANGsk{
\Question{
Riešte neurčitý integrál
\[ \int\mathrm{cotg}^2x\,\mathrm{d}x \]
na intervale \( (\pi; 2\pi) \).}
{\( -\mathrm{cotg}\,x-x+c \), \( c\in\mathbb{R} \)}{\( \mathrm{cotg}\,x-x+c \), \( c\in\mathbb{R} \)}{\( \mathrm{tg}\,x-x+c \), \( c\in\mathbb{R} \)}{\( -\mathrm{tg}^2x+c \), \( c\in\mathbb{R} \)}{}{}}
\LANGes{
\Question{
Resuelve la integral indefinida 
\[ \int\mathrm{cotg}^2x\,\mathrm{d}x \]
en el rango \( (\pi; 2\pi) \).}
{\( -\mathrm{cotg}\,x-x+c \), \( c\in\mathbb{R} \)}{\( \mathrm{cotg}\,x-x+c \), \( c\in\mathbb{R} \)}{\( \mathrm{tg}\,x-x+c \), \( c\in\mathbb{R} \)}{\( -\mathrm{tg}^2x+c \), \( c\in\mathbb{R} \)}{}{}}
\End

\Data\ID{1003107908}
\Updated{2020-07-15 03:07:19}
\Level{C}
\SubArea{Primitive function}
\LANGen{
\Question{
Solve the indefinite integral
\[ \int\frac{\cos2x}{\sin^2x}\mathrm{d}x \]
in the range \( \left(\pi; \frac32\pi\right) \).}
{\( -\mathrm{cotg}\,x-2x+c \), \( c\in\mathbb{R} \)}{\( \mathrm{cotg}\,x-2x+c \), \( c\in\mathbb{R} \)}{\( -\mathrm{tg}\,x-2x+c \), \( c\in\mathbb{R} \)}{\( -\mathrm{cotg}\,x+c \), \( c\in\mathbb{R} \)}{}{}}
\LANGpl{
\Question{
Oblicz całkę nieznaczoną
\[ \int\frac{\cos2x}{\sin^2x}\mathrm{d}x \]
w przedziale \( \left(\pi; \frac32\pi\right) \).}
{\( -\mathrm{cotg}\,x-2x+c \), \( c\in\mathbb{R} \)}{\( \mathrm{cotg}\,x-2x+c \), \( c\in\mathbb{R} \)}{\( -\mathrm{tg}\,x-2x+c \), \( c\in\mathbb{R} \)}{\( -\mathrm{cotg}\,x+c \), \( c\in\mathbb{R} \)}{}{}}
\LANGcs{
\Question{
Řešte neurčitý integrál
\[ \int\frac{\cos2x}{\sin^2x}\mathrm{d}x \]
na intervalu \( \left(\pi; \frac32\pi\right) \).}
{\( -\mathrm{cotg}\,x-2x+c \), \( c\in\mathbb{R} \)}{\( \mathrm{cotg}\,x-2x+c \), \( c\in\mathbb{R} \)}{\( -\mathrm{tg}\,x-2x+c \), \( c\in\mathbb{R} \)}{\( -\mathrm{cotg}\,x+c \), \( c\in\mathbb{R} \)}{}{}}
\LANGsk{
\Question{
Riešte neurčitý integrál
\[ \int\frac{\cos2x}{\sin^2x}\mathrm{d}x \]
na intervale \( \left(\pi; \frac32\pi\right) \).}
{\( -\mathrm{cotg}\,x-2x+c \), \( c\in\mathbb{R} \)}{\( \mathrm{cotg}\,x-2x+c \), \( c\in\mathbb{R} \)}{\( -\mathrm{tg}\,x-2x+c \), \( c\in\mathbb{R} \)}{\( -\mathrm{cotg}\,x+c \), \( c\in\mathbb{R} \)}{}{}}
\LANGes{
\Question{
Resuelve la integral indefinida 
\[ \int\frac{\cos2x}{\sin^2x}\mathrm{d}x \]
en el rango \( \left(\pi; \frac32\pi\right) \).}
{\( -\mathrm{cotg}\,x-2x+c \), \( c\in\mathbb{R} \)}{\( \mathrm{cotg}\,x-2x+c \), \( c\in\mathbb{R} \)}{\( -\mathrm{tg}\,x-2x+c \), \( c\in\mathbb{R} \)}{\( -\mathrm{cotg}\,x+c \), \( c\in\mathbb{R} \)}{}{}}
\End

\Data\ID{1003107909}
\Updated{2020-07-15 03:07:19}
\Level{C}
\SubArea{Primitive function}
\LANGen{
\Question{
Solve the indefinite integral 
\[ \int\frac{\sqrt[5]{\ln x}}x\mathrm{d}x \]
in the range \( (0;\infty) \).}
{\( \frac56\ln x\cdot\sqrt[5]{\ln x}+c \), \( c\in\mathbb{R} \)}{\( \frac56\ln x\cdot\sqrt[6]{\ln x}+c \), \( c\in\mathbb{R} \)}{\( \frac65 \ln x\cdot\sqrt[5]{\ln x}+c \), \( c\in\mathbb{R} \)}{\( \frac65\ln x\cdot\sqrt[6]{\ln x}+c \), \( c\in\mathbb{R} \)}{}{}}
\LANGpl{
\Question{
Oblicz całkę nieznaczoną
\[ \int\frac{\sqrt[5]{\ln x}}x\mathrm{d}x \]
w przedziale \( (0;\infty) \).}
{\( \frac56\ln x\cdot\sqrt[5]{\ln x}+c \), \( c\in\mathbb{R} \)}{\( \frac56\ln x\cdot\sqrt[6]{\ln x}+c \), \( c\in\mathbb{R} \)}{\( \frac65 \ln x\cdot\sqrt[5]{\ln x}+c \), \( c\in\mathbb{R} \)}{\( \frac65\ln x\cdot\sqrt[6]{\ln x}+c \), \( c\in\mathbb{R} \)}{}{}}
\LANGcs{
\Question{
Řešte neurčitý integrál
\[ \int\frac{\sqrt[5]{\ln x}}x\mathrm{d}x \]
na intervalu \( (0;\infty) \).}
{\( \frac56\ln x\cdot\sqrt[5]{\ln x}+c \), \( c\in\mathbb{R} \)}{\( \frac56\ln x\cdot\sqrt[6]{\ln x}+c \), \( c\in\mathbb{R} \)}{\( \frac65 \ln x\cdot\sqrt[5]{\ln x}+c \), \( c\in\mathbb{R} \)}{\( \frac65\ln x\cdot\sqrt[6]{\ln x}+c \), \( c\in\mathbb{R} \)}{}{}}
\LANGsk{
\Question{
Riešte neurčitý integrál
\[ \int\frac{\sqrt[5]{\ln x}}x\mathrm{d}x \]
na intervale \( (0;\infty) \).}
{\( \frac56\ln x\cdot\sqrt[5]{\ln x}+c \), \( c\in\mathbb{R} \)}{\( \frac56\ln x\cdot\sqrt[6]{\ln x}+c \), \( c\in\mathbb{R} \)}{\( \frac65 \ln x\cdot\sqrt[5]{\ln x}+c \), \( c\in\mathbb{R} \)}{\( \frac65\ln x\cdot\sqrt[6]{\ln x}+c \), \( c\in\mathbb{R} \)}{}{}}
\LANGes{
\Question{
Resuelve la integral indefinida 
\[ \int\frac{\sqrt[5]{\ln x}}x\mathrm{d}x \]
en el intervalo \( (0;\infty) \).}
{\( \frac56\ln x\cdot\sqrt[5]{\ln x}+c \), \( c\in\mathbb{R} \)}{\( \frac56\ln x\cdot\sqrt[6]{\ln x}+c \), \( c\in\mathbb{R} \)}{\( \frac65 \ln x\cdot\sqrt[5]{\ln x}+c \), \( c\in\mathbb{R} \)}{\( \frac65\ln x\cdot\sqrt[6]{\ln x}+c \), \( c\in\mathbb{R} \)}{}{}}
\End

\Data\ID{1003107910}
\Updated{2022-04-23 05:04:02}
\Level{C}
\SubArea{Primitive function}
\LANGen{
\Question{
Solve the indefinite integral
\[  \int\mathrm{e}^{\sin x}\cos x\,\mathrm{d}x \]
of a real-valued function.}
{\( \mathrm{e}^{\sin x} +c \), \( c\in\mathbb{R} \)}{\( -\sin x\cdot\mathrm{e}^{\sin x} +c \), \( c\in\mathbb{R} \)}{\( \mathrm{e}^{\cos x} +c \), \( c\in\mathbb{R} \)}{\( \mathrm{e}^{\sin x}\cdot\cos x+c \), \( c\in\mathbb{R} \)}{}{}}
\LANGpl{
\Question{
Oblicz całkę nieznaczoną
\[  \int\mathrm{e}^{\sin x}\cos x\,\mathrm{d}x \]
funkcji rzeczywistej.}
{\( \mathrm{e}^{\sin x} +c \), \( c\in\mathbb{R} \)}{\( -\sin x\cdot\mathrm{e}^{\sin x} +c \), \( c\in\mathbb{R} \)}{\( \mathrm{e}^{\cos x} +c \), \( c\in\mathbb{R} \)}{\( \mathrm{e}^{\sin x}\cdot\cos x+c \), \( c\in\mathbb{R} \)}{}{}}
\LANGcs{
\Question{
Řešte neurčitý integrál
\[  \int\mathrm{e}^{\sin x}\cos x\,\mathrm{d}x \]
v oboru reálných čísel.}
{\( \mathrm{e}^{\sin x} +c \), \( c\in\mathbb{R} \)}{\( -\sin x\cdot\mathrm{e}^{\sin x} +c \), \( c\in\mathbb{R} \)}{\( \mathrm{e}^{\cos x} +c \), \( c\in\mathbb{R} \)}{\( \mathrm{e}^{\sin x}\cdot\cos x+c \), \( c\in\mathbb{R} \)}{}{}}
\LANGsk{
\Question{
Riešte neurčitý integrál
\[  \int\mathrm{e}^{\sin x}\cos x\,\mathrm{d}x \]
v obore reálnych čísel.}
{\( \mathrm{e}^{\sin x} +c \), \( c\in\mathbb{R} \)}{\( -\sin x\cdot\mathrm{e}^{\sin x} +c \), \( c\in\mathbb{R} \)}{\( \mathrm{e}^{\cos x} +c \), \( c\in\mathbb{R} \)}{\( \mathrm{e}^{\sin x}\cdot\cos x+c \), \( c\in\mathbb{R} \)}{}{}}
\LANGes{
\Question{
Resuelve la integral indefinida 
\[  \int\mathrm{e}^{\sin x}\cos x\,\mathrm{d}x \]
en el conjunto de números reales.}
{\( \mathrm{e}^{\sin x} +c \), \( c\in\mathbb{R} \)}{\( -\sin x\cdot\mathrm{e}^{\sin x} +c \), \( c\in\mathbb{R} \)}{\( \mathrm{e}^{\cos x} +c \), \( c\in\mathbb{R} \)}{\( \mathrm{e}^{\sin x}\cdot\cos x+c \), \( c\in\mathbb{R} \)}{}{}}
\End

\Data\ID{1003107911}
\Updated{2020-07-15 03:07:19}
\Level{C}
\SubArea{Primitive function}
\LANGen{
\Question{
Solve the indefinite integral
\[ \int\sin\sqrt x\,\mathrm{d}x \]
in the range \( (0;\infty) \).}
{\( -2\sqrt x\cos\sqrt x+2\sin\sqrt x+c \), \( c\in\mathbb{R} \)}{\( 2\sqrt x\cos\sqrt x+2\sin\sqrt x+c \), \( c\in\mathbb{R} \)}{\( 2\sqrt x\cos\sqrt x-2\sin\sqrt x+c \), \( c\in\mathbb{R} \)}{\( -\cos\sqrt x \), \( c\in\mathbb{R} \)}{}{}}
\LANGpl{
\Question{
Oblicz całkę nieznaczoną
\[ \int\sin\sqrt x\,\mathrm{d}x \]
w przedziale \( (0;\infty) \).}
{\( -2\sqrt x\cos\sqrt x+2\sin\sqrt x+c \), \( c\in\mathbb{R} \)}{\( 2\sqrt x\cos\sqrt x+2\sin\sqrt x+c \), \( c\in\mathbb{R} \)}{\( 2\sqrt x\cos\sqrt x-2\sin\sqrt x+c \), \( c\in\mathbb{R} \)}{\( -\cos\sqrt x \), \( c\in\mathbb{R} \)}{}{}}
\LANGcs{
\Question{
Řešte neurčitý integrál
\[ \int\sin\sqrt x\,\mathrm{d}x \]
na intervalu \( (0;\infty) \).}
{\( -2\sqrt x\cos\sqrt x+2\sin\sqrt x+c \), \( c\in\mathbb{R} \)}{\( 2\sqrt x\cos\sqrt x+2\sin\sqrt x+c \), \( c\in\mathbb{R} \)}{\( 2\sqrt x\cos\sqrt x-2\sin\sqrt x+c \), \( c\in\mathbb{R} \)}{\( -\cos\sqrt x \), \( c\in\mathbb{R} \)}{}{}}
\LANGsk{
\Question{
Riešte neurčitý integrál
\[ \int\sin\sqrt x\,\mathrm{d}x \]
na intervale \( (0;\infty) \).}
{\( -2\sqrt x\cos\sqrt x+2\sin\sqrt x+c \), \( c\in\mathbb{R} \)}{\( 2\sqrt x\cos\sqrt x+2\sin\sqrt x+c \), \( c\in\mathbb{R} \)}{\( 2\sqrt x\cos\sqrt x-2\sin\sqrt x+c \), \( c\in\mathbb{R} \)}{\( -\cos\sqrt x \), \( c\in\mathbb{R} \)}{}{}}
\LANGes{
\Question{
Resuelve la integral indefinida
\[ \int\sin\sqrt x\,\mathrm{d}x \]
en el rango \( (0;\infty) \).}
{\( -2\sqrt x\cos\sqrt x+2\sin\sqrt x+c \), \( c\in\mathbb{R} \)}{\( 2\sqrt x\cos\sqrt x+2\sin\sqrt x+c \), \( c\in\mathbb{R} \)}{\( 2\sqrt x\cos\sqrt x-2\sin\sqrt x+c \), \( c\in\mathbb{R} \)}{\( -\cos\sqrt x \), \( c\in\mathbb{R} \)}{}{}}
\End

\Data\ID{1003107912}
\Updated{2020-07-15 03:07:19}
\Level{C}
\SubArea{Primitive function}
\LANGen{
\Question{
Which method is the most effective to solve the indefinite integral 
\[ \int\frac{\mathrm{d}x}{x\ln x} \]
in the range \( (1;\infty) \)?}
{By substitution, \( a=\ln x \).}{By parts integration, when we let \( u(x)=\frac1x \), where \( u(x) \) is the integrated function, and we let \( v'(x)=\ln x \), where \( v'(x) \) is the differentiated function.}{By substitution, \( a=\frac1x \).}{By factorization into \( \int\frac1x\mathrm{d}x\cdot\int\frac1{\ln x}\mathrm{d}x \).}{}{}}
\LANGpl{
\Question{
Która z podanych metod jest najbardziej efektywna w obliczaniu całki nieoznaczonej
\[ \int\frac{\mathrm{d}x}{x\ln x} \]
w przedziale \( (1;\infty) \)?}
{Przez podstawienie, \( a=\ln x \).}{Przez całki częściowe,  \( u(x)=\frac1x \), gdzie \( u(x) \) jest funkcją układu scalonego, a \( v'(x)=\ln x \), gdzie \( v'(x) \) jest funkcją różniczkowalną.}{Przez podstawienie, \( a=\frac1x \).}{Rozkład na czynniki  \( \int\frac1x\mathrm{d}x\cdot\int\frac1{\ln x}\mathrm{d}x \).}{}{}}
\LANGcs{
\Question{
Kterou metodou lze nejvýhodněji řešit neurčitý integrál
\[ \int\frac{\mathrm{d}x}{x\ln x} \]
na intervalu \( (1;\infty) \)?}
{Substitucí \( a=\ln x \).}{Per partes, jako nederivovanou funkci volíme \( u(x)=\frac1x \), jako derivovanou funkci volíme \( v'(x)=\ln x \).}{Substitucí \( a=\frac1x \).}{Rozložením na součin \( \int\frac1x\mathrm{d}x\cdot\int\frac1{\ln x}\mathrm{d}x \).}{}{}}
\LANGsk{
\Question{
Ktorou metódou je najvýhodnejšie riešiť neurčitý integrál
\[ \int\frac{\mathrm{d}x}{x\ln x} \]
na intervale \( (1;\infty) \)?}
{Substitúciou \( a=\ln x \).}{Per partes, ako nederivovanú funkciu volíme \( u(x)=\frac1x \), ako derivovanú funkciu volíme \( v'(x)=\ln x \).}{Substitúciou \( a=\frac1x \).}{Rozložením na súčin \( \int\frac1x\mathrm{d}x\cdot\int\frac1{\ln x}\mathrm{d}x \).}{}{}}
\LANGes{
\Question{
¿Cuál de los métodos es el más adecuado para resolver la integral indefinida
\[ \int\frac{\mathrm{d}x}{x\ln x} \]
en el rango \( (1;\infty) \)?}
{Por sustitución, \( a=\ln x \).}{Por integración por partes donde \( u(x)=\frac1x \) es la función integrada, y donde \( v'(x)=\ln x \) es la función derivada.}{Por sustitución, \( a=\frac1x \).}{Por factorización de \( \int\frac1x\mathrm{d}x\cdot\int\frac1{\ln x}\mathrm{d}x \).}{}{}}
\End

\Data\ID{1003107913}
\Updated{2020-07-15 03:07:19}
\Level{C}
\SubArea{Primitive function}
\LANGen{
\Question{
Which method is the most effective to solve the indefinite integral 
\[ \int\sin(\ln x)\mathrm{d}x \]
in the range \( (0;\infty) \)?}
{By parts integration, when we let \( u(x)=\sin(\ln x) \), where \( u(x) \) is the integrated function, and we let \( v'(x)=1 \), where \( v'(x) \) is the differentiated function.}{By substitution, \( a=\sin x \).}{By parts integration, when we let \( u(x)=\ln x \), where \( u(x) \) is the integrated function, and we let \( v'(x)=\sin x \), where \( v'(x) \) is the differentiated function.}{By substitution, \( t=\sin(\ln x) \).}{}{}}
\LANGpl{
\Question{
Która z podanych metod jest najbardziej efektywna w obliczaniu całki nieoznaczonej 
\[ \int\sin(\ln x)\mathrm{d}x \]
w przedziale \( (0;\infty) \)?}
{Przez całki częściowe,  \( u(x)=\sin(\ln x) \), gdzie \( u(x) \) jest funkcją układu scalonego, a \( v'(x)=1 \), gdzie \( v'(x) \) jest funkcją różniczkowalną.}{Przez podstawienie, \( a=\sin x \).}{Przez całki częściowe, \( u(x)=\ln x \), gdzie \( u(x) \) jest funkcją układu scalonego, a \( v'(x)=\sin x \), gdzie \( v'(x) \) jest funkcją różniczkowalną.}{Przez podstawienie, \( t=\sin(\ln x) \).}{}{}}
\LANGcs{
\Question{
Kterou metodou lze nejvýhodněji řešit neurčitý integrál
\[ \int\sin(\ln x)\mathrm{d}x \]
na intervalu \( (0;\infty) \)?}
{Per partes, jako nederivovanou funkci volíme \( u(x)=\sin(\ln x) \), jako derivovanou funkci volíme \( v'(x)=1 \).}{Substitucí \( a=\sin x \).}{Per partes, jako nederivovanou funkci volíme \( u(x)=\ln x \), jako derivovanou funkci volíme \( v'(x)=\sin x \).}{Substitucí \( t=\sin(\ln x) \).}{}{}}
\LANGsk{
\Question{
Ktorou metódou je najvýhodnejšie riešiť neurčitý integrál
\[ \int\sin(\ln x)\mathrm{d}x \]
na intervale \( (0;\infty) \)?}
{Per partes, ako nederivovanú funkciu volíme \( u(x)=\sin(\ln x) \), ako derivovanú funkciu volíme \( v'(x)=1 \).}{Substitúciou \( a=\sin x \).}{Per partes, ako nederivovanú funkciu volíme \( u(x)=\ln x \), ako derivovanú funkciu volíme \( v'(x)=\sin x \).}{Substitúciou \( t=\sin(\ln x) \).}{}{}}
\LANGes{
\Question{
¿Cuál de los métodos es el más adecuado para resolver la integral indefinida  
\[ \int\sin(\ln x)\mathrm{d}x \]
en el rango \( (0;\infty) \)?}
{Por integración por partes donde \( u(x)=\sin(\ln x) \) es la función integrada, y donde \( v'(x)=1 \) es la función derivada.}{Por sustitución, \( a=\sin x \).}{Por integración por partes donde \( u(x)=\ln x \) es la función integrada, y donde \( v'(x)=\sin x \) es la función derivada.}{Por sustitución, \( t=\sin(\ln x) \).}{}{}}
\End

\Data\ID{2010000304}
\Updated{2022-08-25 10:08:55}
\Level{C}
\SubArea{Primitive function}
\LANGen{
\Question{
Solve the indefinite integral
\[  \int\mathrm{e}^{\cos x}\sin x\,\mathrm{d}x \]
of a real-valued function.}
{\( -\mathrm{e}^{\cos x} +c \), \( c\in\mathbb{R} \)}{\(- \mathrm{e}^{\cos x}\cdot\cos x+c \), \( c\in\mathbb{R} \)}{\( \mathrm{e}^{\sin x}\cdot\cos x+c \), \( c\in\mathbb{R} \)}{\( \mathrm{e}^{\cos x}\cdot\sin x+c \), \( c\in\mathbb{R} \)}{}{}}
\LANGpl{
\Question{
Oblicz całkę nieoznaczoną 
\[  \int\mathrm{e}^{\cos x}\sin x\,\mathrm{d}x \]
funkcji o wartościach rzeczywistych.}
{\( -\mathrm{e}^{\cos x} +c \), \( c\in\mathbb{R} \)}{\(- \mathrm{e}^{\cos x}\cdot\cos x+c \), \( c\in\mathbb{R} \)}{\( \mathrm{e}^{\sin x}\cdot\cos x+c \), \( c\in\mathbb{R} \)}{\( \mathrm{e}^{\cos x}\cdot\sin x+c \), \( c\in\mathbb{R} \)}{}{}}
\LANGcs{
\Question{
Řešte neurčitý integrál 
\[  \int\mathrm{e}^{\cos x}\cdot\sin x\,\mathrm{d}x \]
v oboru reálných čísel.}
{\( -\mathrm{e}^{\cos x} +c \), \( c\in\mathbb{R} \)}{\(- \mathrm{e}^{\cos x}\cdot\cos x+c \), \( c\in\mathbb{R} \)}{\( \mathrm{e}^{\sin x}\cdot\cos x+c \), \( c\in\mathbb{R} \)}{\( \mathrm{e}^{\cos x}\cdot\sin x+c \), \( c\in\mathbb{R} \)}{}{}}
\LANGsk{
\Question{
Riešte neurčitý integrál
\[  \int\mathrm{e}^{\cos x}\sin x\,\mathrm{d}x \]
v obore reálných čísel.}
{\( -\mathrm{e}^{\cos x} +c \), \( c\in\mathbb{R} \)}{\(- \mathrm{e}^{\cos x}\cdot\cos x+c \), \( c\in\mathbb{R} \)}{\( \mathrm{e}^{\sin x}\cdot\cos x+c \), \( c\in\mathbb{R} \)}{\( \mathrm{e}^{\cos x}\cdot\sin x+c \), \( c\in\mathbb{R} \)}{}{}}
\LANGes{
\Question{
Resuelve la integral indefinida 
\[  \int\mathrm{e}^{\cos x}\sin x\,\mathrm{d}x \]
en el conjunto de números reales.}
{\( -\mathrm{e}^{\cos x} +c \), \( c\in\mathbb{R} \)}{\(- \mathrm{e}^{\cos x}\cdot\cos x+c \), \( c\in\mathbb{R} \)}{\( \mathrm{e}^{\sin x}\cdot\cos x+c \), \( c\in\mathbb{R} \)}{\( \mathrm{e}^{\cos x}\cdot\sin x+c \), \( c\in\mathbb{R} \)}{}{}}
\End

\Data\ID{2010000305}
\Updated{2022-08-25 10:08:55}
\Level{C}
\SubArea{Primitive function}
\LANGen{
\Question{
Evaluate the following integral on the interval \((0;+\infty)\). 
\[ 
 \int \log_2 x\, \mathrm{d}x
\]}
{\(x\log_2x -\frac{x} {\ln 2}+ c,\ c\in \mathbb{R}\)}{\(\log_2 x -\frac{x} {\ln 2}+ c,\ c\in \mathbb{R}\)}{\(x\log_2 x -x+ c,\ c\in \mathbb{R}\)}{\(x\log_2 x +\frac{x} {\ln 2}+ c,\ c\in \mathbb{R}\)}{}{}}
\LANGpl{
\Question{
Oblicz następującą całkę na przedziale  \((0;+\infty)\). 
\[ 
 \int \log_2 x\, \mathrm{d}x
\]}
{\(x\log_2x -\frac{x} {\ln 2}+ c,\ c\in \mathbb{R}\)}{\(\log_2 x -\frac{x} {\ln 2}+ c,\ c\in \mathbb{R}\)}{\(x\log_2 x -x+ c,\ c\in \mathbb{R}\)}{\(x\log_2 x +\frac{x} {\ln 2}+ c,\ c\in \mathbb{R}\)}{}{}}
\LANGcs{
\Question{
Vypočtěte 
\[ 
 \int \log_2 x\, \mathrm{d}x
\] na intervalu \((0;+\infty)\).}
{\(x\log_2x -\frac{x} {\ln 2}+ c,\ c\in \mathbb{R}\)}{\(\log_2 x -\frac{x} {\ln 2}+ c,\ c\in \mathbb{R}\)}{\(x\log_2 x -x+ c,\ c\in \mathbb{R}\)}{\(x\log_2 x +\frac{x} {\ln 2}+ c,\ c\in \mathbb{R}\)}{}{}}
\LANGsk{
\Question{
Vypočítajte 
\[ 
 \int \log_2 x\, \mathrm{d}x
\]
na intervale \((0;+\infty)\).}
{\(x\log_2x -\frac{x} {\ln 2}+ c,\ c\in \mathbb{R}\)}{\(\log_2 x -\frac{x} {\ln 2}+ c,\ c\in \mathbb{R}\)}{\(x\log_2 x -x+ c,\ c\in \mathbb{R}\)}{\(x\log_2 x +\frac{x} {\ln 2}+ c,\ c\in \mathbb{R}\)}{}{}}
\LANGes{
\Question{
Evalúa la siguiente integral en el intervalo \((0;+\infty)\). 
\[ 
 \int \log_2 x\, \mathrm{d}x
\]}
{\(x\log_2x -\frac{x} {\ln 2}+ c,\ c\in \mathbb{R}\)}{\(\log_2 x -\frac{x} {\ln 2}+ c,\ c\in \mathbb{R}\)}{\(x\log_2 x -x+ c,\ c\in \mathbb{R}\)}{\(x\log_2 x +\frac{x} {\ln 2}+ c,\ c\in \mathbb{R}\)}{}{}}
\End

\Data\ID{2010000306}
\Updated{2022-08-25 10:08:55}
\Level{C}
\SubArea{Primitive function}
\LANGen{
\Question{
Evaluate the following integral on the interval \((0;+\infty)\). 
\[ 
 \int x^{3}\ln x\, \mathrm{d}x
\]}
{\(\frac{x^4}{4}\ln x -\frac{x^4} 
{16}+ c,\ c\in \mathbb{R}\)}{\(\frac{x^3}{3}\ln x -\frac{x^3} 
{9}+ c,\ c\in \mathbb{R}\)}{\(\frac{x^2}{2}\ln x -\frac{x^2} 
{4}+ c,\ c\in \mathbb{R}\)}{\(x\ln x -x+ c,\ c\in \mathbb{R}\)}{}{}}
\LANGpl{
\Question{
Oblicz następującą całkę na przedziale \((0;+\infty)\). 
\[ 
 \int x^{3}\ln x\, \mathrm{d}x
\]}
{\(\frac{x^4}{4}\ln x -\frac{x^4} 
{16}+ c,\ c\in \mathbb{R}\)}{\(\frac{x^3}{3}\ln x -\frac{x^3} 
{9}+ c,\ c\in \mathbb{R}\)}{\(\frac{x^2}{2}\ln x -\frac{x^2} 
{4}+ c,\ c\in \mathbb{R}\)}{\(x\ln x -x+ c,\ c\in \mathbb{R}\)}{}{}}
\LANGcs{
\Question{
Vypočtěte 
\[ 
 \int x^{3}\ln x\, \mathrm{d}x
\] na intervalu \((0;+\infty)\).}
{\(\frac{x^4}{4}\ln x -\frac{x^4} 
{16}+ c,\ c\in \mathbb{R}\)}{\(\frac{x^3}{3}\ln x -\frac{x^3} 
{9}+ c,\ c\in \mathbb{R}\)}{\(\frac{x^2}{2}\ln x -\frac{x^2} 
{4}+ c,\ c\in \mathbb{R}\)}{\(x\ln x -x+ c,\ c\in \mathbb{R}\)}{}{}}
\LANGsk{
\Question{
Vypočítajte
\[ 
 \int x^{3}\ln x\, \mathrm{d}x
\]
na intervale \((0;+\infty)\).}
{\(\frac{x^4}{4}\ln x -\frac{x^4} 
{16}+ c,\ c\in \mathbb{R}\)}{\(\frac{x^3}{3}\ln x -\frac{x^3} 
{9}+ c,\ c\in \mathbb{R}\)}{\(\frac{x^2}{2}\ln x -\frac{x^2} 
{4}+ c,\ c\in \mathbb{R}\)}{\(x\ln x -x+ c,\ c\in \mathbb{R}\)}{}{}}
\LANGes{
\Question{
Evalúa la siguiente integral en el intervalo \((0;+\infty)\). 
\[ 
 \int x^{3}\ln x\, \mathrm{d}x
\]}
{\(\frac{x^4}{4}\ln x -\frac{x^4} 
{16}+ c,\ c\in \mathbb{R}\)}{\(\frac{x^3}{3}\ln x -\frac{x^3} 
{9}+ c,\ c\in \mathbb{R}\)}{\(\frac{x^2}{2}\ln x -\frac{x^2} 
{4}+ c,\ c\in \mathbb{R}\)}{\(x\ln x -x+ c,\ c\in \mathbb{R}\)}{}{}}
\End

\Data\ID{2010000307}
\Updated{2022-08-25 10:08:55}
\Level{C}
\SubArea{Primitive function}
\LANGen{
\Question{
Evaluate the following integral on \(\mathbb{R}\). 
\[ 
 \int \sin^{2}x\cos x\, \mathrm{d}x
\]}
{\(\frac{\sin^{3}x} 
 {3} + c,\ c\in \mathbb{R}\)}{\(-\frac{\cos ^{3}x\sin x} 
 {3} + c,\ c\in \mathbb{R}\)}{\(2\sin x + c,\ c\in \mathbb{R}\)}{}{}{}}
\LANGpl{
\Question{
Oblicz następującą całkę na przedziale \(\mathbb{R}\). 
\[ 
 \int \sin^{2}x\cos x\, \mathrm{d}x
\]}
{\(\frac{\sin^{3}x} 
 {3} + c,\ c\in \mathbb{R}\)}{\(-\frac{\cos ^{3}x\sin x} 
 {3} + c,\ c\in \mathbb{R}\)}{\(2\sin x + c,\ c\in \mathbb{R}\)}{}{}{}}
\LANGcs{
\Question{
Vypočtěte 
\[ 
 \int \sin^{2}x\cos x\, \mathrm{d}x
\] na \(\mathbb{R}\).}
{\(\frac{\sin^{3}x} 
 {3} + c,\ c\in \mathbb{R}\)}{\(-\frac{\cos ^{3}x\sin x} 
 {3} + c,\ c\in \mathbb{R}\)}{\(2\sin x + c,\ c\in \mathbb{R}\)}{}{}{}}
\LANGsk{
\Question{
Riešte neurčitý integrál 
\[ 
 \int \sin^{2}x\cos x\, \mathrm{d}x
\]
v \(\mathbb{R}\).}
{\(\frac{\sin^{3}x} 
 {3} + c,\ c\in \mathbb{R}\)}{\(-\frac{\cos ^{3}x\sin x} 
 {3} + c,\ c\in \mathbb{R}\)}{\(2\sin x + c,\ c\in \mathbb{R}\)}{}{}{}}
\LANGes{
\Question{
Evalúa la siguiente integral en \(\mathbb{R}\). 
\[ 
 \int \sin^{2}x\cos x\, \mathrm{d}x
\]}
{\(\frac{\sin^{3}x} 
 {3} + c,\ c\in \mathbb{R}\)}{\(-\frac{\cos ^{3}x\sin x} 
 {3} + c,\ c\in \mathbb{R}\)}{\(2\sin x + c,\ c\in \mathbb{R}\)}{}{}{}}
\End

\Data\ID{9000066001}
\Updated{2020-07-15 03:07:37}
\Level{C}
\SubArea{Primitive function}
\LANGen{
\Question{
Identify the formula for integration by parts.}
{\(\int u(x)v'(x)\, \mathrm{d}x = u(x)v(x) -\int u'(x)v(x)\, \mathrm{d}x\)}{\(\int u(x)v(x)\, \mathrm{d}x = u'(x)v'(x) -\int u'(x)v(x)\, \mathrm{d}x\)}{\(\int u'(x)v'(x)\, \mathrm{d}x = u(x)v(x) -\int u'(x)v(x)\, \mathrm{d}x\)}{\(\int u(x)v'(x)\, \mathrm{d}x = u(x)v(x) +\int u'(x)v(x)\, \mathrm{d}x\)}{}{}}
\LANGpl{
\Question{
Wskaż wzór na całkowanie przez części.}
{\(\int u(x)v'(x)\, \mathrm{d}x = u(x)v(x) -\int u'(x)v(x)\, \mathrm{d}x\)}{\(\int u(x)v(x)\, \mathrm{d}x = u'(x)v'(x) -\int u'(x)v(x)\, \mathrm{d}x\)}{\(\int u'(x)v'(x)\, \mathrm{d}x = u(x)v(x) -\int u'(x)v(x)\, \mathrm{d}x\)}{\(\int u(x)v'(x)\, \mathrm{d}x = u(x)v(x) +\int u'(x)v(x)\, \mathrm{d}x\)}{}{}}
\LANGcs{
\Question{
Vyberte, který z uvedených vztahů je symbolickým zápisem integrační
metody per partes.}
{\(\int u(x)v'(x)\, \mathrm{d}x = u(x)v(x) -\int u'(x)v(x)\, \mathrm{d}x\)}{\(\int u(x)v(x)\, \mathrm{d}x = u'(x)v'(x) -\int u'(x)v(x)\, \mathrm{d}x\)}{\(\int u'(x)v'(x)\, \mathrm{d}x = u(x)v(x) -\int u'(x)v(x)\, \mathrm{d}x\)}{\(\int u(x)v'(x)\, \mathrm{d}x = u(x)v(x) +\int u'(x)v(x)\, \mathrm{d}x\)}{}{}}
\LANGsk{
\Question{
Vyberte, ktorý z uvedených vzťahov je symbolickým zápisom integračnej
metódy per partes.}
{\(\int u(x)v'(x)\, \mathrm{d}x = u(x)v(x) -\int u'(x)v(x)\, \mathrm{d}x\)}{\(\int u(x)v(x)\, \mathrm{d}x = u'(x)v'(x) -\int u'(x)v(x)\, \mathrm{d}x\)}{\(\int u'(x)v'(x)\, \mathrm{d}x = u(x)v(x) -\int u'(x)v(x)\, \mathrm{d}x\)}{\(\int u(x)v'(x)\, \mathrm{d}x = u(x)v(x) +\int u'(x)v(x)\, \mathrm{d}x\)}{}{}}
\LANGes{
\Question{
Identifica la fórmula de integración por partes.}
{\(\int u(x)v'(x)\, \mathrm{d}x = u(x)v(x) -\int u'(x)v(x)\, \mathrm{d}x\)}{\(\int u(x)v(x)\, \mathrm{d}x = u'(x)v'(x) -\int u'(x)v(x)\, \mathrm{d}x\)}{\(\int u'(x)v'(x)\, \mathrm{d}x = u(x)v(x) -\int u'(x)v(x)\, \mathrm{d}x\)}{\(\int u(x)v'(x)\, \mathrm{d}x = u(x)v(x) +\int u'(x)v(x)\, \mathrm{d}x\)}{}{}}
\End

\Data\ID{9000066004}
\Updated{2020-11-12 01:11:01}
\Level{C}
\SubArea{Primitive function}
\LANGen{
\Question{
Evaluate the following integral on \(\mathbb{R}\). 
\[ 
 \int x^{2}\sin x\, \mathrm{d}x
\]}
{\(- x^{2}\cos x + 2x\sin x + 2\cos x + c,\ c\in \mathbb{R}\)}{\(x^{2}\cos x - 2x\sin x - 2\cos x + c,\ c\in \mathbb{R}\)}{\(\frac{1}
{3}x^{3}\cos x + c,\ c\in \mathbb{R}\)}{\(\frac{1}
{3}x^{3} -\cos x + c,\ c\in \mathbb{R}\)}{}{}}
\LANGpl{
\Question{
Wyznacz całkę.
\[ 
 \int x^{2}\sin x\, \mathrm{d}x
\]}
{\(- x^{2}\cos x + 2x\sin x + 2\cos x + c,\ c\in \mathbb{R}\)}{\(x^{2}\cos x - 2x\sin x - 2\cos x + c,\ c\in \mathbb{R}\)}{\(\frac{1}
{3}x^{3}\cos x + c,\ c\in \mathbb{R}\)}{\(\frac{1}
{3}x^{3} -\cos x + c,\ c\in \mathbb{R}\)}{}{}}
\LANGcs{
\Question{
Vypočtěte \(\int x^{2}\sin x\, \mathrm{d}x\) na \(\mathbb{R}\).}
{\(- x^{2}\cos x + 2x\sin x + 2\cos x + c,\ c\in\mathbb{R}\)}{\(x^{2}\cos x - 2x\sin x - 2\cos x + c,\ c\in\mathbb{R}\)}{\(\frac{1}
{3}x^{3}\cos x + c,\ c\in\mathbb{R}\)}{\(\frac{1}
{3}x^{3} -\cos x + c,\ c\in\mathbb{R}\)}{}{}}
\LANGsk{
\Question{
Vypočítajte \(\int x^{2}\sin x\, \mathrm{d}x\) na \(\mathbb{R}\).}
{\(- x^{2}\cos x + 2x\sin x + 2\cos x + c,\ c\in\mathbb{R}\)}{\(x^{2}\cos x - 2x\sin x - 2\cos x + c,\ c\in\mathbb{R}\)}{\(\frac{1}
{3}x^{3}\cos x + c,\ c\in\mathbb{R}\)}{\(\frac{1}
{3}x^{3} -\cos x + c,\ c\in\mathbb{R}\)}{}{}}
\LANGes{
\Question{
Resuelve la siguiente integral en \(\mathbb{R}\).  
\[ 
 \int x^{2}\sin x\, \mathrm{d}x
\]}
{\(- x^{2}\cos x + 2x\sin x + 2\cos x + c,\ c\in \mathbb{R}\)}{\(x^{2}\cos x - 2x\sin x - 2\cos x + c,\ c\in \mathbb{R}\)}{\(\frac{1}
{3}x^{3}\cos x + c,\ c\in \mathbb{R}\)}{\(\frac{1}
{3}x^{3} -\cos x + c,\ c\in \mathbb{R}\)}{}{}}
\End

\Data\ID{9000066005}
\Updated{2022-04-23 05:04:02}
\Level{C}
\SubArea{Primitive function}
\LANGen{
\Question{
Evaluate the following integral on the interval \((0;+\infty)\). 
\[ 
 \int \ln x\, \mathrm{d}x
\]}
{\(x\ln x - x + c,\ c\in \mathbb{R}\)}{\(\ln x - x + c,\ c\in \mathbb{R}\)}{\(x\ln x + x + c,\ c\in \mathbb{R}\)}{\(x\ln x -\frac{1} 
{2}x^{2} + c,\ c\in \mathbb{R}\)}{}{}}
\LANGpl{
\Question{
Wyznacz całkę na przedziale  \((0;+\infty)\).
\[ 
 \int \ln x\, \mathrm{d}x
\]}
{\(x\ln x - x + c,\ c\in \mathbb{R}\)}{\(\ln x - x + c,\ c\in \mathbb{R}\)}{\(x\ln x + x + c,\ c\in \mathbb{R}\)}{\(x\ln x -\frac{1} 
{2}x^{2} + c,\ c\in \mathbb{R}\)}{}{}}
\LANGcs{
\Question{
Vypočtěte \(\int \ln x\, \mathrm{d}x\) na intervalu \((0;+\infty)\).}
{\(x\ln x - x + c,\ c\in\mathbb{R}\)}{\(\ln x - x + c,\ c\in\mathbb{R}\)}{\(x\ln x + x + c,\ c\in\mathbb{R}\)}{\(x\ln x -\frac{1} 
{2}x^{2} + c,\ c\in\mathbb{R}\)}{}{}}
\LANGsk{
\Question{
Vypočítajte \(\int \ln x\, \mathrm{d}x\) na intervalu \((0;+\infty)\).}
{\(x\ln x - x + c,\ c\in\mathbb{R}\)}{\(\ln x - x + c,\ c\in\mathbb{R}\)}{\(x\ln x + x + c,\ c\in\mathbb{R}\)}{\(x\ln x -\frac{1} 
{2}x^{2} + c,\ c\in\mathbb{R}\)}{}{}}
\LANGes{
\Question{
Resuelve la siguiente integral en el intervalo \((0;+\infty)\). 
\[ 
 \int \ln x\, \mathrm{d}x
\]}
{\(x\ln x - x + c,\ c\in \mathbb{R}\)}{\(\ln x - x + c,\ c\in \mathbb{R}\)}{\(x\ln x + x + c,\ c\in \mathbb{R}\)}{\(x\ln x -\frac{1} 
{2}x^{2} + c,\ c\in \mathbb{R}\)}{}{}}
\End

\Data\ID{9000066006}
\Updated{2020-11-12 01:11:12}
\Level{C}
\SubArea{Primitive function}
\LANGen{
\Question{
Evaluate the following integral on the interval \((0;+\infty)\). 
\[ 
 \int x\ln x\, \mathrm{d}x
\]}
{\(\frac{1}
{2}x^{2}\ln x -\frac{1} 
{4}x^{2} + c,\ c\in \mathbb{R}\)}{\(x\ln x -\frac{1} 
{2}x^{2} + c,\ c\in \mathbb{R}\)}{\(x\ln x - x + c,\ c\in \mathbb{R}\)}{\(\frac{1}
{2}x^{2} + \frac{1} 
{|x|} + c,\ c\in \mathbb{R}\)}{}{}}
\LANGpl{
\Question{
Wyznacz całkę na przedziale \((0;+\infty)\).
\[ 
 \int x\ln x\, \mathrm{d}x
\]}
{\(\frac{1}
{2}x^{2}\ln x -\frac{1} 
{4}x^{2} + c,\ c\in \mathbb{R}\)}{\(x\ln x -\frac{1} 
{2}x^{2} + c,\ c\in \mathbb{R}\)}{\(x\ln x - x + c,\ c\in \mathbb{R}\)}{\(\frac{1}
{2}x^{2} + \frac{1} 
{|x|} + c,\ c\in \mathbb{R}\)}{}{}}
\LANGcs{
\Question{
Vypočtěte \(\int x\ln x\, \mathrm{d}x\) na intervalu \((0;+\infty)\).}
{\(\frac{1}
{2}x^{2}\ln x -\frac{1} 
{4}x^{2} + c,\ c\in\mathbb{R}\)}{\(x\ln x -\frac{1} 
{2}x^{2} + c,\ c\in\mathbb{R}\)}{\(x\ln x - x + c,\ c\in\mathbb{R}\)}{\(\frac{1}
{2}x^{2} + \frac{1} 
{|x|} + c,\ c\in\mathbb{R}\)}{}{}}
\LANGsk{
\Question{
Vypočítajte \(\int x\ln x\, \mathrm{d}x\) na intervalu \((0;+\infty)\).}
{\(\frac{1}
{2}x^{2}\ln x -\frac{1} 
{4}x^{2} + c,\ c\in\mathbb{R}\)}{\(x\ln x -\frac{1} 
{2}x^{2} + c,\ c\in\mathbb{R}\)}{\(x\ln x - x + c,\ c\in\mathbb{R}\)}{\(\frac{1}
{2}x^{2} + \frac{1} 
{|x|} + c,\ c\in\mathbb{R}\)}{}{}}
\LANGes{
\Question{
Resuelve la siguiente integral en el intervalo \((0;+\infty)\).  
\[ 
 \int x\ln x\, \mathrm{d}x
\]}
{\(\frac{1}
{2}x^{2}\ln x -\frac{1} 
{4}x^{2} + c,\ c\in \mathbb{R}\)}{\(x\ln x -\frac{1} 
{2}x^{2} + c,\ c\in \mathbb{R}\)}{\(x\ln x - x + c,\ c\in \mathbb{R}\)}{\(\frac{1}
{2}x^{2} + \frac{1} 
{|x|} + c,\ c\in \mathbb{R}\)}{}{}}
\End

\Data\ID{9000066007}
\Updated{2022-04-23 05:04:02}
\Level{C}
\SubArea{Primitive function}
\LANGen{
\Question{
Evaluate the following integral on the interval \((0;+\infty)\). 
\[ 
 \int x^{2}\ln x\, \mathrm{d}x
\]}
{\(\frac{1}
{3}x^{3}\ln x -\frac{1} 
{9}x^{3} + c,\ c\in \mathbb{R}\)}{\(\frac{1}
{2}x^{2}\ln x -\frac{1} 
{4}x^{2} + c,\ c\in \mathbb{R}\)}{\(x\ln x -\frac{1} 
{2}x^{2} + c,\ c\in \mathbb{R}\)}{\(x\ln x - x + c,\ c\in \mathbb{R}\)}{}{}}
\LANGpl{
\Question{
Wyznacz całkę na przedziale  \((0;+\infty)\).
\[ 
 \int x^{2}\ln x\, \mathrm{d}x
\]}
{\(\frac{1}
{3}x^{3}\ln x -\frac{1} 
{9}x^{3} + c,\ c\in \mathbb{R}\)}{\(\frac{1}
{2}x^{2}\ln x -\frac{1} 
{4}x^{2} + c,\ c\in \mathbb{R}\)}{\(x\ln x -\frac{1} 
{2}x^{2} + c,\ c\in \mathbb{R}\)}{\(x\ln x - x + c,\ c\in \mathbb{R}\)}{}{}}
\LANGcs{
\Question{
Vypočtěte \(\int x^{2}\ln x\, \mathrm{d}x\) na intervalu \((0;+\infty)\).}
{\(\frac{1}
{3}x^{3}\ln x -\frac{1} 
{9}x^{3} + c,\ c\in\mathbb{R}\)}{\(\frac{1}
{2}x^{2}\ln x -\frac{1} 
{4}x^{2} + c,\ c\in\mathbb{R}\)}{\(x\ln x -\frac{1} 
{2}x^{2} + c,\ c\in\mathbb{R}\)}{\(x\ln x - x + c,\ c\in\mathbb{R}\)}{}{}}
\LANGsk{
\Question{
Vypočítajte \(\int x^{2}\ln x\, \mathrm{d}x\) na intervale \((0;+\infty)\).}
{\(\frac{1}
{3}x^{3}\ln x -\frac{1} 
{9}x^{3} + c,\ c\in\mathbb{R}\)}{\(\frac{1}
{2}x^{2}\ln x -\frac{1} 
{4}x^{2} + c,\ c\in\mathbb{R}\)}{\(x\ln x -\frac{1} 
{2}x^{2} + c,\ c\in\mathbb{R}\)}{\(x\ln x - x + c,\ c\in\mathbb{R}\)}{}{}}
\LANGes{
\Question{
Resuelve la siguiente integral en el intervalo \((0;+\infty)\).  
\[ 
 \int x^{2}\ln x\, \mathrm{d}x
\]}
{\(\frac{1}
{3}x^{3}\ln x -\frac{1} 
{9}x^{3} + c,\ c\in \mathbb{R}\)}{\(\frac{1}
{2}x^{2}\ln x -\frac{1} 
{4}x^{2} + c,\ c\in \mathbb{R}\)}{\(x\ln x -\frac{1} 
{2}x^{2} + c,\ c\in \mathbb{R}\)}{\(x\ln x - x + c,\ c\in \mathbb{R}\)}{}{}}
\End

\Data\ID{9000066009}
\Updated{2020-11-12 01:11:36}
\Level{C}
\SubArea{Primitive function}
\LANGen{
\Question{
Evaluate the following integral on \(\mathbb{R}\). 
\[ 
 \int x^{2}\mathrm{e}^{x}\, \mathrm{d}x
\]}
{\(x^{2}\mathrm{e}^{x} - 2x\mathrm{e}^{x} + 2\mathrm{e}^{x} + c,\ c\in \mathbb{R}\)}{\(x^{2}\mathrm{e}^{x} + 2x\mathrm{e}^{x} - 2\mathrm{e}^{x} + c,\ c\in \mathbb{R}\)}{\(\frac{1}
{3}x^{3}\mathrm{e}^{x} -\frac{1} 
{2}x^{2}\mathrm{e}^{x} + 2\mathrm{e}^{x} + c,\ c\in \mathbb{R}\)}{\(\frac{1}
{3}x^{3}\mathrm{e}^{x} + \frac{1} 
{2}x^{2}\mathrm{e}^{x} - 2\mathrm{e}^{x} + c,\ c\in \mathbb{R}\)}{}{}}
\LANGpl{
\Question{
Wyznacz całkę.
\[ 
 \int x^{2}\mathrm{e}^{x}\, \mathrm{d}x
\]}
{\(x^{2}\mathrm{e}^{x} - 2x\mathrm{e}^{x} + 2\mathrm{e}^{x} + c,\ c\in \mathbb{R}\)}{\(x^{2}\mathrm{e}^{x} + 2x\mathrm{e}^{x} - 2\mathrm{e}^{x} + c,\ c\in \mathbb{R}\)}{\(\frac{1}
{3}x^{3}\mathrm{e}^{x} -\frac{1} 
{2}x^{2}\mathrm{e}^{x} + 2\mathrm{e}^{x} + c,\ c\in \mathbb{R}\)}{\(\frac{1}
{3}x^{3}\mathrm{e}^{x} + \frac{1} 
{2}x^{2}\mathrm{e}^{x} - 2\mathrm{e}^{x} + c,\ c\in \mathbb{R}\)}{}{}}
\LANGcs{
\Question{
Vypočtěte \(\int x^{2}\mathrm{e}^{x}\, \mathrm{d}x\) na \(\mathbb{R}\).}
{\(x^{2}\mathrm{e}^{x} - 2x\mathrm{e}^{x} + 2\mathrm{e}^{x} + c,\ c\in\mathbb{R}\)}{\(x^{2}\mathrm{e}^{x} + 2x\mathrm{e}^{x} - 2\mathrm{e}^{x} + c,\ c\in\mathbb{R}\)}{\(\frac{1}
{3}x^{3}\mathrm{e}^{x} -\frac{1} 
{2}x^{2}\mathrm{e}^{x} + 2\mathrm{e}^{x} + c,\ c\in\mathbb{R}\)}{\(\frac{1}
{3}x^{3}\mathrm{e}^{x} + \frac{1} 
{2}x^{2}\mathrm{e}^{x} - 2\mathrm{e}^{x} + c,\ c\in\mathbb{R}\)}{}{}}
\LANGsk{
\Question{
Vypočítajte \(\int x^{2}\mathrm{e}^{x}\, \mathrm{d}x\) na \(\mathbb{R}\).}
{\(x^{2}\mathrm{e}^{x} - 2x\mathrm{e}^{x} + 2\mathrm{e}^{x} + c,\ c\in\mathbb{R}\)}{\(x^{2}\mathrm{e}^{x} + 2x\mathrm{e}^{x} - 2\mathrm{e}^{x} + c,\ c\in\mathbb{R}\)}{\(\frac{1}
{3}x^{3}\mathrm{e}^{x} -\frac{1} 
{2}x^{2}\mathrm{e}^{x} + 2\mathrm{e}^{x} + c,\ c\in\mathbb{R}\)}{\(\frac{1}
{3}x^{3}\mathrm{e}^{x} + \frac{1} 
{2}x^{2}\mathrm{e}^{x} - 2\mathrm{e}^{x} + c,\ c\in\mathbb{R}\)}{}{}}
\LANGes{
\Question{
Resuelve la siguiente integral en \(\mathbb{R}\). 
\[ 
 \int x^{2}\mathrm{e}^{x}\, \mathrm{d}x
\]}
{\(x^{2}\mathrm{e}^{x} - 2x\mathrm{e}^{x} + 2\mathrm{e}^{x} + c,\ c\in \mathbb{R}\)}{\(x^{2}\mathrm{e}^{x} + 2x\mathrm{e}^{x} - 2\mathrm{e}^{x} + c,\ c\in \mathbb{R}\)}{\(\frac{1}
{3}x^{3}\mathrm{e}^{x} -\frac{1} 
{2}x^{2}\mathrm{e}^{x} + 2\mathrm{e}^{x} + c,\ c\in \mathbb{R}\)}{\(\frac{1}
{3}x^{3}\mathrm{e}^{x} + \frac{1} 
{2}x^{2}\mathrm{e}^{x} - 2\mathrm{e}^{x} + c,\ c\in \mathbb{R}\)}{}{}}
\End

\Data\ID{9000066010}
\Updated{2020-11-12 01:11:03}
\Level{C}
\SubArea{Primitive function}
\LANGen{
\Question{
Evaluate the following integral on \(\mathbb{R}\). 
\[ 
 \int \mathrm{e}^{2x}\, \mathrm{d}x
\]}
{\(\frac{1}
{2}\mathrm{e}^{2x} + c,\ c\in \mathbb{R}\)}{\(\frac{1}
{3}\mathrm{e}^{3x} + c,\ c\in \mathbb{R}\)}{\(\mathrm{e}^{2x} -\mathrm{e}^{x} + c,\ c\in \mathbb{R}\)}{\(2\mathrm{e}^{2x} + c,\ c\in \mathbb{R}\)}{}{}}
\LANGpl{
\Question{
Wyznacz całkę.
\[ 
 \int \mathrm{e}^{2x}\, \mathrm{d}x
\]}
{\(\frac{1}
{2}\mathrm{e}^{2x} + c,\ c\in \mathbb{R}\)}{\(\frac{1}
{3}\mathrm{e}^{3x} + c,\ c\in \mathbb{R}\)}{\(\mathrm{e}^{2x} -\mathrm{e}^{x} + c,\ c\in \mathbb{R}\)}{\(2\mathrm{e}^{2x} + c,\ c\in \mathbb{R}\)}{}{}}
\LANGcs{
\Question{
Vypočtěte \(\int \mathrm{e}^{2x}\, \mathrm{d}x\) na \(\mathbb{R}\).}
{\(\frac{1}
{2}\mathrm{e}^{2x} + c,\ c\in\mathbb{R}\)}{\(\frac{1}
{3}\mathrm{e}^{3x} + c,\ c\in\mathbb{R}\)}{\(\mathrm{e}^{2x} -\mathrm{e}^{x} + c,\ c\in\mathbb{R}\)}{\(2\mathrm{e}^{2x} + c,\ c\in\mathbb{R}\)}{}{}}
\LANGsk{
\Question{
Vypočítajte \(\int \mathrm{e}^{2x}\, \mathrm{d}x\) na \(\mathbb{R}\).}
{\(\frac{1}
{2}\mathrm{e}^{2x} + c,\ c\in\mathbb{R}\)}{\(\frac{1}
{3}\mathrm{e}^{3x} + c,\ c\in\mathbb{R}\)}{\(\mathrm{e}^{2x} -\mathrm{e}^{x} + c,\ c\in\mathbb{R}\)}{\(2\mathrm{e}^{2x} + c,\ c\in\mathbb{R}\)}{}{}}
\LANGes{
\Question{
Resuelve la siguiente integral en \(\mathbb{R}\). 
\[ 
 \int \mathrm{e}^{2x}\, \mathrm{d}x
\]}
{\(\frac{1}
{2}\mathrm{e}^{2x} + c,\ c\in \mathbb{R}\)}{\(\frac{1}
{3}\mathrm{e}^{3x} + c,\ c\in \mathbb{R}\)}{\(\mathrm{e}^{2x} -\mathrm{e}^{x} + c,\ c\in \mathbb{R}\)}{\(2\mathrm{e}^{2x} + c,\ c\in \mathbb{R}\)}{}{}}
\End

\Data\ID{9000071208}
\Updated{2020-07-15 03:07:37}
\Level{C}
\SubArea{Primitive function}
\LANGen{
\Question{
Evaluate the following integral on the interval \((0;+\infty)\). 
\[ 
 \int x^{2}\ln x\, \mathrm{d}x
\]}
{\(\frac{x^{3}} 
 {3} \left (\ln x -\frac{1} 
{3}\right ) + c,\ c\in \mathbb{R}\)}{\(\frac{x^{2}} 
 {3} + c,\ c\in \mathbb{R}\)}{\(x^{2}\left (\frac{x\ln x}
 {3} -\frac{1} 
{2}\right ) + c,\ c\in \mathbb{R}\)}{}{}{}}
\LANGpl{
\Question{
Wyznacz całkę na przedziale  \((0;+\infty)\). 
\[ 
 \int x^{2}\ln x\, \mathrm{d}x
\]}
{\(\frac{x^{3}} 
 {3} \left (\ln x -\frac{1} 
{3}\right ) + c,\ c\in \mathbb{R}\)}{\(\frac{x^{2}} 
 {3} + c,\ c\in \mathbb{R}\)}{\(x^{2}\left (\frac{x\ln x}
 {3} -\frac{1} 
{2}\right ) + c,\ c\in \mathbb{R}\)}{}{}{}}
\LANGcs{
\Question{
Vypočtěte \(\int x^{2}\ln x\, \mathrm{d}x\) na intervalu \((0;+\infty)\).}
{\(\frac{x^{3}} 
 {3} \left (\ln x -\frac{1} 
{3}\right ) + c,\ c\in \mathbb{R}\)}{\(\frac{x^{2}} 
 {3} + c,\ c\in \mathbb{R}\)}{\(x^{2}\left (\frac{x\ln x}
 {3} -\frac{1} 
{2}\right ) + c,\ c\in \mathbb{R}\)}{}{}{}}
\LANGsk{
\Question{
Vypočítajte \(\int x^{2}\ln x\, \mathrm{d}x\) na intervale \((0;+\infty)\).}
{\(\frac{x^{3}} 
 {3} \left (\ln x -\frac{1} 
{3}\right ) + c,\ c\in \mathbb{R}\)}{\(\frac{x^{2}} 
 {3} + c,\ c\in \mathbb{R}\)}{\(x^{2}\left (\frac{x\ln x}
 {3} -\frac{1} 
{2}\right ) + c,\ c\in \mathbb{R}\)}{}{}{}}
\LANGes{
\Question{
Evalúa la siguiente integral en el intervalo \((0;+\infty)\). 
\[ 
 \int x^{2}\ln x\, \mathrm{d}x
\]}
{\(\frac{x^{3}} 
 {3} \left (\ln x -\frac{1} 
{3}\right ) + c,\ c\in \mathbb{R}\)}{\(\frac{x^{2}} 
 {3} + c,\ c\in \mathbb{R}\)}{\(x^{2}\left (\frac{x\ln x}
 {3} -\frac{1} 
{2}\right ) + c,\ c\in \mathbb{R}\)}{}{}{}}
\End

\Data\ID{9000071209}
\Updated{2022-04-23 05:04:02}
\Level{C}
\SubArea{Primitive function}
\LANGen{
\Question{
Evaluate the following integral on \(\mathbb{R}\). 
\[ 
 \int \cos ^{3}x\sin x\, \mathrm{d}x
\]}
{\(-\frac{\cos ^{4}x} 
 {4} + c,\ c\in \mathbb{R}\)}{\(-\frac{\sin ^{4}x\cos x} 
 {4} + c,\ c\in \mathbb{R}\)}{\(- 3\cos ^{2}x + c,\ c\in \mathbb{R}\)}{}{}{}}
\LANGpl{
\Question{
Wyznacz całkę.
\[ 
 \int \cos ^{3}x\sin x\, \mathrm{d}x
\]}
{\(-\frac{\cos ^{4}x} 
 {4} + c,\ c\in \mathbb{R}\)}{\(-\frac{\sin ^{4}x\cos x} 
 {4} + c,\ c\in \mathbb{R}\)}{\(- 3\cos ^{2}x + c,\ c\in \mathbb{R}\)}{}{}{}}
\LANGcs{
\Question{
Vypočtěte \(\int \cos ^{3}x\sin x\, \mathrm{d}x\) na \(\mathbb{R}\).}
{\(-\frac{\cos ^{4}x} 
 {4} + c,\ c\in \mathbb{R}\)}{\(-\frac{\sin ^{4}x\cos x} 
 {4} + c,\ c\in \mathbb{R}\)}{\(- 3\cos ^{2}x + c,\ c\in \mathbb{R}\)}{}{}{}}
\LANGsk{
\Question{
Vypočítajte \(\int \cos ^{3}x\sin x\, \mathrm{d}x\) na \(\mathbb{R}\).}
{\(-\frac{\cos ^{4}x} 
 {4} + c,\ c\in \mathbb{R}\)}{\(-\frac{\sin ^{4}x\cos x} 
 {4} + c,\ c\in \mathbb{R}\)}{\(- 3\cos ^{2}x + c,\ c\in \mathbb{R}\)}{}{}{}}
\LANGes{
\Question{
Evalúa la siguiente integral en \(\mathbb{R}\). 
\[ 
 \int \cos ^{3}x\sin x\, \mathrm{d}x
\]}
{\(-\frac{\cos ^{4}x} 
 {4} + c,\ c\in \mathbb{R}\)}{\(-\frac{\sin ^{4}x\cos x} 
 {4} + c,\ c\in \mathbb{R}\)}{\(- 3\cos ^{2}x + c,\ c\in \mathbb{R}\)}{}{}{}}
\End

\Data\ID{9000071210}
\Updated{2020-07-15 03:07:37}
\Level{C}
\SubArea{Primitive function}
\LANGen{
\Question{
Evaluate the following integral on \(\mathbb{R}\). 
\[ 
 \int (x + 2)\cos x\, \mathrm{d}x
\]}
{\((x + 2)\sin x +\cos x + c,\ c\in \mathbb{R}\)}{\(-\left (\frac{x^{2}} 
 {2} + 2x\right )\sin x + c,\ c\in \mathbb{R}\)}{\((x + 2)\sin x -\sin x + c,\ c\in \mathbb{R}\)}{}{}{}}
\LANGpl{
\Question{
Wyznacz całkę.
\[ 
 \int (x + 2)\cos x\, \mathrm{d}x
\]}
{\((x + 2)\sin x +\cos x + c,\ c\in \mathbb{R}\)}{\(-\left (\frac{x^{2}} 
 {2} + 2x\right )\sin x + c,\ c\in \mathbb{R}\)}{\((x + 2)\sin x -\sin x + c,\ c\in \mathbb{R}\)}{}{}{}}
\LANGcs{
\Question{
Vypočtěte \(\int (x + 2)\cos x\, \mathrm{d}x\) na \(\mathbb{R}\).}
{\((x + 2)\sin x +\cos x + c,\ c\in \mathbb{R}\)}{\(-\left (\frac{x^{2}} 
 {2} + 2x\right )\sin x + c,\ c\in \mathbb{R}\)}{\((x + 2)\sin x -\sin x + c,\ c\in \mathbb{R}\)}{}{}{}}
\LANGsk{
\Question{
Vypočítajte \(\int (x + 2)\cos x\, \mathrm{d}x\) na \(\mathbb{R}\).}
{\((x + 2)\sin x +\cos x + c,\ c\in \mathbb{R}\)}{\(-\left (\frac{x^{2}} 
 {2} + 2x\right )\sin x + c,\ c\in \mathbb{R}\)}{\((x + 2)\sin x -\sin x + c,\ c\in \mathbb{R}\)}{}{}{}}
\LANGes{
\Question{
Evalúa la siguiente integral en \(\mathbb{R}\). 
\[ 
 \int (x + 2)\cos x\, \mathrm{d}x
\]}
{\((x + 2)\sin x +\cos x + c,\ c\in \mathbb{R}\)}{\(-\left (\frac{x^{2}} 
 {2} + 2x\right )\sin x + c,\ c\in \mathbb{R}\)}{\((x + 2)\sin x -\sin x + c,\ c\in \mathbb{R}\)}{}{}{}}
\End

\Data\ID{9000150104}
\Updated{2020-07-15 03:07:37}
\Level{C}
\SubArea{Primitive function}
\LANGen{
\Question{
Evaluate the following integral on \(\mathbb{R}\). 
\[ 
 \int \cos x\cdot \left (-3 +\sin x\right )^{5}\, \mathrm{d}x
\]}
{\(\frac{\left (-3+\sin x\right )^{6}} 
 {6} + c\text{, }c\in \mathbb{R}\)}{\(6\left (-3 +\sin x\right )^{6} + c,\ c\in \mathbb{R}\)}{\(\frac{\left (-3+\cos x\right )^{6}} 
 {6} + c,\ c\in \mathbb{R}\)}{\(6\left (-3 +\cos x\right )^{6} + c,\ c\in \mathbb{R}\)}{}{}}
\LANGpl{
\Question{
Wyznacz całkę.
\[ 
 \int \cos x\cdot \left (-3 +\sin x\right )^{5}\, \mathrm{d}x
\]}
{\(\frac{\left (-3+\sin x\right )^{6}} 
 {6} + c\text{, }c\in \mathbb{R}\)}{\(6\left (-3 +\sin x\right )^{6} + c,\ c\in \mathbb{R}\)}{\(\frac{\left (-3+\cos x\right )^{6}} 
 {6} + c,\ c\in \mathbb{R}\)}{\(6\left (-3 +\cos x\right )^{6} + c,\ c\in \mathbb{R}\)}{}{}}
\LANGcs{
\Question{
Vypočtěte \(\int \cos x\cdot \left (-3 +\sin x\right )^{5}\, \mathrm{d}x\) na \(\mathbb{R}\).}
{\(\frac{\left (-3+\sin x\right )^{6}} 
 {6} + c,\ c\in \mathbb{R}\)}{\(6\left (-3 +\sin x\right )^{6} + c,\ c\in \mathbb{R}\)}{\(\frac{\left (-3+\cos x\right )^{6}} 
 {6} + c,\ c\in \mathbb{R}\)}{\(6\left (-3 +\cos x\right )^{6} + c,\ c\in \mathbb{R}\)}{}{}}
\LANGsk{
\Question{
Vypočítajte \(\int \cos x\cdot \left (-3 +\sin x\right )^{5}\, \mathrm{d}x\) na \(\mathbb{R}\).}
{\(\frac{\left (-3+\sin x\right )^{6}} 
 {6} + c,\ c\in \mathbb{R}\)}{\(6\left (-3 +\sin x\right )^{6} + c,\ c\in \mathbb{R}\)}{\(\frac{\left (-3+\cos x\right )^{6}} 
 {6} + c,\ c\in \mathbb{R}\)}{\(6\left (-3 +\cos x\right )^{6} + c,\ c\in \mathbb{R}\)}{}{}}
\LANGes{
\Question{
Evalúa la siguiente integral en \(\mathbb{R}\). 
\[ 
 \int \cos x\cdot \left (-3 +\sin x\right )^{5}\, \mathrm{d}x
\]}
{\(\frac{\left (-3+\sin x\right )^{6}} 
 {6} + c\text{, }c\in \mathbb{R}\)}{\(6\left (-3 +\sin x\right )^{6} + c,\ c\in \mathbb{R}\)}{\(\frac{\left (-3+\cos x\right )^{6}} 
 {6} + c,\ c\in \mathbb{R}\)}{\(6\left (-3 +\cos x\right )^{6} + c,\ c\in \mathbb{R}\)}{}{}}
\End
